% Gerado por tools/tei2latex.py a partir de revisado.xml. Regenerar, não editar.

%%>>> parte-2

\part*{PARTE QUARTA\\[0.3em]\large THEODICEA}
\addcontentsline{toc}{part}{PARTE QUARTA --- THEODICEA}

\pg{419} PROLOGO

1. \textbf{Theodicea}. --- \emph{Theodicea} é a parte da Metaphysica, que \emph{trata de} \textsc{Deus}, \emph{emquanto é conhecido pela luz da razão}. --- A \emph{Theodicea} costuma chamar-se tambem \emph{Theologia natural} \footnote{A palavra \emph{theodicea} (\emph{defeza de} Deus), embora pela sua etymologia possa denotar uma parte do tratado de \textsc{Deus} (e neste sentido foi empregada por \emph{Leibnitz}, que a formou), todavia hoje toma-se para significar \emph{toda} a sciencia que trata de \textsc{Deus}, servindo-se da luz natural da razão.}.

2. \textbf{Objecto da Theodicea}. --- O objecto da \emph{Theodicea} é \textsc{Deus}, \emph{emquanto conhecido pela luz da razão}. \textsc{Deus} é o Ser necessario, infinitamente perfeito, creador do mundo. --- A \emph{Theodicea}, pois, distingue-se da Theologia \emph{revelada}; porque esta trata de \textsc{Deus}, \emph{emquanto conhecido pela luz da} Fé \footnote{Lembramos que o criterio para distinguir as sciencias não só se baséa no \emph{objecto}, que attingem, mas tambem no \emph{modo}, por que o attingem. D'ahi a distincção entre a Theologia \emph{revelada} e a \emph{natural}.}.

3. \textbf{Excellencia da Theodicea}. --- A excellencia das sciencias \emph{especulativas} depende principalmente da excellencia dos seus objectos. Ora o objecto da \emph{Theodicea} é o mais excellente entre os objectos de todas as outras sciencias \emph{racionaes}; porque \textsc{Deus} é o Ser infinitamente perfeito. Logo a \emph{Theodicea} é a mais excellente entre todas as sciencias \emph{racionaes} \footnote{A excellencia e sublimidade da \emph{Theodicea} é causa da sua difficuldade. \textsc{Deus}, sendo o ente infinito, excede infinitamente a capacidade da nossa intelligencia, que é essencialmente finita. Sabemos que Deus existe, que é o principio e o-fim de todas as coisas, que contém todas as perfeições possiveis e immaginaveis num grau purissimo e infinito; mas esta idea não representa \emph{adequadamente} a essencia de \textsc{Deus}. Todavia um conhecimento\pgn{420}\emph{inadequado} de \textsc{Deus} aperfeiçoa e ennobrece a nossa intelligencia, mais do que o conhecimento \emph{adequado} de todos os seres creados.\par D'ahi se vê que a \emph{Theodicea} é altissima \emph{sabedoria}. Porquanto, a \emph{sabedoria} consiste no conhecimento da causa \emph{absolutamente suprema} dos entes. Uma causa póde ser \emph{relativa} ou \emph{absolutamente suprema}. E' \emph{relativamente} suprema, quando é a mais alta só num certo genero de causas; pois a \emph{causa} divide-se em \emph{material}, \emph{formal}, \emph{efficiente} e \emph{final}. E' \emph{absolutamente} suprema, quando é mais alta do que toda e qualquer outra causa, de modo que todas as outras causas dependem d'ella. Por isso, se o conhecimento das causas \emph{relativamente} supremas é \emph{sabedoria}, o conhecimento da causa \emph{absolutamente} suprema é \emph{altissima sabedoria}. Ora a \emph{Theodicea} trata da causa \emph{absolutamente} suprema, que é \textsc{Deus}. Logo a \emph{Theodicea} é \emph{altissima sabedoria}.\par Dizemos que a \emph{Theodicea} é mais excellente que todas as outras sciencias \emph{racionaes}, pois a \emph{Sagrada Theologia}, que é mais elevada do que a \emph{Theodicea}, não se guia pela luz da \emph{razão}, mas pela luz da Fé. --- A S. Theologia é mais elevada que a \emph{Theodicea} --- pela \emph{certeza} dos \emph{conhecimentos}, pela \emph{extensão} do \emph{objecto}, e pela \emph{excellencia} do \emph{fim}, que é a visão intuitiva de \textsc{Deus}.}.

\pg{420} 4. \textbf{Utilidade da Theodicea}. --- O estudo da \emph{Theodicea} é summamente util para o \emph{individuo} e para a \emph{sociedade}.

a) E' util para o \emph{individuo}. Por quanto, a nossa \emph{intelligencia}, conhecendo \textsc{Deus}, descobre a Verdade infinita, de que emanam todas as verdades, e a Causa primeira, que produziu todas as creaturas; e a nossa vontade, illuminada e guiada por esse conhecimento, une-se a \textsc{Deus} por amor, e nesta união encontra o descanço nas suas tendencias, a paz nas luctas, a força nas difficuldades, o allivio nas dores e miserias da vida.

b) E' util para a \emph{sociedade}. Na verdade, sendo a sociedade a collecção dos individuos, é claro que, se estes conhecerem e amarem a \textsc{Deus}, aquella terá tranquillidade e socego, e, fortalecida pela harmonia e honestidade dos cidadãos, progredirá num caminho de civilisação e de gloria não caduca \footnote{\textsc{Deus} é o principio e o fim de todas as coisas. Logo ninguem póde prescindir d'Elle. --- Se prescindirmos de \textsc{Deus} na vida \emph{particular}, encontraremos o erro, o vicio, e, por consequencia, o remorso. --- Se prescindirmos de \textsc{Deus} na vida \emph{social}, teremos o despotismo, o roubo, a anarchia. --- Se prescindirmos de \textsc{Deus} nas sciencias, os nossos conhecimentos serão \emph{imperfeitos}, porque não baseados na razão suprema de todas as coisas. Com razão \emph{Bacon} dizia: « Quem considerar a dependencia, a serie e o encadeamento das causas, facilmente se convencerá de que, segundo a mythologia dos poetas, o primeiro annel da cadêa natural está preso aos pés do throno de Jupiter » (\emph{De dign. et augm. scient.}, l. I, n. 5).\par Por isso a \emph{Theodicea} tem um caracter essencialmente real e scientifico, porque encerra a explicação de todos os problemas, a resolução de todas as incognitas, o primeiro e mais forte elo da grande cadêa dos seres. --- Aos \emph{materialistas} e \emph{positivistas}, que excluem do quadro das sciencias a \emph{Theodicea}, offerecemos um magnifico trecho de um escriptor muito conhecido d'elles --- \emph{Vacherot}. E' o seguinte: « O que é esta \emph{Theologia} aristotelica, senão o remate\pgn{421}de uma \emph{psychologia} e de uma \emph{physica}, fundadas na experiencia e na analyse da realidade? Onde está o methodo \emph{a-priori}? Onde está a mera especulação neste vasto systema, cuja base é tão solida como elevado é o seu vertice? Se alli a natureza é explicada pela alma, e a alma por \textsc{Deus}, quem levou o espirito a este supremo principio de explicação, senão a physica primeiramente, e depois a psychologia? Isto manifesta a razão por que a doutrina de Aristoteles ficou firme sobre as ruinas de tantas especulações antigas e modernas, como um d'esses monumentos cuja base solida e profunda desafia a acção do tempo e os golpes da tempestade.... \textsc{Deus} é a maior palavra das linguas humanas. Nenhuma poude esquecel-a. Todas, ainda as mais barbaras, celebraram-na com maior ou menor precisão, profundidade e pureza. E' o problema por excellencia das mais sublimes philosophias » (\emph{Le Nouveau Spiritualisme}, pags. 164 e 288).}.

\pg{421} 5. \textbf{Divisão da Theodicea}. --- Dividimos a \emph{Theodicea} em \emph{quatro} capitulos. No \emph{primeiro} capitulo tratamos da \emph{existencia} de \textsc{Deus}; --- no \emph{segundo}, da sua \emph{essencia} e dos seus \emph{attributos}; --- no \emph{terceiro}, das suas \emph{faculdades e operações immanentes}; --- no \emph{quarto}, das suas \emph{operações transeuntes} \footnote{Ninguem se admire se o nosso tratado da \emph{Theodicea} é relativamente breve. Escrevemos sobretudo para estudantes, que se destinam ao estado ecclesiastico e que devem applicar-se, mais tarde, a um estudo completo da S. Theologia. Por isso, para ganhar tempo e evitar repetições, julgamos conveniente apresentar em resumo as principaes verdades, que se referem a Deus e se conhecem pela luz natural da razão, e depois tratar, com o devido desenvolvimento, as outras partes da Philosophia racional.}.

%%>>> cap-6

\chapter*{CAPITULO PRIMEIRO\\[0.3em]\large EXISTENCIA DE DEUS}
\addcontentsline{toc}{chapter}{CAPITULO PRIMEIRO --- EXISTENCIA DE DEUS}

{\small\noindent\textsc{Summario}: \textsc{Summario:} --- Demonstrabilidade da existencia de \textsc{Deus}. --- Demonstração da existencia de \textsc{Deus}. --- A existencia de \textsc{Deus} e o atheismo.\par}

%%>>> art-36

\section*{ARTIGO I.\\ Demonstrabilidade da existencia de Deus.}
\addcontentsline{toc}{section}{ARTIGO I. --- Demonstrabilidade da existencia de Deus.}

6. A \textbf{existencia de Deus} póde ser \textbf{demonstrada pela razão humana}. --- A razão humana póde deduzir da existencia dos effeitos a existencia da causa. Ora todos os seres sensiveis são effeitos. Logo a razão póde deduzir da existencia dos seres sensiveis a existencia da sua causa, que é \textsc{Deus} \footnote{Os systemas, que contestam a nossa proposição são --- o \emph{tradicionalismo}, admittindo que a existencia de \textsc{Deus} só pode ser conhecida pela\pgn{422}\emph{revelação} ou pela \emph{tradição}, --- o \emph{Kantismo}, sustentando que percebemos os \emph{phenomenos} ou as \emph{manifestações}, mas não os \emph{noumenos} ou as \emph{substancias}, --- o \emph{positivismo} e o \emph{sensualismo}, affirmando que a razão humana só attinge objectos sensiveis, --- o \emph{agnosticismo}, ensinando que toda a nossa sciencia se resume e concentra na esphera dos phenomenos e que por isso o \emph{ultraphenomenico} é \emph{incognoscivel}, --- o \emph{immanentismo}, negando á razão a possibilidade de demonstrar a existencia de \textsc{Deus} e attribuindo-a á \emph{consciencia}.\par Por se sufficientemente conhecida a natureza dos primeiros systemas apontados, falta-nos agora explicar a natureza do ultimo, que é o \emph{immanentismo}. Segundo as pretenções d'este systema, a razão não pode demonstrar a existencia de \textsc{Deus}, como não pode demonstrar nenhuma verdade religiosa e moral, porque esta demonstração não pode ser verificada pelos dados da experiencia. As verdades religiosas e moraes não se provam com os raciocinios, como se provam as verdades mathematicas; devem brotar do fundo da nossa alma, do intimo da nossa consciencia. Se, pois, quizermos encontral-as, devemos reflectir sobre os nossos actos, investigar as exigencias do nosso espirito, sondar as necessidades do nosso coração, escutar a voz da consciencia, onde residem os dictames religiosos e moraes. Para esse trabalho é indispensavel uma vontade bôa e recta. Por quanto, os sentimentos religiosos e moraes, embora se encontrem em todos os homens, mesmo nos incredulos, comtudo só nos homens de bôa vontade e de vida irreprehensivel se encontram num estado \emph{explicito}, \emph{claro} e \emph{consciente}. --- Esta é a chamada \emph{philosophia nova}, ou philosophia da \emph{immanencia}, da \emph{acção}, do \emph{voluntarismo}, do \emph{pragmatismo}, que tem por fundadores \emph{Renouvier}, \emph{Boutroux}, \emph{Bergson}, \emph{Mach}, \emph{William James}. E' uma emanação da doutrina de \emph{Kant}, que, depois de ter negado todo o valor á razão \emph{especulativa}, quando esta investiga o \emph{sobresensivel}, ensina que a razão \emph{pratica} deve admittir a existencia de mesmo \emph{sobresensivel}, como postulado ou presuposto necessario para a ordem moral. --- Esta theoria tem sido perfilhada pelos \emph{modernistas}. Dizem elles que a existencia de \textsc{Deus} não pode ser demonstrada por meio dos factos \emph{externos}, isto é, pela existencia e natureza dos seres \emph{creados}, mas só por meio dos factos \emph{internos} e \emph{subjectivos}. A razão não pode demonstrar a existencia de \textsc{Deus}, ou não o pode fazer sem a cooperação da vontade, do sentimento e dos outros factores \emph{subjectivos}, que \emph{supprem a deficencia dos argumentos racionaes}.\par Tendo refutado na \emph{Anthropologia} o \emph{tradicionalismo}, o \emph{Kantismo}, o \emph{positivismo} e o \emph{sensualismo}, resta-nos refutar o \emph{agnosticismo} e o \emph{immanentismo}.\par A pretenção do \emph{agnosticismo}, que ensina ser \emph{incognoscivel} para a razão a existencia de \textsc{Deus}, é uma pretenção absurda. Que as sciencias \emph{experimentaes} não reconheçam senão os \emph{phenomenos}, e que para ellas o \emph{sobresensivel} seja \emph{incognoscivel}, concedemol-o de bom grado. Mas seguir-se-ha d'ahi que, alem do sensivel, não existe mais nada, e que, atraz do phenomeno, não existe o \emph{noumeno}? Não, com certeza. Se o sobresensivel não é objecto das sciencias experimentaes, é objecto de outra sciencia mais elevada. Os adversarios imitam nisso o celebre D. Ferrante, nos \emph{promessos esposos} de Manzoni, que negava a existencia da peste em Milão, só porque esta não cabia em nenhuma das suas categorias! Se a existencia de \textsc{Deus} não é perceptivel pelos\pgn{423}sentidos, é perceptivel pela razão. E a razão não prescinde dos dados da experiencia sensivel, mas funda-se nelles, explica-os e completa-os. Por quanto, no fundo de cada facto sensivel está sempre uma questão \emph{metaphysica}, que diz respeito á \emph{origem} e á \emph{natureza} do proprio facto e que só a razão resolve. --- A Metaphysica, repetimos, não é um tecido de syllogismos forjados \emph{a-priori}; é a investigação das causas \emph{ultimas}, de que dependem os factos. E, como essas causas devem ser e são reaes, porque reaes são os factos, tambem real deve ser e é a Metaphysica, que tem por seu objecto as causas \emph{ultimas} dos entes. --- Concluimos que a existencia de \textsc{Deus}, se é \emph{incognoscivel} pelos sentidos, não é \emph{incognoscivel} pela razão.\par Não menos absurda é a pretenção dos \emph{immanentistas}. Dizem estes --- que \textsc{Deus} não se conhece pela razão, mas pela \emph{consciencia}, --- que é na consciencia que se manifesta o divino, --- que é no intimo da alma que se encontra \textsc{Deus}. --- Mas, como a consciencia é \emph{psychologica} ou \emph{moral}, conforme refere a existencia das nossas operações intellectuaes, ou indica a \emph{conveniencia} d'essas operações com as regras dos costumes, d'ahi a necessidade de perguntarmos aos \emph{immanentistas} se a consciencia, de que fallam, é a \emph{psychologica}, ou é a \emph{moral}? --- Se entendem fallar da consciencia \emph{psychologica}, enganam-se redondamente. Esta consciencia revela-nos as nossas operações, as nossas necessidades e desejos, a existencia do \emph{eu}, a realidade nossa, mas não a realidade de \textsc{Deus}. Diz-nos, é certo, a consciencia que temos a aspiração, o desejo do infinito, de \textsc{Deus}; mas a \emph{aspiração} do infinito não é a \emph{presença} do infinito. Se o infinito se sentisse na nossa consciencia, não haveria aspiração ou desejo; o que se possue, não se deseja. O infinito está fóra e sobre nós; o divino não é \emph{immanente}, é \emph{transcendente}. Fundando-se nestas aspirações ou desejos, a razão, pelo principio de \emph{causalidade}, eleva-se a \textsc{Deus}... Admittimos que \textsc{Deus} está em nós; mas está como causa creadora e conservadora do nosso ser; e esta presença de \textsc{Deus}, como causa do nosso ser, não se percebe pela consciencia, mas pelo raciocinio, emquanto os seres finitos suppõem e exigem a existencia de um ser infinito... Se teimarmos em procurar a \textsc{Deus} \emph{unicamente} com os dados da consciencia, nunca o encontraremos, a não ser que se admitta --- ou que o ser divino se identifica com o nosso ser, como dizem os \emph{pantheistas}, --- ou que o ser divino, embora distincto do nosso, está sempre presente á nossa intelligencia, como objecto de percepção immediata e directa, como affirmam os \emph{ontologistas}, --- ou que o ser divino é representado e unido á nossa intelligencia por uma idéa ou especie \emph{innata}, como querem os \emph{innatistas}. --- Não menos grosseiro é o engano dos \emph{immanentistas}, se ou quando dizem --- que é a consciencia \emph{moral} que revela o \emph{divino}, --- que Deus vive em nós, como suprema lei moral, --- que a consciencia \emph{moral} é o echo da voz de \textsc{Deus}, é o \emph{proprio Deus que falla}. Por quanto, se sentimos que os dictames da lei moral nos \emph{obrigam}, sentimos tambem que essa \emph{obrigação} não deriva da nossa vontade, mas nos é imposta por um superior, que está sobre e fóra de nós; pois ninguem pode ser subdito e superior de si mesmo. Logo da existencia da lei moral podemos e devemos deduzir a existencia de um legislador supremo... Alem d'isso, a consciencia, sendo a regra \emph{proxima} dos nossos actos, é o \emph{echo} da voz de\pgn{424}\textsc{Deus}, que é a regra \emph{ultima}. E' o echo da voz de \textsc{Deus}; mas não é o proprio \textsc{Deus}; porque o echo de uma voz distingue-se da propria voz, e da pessoa que a emitte. O echo da lei moral só pode servir de base para a conclusão da existencia de um supremo legislador. --- Não podemos tambem negar que a vontade tem o seu logar na demonstração da existencia de \textsc{Deus}. Mas é preciso determinar com exactidão o papel da mesma faculdade nesta demonstração. A vontade pode entrar --- ou para levar a intelligencia a \emph{considerar} os argumentos, que provam a existencia de \textsc{Deus}, --- ou para levar a intelligencia a \emph{admittir cegamente} a existencia de \textsc{Deus}. Na primeira hypothese, o influxo da vontade na intelligencia é apenas \emph{extrinseco} e \emph{mediato}, e isto verifica-se em todos os conhecimentos; na segunda, o influxo é \emph{intrinseco} e \emph{immediato}, e isto verifica-se no acto de Fé, que não é determinado pela \emph{evidencia} da verdade proposta, mas pelo imperio da vontade. Ora é claro que a vontade entra na demonstração da existencia de \textsc{Deus} --- não emquanto leva a intelligencia a \emph{admittir cegamente} essa existencia, mas emquanto a determina a considerar os argumentos, que \emph{provam} tal verdade; de modo que o assenso da intelligencia é devido á força dos argumentos, e não ao imperio da vontade. Esta nunca poderá fazer com que seja verdadeira uma proposição, que a intelligencia julgou falsa. Dizer, pois, com os \emph{immanentistas}, --- que o sentimento deve substituir a razão, --- que a emoção é um conhecimento, --- que a disposição subjectiva deve ter e tem o logar da verdade objectiva, é blasonar um erro tão absurdo quão pueril.\par A Fé confirma as conclusões da razão. O Concilio Vaticano emanou a seguinte definição: « Si quis dixerit Deum unum et verum, Creatorem et Dominum nostrum, \emph{per ea quae facta sunt, naturali rationis lumine certo cognosci non posse}, anathema sit ». --- O Concilio baseou a sua definição no Livro da \emph{Sabedoria} (\emph{c}. 13, vv. 1-5) e na Epistola de S. Paulo aos Romanos (c. 1, vv. 20-21). --- Estas auctoridades devem ser recebidas, ao menos pelos \emph{tradicionalistas}, que reconhecem a \emph{revelação} como unico criterio de verdade.}.

7. A \textbf{existencia de Deus deve ser demonstrada}. --- A existencia de um ente deve ser demonstrada, quando ella não é percebida \emph{directa} e \emph{immediatamente} --- nem pelas faculdades \emph{sensitivas}, --- \pg{422} \pg{423}  nem pela \emph{intelligencia}. Ora a existencia de \textsc{Deus} não é percebida \emph{directa} e \emph{immediatamente} --- nem pelas faculdades \emph{sensitivas}, porque \textsc{Deus} não é um ente sensivel, --- nem pela \emph{intelligencia}, \pg{424}  pois, como demonstrámos na \emph{Anthropologia}, a visão \emph{intuitiva} de \textsc{Deus} excede a natural capacidade da nossa intelligencia. Logo a existencia de \textsc{Deus} deve ser demonstrada \footnote{Os que negam a necessidade de demonstrar a existencia de \textsc{Deus} são --- a) os \emph{sentimentalistas}, sustentando que conhecemos \emph{immediatamente} a \textsc{DEus} por meio de um sentido arcano e divino. --- b) os \emph{ontologistas}, admittindo uma \emph{intuição immediata} ou \emph{directa} de \textsc{Deus}.\par Estas duas opiniões são falsas.\par \emph{a}) E' falsa a dos \emph{sentimentalistas}; pois a existencia d'esse sentido arcano e divino é contraria ao testemunho da consciencia.\par b) E' falsa a dos \emph{ontologistas}; pois o \emph{ontologismo}, como vimos, é falso.\par Aqui apenas accrescentaremos que a S. Egreja, Mestra infallivel de verdade, declarou suspeitas e perigosas, entre outras, estas proposições: « Immediata Dei cognitio, habitualis saltem, intellectui humano essentialis est, ita ut sine ea nihil cognoscere possit: siquidem est ipsum lumen intellectuale ». --- « Esse illud, quod in omnibus et sine quo nihil intelligimus, est esse divinum ». (Decr. S. Officii, die 18 Septemb. a. 1861).}.

8. \textbf{A} \textbf{existencia} de \textbf{Deus} não se demonstra \textbf{a-priori}. --- A demonstração diz-se \emph{a-priori}, quando a existencia de um ser se prova \pg{425}  pela sua \emph{causa efficiente}. Ora a existencia de \textsc{Deus} não se prova pela sua \emph{causa efficiente}; porque \textsc{Deus} é a Causa primeira, e por isso absolutamente independente de toda e qualquer causa. Logo a existencia de \textsc{Deus} não se demonstra \emph{a-priori}.

9. A \textbf{existencia} de \textbf{Deus} não \textbf{se demonstra a-simultaneo}. --- A demonstração é \emph{a-simultaneo}, quando a existencia de um ente se prova pela analyse da idea ou noção do proprio ente. Ora este processo, applicado á existencia de \textsc{Deus}, é vicioso; visto que, para uma coisa existir na realidade, não basta que exista no pensamento, pois deduzir da ordem \emph{ideal} a ordem \emph{real} é tirar uma consequencia mais extensa que as premissas. Logo a existencia de \textsc{Deus} não se demonstra \emph{a-simultaneo} \footnote{Esta proposição --- \textsc{Deus} \emph{existe} --- é \emph{immediatamente} evidente \emph{em si mesma}, mas não é \emph{immediatamente} evidente \emph{em relação a nós}.\par a) E' \emph{immediatamente} evidente \emph{em si mesma}. Uma proposição é \emph{immediatamente} evidente \emph{em si mesma}, quando o predicado convém \emph{essencialmente} ao sujeito. Ora a \emph{existencia} convém \emph{essencialmente} a \textsc{Deus}; porque \textsc{Deus} é o ser \emph{necessario}, e não póde deixar de existir. Logo a proposição --- \textsc{Deus} \emph{existe} --- é \emph{immediatamente} evidente em si mesma.\par \emph{b}) Não é \emph{immediatamente} evidente \emph{em relação a nós}. Uma proposição é \emph{immediatamente} evidente em relação a nós, quando directamente e sem necessidade de raciocinio conhecemos a essencia do sujeito e do predicado; assim a proposição --- \emph{o todo é maior que a parte} --- é \emph{immediatamente} evidente \emph{em relação a nós}, pois basta enunciar as palavras \emph{todo} e \emph{parte}, para logo percebermos que o \emph{todo} é maior que a \emph{parte}. Ora a essencia de \textsc{Deus}, não sendo por nós conhecida senão depois de um raciocinio, e ainda assim muito imperfeitamente, não se nos manifesta logo á primeira vista como exigindo a existencia. Logo a proposição --- \textsc{Deus} \emph{existe} --- não é \emph{immediatamente} evidente \emph{em relação a nós}. --- Cf. S. Thom., --- \emph{Sum}. \emph{Th}., p. I, q. 2, a. 1; --- c. \emph{Gentes}, l. I, cc. 10, 11; 1 \emph{Sent}. Dist. 3, q. 1, a. 2; \emph{De verit}., q. 10, a. 12.\par Os mais celebres argumentos \emph{a-simultaneo} em favor da \emph{existencia} de \textsc{Deus} são os de \emph{Santo Anselmo}, \emph{Descartes} e \emph{Leibnitz}.\par a) \emph{Santo Anselmo} raciocina d'este modo: Uma proposição é \emph{immediatamente} evidente \emph{em relação a nós}, quando a sua verdade se nos manifesta pela simples enunciação dos termos. Ora a verdade d'esta proposição --- \textsc{Deus} \emph{existe} --- torna-se-nos manifesta pela simples enunciação dos termos. Com effeito, quando se profere o nome de \textsc{Deus}, immediatamente entendemos que \textsc{Deus} é o Ser maior que se póde imaginar, e que por isso deve existir na \emph{realidade}; porque, se existisse unicamente na intelligencia, não seria o Ser maior que se póde imaginar, pois superior a Elle seria o ser, que existisse na intelligencia e tambem na \emph{realidade}.\par Critica. Ainda que todos admittissem que \textsc{Deus} é o ser maior que se póde imaginar, todavia d'ahi só poderia concluir-se que \textsc{Deus} existe na \emph{intelligencia}, mas não na \emph{realidade}; porque, como dissemos, deduzir da ordem \emph{ideal} a ordem \emph{real} é tirar uma consequencia, que não está incluida nas premissas\pgn{426}. --- Para provar que \textsc{Deus} existe tambem na \emph{realidade}, deve demonstrar-se que existe na \emph{realidade} um ser, que é o maior que se póde imaginar. --- Cf. S. Thom., c. \emph{Gentes}, 1. I, c. 11.\par b) \emph{Descartes} diz: Todo o attribuito, que está contido na idea \emph{clara} e \emph{distincta} de um sujeito, deve affirmar-se do mesmo sujeito. Ora na idea \emph{clara} e \emph{distincta} de \textsc{Deus} está contida a existencia \emph{real}; porque \textsc{Deus} é infinitamente perfeito, e a existencia \emph{real} é uma perfeição. Logo deve affirmar-se de \textsc{Deus} a existencia \emph{real}, isto é, \textsc{Deus} existe \emph{realmente}.\par Critica. O argumento de \emph{Descartes} tem o mesmo defeito que o de \emph{Santo Anselmo}; pois passa da ordem \emph{ideal} para a \emph{real}. Na verdade, tudo o que está contido na idea \emph{clara} e \emph{distincta} de um sujeito deve attribuir-se ao sujeito na mesma \emph{ordem}, em que está o sujeito; de modo que, se o sujeito estiver na ordem \emph{real}, o predicado tambem estará na ordem \emph{real}, --- mas, se o sujeito estiver na ordem \emph{ideal}, o predicado tambem existirá necessariamente na ordem \emph{ideal}. Logo na \emph{idea} de \textsc{Deus} está contida a existencia \emph{ideal}, e não a \emph{real}. --- Lembramos que, para \emph{Descartes}, a idea do \emph{infinito é innata} (o que é falso), e que a verdade da proposição maior do seu syllogismo se baséa na veracidade do proprio \textsc{Deus} (o que constitue um \emph{circulo vicioso}).\par c) \emph{Leibnitz} apresenta este argumento: O ente infinitamente perfeito é possivel. Se é possivel, existe na \emph{realidade}; porque a existencia \emph{real}, por ser uma perfeição, pertence necessariamente á essencia do Ente infinitamente perfeito. Logo o Ente infinitamente perfeito existe na \emph{realidade}.\par Critica. Este argumento é inconcludente pelo mesmo defeito, que apontámos nos dois antecedentes. Do facto de \emph{concebermos}, como \emph{possivel}, o Ente infinitamente perfeito, cuja essencia exige a existencia, só resulta que não podemos \emph{conceber}, como \emph{possivel}, esse ente, sem que o \emph{concebamos} como existente; mas não se segue que a realidade deva corresponder ao que \emph{concebemos} na intelligencia.}.

10. A \textbf{existencia} de \textbf{Deus demonstra-se a-posteriori}. --- A demonstração \emph{a-posteriori} faz-se, quando, pela existencia do effeito, \pg{426}  se prova a existencia da causa respectiva. Ora a existencia de \textsc{Deus} demonstra-se pela existencia das creaturas; porque estas, por serem \emph{effeitos}, suppõem uma causa proporcionada, que é \textsc{Deus}. Logo a existencia de \textsc{Deus} demonstra-se \emph{a-posteriori} \footnote{Cf. S. Thom., --- \emph{Sum}. \emph{Th}., p. I, q. 2, a. 2; --- \emph{c}. \emph{Gentes}, l. I, c. 12.\par A demonstração da existencia de \textsc{Deus} basea-se no \emph{principio de causalidade}. E' por isso que os \emph{materialistas}, os \emph{positivistas} e todos os \emph{incredulos}, no intuito de negar esta verdade fundamental, ou, ao menos, de prescindir d'ella, atacam a noção de \emph{causa} e esse principio de \emph{causalidade}. --- Não repetimos tudo o que na \emph{Ontologia} escrevemos em defeza do dito principio. Limitamo-nos a observar que o principio de \emph{causalidade} é uma verdade, que a intelligencia descobre e formula á vista das mudanças, dos phenomenos, dos factos, que se succedem, uns aos outros, no mundo visivel. Ora quando um phenomeno começa, percebe-se immediatamente que está relacionado com um outro, de que deriva. Esta relação não é apenas de \emph{successão}, mas de \emph{dependencia}. Tal dependencia é uma coisa real, objectiva, existente nos proprios seres creados; a intelligencia \emph{descobre-a}, mas \emph{não a inventa}. Por isso o principio de \emph{causalidade} tem um valor real e objectivo; porque, embora este juizo --- \emph{todo o effeito tem a sua causa} --- seja um producto da razão, comtudo a \emph{razão não o produz senão fundada nos dados da experiencia sensivel}. A experiencia apresenta os factos; a razão descobre o nexo, que os liga. ---\pgn{427}Advirta-se que o principio de \emph{causalidade} não se applica a todo o ser que \emph{existe}, mas a todo o ser que \emph{começa} a existir, e por isso a Causa \emph{primeira} não tem causa.\par Os \emph{Kantistas} não negam o principio de \emph{causalidade}, mas dizem que elle só tem valor, quando se trata de dois phenomenos, ambos sujeitos á observação sensivel, porque, em tal caso, podemos averiguar a dependencia de um phenomeno do outro, --- mas não tem valor, quando se trata de duas realidades, uma das quaes não está sujeita á experiencia, porque, em tal caso, não podemos conhecer a \emph{dependencia} ou a \emph{origem} de uma realidade da outra, e por isso nem a razão de \emph{causa} e de \emph{effeito}; o que se verifica no nosso caso, porque, se o mundo está sujeito á experiencia, \textsc{Deus} escapa a toda a experiencia. Por isso não podemos basear-nos no principio de \emph{causalidade}, para deduzir da existencia do mundo a existencia de \textsc{Deus}.\par Respondemos que, para o principio de \emph{causalidade} ter todo o seu valor, não é necessario que ambas as realidades, que se dizem proceder uma da outra, estejam sujeitas á experiencia sensivel. Basta que um phenomeno, um facto, presente aos nossos sentidos, manifeste um caracter de mutabilidade e de contingencia, para que possamos e devamos concluir que esse phenomeno \emph{depende} necessariamente de um antecedente. Pouco importa, para o caso, que esse antecedente seja visivel ou invisivel, natural ou sobrenatural; elle existe, deve existir, porque existe o phenomeno, que não tem em si a razão sufficiente da sua existencia. Ninguem viu a força da \emph{gravitação}; e, comtudo, todos a admittem, porque se revela na queda dos graves. Ninguem viu o talento militar de Napoleão; mas todos o admittem, porque se manifestou nos seus grandes planos de guerra. Deus é invisivel na sua essencia, mas revela-se nas coisas creadas, que suppõem e exigem uma força infinita. Portanto, ou devemos negar o principio de \emph{causalidade}, e temos a completa derrocada de todas as sciencias, --- ou devemos admittir que, em força d'esse principio, podemos e devemos elevar-nos ao conhecimento da existencia de \textsc{Deus}. --- Nem se diga que uma percepção mais profunda e completa dos phenomenos naturaes pode fazer com que se prescinda de uma causa primeira. Os phenomenos naturaes, ainda que podessem ser percebidos em toda a sua comprehensão, nunca podem perder os \emph{essenciaes} caracteres de \emph{contingencia} e de \emph{mutabilidade}, e por isso da sua \emph{dependencia} de uma causa primeira. Pelo contrario: quanto mais se conhece o \emph{movimento}, a \emph{harmonia}, a \emph{finalidade} dos seres visiveis, tanto mais necessaria apparece a existencia de um supremo e omnipotente Ordenador.\par Objectam: Todos os nossos conhecimentos são \emph{relativos} ao mundo \emph{visivel} e á intelligencia \emph{finita} do homem. Logo não podemos conhecer a \textsc{Deus}, que se suppõe \emph{invisivel} e \emph{infinito}.\par Resposta. E' certo que todos os nossos conhecimentos, \emph{que adquirimos}\pgn{428}\emph{de um modo directo e immediato}, se referem ao mundo visivel. Mas d'ahi não se segue que a nossa razão não possa, nem de um modo \emph{indirecto} e \emph{mediato}, adquirir outros conhecimentos, relativos ao mundo invisivel. Os caracteres de contingencia e de mutabilidade, que observamos no mundo visivel, levam-nos necessariamente a admittir a existencia de um Ente invisivel, necessario e immutavel. --- Como tambem é certo que todos os conhecimentos são proporcionados á capacidade da intelligencia cognoscente e ao modo por que os objectos se tornam presentes á mesma intelligencia, de modo que só apprehendemos as notas do objecto, que actuaram na faculdade, ficando incognosciveis as que escaparam ao alcance d'ella. Mas esta \emph{relatividade}, se prova que não podemos \emph{comprehender} a Deus, porque a intelligibilidade de Deus excede infinitamente a capacidade da intelligencia creada, não prova que não possamos, \emph{de algum modo}, conhecer o omnipotente Creador de todas as coisas.\par Devemos, porém, advertir que, para os sequazes do \emph{agnosticismo}, \emph{relatividade} significa \emph{subjectivismo}; e por isso quando dizem que os nossos conhecimentos são \emph{relativos}, entendem que são \emph{subjectivos}, em quanto nos fazem conhecer as coisas não como são em si mesmas, mas como são no nosso espirito, ás quaes nada corresponde na realidade. Não nos detemos na refutação d'este erro, a qual já está feita. Só accrescentaremos que, se todos os nossos conhecimentos são \emph{subjectivos}, são \emph{subjectivos} não só os que se referem ao mundo invisivel, mas tambem os que se referem ao visivel, e por isso, se não existe \textsc{Deus}, nem existe coisa alguma no mundo. --- Nem se diga que os conhecimentos relativos ao mundo \emph{visivel} são \emph{objectivos}, porque podem ser verificados pela experiencia, pelos factos: o que se não pode fazer com os conhecimentos relativos ao mundo \emph{invisivel}, a Deus. Negamos que os nossos conhecimentos relativos a Deus não se possam verificar pelos factos. Na verdade, como provamos nós que \textsc{Deus} existe? Provamo-lo pela origem e natureza do mundo, do movimento, da vida, da ordem, da finalidade; porque a origem e a natureza d'estas coisas suppõem a existencia de uma causa extramundana. Ora o mundo, o movimento, a vida, a ordem, a finalidade, são factos, e factos grandiosos, admiraveis, que actuam nos nossos sentidos e despertam a admiração na nossa intelligencia. Portanto o conhecimento de Deus é tão \emph{objectivo}, como \emph{objectivos} são os factos, que nos servem de degrau para nos elevarmos a Elle.\par Continuam: A existencia de Deus não pode ser demonstrada pelos effeitos. Na verdade, entre a causa e o effeito deve existir uma proporção; ora nenhuma proporção pode existir entre as creaturas, que são seres finitos, e Deus, que é o Ser infinito.\par Resposta. Demonstra-se a causa pelos effeitos por dois modos. Se os effeitos possuem a mesma natureza da causa (o que se verifica nas causas univocas), então não só levam ao conhecimento da existencia, mas tambem da \emph{natureza} da causa, de modo que adquirimos um conhecimento \emph{perfeito} da causa. Se os effeitos não têm a mesma natureza da causa (o que se verifica nas causas \emph{equivocas}), então só levam ao conhecimento da existencia da causa, e não da \emph{natureza}, e por isso o conhecimento da causa é \emph{imperfeito}\pgn{429}, e o grau d'esta imperfeição é proporcionado ao grau de dessemelhança entre os effeitos e a causa. Por isso, \emph{todo} o effeito tem com a sua causa uma proporção de \emph{ordem}, emquanto é produzido por ella; mas nem \emph{todo} o effeito tem com a causa uma proporção de \emph{natureza}. Appliquemos estes principios ao nosso caso. Os seres visiveis, sendo finitos e contingentes, não exprimem toda a virtude da sua causa, não possuem a mesma \emph{natureza} de \textsc{Deus}, e por isso não nos levam ao conhecimento \emph{adequado} da \emph{natureza} de \textsc{Deus}. Mas, sendo \emph{effeitos}, e effeitos que só podem derivar de um Ente necessario e omnipotente, levam certamente ao conhecimento da \emph{existencia} d'esse Ente. --- Todavia, embora o conhecimento, que temos da \emph{natureza} de \textsc{Deus}, não seja \emph{adequado}, mas \emph{imperfeito}, comtudo é \emph{positivo}; pois as coisas creadas são semelhanças \emph{positivas}, embora \emph{imperfeitas}, das divinas perfeições, visto que todo o agente tende a imprimir no effeito a propria \emph{semelhança}.}.

%%>>> art-37

\section*{Artigo II.\\ Demonstração da existencia de Deus.}
\addcontentsline{toc}{section}{Artigo II. --- Demonstração da existencia de Deus.}

11. \textbf{Especies de argumentos para} a \textbf{demonstração da existencia de Deus}. --- A existencia de \textsc{Deus} demonstra-se \emph{a-posteriori}. Ora os argumentos, que demonstram \emph{a-posteriori} a existencia de \pg{427} \pg{428}  \textsc{Deus}, são de tres especies: \emph{metaphysicos}, \emph{physicos} e \emph{moraes}. Os \emph{metaphysicos} basêam-se na propria \emph{natureza} das coisas, --- os \emph{physicos}, na \emph{ordem} que existe no mundo e, em particular, nas \pg{429}  \emph{operações} dos agentes \emph{irracionaes}, --- os \emph{moraes}, na existencia da lei \emph{moral} e no \emph{consenso} de todos os povos \footnote{Na classificação dos argumentos, que demonstram a existencia de \textsc{Deus}, em \emph{methaphysicos}, \emph{physicos}, e \emph{moraes}, conforme se baséam na \emph{essencia} das coisas, ou na \emph{ordem} existente no mundo e nas \emph{operações} dos agentes \emph{irracionaes}, ou na lei \emph{moral} e no \emph{consenso} dos povos, seguimos a doutrina dos mais sensatos auctores.\par Alguns dividem esses argumentos em \emph{cosmologicos}, \emph{psychologicos} e \emph{moraes}; conforme os factos, em que se fundam, são \emph{cosmologicos} (baseados nas leis \emph{physicas}), \emph{psychologicos} (baseados nas leis \emph{intellectuaes}), ou \emph{moraes} (baseados nas leis \emph{moraes}). --- Esta divisão tambem é acceitavel.\par \emph{Kant}, dizendo que o espirito humano se eleva a \textsc{DEus} por meio de um d'estes tres principios: de \emph{identidade}, de \emph{causalidade} ou de \emph{finalidade}, reduz as provas a tres classes: \emph{ontologicas}, \emph{cosmologicas} e \emph{physico-teleologicas}, conforme se baseam no principio de \emph{identidade}, de \emph{causalidade} ou de \emph{finalidade}. --- A doutrina de \emph{Kant} é inacceitavel. O principio, em que se baseiam todas as provas em favor da existencia de \textsc{Deus}, é um só: o \emph{principio de causalidade}, e a distincção das provas deriva da natureza dos \emph{factos}, aos quaes o dito principio se applica. O principio de \emph{identidade} não prova coisa alguma, como o proprio \emph{Kant} confessa; o principio de \emph{finalidade}, longe de ser a base do argumento \emph{teleologico}, é apenas um corollario d'esse argumento.\par A existencia de \textsc{DEus} é uma verdade tão necessaria e fundamental, que recorre em todos os grandes problemas da philosophia. Nem admira. Não encontrando nos agentes finitos a causa \emph{sufficiente} dos seres, a razão deve elevar-se a uma causa \emph{suprema} e \emph{infinita}. \emph{A origem} do \emph{mundo} (\emph{Cosm}., \emph{ger}.\pgn{430}Cap. I, art. V), á \emph{finalidade} nas operações dos seres irracionaes (\emph{Cosm}. \emph{ger}., Cap. II, art. III), a ordem \emph{physica} com as suas \emph{leis} (\emph{Cosm}. \emph{ger}., Cap. II, art. IV e V), a \emph{origem} e o \emph{desenvolvimento da vida} (\emph{Cosm}. \emph{esp}., Cap. II, art. III), á \emph{origem} da \emph{alma} e do \emph{corpo} humano (\emph{Anthr}., Sec. IV, Cap. I), tudo o que nos cerca, com os seus caracteres de \emph{contingencia} e de \emph{mutabilidade}, proclama a existencia de uma Causa infinitamente sabia, poderosa e boa --- a existencia de \textsc{DEus}.\par Escutámos esta voz maravilhosa da creação e podemos dizer, com immensa satisfação da nossa alma, que procurámos sempre pôr em relevo a racionabilidade e a necessidade da existencia de um Ente supremo, e mostrar a fatuidade, os erros e as contradicções dos adversarios.\par Mas isso não obsta a que façamos aqui, como no logar proprio, uma synthese do que dissemos em todo este trabalho, indicando os principaes argumentos em defeza da \emph{existencia} de Deus. --- Cf. \emph{Sum. Th.}, p. I, q. 2, a. 3; --- \emph{c. Gentes}, l. I, c. 13.}.

\textbf{12. Argumentos metaphysicos}. --- Os argumentos \emph{metaphysicos}, em favor da \emph{existencia} de \textsc{Deus}, baséam-se em que, --- a) se no mundo existe o movimento ou mudança, deve existir um motor \emph{immovel}, --- \emph{b}) se existem as causas \emph{segundas}, deve existir a causa \emph{primeira}, --- c) se existem seres \emph{contingentes}, deve existir um ser \emph{necessario}, --- \emph{d}) se existem varios e diversos graus de perfeição, deve existir um ser infinitamente perfeito, --- e) se existem verdades \emph{absolutas} e \emph{necessarias}, deve existir uma intelligencia \emph{absoluta} e \emph{necessaria}.

\pg{430} a) As substancias creadas estão sujeitas a continuos movimentos ou mudanças. Ora os movimentos, sendo passagens da \emph{potencia} para o \emph{acto}, devem ser determinados --- não pelas proprias substancias que se movem, pois ninguem dá o que não tem, --- mas por um motor distincto, pois o ente, dotado apenas de \emph{potencia}, não adquire o \emph{acto}, se lh'o não communica um outro ente, que já o possue. Este motor --- ou é \emph{movido}, --- ou é \emph{immovel}. Se é \emph{immovel}, é \textsc{Deus}. Se é \emph{movido}, subiremos de motor em motor, até encontrarmos necessariamente um motor \emph{immovel}; porque, se não existisse este motor \emph{immovel}, não poderiam existir os motores \emph{moveis}, que não têm em si a razão do seu movimento, e por isso não haveria movimento nas substancias creadas. Logo existe \textsc{Deus} \footnote{Aqui a palavra \emph{movimento} toma-se em sentido mais lato, emquanto exprime não só o \emph{movimento local}, mas toda e qualquer \emph{mudança} ou \emph{passagem da potencia para o acto}. D'este modo o argumento exposto equivale a este: O mundo é \emph{mutavel}. Ora todo o ente \emph{mutavel} depende do ente \emph{immutavel}. Logo existe o ente \emph{immutavel}. O ente \emph{immutavel} é \textsc{Deus}. Logo existe \textsc{Deus}. --- Ora, encontrando-se a \emph{mudança} não só nos corpos mas tambem nos espiritos, o argumento póde desdobrar-se; pois, baseando-nos tanto no movimento dos \emph{corpos} como no dos \emph{espiritos}, podemos elevar-nos ao descobrimento de uma causa, que move tudo, sendo ella \emph{immovel} ou \emph{immutavel}.\par Eis os dois argumentos:\par a) A experiencia attesta que as substancias materiaes, por exemplo, os planetas, estão sujeitas a continuos movimentos. Ora todo o corpo que se move, é movido por um agente; pois o corpo, sendo por si indifferente tanto para a quietação como para o movimento, não póde determinar-se por si ao movimento. Logo ou devemos admittir uma serie infinita de motores \emph{moveis}, ou deve existir um motor \emph{immovel}, que seja a causa de todo o movimento dos corpos. Ora uma serie infinita de motores \emph{moveis} repugna; porque numa serie \emph{infinita} de motores \emph{moveis} não haveria um \emph{primeiro} motor, e assim não existiria uma causa proporcionada ou adequada de todos os movimentos\pgn{431}nos corpos. Logo deve existir um motor \emph{immovel}. Este motor \emph{immovel} deve ser infinitamente perfeito; pois, se o não fosse, se lhe faltasse alguma perfeição, poderia tender para ella, e assim passaria da potencia ao acto, e não seria \emph{immovel}. O ser infinitamente perfeito é \textsc{Deus}. Logo existe \textsc{Deus}.\par Este argumento, muito popular, pode desenvolver-se \emph{scientificamente} da maneira seguinte. Existe no mundo o movimento. Para a explicação da existencia do movimento adduzem-se duas hypotheses: a \emph{mechanica} e a \emph{dynamica}. Ora nenhuma d'estas hypotheses pode, na explicação do movimento, prescindir da intervenção de um primeiro Motor. --- Não pode prescindir a hypothese \emph{mechanica}. Na verdade, para os sequazes d'esta hypothese, a materia é \emph{meramente passiva}, de modo que os corpos são movidos por um impulso extrinseco, mas elles não movem, e só communicam aos outros corpos o movimento recebido. Se todos os corpos, de que se compõe o universo, são meramente passivos, é necessario que a força impulsiva do movimento venha de fóra, e proceda de um primeiro Motor, que não seja movido. Nem se diga que o movimento da materia teve um começo \emph{necessario}, ou que é \emph{eterno}; porque um \emph{começo necessario} é uma contradicção (o que começa é necessariamente contingente), e um movimento \emph{eterno} da materia, embora não envolvesse repugnancia, não deixaria de exigir um primeiro motor, desde que se supponha a materia meramente \emph{passiva}. --- Não pode prescindir a hypothese \emph{dynamica}. Com effeito, para os sequazes d'esta hypothese, a materia não é meramente passiva, mas é essencialmente dotada de actividade, em virtude da qual um corpo pode, em certas circumstancias, mover outro corpo, \emph{attrahendo-o}. E' certo que o mundo inteiro se move pela \emph{lei da attracção}; mas a \emph{attracção} não é o \emph{primeiro} motor, nem é o motor \emph{sufficiente} do nosso systema solar. Não é o \emph{primeiro} motor; porquanto, a força de attracção, de que é dotada a materia, não está sempre e necessariamente \emph{em acto}, mas, pelo contrario, tende para o repouso das partes numa unica massa, e por isso é indifferente para o acto e para a potencia, e, não contendo em si a razão da passagem da potencia para o acto, exige e suppõe uma causa que a determine e mova. Não é o motor \emph{sufficiente} do nosso systema solar; porquanto este systema deve ser mantido em equilibrio e regulado nos seus movimentos por duas forças --- uma \emph{centripeta} (sem a qual não haveria unidade no systema) e outra \emph{centrifuga} (sem a qual todos os corpos confundir-se-hiam numa unica massa), e se a força \emph{centripeta}, da qual deriva a attracção, pode dizer-se essencial á materia, não se pode dizer essencial a força \emph{centrifuga}, a qual por isso deve ser accrescentada e unida por um motor extranho e superior á propria materia.\par Dizem que este argumento, se prova a existencia de um Ente superior ao mundo e causa de todo o movimento, não prova a existencia de um Deus tão perfeito, como o concebem os espiritualistas. --- Respondemos que o nosso argumento, embora não provasse senão a \emph{existencia} de um Ser supremo, já teria alcançado o seu fim, que é o de demonstrar a mesma \emph{existencia}, e não a \emph{natureza} do Ente supremo. São duas questões distinctas. --- Todavia devemos admittir que este argumento, que prova a necessidade de um \emph{primeiro Motor}, contém, \emph{em germe}, a razão de todos os attributos da Divindade.\pgn{432}Na verdade, se a materia, por ser indifferente ao repouso e ao movimento, tem a razão do seu movimento não em si mas numa causa exterior, segue-se que tal causa existe fóra da materia, e por isso é immaterial, intelligente, eterna, necessaria, e tambem infinita e unica. Porquanto, o primeiro Motor, não podendo ser movido, é \emph{todo acto}, sem mistura de potencialidade, é acto \emph{purissimo}, isto é, infinitamente perfeito e é unico.\par Os melhores naturalistas e philosophos modernos reconhecem esta verdade. Citaremos apenas \emph{Couchy}, \emph{Clarke}, \emph{B. S. Hilaire}, \emph{P. Janet}. No fim, ouviremos um insuspeito, \emph{Rousseau}.\par \emph{Couchy} diz: « O unico ser, que póde produzir as forças physicas, é o Ente necessario. A força é a expressão da sua vontade. As diversas forças do equilibrio e movimento são causas secundarias. A gravitação universal, a attracção dos corpos na razão inversa do quadrato das distancias, que os separam, é uma lei que exprime a vontade divina » (\emph{Comptes rendus hebd. de l'Acad. des sciences}, t. XXI, Paris, 1845).\par \emph{Clarke}, discipulo de \emph{Newton}, fallando da \emph{attracção}, de que depende todo o movimento dos corpos celestes, escreve: « Ella póde ser o effeito de um impulso, mas não póde ser material, é antes o effeito de uma causa immaterial » (\emph{Physica de Rohault}, trad. por \emph{Clarke}, II, c. XI).\par \emph{Barthélemy Saint-Hilaire} diz: « Concebe-se que a materia siga hoje as leis infalliveis, que a regem, e conserve uma regularidade eterna; mas foi necessario um impulso primitivo, que ordenou tudo para a inexhaurivel serie dos tempos. Até as anomalias, que offerece o admiravel systema dos ceus, attestam a presença immanente e indefectivel d'Aquelle que os fez. O sobrenatural está em tudo... Só é necessario que a sciencia se resigne a resolver certos problemas por outros meios, que não sejam os de uma experiencia impossivel; e o da origem de todas as coisas é um d'esses problemas, a que só se renuncia por timidez, pensando que se pratica uma reserva prudente. \emph{A questão da origem é inevitavel}, e de nada valeria querer illudil-a ». (\emph{Journal des Savants}, 1862, pag. 607).\par \emph{P. Janet} diz: « A experiencia não mostra senão que um corpo é movido por outro, e este por outro, e isto ao infinito. Ora, se fóra d'esta cadêa não existe uma causa motora, teremos uma serie de movimentos sem causa e sem razão » (\emph{Le matérialisme contemporain}, 1875, pag. 58).\par \emph{Rousseau} escreve: « Vejo a materia ora em movimento e ora em repouso; d'ahi concluo que nem o repouso nem o movimento lhe são essenciaes. Mas o movimento, sendo uma acção, mostra ser effeito de uma causa... Conceber a materia productora do movimento é conceber um effeito destituido de causa... Não é evidente que, se o movimento fosse essencial á materia, um não poderia separar-se da outra, mas estaria sempre unido com ella, conservando-se sempre no mesmo grau, sempre o mesmo em cada uma das partes, não poderia augmentar nem diminuir, e não se poderia conceber a materia em repouso?!... Quando me dizem que o movimento não é \emph{essencial} á materia, mas \emph{necessario}, pretendem embaraçar-me com palavras, que poderiam facilmente confutar-se, se tivessem algum sentido. Porque o movimento da materia --- ou deriva da propria materia, e então lhe é essen-\pgn{433}cial, --- ou deriva de uma causa extrinseca, e já não é necessario á propria materia, senão emquanto a causa motora actua nella » (\emph{Emilio}, t. III).\par Na \emph{Cosmologia geral} tambem desenvolvemos este ponto.\par b) No mundo não só existe o movimento \emph{local}, mas tambem o movimento \emph{espiritual}; e este, como aquelle, suppõe e exige necessariamente a existencia de um primeiro motor \emph{immovel}. Na verdade, se reflectirmos em nós mesmos, percebemos que a nossa alma tambem se move, emquanto passa continuamente da potencia para o acto. O acto é um complemento, uma addição que aperfeiçõa a potencia. Ora este acto não pode ser e não é extrahido do fundo da potencia, porque ninguem dá o que não tem; mas deve provir de um agente extrinseco, de um motor, que possua \emph{actualmente} o que falta á potencia e que seja capaz de lh'o communicar. Este motor é a propria alma, que se move a si mesma emquanto uma faculdade, que está \emph{em acto}, move uma outra, que está \emph{em potencia}. Todavia a alma não pode ser o \emph{primeiro} motor. O primeiro acto da faculdade, que move todas as outras, não pode ser o effeito do impulso das outras faculdades, que estão em potencia, mas deve ser o effeito de um motor distincto e superior á propria faculdade. Este motor --- ou é o primeiro Motor, --- ou é um motor intermedio, movido, finalmente, por esse primeiro Motor \emph{immovel}.\par Este argumento, que se refere ao movimento da alma humana em geral, conserva toda a sua força, quando se refere ao movimento da \emph{intelligencia} e da \emph{vontade}.\par a) Dá-se o movimento na \emph{intelligencia}. Ainda que a intelligencia mova a vontade, e esta mova aquella, comtudo, por ser absurdo o processo ao infinito, devemos admittir que o primeiro acto da intelligencia precedeu o primeiro acto da vontade. Diz S. Thomaz: « Non oportet procedere in infinitum, sed statur in intellectu, sicut in primo. Omnem enim voluntatis motum necesse est quod praecedat apprehensio; sed non omnem apprehensionem praecedit motus voluntatis » (\emph{Sum. Th.}, p. I, q. 82, a. 4, ad 3). Mas de que modo passa a intelligencia da potencia para o acto? Espontaneamente? Não; porque faltaria a razão sufficiente da passagem. Pela acção do objecto intelligivel? Tambem não; porque, sendo este objecto formado pela intelligencia agente, a acção do objecto é \emph{posterior} á da intelligencia. A causa motora deve ser superior ao objecto e a sujeito, ao intelligivel e ao intelligente; e esta causa é Deus. Aristoteles, ainda que pagão, exprime esta verdade nas seguintes magnificas palavras: « Quod est principium motus in anima? Patet vero, quemadmodum in universo est Deus, sic etiam in anima; moventur enim omnia ab inexistente in nobis Numine. Intellectus autem principium non intellectus sed praestantius quidpiam est. Quid autem et scientia praestantius est et ratione, nisi Deus? (\emph{Moral. Eudem.}, l. 7, c. 14). D'onde se vê a verdade do aphorismo: \emph{Eu penso; logo Deus existe.}\par \emph{b}) Não só o acto da intelligencia, mas tambem o da \emph{vontade} exige e suppõe a existencia de um primeiro Motor. A vontade é movida necessariamente pelo bem absoluto. Esse movimento provoca uma tendencia \emph{universal, necessaria, indestructivel}. Ora um effeito, dotado de semelhantes qualidades, como é essa tendencia, não pode derivar de uma causa imaginaria ou phantastica, mas deve derivar de uma causa real e objectiva, que deve\pgn{434}identificar-se com o proprio bem absoluto; porque todo o effeito suppõe uma causa proporcionada. Ora o bem absoluto é Deus. Logo Deus existe.\par Como se vê, o principio --- \emph{tudo o que se move é movido por um motor distincto e, em ultima analyse, immovel} --- é muito fecundo e serve para a intelligencia e resolução de innumeraveis problemas.}.

\emph{b}) No mundo existem muitas e diversas \emph{causas efficientes}. Estas causas estão por tal modo subordinadas entre si, que formam uma serie, na qual uma depende de outra, a inferior \pg{431} \pg{432}  é \emph{effeito} da superior. Ora, nesta serie de \emph{causas efficientes}, deve existir uma causa \emph{primeira}, que seja só \emph{causa} e não \emph{effeito}, que seja causa de todas as causas; visto que, se não existisse uma \pg{433}  causa \emph{primeira}, não poderiam existir, como existem, as causas \emph{segundas}, pois as causas \emph{segundas} são \emph{effeitos}, e todo o effeito \pg{434}  exige necessariamente uma causa proporcionada. A causa \emph{primeira} é \textsc{Deus}. Logo existe \textsc{Deus} \footnote{Os \emph{materialistas}, para destruirem a força d'este argumento, recorrem a uma serie \emph{infinita} de causas \emph{efficientes}, das quaes uma produza outra, sem sermos obrigados a admittir uma \emph{primeira} causa \emph{efficiente}.\par O recurso é vão. Porquanto, essa serie \emph{infinita}, embora fosse \emph{possivel}. estaria disposta ou em linha \emph{recta}, ou em linha \emph{circular}. Ambas as hypotheses são absurdas. --- E' absurda a \emph{primeira} ; porque, se estivesse disposta em linha \emph{recta}, de modo que a causa inferior fosse produzida pela superior, essa serie não teria uma causa \emph{primeira} (porque o infinito não tem principio). Não tendo uma causa primeira, não teria a \emph{segunda}, nem a \emph{terceira}, nem causa alguma; visto que numa serie de causas, dispostas em linha \emph{recta}, a inferior não póde existir sem a superior. Logo deve admittir-se, fóra d'esta serie, uma causa primeira. --- E' absurda a segunda hypothese. Na verdade, se a serie \emph{infinita} das causas \emph{efficientes} estivesse disposta em linha \emph{circular}, um e o mesmo ente seria causa e effeito de outro ente. Na verdade, supponhamos uma linha \emph{circular}, formada por a, b, c, d... E' evidente que, sendo a causa de b, b de c, c de d, d de a, d seria causa e effeito (ainda que remoto) de a; o que é absurdo.\par Os adversarios dizem que a \emph{primeira} causa \emph{efficiente} é a \emph{materia}.\par Esta asserção tambem é absurda. Por quanto, como provámos, a materia não é nem póde ser eterna, mas foi produzida, e por isso não póde ser a causa \emph{primeira} de todas as coisas.\par Note-se a differença entre este e o precedente argumento. Aqui, \textsc{Deus} considera-se como a causa \emph{efficiente} de tudo o que existe; além, \textsc{Deus} considerou-se tambem como causa \emph{final} de todas as coisas, emquanto tudo tende para \textsc{Deus}, como Ente infinitamente intelligivel e amavel.}.

c) Os entes visiveis são todos \emph{contingentes}, isto é, existem mas podiam deixar de existir. Se são \emph{contingentes}, foram produzidos por outro ente; porque o \emph{contingente}, podendo existir e deixar de existir, não contém em si a razão sufficiente da propria existencia. Este outro ente, que produziu o \emph{contingente}, ou é tambem \emph{contingente}, ou é \emph{necessario}. Se é \emph{necessario}, é \textsc{Deus}. Se é \emph{contingente}, foi produzido por outro, e este por outro, até se encontrar inevitavelmente o ente \emph{necessario} ; porque, se não existisse o ente \emph{necessario}, não poderia existir nenhum ente \emph{contingente}, pois repugna uma serie de entes \emph{dependentes} sem um \emph{primeiro}, de que todos dependam. Logo existe \textsc{Deus} \footnote{Em poucas palavras: O mundo é \emph{contingente}. Ora o ente \emph{contingente} depende do ente \emph{necessario}. Logo existe um ente \emph{necessario}. O ente necessario é \textsc{Deus}. Logo existe \textsc{Deus}.\par \emph{Vazherot} afirma muitas vezes (sem nunca o demonstrar) que o mundo,\pgn{435}se é \emph{contingente} na \emph{apparencia}, \emph{é necessario} e \emph{absoluto} na \emph{realidade}. --- Esta affirmação é absurda. O mundo, como demonstrámos na \emph{Cosmologia geral}, \emph{é essencialmente contingente} no seu conjuncto e em cada uma das suas partes.\par \emph{Kant} diz que a prova, deduzida da idea do ser \emph{necessario}, não tem valor; porque, como a prova \emph{ontologica}, passa da ordem \emph{ideal} para a \emph{real}, apoiando-se num principio meramente \emph{subjectivo}, que é o de \emph{causalidade}.\par Rejeitamos com \emph{Kant} a prova \emph{ontologica} acerca da existencia de \textsc{Deus}, mas negamos que a prova \emph{cosmologica} se reduza á \emph{ontologica}. Na verdade, a prova \emph{cosmologica} não consiste, como pensa \emph{Kant}, em deduzir da \emph{idea} de um ser \emph{necessario} a sua existencia \emph{real}, mas em deduzir da existencia \emph{real} dos seres \emph{contingentes} a existencia \emph{real} do ser \emph{necessario}, que é \textsc{Deus}. --- Além d'isso, negamos que o principio de \emph{causalidade} seja meramente \emph{subjectivo}; elle é tambem, e primeiramente, \emph{objectivo}, porque não é lei do nosso espirito, senão porque é lei das coisas.}.

\pg{435} d) Os entes creados, embora possuam todos uma perfeição \emph{finita}, não a possuem todos no mesmo grau, mas um é \emph{mais} perfeito, e outro o é \emph{menos}. Isto mostra que a perfeição das creaturas é \emph{participada}; porque, se o não fosse, não haveria razão por que numa creatura a perfeição fosse \emph{maior}, e noutra \emph{menor}. Sendo \emph{participada}, a perfeição das creaturas deve derivar de um ente, e de um ente, que a possua pela \emph{sua propria essencia}. Ora a perfeição, que um ente possue pela sua \emph{propria essencia}, deve ser \emph{infinita}; porque, se fosse \emph{finita}, seria \emph{participada}. Logo existe um ente \emph{infinitamente} perfeito, de que derivaram, como de um oceano immenso, as muitas e diversas perfeições das creaturas. O ente \emph{infinitamente} perfeito é \textsc{Deus}. Logo existe \textsc{Deus} \footnote{Em poucas palavras: Os entes creados são todos \emph{finitos}. Ora a existencia dos entes \emph{finitos} suppõe necessariamente a existencia do ente \emph{infinito}. Logo existe o ente \emph{infinito}. O ente \emph{infinito} é \textsc{Deus}. Logo existe \textsc{Deus}.\par Objectam: A diversidade das perfeições, que se encontram nos seres finitos, não demonstra a existencia de um ente infinito.\par A objecção já está refutada no texto. Entre os seres finitos, um é \emph{mais} ou \emph{menos} perfeito que o outro. Ora uma coisa diz-se \emph{mais} ou \emph{menos} perfeita, emquanto \emph{mais} ou \emph{menos} se approxima de um ente \emph{absolutamente perfeito}: assim uma photographia diz-se \emph{mais} ou \emph{menos} perfeita, emquanto \emph{mais} ou menos se approxima do original, que é uma coisa \emph{absoluta}. Logo deve existir um ente \emph{absoluta} e \emph{infinitamente} perfeito, que seja a causa das \emph{varias} e \emph{limitadas} perfeições, que se encontram nas creaturas. --- O \emph{relativo}, que se não refira ao \emph{absoluto}, é uma contradicção.}.

e) A intelligencia humana conhece verdades \emph{absolutas} e \emph{necessarias}, tanto da ordem \emph{especulativa} como da \emph{pratica}. Taes verdades devem ter uma \emph{razão sufficiente}. Esta \emph{razão sufficiente} não se encontra nos \emph{seres creados}, nem na propria \emph{intelligencia humana}. Não se encontra nos \emph{seres creados}; pois estes, sendo \emph{relativos} e \emph{contingentes}, não podem produzir o \emph{absoluto} e o \pg{436}  \emph{necessario}. Nem se encontra na \emph{intelligencia creada}; pois esta não forma as verdades \emph{absolutas} e \emph{necessarias}, descobre-as apenas e recebe-as como regra do seu pensamento. Logo a \emph{razão sufficiente} d'aquellas verdades deve encontrar-se numa \emph{intelligencia absoluta} e \emph{necessaria}; pois ellas são entidades \emph{ideaes}, \emph{absolutas} e \emph{necessarias}. A intelligencia \emph{absoluta} e \emph{necessaria} é a de \textsc{Deus}. Logo existe \textsc{Deus} \footnote{E' certo que, além da ordem \emph{physica}, composta de substancias materiaes, existe uma ordem \emph{logica}, composta de verdades \emph{necessarias}, \emph{absolutas}; por quanto estas verdades não podem ser formadas nem mudadas pela nossa inteligencia. Esta faculdade limita-se a descobril-as, e descobre-as como independentes de toda a experiencia e de todo o pensamento creado. Logo ou devemos negar a verdade do principio \emph{de} razão \emph{sufficiente}, ou deve existir uma intelligencia \emph{necessaria} e \emph{absoluta}, em que estas verdades encontram o seu fundamento. Ora o principio \emph{de razão sufficiente} é verdadeiro. Logo deve existir uma intelligencia \emph{absoluta} e \emph{necessaria}, que é a de \textsc{Deus}.\par Nem se diga que, neste argumento, passamos da ordem \emph{ideal} para a \emph{real}. Por quanto, as verdades \emph{absolutas} e \emph{necessarias} são \emph{objectivas}; pois, como dissemos, ellas são verdades independentemente do nosso pensamento, não podemos formal-as nem mudal-as, são sempre o que são, embora não pensemos nellas e as contestemos. Sendo \emph{objectivas}, devem encontrar na \emph{realidade} um fundamento. Este fundamento não se encontra nas coisas mutaveis, nem na intelligencia creada. Logo deve encontrar-se num Ente \emph{necessario} e \emph{eterno}, em \textsc{Deus}.}.

13. \textbf{Argumentos physicos}. --- Os argumentos \emph{physicos} basêam-se em que, se existe uma ordem maravilhosa no mundo e nas operações dos agentes \emph{irracionaes}, deve existir um ordenador infinitamente sabio e poderoso.

a) Uma ordem maravilhosa e constante existe entre as muitas e diversas substancias, que compõem o universo. Esta ordem exige uma causa infinitamente sabia e poderosa; pois só essa causa, --- por ser \emph{infinitamente sabia}, podia conhecer a natureza de todas as substancias e o posto de cada uma d'ellas no plano da creação, e escolher as leis mais aptas para a constancia e regularidade da mesma ordem; --- e, por ser \emph{infinitamente poderosa}, podia executar esse plano. A causa infinitamente \emph{sabia} e \emph{poderosa} é \textsc{Deus}. Logo existe \textsc{Deus} \footnote{Não insistimos neste ponto; porque na \emph{Cosmologia geral} provámos que a ordem \emph{physica}, com as suas \emph{leis}, é effeito de uma intelligencia \emph{infinita}.\par \emph{Kant} diz: A ordem \emph{physica} é effeito de uma intelligencia, mas não de uma intelligencia \emph{infinita}. Logo não demonstra a existencia de \textsc{Deus}.\par O que dizemos no texto mostra a falsidade da affirmação de \emph{Kant}. --- Além d'isso, a intelligencia, que estabeleceu a ordem no mundo, só podia\pgn{437}ser a do Ente, que \emph{criou} o proprio mundo. Ora o ente, que \emph{criou} o mundo, deve ter um poder \emph{infinito}. O ente, dotado de poder \emph{infinito}, deve ter uma natureza \emph{infinita} e por isso uma intelligencia \emph{infinita}.\par Os materialistas objectam: Esta ordem e regularidade, que existe em todo o universo e em cada substancia, não é effeito de uma causa intelligente e deliberada, mas é obra do \emph{acaso}.\par Resposta. O \emph{acaso}, sendo \emph{privação} de intelligencia e deliberação, é \emph{nada}. Ora é absurdo dizer que o \emph{nada} produziu a \emph{realidade}. Logo a ordem, que existe no mundo, sendo uma \emph{realidade}, não foi produzida pelo \emph{acaso}, mas por um Ente real, infinitamente sabio e poderoso.}.

b) Não só a ordem, que existe no universo, mas tambem a que existe nas operações de cada agente \emph{irracional}, proclama \pg{437}  a existencia de uma Intelligencia soberana e infinita. Na verdade, todo o agente \emph{irracional}, nas suas operações, tende para um fim, empregando para isso os meios mais efficazes, e executando um plano com maravilhosa exactidão, constancia e uniformidade. Este fim, esta escolha de meios e este plano, devem ser effeitos de uma intelligencia \emph{infinita}; porque só uma intelligencia \emph{infinita} podia determinar a cada agente \emph{irracional} um fim proporcionado, modelar planos complicados e ao mesmo tempo perfeitos, e, para os realizar, servir-se da forças cegas da natureza. A intelligencia \emph{infinita} é propria de \textsc{Deus}. Logo existe \textsc{Deus} \footnote{Tambem este ponto ácerca da \emph{finalidade} nas operações dos agentes \emph{irracionaes} foi sufficientemente tratado na \emph{Cosmologia geral}.}:

14. \textbf{Argumentos moraes}. --- Os argumentos \emph{moraes} fundam-se em que, --- a) se existe uma lei \emph{moral}, deve existir um legislador, --- b) se todos os povos admittiram a existencia de um Ente supremo, Este deve existir.

a) A consciencia attesta-nos que temos gravada no nosso coração uma lei moral, á qual não nos podemos subtrahir, e que o cumprimento d'essa lei produz tranquillidade, assim como a sua transgressão causa remorso. Ora a existencia d'esta lei, como a respectiva sancção, suppõe a existencia de um legislador. Este legislador não póde ser o proprio homem; aliás seria inexplicavel o sentimento da nossa dependencia de um poder superior. Logo deve ser um Ente superior ao homem e dotado de virtude ou poder \emph{infinito}; porque só um ente infinitamente poderoso podia gravar no coração dos homens essa lei. O Ente superior, de que o homem depende e a que deve obedecer, é \textsc{Deus}. Logo existe \textsc{Deus} \footnote{Este argumento, apresentado tambem por \emph{Kant} para a demonstração da existencia de \textsc{Deus}, será mais desenvolvido na \emph{Moral}.\par Aqui apenas advertimos que a existencia da \emph{lei moral} considera-se como um \emph{facto experimental}, attestado pela consciencia. Provado o facto da existencia da \emph{lei}, legitimamente deduz-se a existencia de um \emph{legislador}.}.

\pg{438} b) Todos os homens, em todos os tempos e em todos os logares, reconheceram e admittiram a existencia de um Ente supremo. Este consenso unanime, constante e universal, não póde ser effeito do \emph{erro}; pois o \emph{erro} é causa de discordias, é mutavel e particular. Logo é effeito da \emph{verdade}; porque só a \emph{verdade} é unanime, constante e universal, e só ella pode levar os homens a mortificar as paixões e renunciar á propria vontade, para se sujeitarem á vontade de um Ente supremo. Logo existe \textsc{Deus}\footnote{O consenso do genero humano, relativamente á \emph{existencia} de um Ente supremo, não podia ser mais constante, nem mais universal. Esta crença num Ente supremo podia ser e foi desfigurada por muitos erros; mas conservou sempre o fundo intacto. Muitos povos ignoraram e ignoram ainda a \emph{essencia} de \textsc{Deus}, mas nenhum ignorou ou ignora a sua \emph{existencia.}\par \emph{Cicero} disse: « Não existe povo algum, por inculto e selvagem, que não tenha fé em \textsc{Deus}, embora lhe não conheça a essencia » (\emph{De nat}. \emph{Deor}. I, 17).\par \emph{Quatrefages} escreve: « Obrigado pelo meu ensino a passar revista a todas as raças humanas, procurei o \emph{atheismo} tanto entre os povos mais selvagens, como entre os mais civilisados. \emph{Não o encontrei em parte alguma, a não ser nalgum individuo ou eschola muito limitada,} como se viu na Europa no seculo passado e como se vê ainda no presente » (\emph{L'espèce humaine}, pag. 355, Paris, 1877)... E' noutro logar não se exprime menos claramente: « Pouco a pouco fez-se a luz, e os povos da \emph{Australia} e da \emph{Melanesia}, os \emph{Boschimans}, os \emph{Ottentotes}, os \emph{Cafres}, os \emph{Bechuanas}, foram eliminados do numero dos povos atheus e reconhecidos como religiosos » (\emph{Rapport}, pag. 410).\par De resto, os \emph{Annaes da propagação da Fé} provam, com razões claras e baseadas na experiencia, que povos atheus não passam de um mytho.\par Objectam: A crença de todos os povos na existencia de um Ente supremo derivou da \emph{ignorancia}, ou da astucia dos \emph{sacerdotes}, ou da prepotencia dos \emph{legisladores}, ou do \emph{medo}, ou dos \emph{prejuizos}.\par Resposta: A crença unanime de todos os povos em \textsc{Deus} não derivou --- a) nem da \emph{ignorancia}, --- \emph{b}) nem da astucia dos \emph{sacerdotes}, --- \emph{c}) nem da prepotencia dos \emph{legisladores}, --- \emph{d}) nem do \emph{medo}, --- \emph{e}) nem dos \emph{prejuizos}.\par a) Não derivou da \emph{ignorancia}. Por quanto, o estudo das causas naturaes é o meio mais facil e efficaz para se conhecer a existencia de \textsc{Deus}; visto que essas causas, sendo \emph{secundarias}, \emph{contingentes}, \emph{imperfeitas}, exigem uma causa \emph{primeira}, \emph{necessaria}, \emph{perfeitissima}, que é \textsc{Deus}.\par A experiencia confirma a these. Na verdade, os maiores naturalistas são os mais convencidos e enthusiastas adoradores de \textsc{Deus}. --- Citaremos apenas os testemunhos de \emph{Linneu} e \emph{Kepler}.\par \emph{Linneu} escreve: « Quando acordei, \textsc{Deus} \emph{eterno}, \emph{immenso}, \emph{omnisciente}, \emph{omnipotente}, acabava de passar. Vi-o ao longe, e fiquei profundamente maravilhado. Segui os vestigios de seus passos, atravez das obras da creação, e por toda a parte, mesmo nas mais pequenas coisas, que poder! que sabedoria! que ineffavel perfeição! O sol e todo o systema planetario, immenso, incalculavel, appareceram-me suspensos pelo Primeiro Motor, Causa das\pgn{439}causas, Governador e Conservador do universo. Todas as coisas são, pois, testemunho da sabedoria divina; a belleza, a harmonia, as justas proporções das creaturas proclamam o poder do grande \textsc{Deus} » (\emph{Syst}. \emph{de la Nat}. p. 10).\par \emph{Kepler} conclue a sua obra, intitulada « \emph{Os quatro livros da harmonia celeste} », por estas sublimes palavras: « E agora só me resta levantar as mãos e os olhos para ao ceu, e dirigir ao Auctor de toda a luz a minha humilde oração: O' Vós, que, pelos esplendores sublimes que derramastes sobre toda a natureza, elevais os nossos desejos até á divina luz da vossa graça, sede bemdito! Senhor e Creador de todas as consolações, que experimentei no extasis, a que me elevou a contemplação da obra das vossas mãos, eis acabado este livro, que contém o fructo do meu trabalho e de toda a intelligencia, que me destes. Proclamei deante dos homens toda a grandeza das vossas obras, mostrando-lhes, na medida das minhas forças, a sua immensa perfeição. Esforcei-me por chegar até á verdade, por a conhecer o mais perfeitamente possivel, e, se me escapou alguma coisa indigna de Vós, fazei que a conheça para a eliminar. Ter-me-hei deixado seduzir pela presumpção, na presença da belleza maravilhosa das vossas obras? Terei eu procurado a minha gloria entre os homens, levantando este monumento, que só devia ser consagrado á vossa honra? Ah! se assim foi, tende piedade e misericordia de mim, e concedei-me a graça de que o meu trabalho não produza mal algum, mas concorra para a glorificação do vosso Nome e para a salvação das almas ».\par Os mesmos sentimentos encontram-se nos escriptos dos mais celebres naturalistas, como \emph{Copernico}, \emph{Galileu}, \emph{Newton}, \emph{Jussieu}, \emph{Cuvier}, \emph{Agassiz}, \emph{Biot}, \emph{Buffon}, \emph{Lavoisier}. \emph{Brongniart}, \emph{Blainville}, \emph{Bertholet}, \emph{Gay-Lussac}, \emph{Elie de Beaumont}, \emph{Thénard}, \emph{Réaumur}, \emph{Barande}, \emph{Milne-Edwards}, \emph{Dumas}, \emph{Chevreul}. \emph{Liebig}, \emph{Flourens}, \emph{Quatrefages}, \emph{Pasteur}, etc. Nem admira; porque, como diz \emph{Herschel}, « quanto mais se extende o campo da Physica, tanto mais numerosas e\pgn{440}inconcussas se tornam as demonstrações da existencia eterna de uma intelligencia creadora e omnipotente; os geologos, os astronomos, os naturalistas trouxeram todos a sua pedra para este grande templo da sciencia, --- templo levantado ao proprio \textsc{Deus} ». --- Não nos cançaremos de repetir com \emph{Bacon}; « A pouca sciencia afasta de \textsc{Deus}, a muita conduz a Elle ».\par b) Não derivou da astucia dos \emph{sacerdotes}; pois a instituição dos \emph{sacerdotes é effeito}, e não \emph{causa}, da crença universal na Divindade.\par c) Não derivou da prepotencia dos \emph{legisladores}. Os legisladores basearam as leis na propria crença religiosa dos seus povos. --- Além d'isso, os povos não teriam acceitado, ao menos em toda a parte, uma crença, que é contraria a todas as paixões. --- Finalmente, as leis variam em todos os povos, mas a crença em \textsc{Deus} foi sempre a mesma por toda a parte.\par d) Não derivou do \emph{medo}. Na verdade, esse \emph{medo} ou era conforme com a razão, ou não. Se não era conforme com a razão, não podia produzir a crença em \textsc{Deus}; pois essa crença é universal e constante, e um effeito universal e constante não póde derivar de uma causa particular e mutavel, como é o \emph{medo}. Se era conforme com a razão, esse \emph{medo} já suppunha a crença em \textsc{Deus}; pois ninguem tem \emph{medo} de um ente, que se não conhece como existente. --- A verdade é que o \emph{medo}, longe de produzir a crença em \textsc{Deus}, deu sempre occasião ao \emph{atheismo}, pois os libertinos, \emph{pelo facto de temer o castigo}, esforçam-se por negar a existencia de \textsc{Deus}.\par e) Não derivou dos \emph{prejuizos}; pois os \emph{prejuizos} variam segundo a diversidade dos povos, ao passo que a crença em \textsc{Deus} é universal e constante.}.

%%>>> art-38

\section*{Artigo III.\\ A existencia de Deus e o atheismo.}
\addcontentsline{toc}{section}{Artigo III. --- A existencia de Deus e o atheismo.}

\pg{439} 15. \textbf{Atheismo}. --- \emph{Atheismo} é o erro dos que não reconhecem ou negam a \emph{existencia} de \textsc{Deus}.

16. \textbf{Divisão do atheismo}. --- O \emph{atheismo} divide-se em \emph{pratico} e \emph{theoretico}. --- O \emph{pratico} é o erro dos que, embora affirmem a existencia de \textsc{Deus} com a \emph{razão}, negam-na com os \emph{factos}, pois vivem, como se \textsc{Deus} não existisse. --- O \emph{theoretico} é o erro dos que negam com a \emph{razão} a existencia de \textsc{Deus}.

17. \textbf{Divisão do atheismo theoretico}. --- O atheismo \emph{theoretico} póde ser \emph{negativo} ou \emph{positivo}. --- O \emph{negativo} é o erro dos que não admittem a existencia de \textsc{Deus}, porque o não conhecem. --- O \emph{positivo} é o erro dos que, tendo conhecido a \textsc{Deus}, estão \emph{séria} e \emph{sinceramente} convencidos de que Elle não existe.

\pg{440} 18. \textbf{O atheismo pratico é possivel}. --- O atheismo \emph{pratico} não é uma depravação da \emph{intelligencia}, mas sim da \emph{vontade}. Ora a nossa vontade está sujeita á depravação, e, para se entregar mais livremente ás desordens, póde deixar de pensar em \textsc{Deus} e de Lhe prestar as homenagens da sua piedade e obediencia. --- Este \emph{atheismo} não só é \emph{possivel}, mas existe realmente\footnote{\emph{Bayle} diz: « Os que não servem a \textsc{Deus}, e que antes o offendem com uma vida cheia de crimes e vicios, são uma especie de atheus, é o atheismo dos libertinos e devassos; porque effectivamente vivem, como se não houvesse \textsc{Deus}, sem medo da sua Justiça, sem gratidão para com a sua Bondade, sem respeito para o seu Nome, sem obediencia ás suas leis, e, ainda que podessem tirar \textsc{Deus} do seu throno para o esmagar sob os seus pés e aniquilar para sempre, não viveriam com mais licença e ousadia. São atheus pelas suas acções; conhecem a \textsc{Deus} com a lingua, mas negam-no com os costumes; fallam como crentes, mas operam como incredulos » (\emph{Pensées div}., t. IV, pag. 94).}.

19. \textbf{O atheismo theoretico-negativo não é possivel nos homens adultos}. --- A existencia de \textsc{Deus}, embora nos não seja \emph{immediatamente} evidente, todavia descobre-se logo que começamos a applicar os primeiros principios da razão. Por quanto, sendo natural e facillimo deduzir da existencia dos effeitos a existencia da causa respectiva, é claro que, contemplando o universo, não podemos deixar de reconhecer uma causa, que o \pg{441}  produziu e formou. Logo o atheismo \emph{theoretico-negativo} é impossivel nos homens \emph{adultos} \footnote{A historia contém exemplos de homens que, isolados, desde crianças, de toda a convivencia social, chegaram, só pela força da razão, ao conhecimento de uma causa productora do mundo.}.

20. \textbf{O atheismo theoretico-positivo é impossivel.} --- E' impossivel que o homem esteja \emph{séria} e \emph{sinceramente} convencido da falsidade de uma coisa, attestada pela voz da natureza. Ora a existencia de \textsc{Deus} é attestada pela voz da natureza; pois é admittida por todos os homens, em todos os tempos, em todos os logares, e só o que deriva da natureza se encontra em todos os homens, em todos os tempos, em todos os logares. Logo é impossivel que o homem esteja \emph{séria} e \emph{sinceramente} convencido de que não existe \textsc{Deus}. --- A experiencia confirma a these; pois os que se dizem \emph{atheus} são \emph{scepticos} ou \emph{pantheistas} \footnote{O nosso seculo foi chamado \emph{era do atheismo}. E' certo que alguns infelizes empregam esforços para destruir, se fosse possivel, a propria idea de \textsc{Deus}. \emph{Fouerbach} diz: « Só o homem é o verdadeiro salvador! Só o homem é o nosso \textsc{Deus}, o nosso Juiz, o nosso Redemptor ». \emph{Rénan} escreve: « \textsc{Deus}, a Providencia, a immortalidade, não passam de palavras velhas e pesadas ». A loja maçonica de Liege, em 1865, dizia: « Não só devemos collocar-nos acima de todas as religiões, mas acima de toda a crença num Deus ».\par Mas as pretendidas \emph{convicções} d'esses desvairados não passam do \emph{duvidas}. O espirito nunca poderá sinceramente convencer-se de que \textsc{Deus} não existe e, pelo menos com as suas duvidas, protestará sempre contra as estultas affirmações do coração. Já nos seus tempos \emph{Seneca} dizia: « Mentem os que dizem estar convencidos de que \textsc{Deus} não existe; pois, se ás vezes o dizem, todavia de noite, quando estão sós, começam a duvidar ».\par Estas \emph{duvidas} ácerca da existencia de \textsc{Deus} derivam da \emph{superficialidade} do pensamento e, sobretudo, da \emph{corrupção} da vontade.\par As \emph{duvidas} derivam da \emph{superficialidade} do pensamento. A este respeito, \emph{Caro} escreve: « Ha entre nós uma classe de espiritos scientificamente desorganisados, habeis para comprehender os phenomenos e as leis, nas diversas ordens da realidade, na natureza ou na historia, mas incapazes, por causa de uma inercia systematica, que se toma por uma força, de se elevar á concepção da \emph{causa}, que é a razão ultima das leis naturaes, assim como as leis são a razão dos factos » (\emph{L'idée de Dieu}, 1873, pag. 42).\par Mas as \emph{duvidas} derivam principalmente da \emph{corrupção} da vontade. Não é na \emph{intelligencia}, que o atheu nega a existencia de \textsc{Deus} ou duvida d'ella, mas é no \emph{coração}. O homem, cuja vida é pura e ordenada, nem sequer duvida da existencia de \textsc{Deus}. \emph{Rousseau} dizia: « Conservai a vossa alma num estado de desejar que exista \textsc{Deus}, e nunca duvidareis da sua existencia ». \emph{La Bruyère} escreve: « Desejava ouvir a um homem sabio, moderado, casto e justo, que \textsc{Deus} não existe ou que a nossa alma não é immortal; elle fallaria sem interesse. Mas esse homem não se encontra ».}.

21. \textbf{Consequencias do atheismo.} --- O \emph{atheismo} é a destruição da \emph{sciencia}, da \emph{sociedade} e da \emph{natureza humana}.

\pg{442} a) E' a destruição da \emph{sciencia}. Na verdade, a sciencia é o conhecimento das \emph{causas} dos seres. Ora, se \textsc{Deus} não existisse, os problemas mais fundamentaes seriam inexplicaveis. Porquanto, --- na ordem \emph{physica}, seria inexplicavel a origem do mundo, da vida, da alma humana, a uniformidade e constancia da ordem e das leis physicas; --- na ordem \emph{logica}, faltaria a base da \emph{possibilidade} das coisas, da \emph{eternidade} e \emph{necessidade} das essencias, da \emph{immutabilidade} da verdade; --- na ordem \emph{moral}, seria inexplicavel a existencia da \emph{lei moral}, a virtude não teria o premio, nem o vicio o castigo, e o bem não differiria do mal.

b) E' a destruição da \emph{sociedade}. Porquanto, uma sociedade não póde subsistir, sem que haja fidelidade e respeito nos subditos, desinteresse e justiça no principe, amor e concordia em todos. Mas uma sociedade sem \textsc{Deus} é o reino da anarchia, tendo por direito a força material, e por lei o despotismo \footnote{\emph{Voltaire} diz (\emph{Dict}. \emph{phil}., a. \emph{athée}): « Eu não quereria ser subdito de um principe atheu, que, se encontrasse o seu interesse em me esmigalhar, com certeza esmigalhar-me-hia. Eu não quereria, se fosse soberano, ter cortezãos atheus, pois, se tivessem interesse em me envenenar, deveria tomar, todos os dias, contra-venenos. E', pois, absolutamente necessario, para os principes e para os povos, que a idea de um Ente supremo, creador, governador, remunerador, seja gravada profundamente em todos os espiritos »}.

c) E' a destruição da \emph{natureza humana}. Com effeito, se \textsc{Deus} não existisse, a nossa alma não seria immortal, e assim o homem seria a mais desgraçada das creaturas; porque, depois de ter soffrido tantas dores, depois de ter suspirado tanto pela felicidade sem nunca a encontrar, o \emph{nada} seria a recompensa dos seus soffrimentos e o descanço das suas tendencias \footnote{Estas consequencias do \emph{atheismo} são tambem provas, ainda que \emph{indirectas}, da nossa these. Comtudo, \emph{Stuart Mill} diz que não temos prova alguma, nem \emph{seriamente provavel}, para demonstrar a existencia de \textsc{Deus}!}.

\pg{443} CAPITULO SEGUNDO

\textbf{ESSENCIA} E \textbf{ATTRIBUTOS DE DEUS}

\textsc{SUMMARIO}: --- Essencia de \textsc{Deus}, seu constitutivo e attributos. --- Attributos de \textsc{Deus} negativos. --- Attributos de \textsc{Deus} positivos.

%%>>> art-39

\section*{Artigo I.\\ Essencia de Deus, seu constitutivo e attributos.}
\addcontentsline{toc}{section}{Artigo I. --- Essencia de Deus, seu constitutivo e attributos.}

\textbf{22. Essencia de Deus.} --- \emph{Essencia} de \textsc{Deus} é o que faz com que Elle seja \textsc{Deus}, e se distinga de todos os outros seres.

\textbf{23. Essencia physica e metaphysica de Deus.} --- A essencia \emph{physica} de \textsc{Deus} é constituida por todas as perfeições, de que a Divindade é dotada. A \emph{metaphysica} é a perfeição fundamental, que \emph{se concebe} como raiz e fonte de todas as outras perfeições divinas, e pela qual \textsc{Deus} se distingue das creaturas. --- A essencia \emph{physica} representa \textsc{Deus} como é em si mesmo; a \emph{metaphysica} exprime \textsc{Deus} como é em relação ao nosso modo de entender \footnote{Lembramos que a essencia --- é \emph{physica}, quando os seus elementos constitutivos são \emph{reaes} e \emph{concretos} (assim o \emph{corpo organico} e á \emph{alma racional} constituem a essencia \emph{physica} do homem), --- é \emph{metaphysica}, quando os seus elementos constitutivos são \emph{logicos} e \emph{abstractos}, denotando o \emph{genero} e a \emph{differença} (assim a \emph{animalidade} e a \emph{racionalidade} constituem a essencia \emph{metaphysica} do homem). --- Da essencia brotam as \emph{propriedades}. Estas distinguem-se \emph{realmente} da essencia \emph{physica}, porque não a constituem mas resultam d'ella (assim a intelligencia e a vontade são realmente distinctas da essencia \emph{physica} da nossa alma), --- mas podem não distinguir-se \emph{realmente} da essencia \emph{metaphysica}, porque, ainda que não possam considerar-se como elementos \emph{constitutivos} da essencia, comtudo podem identificar-se com esses elementos (assim a \emph{immortalidade}, embora não possa considerar-se como um elemento constitutivo da essencia \emph{metaphysica} da nossa alma, comtudo não se distingue \emph{realmente} da mesma essencia). E' claro que não podemos fallar da essencia de \textsc{Deus}, como fallamos da essencia das creaturas. Em Deus, a essencia \emph{physica} significa e abrange toda a perfeição, que compete ao Ente infinito, sem sombra de imperfeição, --- e a essencia \emph{metaphysica} exclue toda e qualquer composição. Por isso o \emph{constitutivo} da essencia \emph{metaphysica} de Deus não é uma coisa realmente distincta das outras propriedades, e a sua relação com estas deve, mais ou menos, conceber-se como se concebe a relação do ente, considerado na sua razão communissima, com as suas propriedades \emph{transcendentaes}.\par Dizemos que a essencia \emph{metaphysica} exprime \textsc{Deus} como Elle é em \emph{relação} ao \emph{nosso modo de entender}. Por quanto, considerando os seres creados,\pgn{444}encontramos em cada um d'elles uma perfeição, que é a primeira a ser conhecida, que constitue o mesmo ser numa determinada especie, distinguindo-o de todos os outros, e que é a raiz e a fonte de todas as outras perfeições. E, como subimos ao conhecimento de \textsc{Deus} por meio das creaturas, procuramos encontrar em \textsc{Deus} essa perfeição \emph{primeira}, \emph{caracteristica e radical}; embora na essencia de \textsc{Deus}, como é em si mesma, não se encontre attributo algum, que seja anterior aos outros attributos e de que estes fluam. --- Este modo de entender a essencia de \textsc{Deus} é o resultado da limitação da nossa intelligencia; não podemos entender o infinito senão por partes, e, ainda assim, com muita imperfeição. Todavia, a propria intelligencia emenda ou explica o seu modo de ver; pois que, depois de ter considerado um attributo divino, como fundamento dos outros, prova que, em \textsc{Deus}, não só não ha distincção alguma entre a essencia e os attributos, mas nem sequer entre os mesmos attributos.}.

\pg{444} \textbf{24. A essencia physica de Deus consiste na infinidade.}

a) Em \textsc{Deus}, o ser (ou a existencia) não é recebido nem subsiste na essencia, mas é \emph{subsistente por si}; visto que o ser de \textsc{Deus} não se distingue da sua essencia, e o que não se distingue de uma coisa, não pode ser recebido nem subsistir nella. Ora o ente, cujo ser não é recebido nem subsiste na essencia, é necessariamente infinito; porque o ser é limitado ou finito só porque e quando é recebido numa certa e determinada essencia. Logo Deus é infinito, isto é, a essencia \emph{physica} de \textsc{Deus} consiste na \emph{infinidade} \footnote{Desenvolvamos este argumento. Uma substancia, como dissemos na \emph{Ontologia}, pode \emph{subsistir}, ou \emph{existir em si}, por dois modos: ou porque é a existencia, ou porque \emph{tem} a existencia. No primeiro caso, a existencia (ou o ser) não é recebida na essencia, ou natureza, como na sua potencia proporcionada, e por isso é acto puro e subsistente por si; --- no segundo, a existencia (ou o ser) é recebida na essencia, como acto e principio de subsistencia, e por isso não é acto puro, nem subsiste por si, mas subsiste na essencia, como na sua potencia. Ora, em Deus, a essencia e o ser, ou a existencia, não se distinguem \emph{realmente}, porque o ser não é accidental, mas é essencial á essencia divina; porque esta é \emph{essencialmente} existente. Se o ser não se distingue \emph{realmente} da essencia divina, Deus, assim como é a sua essencia, isto é, a sua Divindade, assim tambem é o seu proprio ser, ou a sua propria existencia. Mas o ser, que não se distingue da essencia, não é recebido na propria essencia; porque o ser, que é recebido na essencia, deve necessariamente distinguir-se d'esta, como o acto se distingue necessariamente da potencia (assim a côr, recebida numa parede, distingue-se da propria parede). O ser, que não é recebido numa essencia, não é limitado ou determinado a uma certa especie; visto que o ser emtanto se limita ou se determina a uma certa especie, emquanto é recebido numa certa potencia ou natureza (assim o ser \emph{humano} é limitado ou determinado a uma certa especie, e não a uma outra, porque é recebido na \emph{essencia} ou natureza \emph{humana}). Portanto, o ser, que não é recebido numa essencia, não é limitado ou circumscripto, mas é illimitado ou infinito. E, como o \emph{ser}, por ser \emph{acto},\pgn{445}indica e é \emph{perfeição}, segue-se que o ser subsistente por si é a perfeição subsistente, illimitada, infinita; de modo que não ha nem pode conceber-se entitade, perfeição, attributo, que esse ser não possua, assim como a alvura, se fosse subsistente em si, havia de conter tudo o que tem razão de alvor. Deus, portanto, sendo o ser subsistente por si, é infinito na sua perfeição. E' o raciocinio do Doutor Angelico, quando diz: « Deus est ipsum esse per se subsistens; ex quo oportet quod totam perfectionem essendi in se contineat. Manifestum est enim quod si aliquod calidum non habeat totam perfectionem calidi, hoc ideo est, quia calor non participatur secundum perfectam rationem; sed, si calor esset per se subsistens, non posset ei aliquid deesse de virtute caloris. Unde cum Deus sit ipsum esse subsistens, nihil de perfectione essendi potest ei deesse. Omnium autem perfectiones pertinent ad perfectionem essendi; secundum hoc enim aliqua perfecta sunt quod aliquo modo esse habeant. Unde sequitur quod nullius rei perfectio Deo desit » (\emph{Sum. Theolog.}, p. I, q. 4, a. 2). --- Mais adiante accrescenta: « Cum esse divinum non sit esse receptum in aliquo; sed ipse sit suum esse subsistens, manifestum est quod ipse Deus sit infinitus et perfectus » (\emph{Ib.}, q. 7, a. 1). E noutro logar: « Deus est infinitus: quod sic videri potest. Esse enim hominis terminatum est ad hominis speciem, quia est receptum in natura speciei humanae, et simile est de esse equi, vel cujuslibet creaturae. Esse autem Dei, quum non sit in aliquo receptum, sed sit esse purum, non limitatur ad aliquem modum perfectionis essendi, sed totum esse in se habet » (\emph{De Pot.}, q. 1, a. 2).}.

\pg{445} b) A essencia \emph{physica} de \textsc{Deus} é constituida por todas as perfeições, de que a Divindade é dotada. Ora as perfeições, de que a Divindade é dotada, são \emph{infinitas}; visto que \textsc{Deus}, sendo \emph{immutavel} e \emph{necessario}, deve possuir todas as possiveis perfeições; pois, se Lhe faltasse alguma, seria, com relação a esta, \emph{mutavel} e \emph{contingente}, e por isso dependeria de uma causa. Logo a essencia \emph{physica} de \textsc{Deus} consiste na \emph{infinidade}.

c) Se as perfeições, que constituem a essencia \emph{physica} de \textsc{Deus}, fossem \emph{finitas} ou \emph{limitadas}, essa limitação deveria provir --- ou de uma causa \emph{extrinseca}, --- ou da indole da propria essencia divina, em que as perfeições são recebidas. Ora não pode provir --- nem de uma causa \emph{extrinseca}, porque \textsc{Deus}, sendo a Causa \emph{primeira}, não depende de outra causa, --- nem da indole da propria essencia divina, porque esta, sendo perfeita por \emph{si mesma}, e não por \emph{participação}, não pode deixar de possuir a plenitude da perfeição. Logo as perfeições, que constituem a essencia \emph{physica} de \textsc{Deus}, são \emph{infinitas}. Logo a essencia \emph{physica} de \textsc{Deus} consiste na \emph{infinidade} \footnote{A \emph{infinidade} de \textsc{Deus} \emph{indirectamente} demonstra-se pelo facto de se não encontrar n'Elle nenhuma das causas, que limitam as creaturas. Na verdade, as creaturas, que são compostas de \emph{materia} e \emph{forma}, são limitadas, porque a \emph{forma} limita a \emph{materia}, quando a determina a constituir uma certa especie e não outra, --- e a \emph{materia} limita a \emph{forma}, tornando-a \emph{sin}\pgn{446}\emph{gular} e \emph{individua}. Ora \textsc{Deus} não pode ser limitado d'este modo: porque, como provaremos d'aqui a pouco, não é composto da \emph{materia} e \emph{forma}. --- As creaturas todas pertencem a um certo \emph{genero} e a uma certa \emph{especie}, e por isso são limitadas, porque não possuem as perfeições proprias dos outros generos e das outras especies. Ora \textsc{Deus}, sendo o ser \emph{necessario}, \emph{immutavel} e \emph{independente}, transcende todos os generos e todas as especies. --- As creaturas todas são compostas de \emph{essencia} e de \emph{ser}, ou \emph{existencia}, e por isso são limitadas, porque o \emph{ser}, sendo recebido numa essencia \emph{finita}, como é a essencia da creatura, não pode deixar de ser \emph{finito}. Ora, como dissemos, o \emph{ser} divino, \emph{existindo} por assim o exigir a sua \emph{essencia}, subsiste por \emph{si} e não é recebido numa essencia, que o limite. --- As creaturas são limitadas pelo \emph{espaço} e pelo \emph{tempo}. \textsc{Deus}, por ser \emph{immenso}, não é limitado pelo \emph{espaço}, e, por ser \emph{eterno}, não é determinado pelo \emph{tempo}. Logo Deus é \emph{infinito}. --- S. Thomaz prova a mesma verdade, dizendo que Deus é \emph{infinito}; porque a razão da limitação num ente deriva ou da \emph{causa} do proprio ente ou da \emph{essencia}, em que o ser é recebido, mas Deus não tem causa e o seu ser não é recebido na essencia, mas, como provámos, é subsistente por si. Eis as suas palavras: « Ipsum esse, absolute consideratum, infinitum est... Si igitur alicujus esse sit finitum, oportet quod limitetur esse illud per aliquid aliud, quod sit aliqualiter causa illius esse, vel \emph{receptivum} ejus. Sed esse divini non potest esse aliqua causa, quia ipse necesse est esse per seipsum: nec esse ejus est receptivum, quum ipse sit suum esse. Ergo esse suum est infinitum, et ipse infinitus » (\emph{C. Gent.}, lib. I, c. 43).\par Os erros ácerca da \emph{infinidade} de \textsc{Deus} podem reduzir-se a tres.\par \emph{a)} Uns escriptores, não entendendo por \emph{infinito} senão a \emph{quantidade}, disseram que \textsc{Deus} não é infinito, porque uma \emph{quantidade infinita} repugna.\par Esta opinião refuta-se, negando o \emph{supposto}. A \emph{infinidade}, que se attribue a \textsc{Deus}, não é de \emph{quantidade}, mas sim de \emph{perfeição}.\par \emph{b)} Outros, como \emph{Locke}, \emph{Hobbes}, \emph{Robinet}, dizem que Deus é \emph{infinito}, não \emph{absoluta} mas \emph{relativamente}; pois faltam-lhe muitas perfeições.\par Esta opinião é absurda; porque \textsc{Deus}, como provámos, é \emph{essencialmente} infinito nas suas adoraveis perfeições.\par \emph{c)} Outros, como \emph{Hegel}, seguido por \emph{Vacherot}, \emph{Renan} e \emph{Taine}, não fazendo distincção entre a idea meramente \emph{subjectiva} do ente em geral e o Ente necessario, nem entre o conceito de \emph{infinidade} e o Infinito real, dizem que \textsc{Deus} não é um ente \emph{real} nem \emph{pessoal}, mas é a \emph{categoria do ideal}.\par Como é evidente, estas affirmações encerram o \emph{pantheismo}. Um \textsc{Deus} ideal, formado por nós, não é \textsc{Deus}. --- Não nos demoraremos na refutação do \emph{systema} de \emph{Hegel}; pois o refutámos na \emph{Cosmologia geral}, tratando do \emph{pantheismo}. Aqui apenas accrescentaremos que as theorias de \emph{Hegel}, servilmente copiadas ou imitadas por muitos escriptores, são tidas, na propria Allemanha, naquella conta que merecem. \emph{Schopenhauer} diz --- que \emph{Hegel} « é uma \emph{cabeça mediocre},\pgn{447}\emph{que por todos os meios quiz passar por um grande philosopho} », --- e que a philosophia de \emph{Hegel}, como a de \emph{Fichte} e de \emph{Schelling}, \emph{é} « \emph{um} minimum \emph{de pensamento diluido in quinhentas paginas de phraseologia nauseabunda} »!...\par \textsc{Deus}, sendo infinito nas suas perfeições, deve possuir todas as perfeições, espalhadas nas creaturas; pois a causa \emph{efficiente} deve possuir toda a perfeição, que se encontra no effeito. --- E' o processo da \emph{causalidade}.\par Mas esta posse é diversa, segundo a natureza das perfeições. Por quanto as perfeições podem ser \emph{absolutas} ou \emph{mixtas}.\par \emph{Absolutas} são as perfeições, que no seu conceito não encerram imperfeição alguma, nem se oppõem a perfeições mais elevadas, como são a \emph{vida}, a \emph{belleza}, a \emph{bondade}, a \emph{sabedoria}, etc. Estas perfeições attribuem-se a \textsc{Deus} no proprio conceito, que exprimem, mas por um modo mais \emph{eminente}, porque em \textsc{Deus} se encontram por um modo infinitamente mais perfeito e sublime do modo, por que se encontram nas creaturas, pois \textsc{Deus} não é \emph{bom}, \emph{santo}, \emph{amavel}, como é o homem. Por isso os nomes, com que exprimimos as perfeições \emph{absolutas} de \textsc{Deus}, não são \emph{univocos} nem \emph{equivocos}, mas \emph{analogos}. (\emph{Sum. Th.}, p. I, q. 13, a. 5). --- E' o processo da \emph{eminencia}.\par \emph{Mixtas} são as perfeições, que, embora no seu conceito exprimam uma entidade positiva, todavia encerram alguma imperfeição e se oppõem a uma perfeição mais elevada; tal é a \emph{corporeidade}, que, embora signifique uma coisa positiva, todavia encerra alguma imperfeição (pois o \emph{corpo}, por ser \emph{composto}, é \emph{dependente}, \emph{potencial}, \emph{finito}, etc.) e se oppõe a uma perfeição mais elevada, que é a \emph{espiritualidade} (pois o \emph{espirito} é mais nobre que o \emph{corpo}). Estas perfeições não se attribuem a \textsc{Deus} no seu proprio conceito, mas só \emph{eminentemente}, emquanto attribuimos a \textsc{Deus} toda a entidade, que essas perfeições encerram, e afastamos d'Elle toda a imperfeição, de que as perfeições estão cercadas. --- E' o processo da \emph{exclusão}.\par D'ahi se vê que o processo da \emph{exclusão} é, para nós, o mais natural, para nos elevarmos ao conhecimento de \textsc{Deus}. Por quanto, devemos afastar de \textsc{Deus} --- não só as perfeições creadas, que no seu proprio conceito indicam alguma imperfeição, --- mas tambem todo o limite de que estão cercadas as perfeições creadas, que no seu conceito não encerram alguma imperfeição. D'este modo mais sabemos o que \textsc{Deus} não é, do que o que Elle é.\par Todavia, este processo, --- além de manifestar a perfeição infinita de \textsc{Deus}, pois o Ente, de que se afasta toda a imperfeição, é perfeitissimo, --- subministra-nos um verdadeiro conhecimento de \textsc{Deus}, pois tanto melhor conheceremos um ente, quanto mais o distinguirmos de todos os outros seres. --- \textsc{Deus}, pois, não é absolutamente \emph{incognoscivel}, como pretende \emph{H. Spencer}.}.

\textbf{25. A essencia metaphysica de Deus consiste na subsistencia do proprio ser.} --- A essencia \emph{metaphysica} de \textsc{Deus} consiste naquella perfeição, --- que se concebe como raiz e fonte de todas as outras perfeições divinas, e pela qual \textsc{Deus} se distingue das creaturas. Ora a \emph{subsistencia do proprio ser} é a perfeição \pg{446} \pg{447}  --- que se concebe como raiz e fonte de todas as outras perfeições divinas, porque, como provámos, Deus possue todas as perfeições, possiveis e imaginaveis, pelo facto de ser o \emph{proprio ser subsistente}, --- e pela qual se distingue das creaturas, visto que Deus é o Ente infinito, porque \emph{é} o proprio ser subsistente, e as creaturas são finitas, porque \emph{têm} o ser recebido numa determinada essencia. Logo a essencia \emph{metaphysica} de \textsc{Deus} consiste na \emph{subsistencia do proprio ser} \footnote{A conclusão é certa, e é uma consequencia dos principios expostos no numero anterior. O ser divino, pelo facto de não ser recebido em nenhuma\pgn{448}essencia ou potencia, não tem limites, encerra uma perfeição absoluta sem sombra de imperfeição. Repetimos con S. Thomaz: « \emph{Ex hoc quod Deus est ipsum esse subsistens, oportet quod totam perfectionem essendi in se contineat} » (\emph{Sum. Th.}, p. I, q. 4, a. 2). Sendo subsistente, o ser divino distingue-se das creaturas, que não \emph{são} o ser, mas \emph{têm} o ser, recebido numa determinada potencia ou essencia e por isso limitado a uma certa especie. Como o ser é \emph{acto}, o ser que não é recebido numa potencia, mas que é acto \emph{puro}, é necessariamente infinito ou illimitado; ao passo que o ser, que é recebido numa potencia, é limitado por ella. Por isso, não só pode dizer-se que o \emph{constitutivo} da essencia divina consiste em que \textsc{Deus} \emph{é o ser subsistente por si}, mas tambem pode dizer-se que esse \emph{constitutivo} consiste em que \textsc{Deus} é \emph{acto purissimo}, \emph{não recebido em nenhuma potencia}. As expressões são synonimas. --- Concluimos --- que a razão formal e intrinseca, pela qual \textsc{Deus} é unico, eterno, infinito, é porque Elle é o proprio ser subsistente, é acto purissimo, cuja essencia é constituida pelo proprio ser, --- e que, pelo contrario, a razão, pela qual as creaturas são multiplicaveis, imperfeitas, finitas, é porque têm o ser recebido na essencia, e por isso são compostas de potencia e de acto, de essencia e de ser.\par Sendo Deus o proprio ser subsistente, o nome mais proprio de Deus é \textsc{JEHOVAH}, que significa: \textsc{Qui EST} --- \textsc{AQUELLE QUE} É. Por quanto, os nomes significam a forma ou essencia das coisas. Ora, em Deus, a essencia identifica-se com o ser, isto é, o ser é a sua propria essencia, de modo que á pergunta: \emph{o que é Deus?} podemos responder: \emph{Deus é o proprio ser} --- \emph{Deus est ipsum esse}. Logo o nome \emph{Qui est} é muito proprio de Deus. Com effeito, de nenhuma creatura podemos dizer que ella é o proprio ser; porque o \emph{ser}, nas \emph{creaturas}, não é a \emph{essencia}, nem é um \emph{predicado essencial}, mas é distincto da essencia, embora seja recebido nella; e por isso os nomes, que significam as coisas creadas, exprimem a essencia, e não e ser, assim \emph{homem} exprime a essencia de um individuo, isto é, \emph{a humanidade}, mas não exprime o \emph{ser} ou a \emph{existencia} d'esse individuo, que poderia não existir. Escutemos S. Thomaz: « Cum in omni quod est, sit considerare \emph{quidditatem suam}, per quam subsistit in natura determinata, et \emph{esse suum}, per quod dicitur de eo quod est in actu, hoc nomen \emph{res} imponitur rei a quidditate sua, hoc nomen \emph{qui est} vel \emph{ens} imponitur ab ipso actu essendi. Cum autem ita sit quod in qualibet re creata, essentia sua differat a suo esse, res illa proprie denominatur a quidditate sua et non ab actu essendi, sicut homo ab humanitate. In Deo autem ipsum esse suum est sua quidditas; et ideo nomen quod sumitur ab esse, proprie nominat ipsum, et est proprium nomen ejus; sicut proprium nomen hominis quod sumitur a quidditate sua » (I \emph{Sent}, Dist. 8, q. 1, a. 1).\par E não pode deixar de ser assim. Quando digo: \emph{Pedro é homem}, attribuo a esse individuo duas coisas: o \emph{ser}, ou a \emph{existencia}, e a \emph{natureza humana}, isto é, attribuo-lhe o ser limitado e determinado a uma certa especie; mas quando digo: \emph{Deus é aquelle que é,} quero affirmar que Deus não é este ou aquelle ente, mas que é o ser, todo o ser, o proprio ser, o ser infinito, de modo que tudo o que tem razão de ser deve encontrar-se e se encontra\pgn{449}em Deus. Eis o sentido do nome \emph{qui est}, que o proprio Deus manifestou a Moyses. Lê-se no \emph{Exodo}; « Ait Moyses ad Deum: Ecce ego vadam ad filios Israel, et dicam eis: Deus patrum vestrorum misit me ad vos. Si dixérint mihi: quod est nomen ejus? quid dicam eis? Dixit Deus ad Moysen: \textsc{Ego sum qui sum}. Ait: Sic dices filiis Israel: \textsc{Qui EST} misit me ad vos » (Cap. III, vv. 13, 14). --- E' este um nome mais universal e menos determinado que os outros; mas por isso mesmo é o mais proprio. Sendo a essencia divina infinitamente superior á intelligencia creada, o nome mais proprio de Deus é o que melhor traduz essa transcendencia e indefinibilidade, e exclue as imperfeitas determinações das creaturas. Diz S. Thomaz: « Quia divina comprehendi a nobis non possunt, convenientius a nobis significantur per nomina communia, quae indefinite aliquid significant, quam per nomina specialia, quae definite rei species exprimunt; unde hoc nomen \emph{qui est}, quod secundum Damascenum significat substantiae pelagus infinitum, convenientissimum nomen dicitur esse » (\emph{De pot.}, q. 10, a. 1 ad 9).\par Nem se diga que o nome \emph{qui est} não é o mais proprio de Deus, porque não o distingue das creaturas, pois Deus \emph{é} e as creaturas tambem \emph{são}. Por quanto, ha uma grande differença entre o \emph{ser} de Deus e o \emph{ser} das creaturas. O ser de Deus, por ser subsistente em si, identifica-se com a essencia, e por isso se torna individuo e distincto de todos os outros entes; ao passo que o ser das creaturas, por não ser subsistente em si, não se identifica com a essencia, mas é recebido e subsiste na essencia, e por isso se torna individuo e distincto por causa da mesma essencia. Portanto o nome \emph{qui est} é exclusivamente proprio de Deus, porque só em Deus o ser é a mesma essencia, e o nome, como dissemos, exprime a essencia. Ouçamos de novo S. Thomaz: « \emph{Ipsum esse Dei distinguitur et individuatur a quolibet alio esse per hoc ipsum, quod est esse per se subsistens, et non adveniens alicui naturae, quae sit aliud ab ipso esse. Omne autem aliud esse, quod non est subsistens, oportet quod individuetur per naturam et substantiam, quae in tali esse subsistit. Et in eis verum est, quod esse hujus est aliud ab esse illius per hoc quod est alterius naturae} » (\emph{De pot.}, q. 7, a. 2 ad 5).\par *\par Todos os philosophos concordam em que em \textsc{Deus}, considerado segundo o nosso modo de entender, existe um attributo constitutivo da essencia divina, mas divergem na determinação d'aquelle attributo.\par Ha quatro outras opiniões a este respeito.\par \emph{a}) Uns dizem que o attributo \emph{constitutivo} da essencia divina é a somma \emph{actual} de todas as perfeições possiveis.\par \emph{b}) Outros, como \emph{Scoto}, sustentam que a essencia divina não é constituida pela somma \emph{actual}, mas pela \emph{exigencia radical} de todas as perfeições.\par \emph{c}) Outros, seguindo \emph{Billuart}, afirmam que o \emph{constitutivo} da essencia divina é a \emph{intelligencia}.\par \emph{d}) Outros finalmente ensinam que o constitutivo da essencia de \textsc{Deus} é a \emph{aseidade} (\emph{ens a se}), isto é, aquella propriedade que faz com que \textsc{Deus} exista \emph{por si}, sem receber de outra causa o ser que possue.\pgn{450}\par Todas estas opiniões são inacceitaveis.\par \emph{a}) E' inacceitavel a \emph{primeira}. Por quanto, a \emph{infinidade} das perfeições constitue a essencia \emph{physica} de \textsc{Deus}, mas não a \emph{metaphysica}; e aqui procuramos o \emph{constitutivo} da essencia \emph{metaphysica} de \textsc{Deus}.\par \emph{b}) E' inacceitavel a \emph{segunda}. Com effeito, aquella \emph{existencia radical}, longe de constituir a essencia de \textsc{Deus}, suppõe-na já constituida.\par \emph{c}) E' inacceitavel a \emph{terceira}. Na verdade, a \emph{intelligencia} não é a \emph{primeira} coisa, que conhecemos de \textsc{Deus}, nem o distingue de todos os seres creados; pois ha muitas creaturas dotadas de intelligencia. --- Além d'isso, não podemos deduzir da \emph{intelligencia} muitos attributos de \textsc{Deus}, como a \emph{eternidade}, a \emph{immensidade}, a \emph{immutabilidade}, etc.\par \emph{d}) E' inacceitavel a \emph{quarta}. Porquanto, a \emph{aseidade}, embora distinga Deus das creaturas, comtudo não representa o intrinseco, proximo e formal principio, que se concebe como constitutivo da essencia divina, e pelo qual Deus se distingue de todas as outras coisas. Na verdade, a independencia (\emph{esse a se}), embora exclua a causa efficiente, todavia não se concebe como uma coisa \emph{intrinseca} ao proprio ente, nem como a raiz de todos os attributos divinos, porque podemos sempre perguntar: porque Deus é \emph{ens a se}? E a resposta deve ser esta: porque \textsc{Deus} é o ser subsistente por si, e o acto purissimo. Por isso, este attributo, embora não seja o primeiro na ordem do \emph{conhecimento}, é o primeiro na ordem da \emph{realidade}, porque é o que se concebe como uma coisa \emph{intrinseca} e como a raiz e a fonte de todos os outros attributos divinos. --- Notamos que alguns philosophos e theologos tomam estas duas expressões --- \emph{Deus est ens per se subsistens} e \emph{Deus est ens a se} --- como equivalentes, e, dizendo que o constitutivo da essencia divina é a \emph{aseidade}, entendem que é a \emph{subsistencia do ser divino}. Mas, rigorosamente fallando, ha uma differença entre as duas expressões.}.

\textbf{26.} \textbf{Especies de attributos divinos, derivados da essencia}. --- Os attributos divinos, derivados da essencia, dividem-se em \emph{quiescentes} e \emph{operativos}. --- Os \emph{quiescentes} são os que se referem á \pg{448} \pg{449}  \emph{essencia}, emquanto é uma \emph{entidade}. Os \emph{operativos} são os que se referem á \emph{essencia}, emquanto é \emph{principio de operações}, e comprehendem as \emph{faculdades}, com as \emph{operações} respectivas.

\pg{450} Os attributos \emph{quiescentes} dividem-se em \emph{negativos} e \emph{positivos}. Os \emph{negativos} são os que \emph{negam} de \textsc{Deus} as imperfeições das creaturas. Os \emph{positivos} são os que \emph{affirmam} de \textsc{Deus} as perfeições creadas, mas isentas de todo o limite.

Neste artigo e no seguinte tratamos dos attributos \emph{quiescentes}, \emph{negativos} e \emph{positivos}; nos capitulos seguintes trataremos dos attributos \emph{operativos} \footnote{A nossa intelligencia, não podendo, por ser finita, comprehender a Essencia divina, que excede, pela sua perfeição infinita, a capacidade de toda a creatura, procura exprimil-a, de algum modo, por muitos e differentes conceitos. A esta pluralidade e diversidade de conceitos corresponde, da parte de Deus, não uma pluralidade ou diversidade real de perfeições, mas uma plena e infinita perfeição, que equivale a todas as perfeições e que as contém de um modo eminente. --- A esta pluralidade e diversidade de conceitos, que exprimem a unica e simplicissima Essencia divina, correspondem varios nomes, que por isso não são \emph{synonymos}, porque, embora todos exprimam a mesma Essencia, comtudo não a exprimem todos sob o mesmo aspecto, mas um exprime-a sob um aspecto e outro sob outro aspecto, e sempre de um modo \emph{deficiente} e \emph{parcial} (Cf. \emph{Sum}. \emph{Th}., p. I, q. 13, a. 4).}.

%%>>> art-40

\section*{Artigo II.\\ Attributos de Deus negativos.}
\addcontentsline{toc}{section}{Artigo II. --- Attributos de Deus negativos.}

\pg{451} 27. Attributos de \textbf{Deus} negativos. --- Attributos \emph{negativos} são, como dissemos, os que \emph{negam} de \textsc{Deus} as imperfeições proprias das creaturas. --- Os principaes são: a \emph{simplicidade}, a \emph{immutabilidade}, a \emph{eternidade}, a \emph{immensidade} e a \emph{unidade} \footnote{Estes attributos são \emph{negativos} em quanto ao \emph{termo}, mas não em quanto á \emph{significação}, porque significam uma perfeição pura e simplicissima, de que a essencia de \textsc{Deus} é dotada. --- A \emph{infinidade} tambem colloca-se entre os attributos \emph{negativos}; mas d'ella fallámos no artigo precedente.}.

28. Simplicidade de \textbf{Deus}. --- A \emph{simplicidade} é o attributo, que exclue de \textsc{Deus} toda a \emph{composição} e \emph{divisibilidade} \footnote{\emph{Simplicidade}, em geral, é a \emph{negação de partes}; porque todos chamam \emph{simples} o ente, que não tem partes. A \emph{simplicidade} é uma especie de \emph{unidade}, á qual accrescenta a \emph{indivisibilidade}. \emph{Uno} é o ente, que é \emph{indiviso}. Mas o ente \emph{uno} pode ter partes, e, pela divisão, tornar-se multiplice; por isso, se exclue a multiplicidade \emph{actual}, não exclue a multiplicidade \emph{potencial}. Pelo contrario, o \emph{simples}, por não ter partes, exclue a multiplicidade \emph{actual} e a \emph{potencial}. --- A \emph{simplicidade}, considerada em si, não representa nem uma perfeição, nem uma imperfeição. Ás vezes o ente \emph{simples} é mais perfeito que o \emph{composto}, ás vezes é menos perfeito. E' mais perfeito, quando representa um \emph{acto} ou uma \emph{forma} completa e subsistente, dotada de toda a perfeição, propria da sua especie; assim o Anjo é mais perfeito que o homem. E' menos perfeito, quando representa uma \emph{parte} de um todo; assim a alma humana é menos perfeita que o composto humano. Diz S. Thomaz: « Simplicitas \emph{per se} non est causa nobilitatis, sed perfectio; unde ubi perfecta bonitas in uno simplici invenitur, simplex est nobilius, quam compositum: quando autem e contrario simplex est imperfectum, tunc compositum est nobilius, quam simplex, sicut homo nobilior est terra » (4 Sent. \emph{Dist}. 11, q. 2, a. 1 ad 1).\par A' \emph{simplicidade} oppõe-se diametralmente a \emph{composição}. --- Esta pode ser \emph{real}, ou \emph{de razão}. E' \emph{real}, quando a coisa é por si composta de partes, independentemente da nossa reflexão. E' \emph{de razão}, quando a nossa intelligencia, pela sua limitação natural, exprime por varios e diversos conceitos uma coisa, que por si é simplicissima. A composição \emph{real} denota imperfeição na coisa, e por isso não pode admittir-se em \textsc{Deus}; a \emph{de razão} denota apenas imperfeição na nossa intelligencia, e é compativel com a simplicidade de \textsc{Deus} (Cf. \emph{C. Gent.}, 1. II, c. 14, --- \emph{De ver.}, q. 5, a. 2 ad 2). --- A composição \emph{real} subdivide-se em \emph{physica}, \emph{metaphysica} e \emph{logica}. A \emph{physica} é propria do ente, que consta de diversos elementos \emph{physicos}, como são a \emph{materia} e a \emph{forma}. A \emph{metaphysica} é propria do ente, que consta de partes \emph{metaphysicas}, como são a \emph{potencia} e o \emph{acto}, a \emph{essencia} e o \emph{supposto}, a \emph{essencia} e a \emph{existencia}, a \emph{essencia} e os \emph{attributos}, a \emph{substancia} e o \emph{accidente}. A \emph{logica} é propria do ente, que consta de elementos \emph{logicos}, como são o \emph{genero} e a \emph{differença}. --- Ninguem se admire se reduzimos a composição \emph{logica á real}. A composição \emph{logica}, se é \emph{de razão} emquanto ao seu conceito \emph{formal}, porque\pgn{452}os elementos são productos da intelligencia, é, comtudo, \emph{real} emquanto ao \emph{fundamento}, pois presuppõe uma composição \emph{real} no ente composto de \emph{genero} e de \emph{differença}, visto que o \emph{genero} se funda na \emph{potencia}, e a \emph{differença} no \emph{acto}. --- Provaremos que \textsc{Deus} é \emph{simplicissimo}, se provarmos que Elle exclue de si toda e qualquer especie de \emph{composição}.}.

\pg{452} 29. \textbf{Deus} é \textbf{simplicissimo}.

a) Todo o ente composto é \emph{finito}; porque resulta de partes, e estas são essencialmente finitas. Ora \textsc{Deus} é \emph{infinito}. Logo \textsc{Deus} não é \emph{composto}, mas \emph{simplicissimo}.

b) O ente \emph{composto} não só é \emph{posterior}, ao menos por \emph{natureza}, ás partes componentes, mas tambem depende de uma causa \emph{extrinseca}, que reune as mesmas partes componentes. Ora \textsc{Deus} é o ente \emph{primeiro}, a causa independente de toda a causa. Logo \textsc{Deus} não é \emph{composto}, mas \emph{simplicissimo} \footnote{Cf. S. Thom., \emph{Sum. Th.}, p. I, q. 3, a. 7; --- c. \emph{Gent.}, l. I, c. 16 e seg., --- \emph{De pot.} q. 7, a. 1. --- Depois de termos demonstrado em geral que em \textsc{Deus} não ha composição alguma, demonstraremos em particular que em \textsc{Deus} não ha composição \emph{physica}, nem \emph{metaphysica}, nem \emph{logica}.\par a) Em \textsc{Deus} não ha composição \emph{physica}. --- Em \textsc{Deus} haveria composição \emph{physica}, se elle fosse \emph{corpo}, ou \emph{forma corporea}, ou \emph{composto} de \emph{materia} e \emph{forma}. Ora \textsc{Deus} não é \emph{corpo}, nem \emph{forma} de corpo, nem \emph{composto}.\par a) \textsc{Deus} não é \emph{corpo}. O corpo não é a mais nobre entre as substancias. Porquanto, o corpo \emph{vivo} é mais nobre do que o \emph{não-vivo}. Ora o corpo \emph{vivo} é mais nobre do que o \emph{não-vivo}, não emquanto é \emph{corpo} (aliás todos os corpos seriam eguaes em \emph{nobreza}), mas emquanto é \emph{vivo}. Logo o principio, de que deriva a \emph{vida} para o corpo, é mais nobre do que o \emph{corpo}, como tal. Ora \textsc{Deus} é o ser nobilissimo. Logo \textsc{Deus} não é \emph{corpo} (\emph{Sum. Th.}, p. I, q. 3, a. 1).\par b) \textsc{Deus} não é \emph{forma} nem \emph{materia} de algum corpo. --- Não é \emph{forma}. Porquanto, a \emph{forma corporea}, além de ser \emph{finita} e por isso \emph{creada} e \emph{contingente}, depende, pelo menos \emph{extrinsecamente} (como a alma humana), da \emph{materia}. Ora \textsc{Deus} é infinito, Causa primeira, Ente necessario e absolutamente independente. Logo não pode ser \emph{forma} de algum corpo. --- Muito menos pode ser \emph{materia}; pois esta é um ser \emph{potencial}, \emph{finito}, \emph{contingente}, e por isso dotada de propriedades contrarias ás de \textsc{Deus} (\emph{Sum. Th.}, p. I, q. 3 a. 8).\par c) \textsc{Deus} não é \emph{composto} de \emph{materia} e \emph{forma}. Com effeito, todo o \emph{composto}, além de ser \emph{finito}, é \emph{dependente} não só da causa \emph{efficiente}, que reune as partes, mas tambem da \emph{forma}, de que recebe as perfeições. Ora \textsc{Deus} não depende de nenhuma causa, por ser a Causa primeira, e possue todas as perfeições não por \emph{participação}, mas por \emph{essencia}. (\emph{Sum. Th.}, p. I, q. 3, a. 2). --- E', pois, absurdo o \emph{anthropomorphismo}, que attribue a \textsc{Deus} a \emph{forma humana}.\par D'isto resulta que \textsc{Deus} é \emph{purissimo espirito}. \emph{Espirito} é uma substancia independente da \emph{materia}. Ora a substancia divina exclue toda e qualquer mistura de \emph{materia} e de \emph{potencia}. Logo \textsc{Deus} é \emph{purissimo espirito}.\par Totavia, do facto de \textsc{Deus} ser \emph{purissimo espirito}, e não ter \emph{corpo}, não póde concluir-se (como concluem os \emph{pantheistas}) que Elle não é \emph{pessoa}. Por quanto, \emph{pessoa} é a \emph{substancia individua de natureza racional}. Ora \textsc{Deus} é \emph{substancia}, subsistente em si e por si, e esta \emph{substancia} é infinitamente \emph{intelligente}. Logo \textsc{Deus} é \emph{pessoa}. --- E não só é \emph{pessoa}, mas possue a \emph{personalidade}\pgn{453} no grau mais elevado, que póde imaginar-se; porque a \emph{substancia} de \textsc{Deus} é \emph{perfeitissima}, e a sua intelligencia é \emph{infinita}. --- A materia não é sempre elemento necessario da \emph{personalidade}; pelo contrario, sendo a \emph{razão} o elemento principal da \emph{pessoa}, é claro que, quanto mais a razão for independente da materia, tanto mais nobre será a \emph{pessoa}.\par b) Em \textsc{Deus} não ha composição \emph{metaphysica}. --- Na verdade, \textsc{Deus} não é composto de \emph{potencia} e \emph{acto}, de \emph{essencia} e \emph{supposto}, de \emph{essencia} e \emph{existencia}, de \emph{essencia} e \emph{attributos}, de \emph{substancia} e \emph{accidentes}.\par a) \textsc{Deus} não é composto de \emph{potencia} e \emph{acto}. Todo o ente, composto de \emph{potencia} e \emph{acto}, é \emph{mutavel}. Ora \textsc{Deus} é \emph{immutavel}. Logo não é composto de \emph{potencia} e \emph{acto}. --- Por isso \textsc{Deus} é \emph{acto purissimo}, contendo \emph{todas} as \emph{possiveis} perfeições num grau \emph{infinito}; porque, se não possuisse alguma perfeição ou fosse capaz de augmento, estaria em \emph{potencia} para aquella perfeição ou augmento de perfeição, e por isso seria \emph{composto} de \emph{potencia} e \emph{acto}.\par b) \textsc{Deus} não é composto de \emph{essencia} e de \emph{supposto}. Um ente é composto de \emph{essencia} e de \emph{supposto}, quando os elementos, que estão contidos neste, differem, na \emph{especie} ou no \emph{numero}, dos que estão contidos naquella; de modo que não se póde dizer que o \emph{supposto} é a \emph{essencia}. Assim o homem é composto de \emph{essencia} e de \emph{supposto}; porque, além de \emph{todos} os elementos \emph{especificos}, que estão contidos na \emph{humanidade} e que são a \emph{animalidade} e a \emph{racionalidade}, o \emph{supposto humano} contém outros elementos, que são as \emph{propriedades individuaes} e que não estão contidas no conceito \emph{especifico} de \emph{humanidade}; de modo que não podemos dizer que \emph{Pedro} é a \emph{humanidade}, mas só que \emph{Pedro} possue a \emph{humanidade}. Ora as \emph{propriedades individuaes} derivam da \emph{materia}. Logo nos \emph{espiritos puros}, em que não ha \emph{materia}, tudo o que está no \emph{supposto}, encontra-se na \emph{essencia}; de modo que esta póde attribuir-se áquelle por \emph{identidade}, dizendo que o \emph{supposto} é a propria \emph{essencia}. \textsc{Deus} é \emph{purissimo espirito}. Logo, em \textsc{Deus}, não ha distincção entre \emph{supposto} e \emph{essencia}, e por isso \textsc{Deus} não é composto de \emph{supposto} e \emph{essencia}; de modo que podemos dizer que \textsc{Deus} é a \emph{Divindade} (\emph{Sum. Th.}, p. I, q. 3, a. 3; --- c. \emph{Gent.} l. I, c. 21).\par c) \textsc{Deus} não é composto de \emph{essencia} e de \emph{existencia}. Um ente, composto de \emph{essencia} e de \emph{existencia}, isto é, cuja \emph{essencia} se distingue \emph{realmente} da \emph{existencia}, não existe \emph{essencialmente}, por assim o exigir a sua essencia, mas recebeu o \emph{ser} de outra causa. Ora \textsc{Deus} existe \emph{essencialmente}, não recebendo de outra causa o ser que possue. Logo \textsc{Deus} não é composto de \emph{essencia} e de \emph{existencia}. --- Além disso, a \emph{essencia} refere-se á \emph{existencia}, como a \emph{potencia} ao \emph{acto}. Ora \textsc{Deus} não é composto de \emph{potencia} e \emph{acto} (\emph{Sum. Th.}, p. I, q. 3, a. 4; --- c. \emph{Gentes}, l. I, c. 22).\par d) \textsc{Deus} não é composto de \emph{essencia} e \emph{attributos}. A Essencia divina, sendo o proprio ser, possue a plenitude da perfeição, e por isso não ha attributo algum, que Ella não contenha. Ora a essencia de \textsc{Deus}, se se distinguisse \emph{realmente} dos seus attributos, aperfeiçoar-se-hia pelos mesmos attributos, e já não possuiria a plenitude da perfeição. Logo em \textsc{Deus} não ha distincção \emph{real} entre a \emph{essencia} e os \emph{attributos}. Logo \textsc{Deus} não é composto de \emph{essencia} e de \emph{attributos}. --- Alem d'isso, quando os attributos são \emph{realmente} distinctos da essencia, esta refere-se áquelles, como a potencia ao acto. Ora a Essencia divina, por ser o proprio ser subsistente, é acto purissimo, sem sombra de\pgn{454}potencialidade. Diz S. Thomaz: « Perfectio divini \emph{esse} non attenditur secundum aliquod \emph{additum} supra ipsum; sed quia ipsum secundum \emph{seipsum} perfectum est. Bonitas igitur Dei non est aliquid additum suae substantiae, \emph{sed} sua \emph{substantia est sua bonitas} » (\emph{C. Gentes}, l. I, c. 38). --- Se, em \textsc{Deus}, os attributos são uma e a mesma coisa com a essencia, é claro que, assim como dizemos que \textsc{Deus} é a \emph{Divindade}, assim tambem devemos dizer que Elle é a \emph{Verdade}, a \emph{Belleza}, a \emph{Bondade}, a \emph{Misericordia}, etc.\par Se, em \textsc{Deus}, os attributos não se distinguem \emph{realmente} da \emph{essencia}, tambem não se distinguem \emph{realmente} entre si, e por isso, na realidade, um identifica-se com outro. E' um corollario da these precedente. Se cada attributo divino se identifica com a Essencia, é evidente que um não se distingue realmente do outro, porque uma coisa não se distingue realmente de si mesma. A Essencia divina encerra, na sua absoluta simplicidade, todas as perfeições possiveis; e só a imperfeição e limitação da nossa intelligencia considera como derivadas da essencia e distinctas entre si as propriedades, que, na realidade, são uma e a mesma coisa. --- E' claro que isto se refere aos attributos \emph{absolutos} de \textsc{Deus}, e não aos que se dizem \emph{relativos ad intra} e que são oppostos entre si, como são o \emph{Paternidade}, a \emph{Filiação} e a \emph{Processão}; porque estes attributos são realmente distinctos e oppostos e constituem as tres Pessoas da Santissima Trindade --- \textsc{Padre}, \textsc{Filho} e \textsc{Espirito Santo}. Diz S. Thomaz: « Omnia in Deo sunt unum \emph{re}, praeter ingenerationem, generationem et processionem, quae constituunt personas re distinctas » (I \emph{Sent.} \emph{Dist.} 2, q. 1, a. 3 ad 6.\par Mas, se os attributos divinos não se distinguem \emph{realmente} um do outro, porque cada um se identifica com a Essencia divina, comtudo distinguem-se \emph{logicamente}, porque o conceito de um não se identifica com o conceito de outro; assim o conceito de \emph{justiça} não é o conceito de \emph{misericordia}. E esta distincção \emph{logica} não é arbitraria, mas tem o seu \emph{fundamento real} na propria Essencia divina. Porquanto, a cada um d'esses diversos conceitos, relativos a attributos diversos, corresponde, na sua plenitude simplicissima, a unica Essencia divina. Se assim não fosse, isto é, se de nenhum modo existisse em \textsc{Deus} o fundamento d'essa distincção, os nomes \emph{absolutos} e \emph{positivos}, que exprimem as divinas perfeições, seriam \emph{synonymos}; o que é falso; pois esses nomes, embora exprimam todos a Essencia divina, exprimem-na comtudo sob differentes aspectos; assim a \emph{justiça} exprime a Essencia divina emquanto dá a cada um o que merece,a \emph{misericordia} exprime a mesma Essencia divina emquanto tem dó dos nossos males. E' a doutrina de S. Thomaz. Diz assim: « In Deo est sapientia, bonitas, et hujusmodi, quorum quodlibet est ipsa divina essentia, et ita sunt unum re. Et quia unumquodque eorum est in Deo secundum verissimam sui rationem, et ratio sapientiae non est ratio bonitatis in quantum hujusmodi, relinquitur quod sunt diversa ratione non tantum ex parte ipsius ratiocinantis, sed ex proprietate ipsius rei; et inde est quod Ipse non est causa rerum omnino aequivoca, cum secundum formam suam producat effectus similes non univoce sed analogice » (I \emph{Sent}. Dist. 2, q. 1, a. 2; cf. \emph{Sum. Th.}, p. I, q. 13, a. 2). E mais adeante: « Pluralitas nominum non tantum est ex parte intellectus nostri, formantis diversas conceptiones de Deo; sed etiam ex parte ipsius Dei, inquantum scilicet est\pgn{455}aliquid in Deo correspondens omnibus istis conceptionibus, scilicet plena et omnimoda ejus perfectio » (I \emph{Sent}. Dist. 2 q. 1, a. 3).\par e) \textsc{Deus} não é composto de \emph{substancia} e \emph{accidentes}. O \emph{accidente} refere-se á \emph{substancia}, como o \emph{acto} á \emph{potencia}. Ora \textsc{Deus} é \emph{acto purissimo}. --- Além d'isso, o \emph{accidente} aperfeiçoa a \emph{substancia}. Ora a \emph{substancia divina}, possuindo a plenitude da perfeição, não pode receber novas perfeições. Logo \textsc{Deus} não é composto de \emph{substancia} e de \emph{accidentes}.\par c) Em \textsc{Deus} não existe composição \emph{logica}, isto é, de \emph{genero} e \emph{differença}. Por quanto o \emph{genero} refere-se á \emph{differença}, como a \emph{potencia} ao \emph{acto}. Logo \textsc{Deus}, por ser \emph{acto purissimo}, não é composto de \emph{genero} e de \emph{differença}. --- Além d'isso, o ente, composto de \emph{genero} e de \emph{differença}, é \emph{finito}, porque só possue a perfeição propria de uma certa especie de entes. Ora \textsc{Deus} é \emph{infinito}. --- Mais, a composição de \emph{genero} e de \emph{differença}, embora seja \emph{intencional}, comtudo suppõe uma composição \emph{real} de \emph{potencia} e de \emph{acto}; pois o \emph{genero} funda-se na \emph{potencia}, e a \emph{differença} no \emph{acto}. Ora em \textsc{Deus} não ha composição \emph{real}, e por isso não pode haver composição \emph{logica}. --- Se \textsc{Deus} podesse estar subordinado a um \emph{genero}, este seria o \emph{ente} ou a \emph{substancia}. Ora não pode ser o \emph{ente}; porque este não se attribue a \textsc{Deus} e ás creaturas em significação \emph{univoca}; visto que, quando se attribue a \textsc{Deus}, \emph{ente} significa o proprio ser subsistente em si, e, quando se attribue ás creaturas, significa o ser recebido num sujeito ou potencia, que é a essencia. Nem pode ser a \emph{substancia}; porque a \emph{substancia}, que é \emph{genero}, não é o proprio ser, aliás toda a substancia seria o proprio ser, isto é, o seu ser não seria produzido por outra causa; e por isso \textsc{Deus}, que é o proprio ser, não pode estar sujeito ao genero de \emph{substancia}. --- Todavia, \emph{analogicamente} e num sentido lato, podemos dizer que \textsc{Deus} é \emph{substancia}, emquanto esta significa um ente \emph{subsistente}. --- D'ahi se vê que não podemos formar uma \emph{rigorosa} definição de \textsc{Deus}, porque uma definição \emph{rigorosa} é composta de \emph{genero} e de \emph{differença}. Mas nem por isso confundimos \textsc{Deus} com as creaturas. \textsc{Deus} distingue-se das creaturas --- não por differenças especificas e componentes, --- mas pela pureza absoluta, pela substancia simplicissima do proprio ser (\emph{Sum. Th.}, p. I, q. 3, a. 5).}.

\pg{453} 30. \textbf{Immutabilidade de Deus}. --- \emph{Immutabilidade} é o attributo, que afasta de \textsc{Deus} toda a \emph{mudança}. --- A \emph{mudança} pode ser

\pg{454} duplice: \emph{physica} ou \emph{moral}. A mudança \emph{physica} consiste na adquisição ou na perda de alguma entidade ou perfeição; a

\pg{455} \emph{moral} consiste na passagem da \emph{vontade} de um para outro acto, de um para outro proposito \footnote{S. Thomaz trata d'isto --- na \emph{Sum. Th.}, p. I, q. 9; --- q. 14, a. 15; q. 19, a. 7; --- \emph{c. Gentes}, l. I, c. 13 e seg.; --- I \emph{Sent.} \emph{Dist.} 8, q. 3, a. 1 e 2.}.

31. \textbf{Deus} é \textbf{immutavel}.

a) \textsc{Deus} é \emph{absolutamente} immutavel. Todo o ente, que se muda de algum modo, é composto de \emph{potencia} e de \emph{acto}. Ora \textsc{Deus} é acto simplicissimo e purissimo, sem sombra de potencialidade. Logo \textsc{Deus} é \emph{absolutamente} immutavel.

b) \textsc{Deus} é \emph{physicamente} immutavel. Todo o ente, que se muda \emph{physicamente}, ou adquire uma perfeição que não tinha, ou perde uma perfeição que tinha. Ora \textsc{Deus}, sendo o Ente \emph{infinito e necessario}, possue todas as perfeições possiveis, e não póde perdel-as, aliás perderia a essencia, com que as perfeições se identificam. Logo \textsc{Deus} é \emph{physicamente} immutavel.

\pg{456} c) \textsc{Deus} é \emph{moralmente} immutavel. Com effeito, a mudança \emph{moral}, ou de \emph{vontade}, deriva de tres causas, --- ou da mudança \emph{physica} do sujeito, que, por isso, começa a desejar uma coisa, pela qual tinha repugnancia, --- ou da mudança do \emph{conhecimento}, emquanto a intelligencia descobre uma causa, que antes ignorava, e assim determina a vontade a mudar de proposito, --- ou da leviandade e inconstancia da propria \emph{vontade}, que, sem motivo sufficiente, passa de um proposito para outro. Ora nenhuma d'estas tres causas pode admittir-se em \textsc{Deus}, que, sendo infinitamente perfeito, não está sujeito a mudanças physicas, nem á ignorancia ou inconstancia. Logo \textsc{Deus} é \emph{moralmente} immutavel \footnote{D'onde resulta que \textsc{Deus}, quando quer ou permitte effeitos diversos ou contrarios entre si, não muda de vontade; porque uma coisa é mudar de vontade e outra coisa é querer a mudança no effeito, e um agente póde, com o mesmo decreto, prescrever que presentemente se faça uma coisa, e que depois se faça o contrario. --- Como tambem, quando criou o mundo, \textsc{Deus} não mudou de vontade, porque não mudou de essencia (visto que em \textsc{Deus}, como provaremos, as \emph{operações} se identificam com a \emph{essencia}), mas, desde a eternidade, decretou que, em um certo tempo, existisse o mundo. Diz S. Thomaz: « Aliud est mutare voluntatem, et aliud est velle aliquarum rerum mutationem. Potest enim aliquis, eâdem voluntate immobiliter permanente, velle quod nunc fiat hoc, et postea fiat contrarium. Sed tunc voluntas mutaretur, si aliquis inciperet velle quod prius non voluit. Quod quidem accidere non potest, nisi praesupposita mutatione vel ex parte cognitionis, vel circa dispositionem substantiae ipsius volentis » (\emph{Sum. Th.}, p. I, q. 19, a. 7). --- Entende-se egualmente como é que \textsc{Deus}, sem alguma mudança na sua vontade, possa agora amar e amanhã deixar de amar uma e a mesma pessoa, que hoje é santa e amanhã é peccadora; porque a mudança não se deu na vontade de \textsc{Deus}, mas só no homem, que antes era digno de amor, e depois se tornou digno de odio. Diz S. Agostinho: « Cum iratus malis dicitur (Deus), et placidus bonis, illi mutantur, non Ipse: sicut lux infirmis oculis aspera, firmis lenis est, ipsorum scilicet mutatione, non sua » (\emph{De Trin.}, l. V, c. 16).\par Objectam: \textsc{Deus} é livre com relação ás creaturas. A liberdade suppõe a indifferença; de modo que o acto, que se diz livre, pode ser e pode não ser realizado. Ora se \textsc{Deus} pode realizar e pode deixar de realizar um acto, é evidente que, quando o realiza, a sua vontade se muda.\par Resposta. Tudo o que \textsc{Deus} quer fóra da sua essencia, não só o quer desde a eternidade, mas tambem com o mesmo simplicissimo acto, com que ama a si mesmo; pois, em \textsc{Deus}, não ha nem pode haver composição ou vicissitude de actos. O acto da divina vontade, sendo simplicissimo, refere-se ao objecto de um modo proporcionado á condição do mesmo objecto. Logo \textsc{Deus} ama a si mesmo como \emph{fim}, e ama as outras coisas como \emph{meios} para o mesmo fim. E, como estes meios não têm uma relação necessaria com o fim, \textsc{Deus} pode querel-os e deixar de querel-os; e, ainda que ame a si mesmo \emph{necessariamente}, comtudo ama ou quer as outras coisas \emph{contingente} e \emph{livremente}, porque assim o exige a diversa qualidade dos objectos. --- Por isso\pgn{457}Deus quer verdadeiramente as creaturas, quando as quer; e não as quer verdadeiramente, quando não as quer; mas o querer e o não querer de Deus não suppõe uma mudança da parte de Deus, mas da parte das creaturas, que se referem \emph{realmente} ao proprio Deus. --- Para a comprehensão d'esta verdade, lembramos que a causa efficiente creada não se muda, porque opera; mas muda-se, porque, operando, passa necessariamente da potencia para o acto; e por isso Deus, que é acto purissimo, não se muda quando opera. --- Em poucas palavras: a liberdade, no seu conceito abstractissimo, exprime indifferença não para com a propria \emph{entidade}, mas para com o \emph{termo}, de modo que a vontade possa querer ou deixar de querer um objecto; mas não exige que este querer ou deixar de querer seja expressão de actos \emph{entitativamente} diversos, se não quando ella se considera na vontade creada, a qual, pela sua volição, passa da potencia para o acto.\par \emph{Renan}, seguido pelos defensores do \emph{progresso indefinido}, distingue um duplice \textsc{Deus}, --- o \textsc{Deus} \emph{ideal} e o \textsc{Deus} \emph{real}. O \textsc{Deus} \emph{ideal} é, segundo esse escriptor, immutavel e eterno; o \textsc{Deus} \emph{real}, que se identifica com a \emph{universalidade das coisas existentes}, \emph{é mutavel} e capaz de progresso indefinido.\par Estas impias e absurdas affirmações, que encerram o mais crasso \emph{pantheismo}, não merecem ser novamente refutadas, pois já o foram no capitulo antecedente, quando provámos que existe um \textsc{Deus} \emph{real} e \emph{verdadeiro}.}.

\pg{457} 32. \textbf{Eternidade de Deus}. --- A \emph{eternidade} é o attributo, que exclue de \textsc{Deus} o \emph{principio}, o \emph{fim}, e toda a \emph{successão} \footnote{\emph{Boecio} dá á seguinte definição da \emph{eternidade}: « A \emph{eternidade é a posse total e perfeita de uma vida sem fim} ». Analysemos esta definição. A \emph{eternidade} --- é \emph{posse}, porque é uma coisa estavel, que se conserva com socego e não póde ser tirada por nenhuma força extranha; --- é posse \emph{total}, porque a duração de \textsc{Deus} é sempre completa e indivisivel, e distingue-se da duração do ser finito, a qual não é \emph{total} mas \emph{successiva}; --- \emph{perfeita}, porque a \emph{eternidade} é o \emph{presente}, que dura sempre, e por isso é diversa do tempo \emph{presente} dos seres finitos, que, embora, por não ter partes, seja \emph{total}, comtudo não é uma coisa absolutamente perfeita (pois não é perfeição possuir uma coisa apenas por um instante, e o tempo \emph{presente} das coisas finitas é um instante); --- é posse da uma \emph{vida}; porque a duração refere-se principalmente ao \emph{acto} dos seres nobilissimos, que é \emph{viver}. --- de uma vida \emph{sem fim}, isto é, que não principiou nem ha de acabar. --- A \emph{posse} sem \emph{principio} e sem \emph{fim} é o elemento \emph{material} da \emph{eternidade}; a \emph{simultaneidade} da posse é o elemento \emph{formal}.\par \emph{A eternidade}, pois, consiste na exclusão do \emph{principio}, do \emph{fim}, e de toda a \emph{successão}, e identifica-se com a propria substancia de \textsc{Deus}, emquanto Esta não teve \emph{principio}, nem terá \emph{fim}, e não está sujeita a \emph{successão}.\par Por isso o \emph{tempo} não é uma parte da \emph{eternidade}, porque a \emph{eternidade} não tem \emph{partes}; mas a \emph{eternidade} corresponde e assiste a todos os seculos, que se vão succedendo, assim como um \emph{centro} immovel corresponde, por meio de linhas, a todos os pontos de uma esphera, que se mova em volta d'elle.\par D'onde se vê o erro --- tanto de \emph{Locke}, quando disse que a \emph{eternidade} é constituida pela infinita multiplicação do \emph{tempo}, --- como de \emph{Clarke}, quando opinou que a \emph{eternidade} é o proprio tempo, considerado em \emph{abstracto}.\par S. Thomaz trata d'este argumento --- na \emph{Summ. Theol.}, p. I, q. 10; --- \emph{c. Gentes}, l. I, c. 15; --- I \emph{Sent.} Dist. 8, q. 2, a. 1, e Dist. 19, q. 2, a. 1 e 2.}

\pg{458} 33. \textbf{Deus} é \textbf{eterno}. --- Um ente e é diz-se \emph{eterno}, quando não tem \emph{principio}, nem \emph{fim}, e não está sujeito a \emph{successão}. Ora \textsc{Deus} não tem \emph{principio} nem \emph{fim}, e não está sujeito a \emph{successão}. --- Não tem \emph{principio} nem \emph{fim}. Na verdade, \textsc{Deus}, sendo o Ente \emph{necessario}, \emph{immutavel} e \emph{independente}, existiu sempre e não póde deixar de existir; pois, se tivesse principiado a existir, quando não existia, ou podesse deixar de existir, seria \emph{contingente}, \emph{mutavel} e \emph{dependente}. --- Não está sujeito a \emph{successão}. Porquanto, a \emph{successão} indica \emph{mudança} e \emph{imperfeição}, e \textsc{Deus} é \emph{immutavel} e infinitamente \emph{perfeito}. Logo \textsc{Deus} é \emph{eterno} \footnote{Dizemos no texto que toda a \emph{successão} indica \emph{imperfeição}; porquanto o ente, que admitte \emph{successão} nos seus estados ou qualidades, adquire perfeições que não tinha, ou perde as que tinha, e por consequencia é finito ou limitado nas suas perfeições.\par \textsc{Deus} não só é \emph{eterno}, mas é a sua \emph{eternidade}; porque a eternidade não é uma coisa distincta de Deus, mas é o seu proprio ser uniforme; e por isso Deus, assim como é a sua essencia, assim tambem é a sua eternidade. Diz S. Thomaz: « Ratio aeternitatis consequitur immutabilitatem, sicut ratio temporis consequitur motum. Unde cum Deus sit maxime immutabilis, sibi maxime competit esse aeternum. Nec solum est aeternus, sed est sua aeternitas; cum tamen nulla alia res sit sua duratio, quia non est suum \emph{esse}. Deus autem est suum esse uniforme. Unde sicut est sua essentia, ita sua aeternitas » (\emph{Sum}. \emph{Th.,} p. I, q. 10, a. 2).\par Como a \emph{eternidade} consiste na perfeita immutabilidade, e todas as creaturas são mutaveis, segue-se que a nenhuma creatura compete a eternidade, no sentido rigoroso da palavra, e por isso só Deus é eterno. --- Todavia costumamos dizer que algumas creaturas participam da eternidade, como participam da immutabilidade no seu ser ou nos seus actos. Mas trata-se sempre de uma eternidade impropria e imperfeita (\emph{Sum}. \emph{Th.,} p. I, q. 10, a. 3).}.

34. \textbf{Immensidade de Deus}. --- \emph{Immensidade} é o attributo, pelo qual \textsc{Deus} está presente \emph{actualmente} em todas as coisas \emph{existentes} e \emph{virtualmente} em todas as coisas \emph{possiveis} \footnote{Quando dizemos que \textsc{Deus} está \emph{virtualmente} presente em todas as coisas \emph{possiveis}, entendemos que \textsc{Deus} tem a \emph{virtude} ou o poder de estar em todas as coisas, que não existem, mas que podem existir.\par E' necessario distinguir a \emph{immensidade} da \emph{ubiquidade}. --- A \emph{immensidade} é a presença divina, \emph{actual} ou \emph{virtual}, em todas as coisas, tanto \emph{existentes} como \emph{possiveis}; a \emph{ubiquidade} é a \emph{actual} presença de \textsc{Deus} nas coisas \emph{existentes}. --- A \emph{immensidade} é uma coisa eterna, absoluta e infinita, porque se extende a todo o possivel, e prescinde da \emph{actual} existencia das coisas; a \emph{ubiquidade} é uma coisa temporanea, relativa, não é rigorosamente infinita, porque suppõe a existencia do mundo, e porque as coisas creadas não são infinitas. --- A \emph{ubiquidade} está contida na \emph{immensidade}. --- Por isso \textsc{Deus},\pgn{459}embora não esteja presente nas coisas finitas, senão depois da creação, todavia é \emph{immenso} desde a eternidade; porque a propria Essencia de \textsc{Deus} exige que Elle possa e deva estar presente em todas as coisas, \emph{existentes} e \emph{possiveis}. --- S. Thomaz trata d'este assumpto --- na \emph{Sum}. \emph{Th.,} p. I, q. 8, --- c. \emph{Gentes}, 1. III, c. 68; --- I \emph{Sent}. Dist. 37, q. 3, a. 1.}.

35. \textbf{Deus} é immenso. --- \textsc{Deus} será \emph{immenso}, se estiver presente --- não só em todas as coisas \emph{creadas} e em todos os logares \emph{existentes}, --- mas tambem em todas as coisas e logares \pg{459}  \emph{possiveis}. Ora \textsc{Deus} --- está presente em todas as coisas \emph{creadas}, porque Elle produziu e conserva todas as creaturas, e um agente não póde operar onde não existe; --- está presente em todos os logares \emph{existentes}, porque todo o logar é uma coisa creada; --- póde estar presente em todas as coisas e logares \emph{possiveis}, porque \textsc{Deus} deve estar em todo o ser, a que se extende o seu poder (pois, em \textsc{Deus}, os attributos identificam-se com a substancia), e o poder de \textsc{Deus} extende-se a todo o \emph{possivel}. Logo \textsc{Deus} é \emph{immenso} \footnote{\textsc{Deus} está em todos os logares --- com o seu \emph{poder}, emquanto tudo Lhe obedece, --- com a sua \emph{presença}, emquanto Elle vê tudo o que se faz e pensa, --- com a sua \emph{essencia}, emquanto produz e conserva o ser das creaturas (Cf. \emph{Sum}. \emph{Th.,} p. I, q. 8, a. 3).\par Pergunta-se: Onde estava \textsc{Deus}, antes da creação do mundo? e agora como está Elle no mundo?\par Antes da creação, \textsc{Deus} estava em si mesmo. E, quando o mundo foi creado, \textsc{Deus} não deixou de estar em si mesmo, nem se mudou pelo facto de estar presente em todas as coisas; assim como não se muda o sol, quando se mudam as coisas, que elle illumina. --- E, agora, \textsc{Deus} está no mundo, não diffundido por toda a parte, como a luz ou o ar, mas de um modo ineffavel, que excede toda a \emph{imaginação}. Uma semelhança, embora infinitamente \emph{inadequada}, da presença de \textsc{Deus} nas coisas é a nossa alma, a qual, ainda que \emph{simples} e \emph{espiritual}, está toda em todo o organismo e toda em cada parte d'elle. Dizemos --- infinitamente \emph{inadequada}, porque a nossa alma é limitada ao nosso corpo, ao passo que \textsc{Deus} não é limitado a este mundo.\par Muitos negaram ou desvirtuaram a \emph{immensidade} de \textsc{Deus}.\par \emph{a}) Uns, como os \emph{Manicheus} e \emph{Descartes}, negaram a \emph{ubiquidade} de \textsc{Deus}; pois aquelles diziam que só as coisas \emph{espirituaes} estão sujeitas ao \emph{poder} de \textsc{Deus}, e \emph{Descartes} affirmou que \textsc{Deus} está em todas as coisas com o seu \emph{poder} e \emph{presença}, mas não com a sua \emph{essencia} ou \emph{substancia}.\par Ambas essas affirmações são falsas. --- E' falsa a dos \emph{Manicheus}; porque o poder de \textsc{Deus} é infinito e extende-se a todas as creaturas, espirituaes e corporeas. --- E' falsa a de \emph{Descartes}; porque, identificando-se em \textsc{Deus} as operações com a essencia, onde existem aquellas, deve existir esta.\par \emph{b}) Outros, seguindo as theorias de \emph{Newton} e \emph{Clarcke}, disseram que a \emph{immensidade divina} é o espaço infinito.\par Esta opinião foi refutada na \emph{Ontologia}.\par \emph{c}) \emph{Spinoza} disse que a \emph{immensidade} não é o attributo de um Ente individuo e pessoal, mas é a \emph{collecção} de todas as coisas.\par Estas affirmações contêm o \emph{pantheismo}, que já foi refutado.\par \emph{d}) \emph{Bayle} diz que \textsc{Deus}, sendo \emph{simplicissimo}, não póde ser \emph{immenso}; porque a \emph{simplicidade} e a \emph{immensidade} são attributos contradictorios.\par Respondemos que a \emph{simplicidade} não só não é contraria á \emph{immensidade},\pgn{460}mas é uma condição indispensavel para ella; porque \textsc{Deus} não seria \emph{immenso}, se não fosse \emph{simplicissimo}, pois o \emph{composto}, resultando de partes necessariamente finitas e limitadas, não póde ser \emph{immenso}.}.

\pg{460} 36. \textbf{Unidade de Deus}. --- \emph{Unidade} é o attributo, pelo qual a Essencia divina não póde dividir-se, multiplicar-se e constituir varios e diversos deuses \footnote{Quando, pois, dizemos que \textsc{Deus} é \emph{uno}, entendemos que é uno não só \emph{especifica} mas tambem \emph{numericamente}, e por isso é \emph{unico}.\par S. Thomaz trata d'este assumpto --- na \emph{Sum Th.}, p. I, q. 11, a. 3 e 4; --- c. \emph{Gentes}, l. I, c. 42; --- \emph{De pot.}, q. 3, a. 6; --- I \emph{Sent}. Dist. 2, a. 1, --- II \emph{Sent}. Dist. 21, q. 1, a. 1.}.

37. \textbf{Deus} é uno.

a) \textsc{Deus} é o Ente infinitamente perfeito. Ora o ente infinitamente perfeito deve ser \emph{uno}. Porquanto, se fossem varios, deveria haver entre elles uma distincção. Esta distincção só poderia dar-se, emquanto um fosse dotado de uma perfeição, que faltasse a outro. Mas, neste caso, aquelle, a quem faltasse uma perfeição, não seria \textsc{Deus}; mas só o seria Aquelle, que fosse infinitamente perfeito. Logo \textsc{Deus} é \emph{uno}.

b) Uma forma multiplica-se, emquanto é recebida em muitos e diversos sujeitos; assim a \emph{humanidade} multiplica-se, emquanto é recebida em muitos e diversos homens. Se, porém, a forma não for recebida em sujeito algum, mas subsistir em si mesma, não poderá multiplicar-se, e existirá um unico individuo d'aquella especie; assim, se a \emph{humanidade} fosse subsistente em si e por si, existiria um unico homem. Ora a Divindade não é recebida em sujeito algum (porque, em \textsc{Deus}, infinitamente \emph{simples}, a \emph{essencia} identifica-se com o \emph{supposto}), mas é subsistente em si e por si. Logo a Divindade não póde multiplicar-se e formar varios suppostos. Logo \textsc{Deus} é \emph{uno} \footnote{Nunca devemos esquecer que a \emph{subsistencia do ser divino} é, segundo o nosso modo de entender, a razão e a raiz de todos os attributos divinos. O ser das creaturas é simples como é simples uma parte ou uma coisa incompleta, porque é uma forma recebida num sujeito e individualisada pela natureza, que subsiste nesse ser, e por isso é multiplicavel e multiplica-se conforme a multidão das creaturas. Pelo contrario, o ser de Deus é simples, como é simples o acto e a perfeição, porque é o ser por si subsistente e por si individualisado e distincto de todo e qualquer outro ser, e por isso não é multiplicavel, mas é \emph{unico}.}.

c) Se existissem varios deuses, \emph{especificamente} eguaes, mas \emph{numericamente} distinctos, seriam necessariamente compostos de \emph{natureza divina}, \emph{commum} a todos, e de \emph{notas} ou \emph{attributos}, \emph{proprios} de cada um (aliás seria impossivel a egualdade \emph{especifica} a e distincção \emph{numerica}). Ora \textsc{Deus}, sendo \emph{simplicissimo}, exclue toda \pg{461}  e qualquer composição. Logo não podem existir varios deuses. Logo \textsc{Deus} é \emph{uno} \footnote{Os systemas ou crenças contrarias á \emph{unidade} de \textsc{Deus} foram o \emph{dualismo} e o \emph{polytheismo}.\par a) O \emph{dualismo} admittia que existem \emph{dois principios infinitos}, \emph{necessarios}, e \emph{supremos} --- um \emph{bom} e outro \emph{mau}, e que do \emph{principio bom} derivam todos os \emph{bens} e do \emph{mau} todos os \emph{males}.\par O \emph{dualismo} foi chamado tambem \emph{Manicheismo}, porque foi \emph{Manes} ou \emph{Manicheu} quem primeiramente o ensinou. --- Nos tempos modernos foi resuscitado pelo sceptico \emph{Bayle}, ensinando que o \emph{dualismo}, embora seja absurdo \emph{a-priori}, comtudo póde provar-se \emph{a-posteriori}.\par O \emph{dualismo} é falso pelas seguintes razões:\par \emph{a}) \emph{Baséa-se num falso fundamento}. Porquanto, admitte que existe um principio inteiramente mau; porque, aliás, não poderia explicar-se a existencia do mal. Ora esta razão é falsa; porque o mal, derivando das causas creadas, é uma consequencia do \emph{bem limitado}.\par \emph{b}) \emph{Um principio inteiramente mau não póde existir}. Porquanto, esse principio --- ou é privado de toda a perfeição e entidade, --- ou falta-lhe apenas a bondade \emph{moral}. --- Na primeira hypothese, um ente inteiramente mau seria \emph{nada} e o \emph{nada} não existe. --- Na segunda, um principio inteiramente mau seria, ao mesmo tempo, infinito e finito. Seria \emph{infinito}, porque o systema suppõe-no tal. Seria \emph{finito}, porque faltar-lhe-hia uma perfeição.\par \emph{c}) \emph{O dualismo perverte a moral}. Com effeito, se todo o \emph{bem} devesse attribuir-se a um principio \emph{bom}, e todo o \emph{mal} a um principio \emph{mau}, o homem não poderia dizer-se e não seria causa do bem nem do mal, e por isso não seria merecedor de premio nem de castigo.\par \emph{d}) \emph{Esta hypothese não explica a existencia do bem e do mal}. Porquanto --- ou estes dois principios seriam eguaes, --- ou o principio \emph{bom} seria mais forte que o \emph{mau}, --- ou este mais forte que aquelle. Na primeira hypothese, as forças dos dois principios paralysar-se-hiam, e assim não existiria, sobre a terra, nem o bem nem o mal. Na segunda, não haveria mal algum. Na terceira, não haveria bem algum. Ora estas consequencias são falsas.\par \emph{Bayle} diz que os dois principios fizeram um pacto, para se não transtornarem mutuamente nas suas obras.\par Mas esta explicação é absurda. Se assim fosse, esses dois principios seriam fracos, e nenhum d'elles seria infinitamente perfeito. --- De mais, esse pacto ou seria bom ou mau. Se fosse bom, não poderia ser acceite pelo principio \emph{mau}; pois este, sendo inteiramente \emph{mau}, não póde fazer nem tolerar bem algum. Se fosse mau, não poderia ser acceite pelo principio \emph{bom}; pois este, segundo o \emph{dualismo}, não póde fazer nem tolerar mal algum.\par b) O \emph{polytheismo} é o erro que admitte varios \emph{deuses}.\par Este erro, que derivou sobre tudo da corrupção da vontade, não foi constante nem universal. --- Não foi \emph{constante}; porque não existia antes do \emph{diluvio}, e porque a primeira crença dos povos foi o \emph{monotheismo} (Cf. Hooke, \emph{Princip}. \emph{rélig}. \emph{natural}. \emph{et révél}.,--- Frank, \emph{Études Orientales}). --- Não foi \emph{universal}; porque, além de muitos philosophos (como \emph{Anaxagoras}, \emph{Socrates}, \emph{Platão}, \emph{Aristoteles}, \emph{Cicero}), alguns povos, como o \emph{d'Israel}, professavam o \emph{monotheismo}, e muitos outros, embora admittissem deuses, comtudo reconheciam \emph{Jupiter}, como pae e senhor de todos os deuses.}.

%%>>> art-41

\section*{ARTIGO III.}
\addcontentsline{toc}{section}{ARTIGO III.}

{\small\noindent\textsc{Summario}: \textsc{Summario}: --- Intelligencia de Deus, sua natureza e objecto. --- Vontade de Deus, sua natureza e objecto. --- Vida divina, sua existencia e natureza.\par}

\pg{462} \textbf{Attributos de Deus positivos.}

38. \textbf{Attributos de Deus positivos.} --- Attributos de \textsc{Deus} \emph{positivos} são os que \emph{affirmam} de \textsc{Deus} as perfeições dos seres creados, mas isentas de toda a imperfeição ou de todo o limite. --- Os principaes são: a \emph{verdade}, a \emph{bondade} e a \emph{belleza} \footnote{Depois dos attributos \emph{negativos} de \textsc{Deus}, devemos considerar os attributos \emph{positivos}. E primeiramente consideraremos os attributos \emph{transcendentaes}, que \emph{analogicamente} convêm a \textsc{Deus} e ás creaturas. Como dissemos na Ontologia, attributos \emph{transcendentaes} são os que convêm a \emph{todo o ente}, emquanto \emph{tal}, e que, embora não sejam realidades distinctas do proprio ente, comtudo são explicações mais claras do mesmo, porque exprimem o que o ente não exprime. --- São quatro: \emph{unidade}, \emph{verdade}, \emph{bondade}, \emph{bellezza}. --- Tendo tratado da \emph{unidade}, resta-nos tratar dos outros attributos.}.

39. \textbf{Verdade de Deus.} --- \emph{Verdade} de \textsc{Deus} é o attributo, que consiste na \emph{conformidade}, ou \emph{adequação}, entre o acto da intelligencia divina e o ser divino \footnote{S. Thomaz trata d' este assumpto --- na \emph{Sum}. \emph{Th}., p. I, q. 16, --- c. \emph{Gentes}, l. I, c. 60 e seg., --- \emph{De verit}., q. I, --- I \emph{Sent}. Dist. 19, q. 5, a. a. 2 e 3.}.

40. \textbf{Deus é verdadeiro.} --- Um ente é \emph{verdadeiro}, quando ha conformidade entre o seu ser e o acto da intelligencia, á qual de algum modo se refere. Ora entre o ser de \textsc{Deus} e o acto da sua intelligencia não só ha \emph{conformidade}, mas absoluta \emph{identidade}; porque, em \textsc{Deus}, a substancia e os actos das faculdades são uma e a mesma coisa. Logo \textsc{Deus} é \emph{verdadeiro} \footnote{A verdade, por ser uma adequação entre a intelligencia e o objecto, deve encontrar-se em \textsc{Deus}, cuja intelligencia é uma e a mesma coisa com a essencia. Na sua essencia, como veremos, \textsc{Deus} conhece todas as coisas. --- Nem se diga que a verdade, consistindo \emph{na composição} e na \emph{divisão}, não pode encontrar-se em \textsc{Deus}, cuja intelligencia é \emph{simplicissima}, como é a sua essencia. O principio --- a \emph{verdade consiste na composição e na divisão} --- verifica-se unicamente com relação á verdade, que é propria da intelligencia \emph{humana}, e não pode applicar-se á intelligencia \emph{divina}. Com effeito, quando se diz que a verdade consiste na adequação entre a intelligencia e o objecto, --- não devemos entender que o \emph{acto} da intelligencia deve por tal modo egualar o \emph{objecto}, que, se este fôr composto e material, tambem composto e material deva ser aquelle (o que é absurdo), --- mas devemos entender que o \emph{verbo} da intelligencia deve por tal forma corresponder ao seu \emph{objecto}, que tudo o que o \emph{verbo} exprime se deve encontrar realmente no \emph{objecto}, o qual pode ser simples ou composto, espiritual ou material. Portanto a \emph{qualidade}, isto é, a \emph{simplicidade} ou a \emph{composição} do \emph{acto intellectual} não entra na razão ou essencia da \emph{verdade}, mas refere-se exclusivamente ao principio intelligente; e por isso, se a intelligencia do homem, por ser\pgn{463}a menos perfeita, não pode alcançar a verdade se não por meio da \emph{composição} e da \emph{divisão} (juizo \emph{affirmativo} ou \emph{negativo}), a intelligencia de \textsc{Deus}, por ser infinitamente perfeita e simplicissima, como a sua essencia, com um acto unico e simplicissimo alcança toda a verdade, quer esta se refira á essencia divina, quer se refira ás coisas creadas, simples ou compostas, espirituaes ou materiaes. (\emph{Cf}. \emph{C}. \emph{Gentes}, l. I. c. 59).}.

\pg{463} 41. \textbf{Deus é a verdade.} --- Consistindo a verdade na conformidade entre o acto da intelligencia e o ser, \textsc{Deus} será a propria \emph{verdade}, se fôr o seu entender e o seu proprio ser. Ora \textsc{Deus} é, effectivamente, o seu entender e o seu proprio ser. Logo \textsc{Deus} é a \emph{verdade}. --- \textsc{Deus} é a verdade \emph{summa} e \emph{infinita}, porque summa e infinitamente perspicaz é a sua intelligencia, e summa e infinitamente intelligivel é o seu ser \footnote{Em \textsc{Deus} não só ha a verdade, mas Elle é a propria verdade. Por quanto, tudo o que está em Deus não é \emph{participado} nem \emph{accidental}, mas é \emph{essencial} e \emph{substancial}, é o proprio ser divino, é o proprio Deus. Logo Deus, assim como é a sua essencia, o seu ser, assim tambem é a sua verdade. Diz S. Thomaz: « De Deo nihil \emph{participative} dici potest, cum sit suum esse, quod nihil participat. Sed veritas est in Deo. Si igitur non dicatur \emph{participative}, oportet quod dicatur \emph{essentialiter}. \textsc{Deus} ergo est sua veritas ». (\emph{C}. \emph{Gentes}, l. I, c. 60). --- A verdade divina é tão pura, que não tolléra sombra de falsidade ou de engano. « A verdade, diz o mesmo Angelico, não é compativel com a falsidade, como a candura não é compativel com a côr preta. Ora Deus não só é verdadeiro, mas é a propria verdade. Logo em Deus não pode haver falsidade » (\emph{Ib}., c. 61).}.

42. \textbf{Deus é a verdade na ordem ontologica e na ordem logica.}

a) \textsc{Deus} É A VERDADE NA ORDEM ONTOLOGICA. --- \textsc{Deus} será a verdade na ordem \emph{ontologica}, se a sua essencia ou substancia se conformar com a sua intelligencia; porque a verdade \emph{ontologica} consiste na conformidade do objecto com a intelligencia. Ora a essencia de \textsc{Deus}, não só se conforma com a sua intelligencia, mas identifica-se com ella, e \textsc{Deus} é a sua essencia. Logo Deus é a verdade na ordem \emph{ontologica}.

b) \textsc{Deus} É A VERDADE NA ORDEM LOGICA. --- \textsc{Deus} será a verdade na ordem \emph{logica}, se a sua intelligencia se conformar com a sua essencia; porque a verdade \emph{logica} consiste na conformidade da intelligencia com o objecto. Ora, como dissemos, a intelligencia de \textsc{Deus} não só se conforma com a sua essencia, mas identifica-se com ella, e \textsc{Deus} é a sua intelligencia. Logo \textsc{Deus} é a verdade na ordem \emph{logica} \footnote{A verdade, como dissemos, pode considerar-se no \emph{objecto} e na \emph{intelligencia}. Considerada no \emph{objecto}, é a conformidade d'este com a intelligencia (\emph{veritas rei}); --- considerada na intelligencia, é a conformidade d'esta com o objecto (\emph{veritas intellectus}). Em \textsc{Deus} ha a verdade do \emph{objecto} ou \emph{ontologica}, e ha a verdade da \emph{intelligencia} ou \emph{logica}. Diz S. Thomaz: « Veritas\pgn{464}invenitur \emph{in intellectu}, secundum quod apprehendit \emph{rem ut est}; et \emph{in} re, secundum quod habet esse \emph{conformabile intellectui}. Hoc autem maxime invenitur in Deo. Nam esse suum non solum est conforme suo intellectui, sed etiam est ipsum suum intelligere..., et ipse est suum esse et intelligere. Unde sequitur quod non solum in ipso sit veritas, sed quod ipse sit ipsa summa veritas » (\emph{Sum}. \emph{Th}., p. I, q. 16, a. 5).}.

\pg{464} \textbf{43. A verdade ontologica de Deus é a causa efficiente da verdade ontologica das creaturas.} --- A verdade \emph{ontologica} de \textsc{Deus} será a causa \emph{efficiente} da verdade \emph{ontologica} das creaturas, se o ser d'estas for effeito e expressão da Essencia divina. Ora o ser das creaturas, tendo sido produzido por \textsc{Deus} e em conformidade com as idéas archetypas, que são imitações diversas da Essencia infinita, é, na sua physica realidade, expressão e effeito da mesma Essencia. Logo a verdade \emph{ontologica} de \textsc{Deus} é a causa \emph{efficiente} da verdade \emph{ontologica} das creaturas \footnote{A Essencia divina é uma simplicissima unidade, mas é imitavel e participavel, por infinitos e diversos modos, pelas creaturas. « Cada creatura, diz S. Thomaz, possue a propria especie, emquanto de algum modo participa da semelhança da Essencia divina ». (\emph{Sum}. \emph{Th}., p. 1, q. 15, a. 2). Por isso cada realidade creada é uma \emph{expressão}, ou \emph{semelhança}, da unica Essencia divina. --- E' uma \emph{expressão}, e tambem um \emph{effeito}. Diz o mesmo santo Doutor: « Um ente, que possue o ser \emph{por essencia}, é necessariamente causa de todos os entes, que o possuem por \emph{participação}... Mas Deus possue o ser pela sua essencia, porque é o proprio ser; e todos os outros entes só o possuem por participação, porque o ente, que é o proprio ser, é unico. Portanto, Deus é causa do ser de todos os outros entes » (\emph{C}. \emph{Gentes}, l. II, c. 15). --- Sendo uma expressão e um effeito da Essencia divina, o ser das creaturas não pode deixar de corresponder ás ideas archetypas da intelligencia divina. Por isso, o ser das creaturas refere-se \emph{necessariamente} á intelligencia divina, emquanto não pode deixar de corresponder ás ideas archetypas, e só \emph{accidentalmente} á intelligencia creada, emquanto \emph{pode} ser conhecido por ella; e a verdade \emph{ontologica} das creaturas é \emph{necessaria} com relação a Deus, e é \emph{contingente} com relação a nós. --- Como tambem essa verdade é \emph{unica} e \emph{multipla}. E' \emph{unica} com relação a Deus: porque unica é a Essencia divina, que pode ser e é diversamente imitada e participada. E' \emph{multipla} com relação ás creaturas; porque, consistindo a verdade \emph{ontologica} na conformidade do ente com a intelligencia, são tantas as verdades, quantas são as entidades.}.

\textbf{44. A verdade logica de Deus é a causa exemplar da verdade logica da intelligencia humana.} --- A verdade \emph{logica} da intelligencia humana depende das coisas, emquanto a nossa intelligencia, para alcançar e possuir essa verdade, deve conformar-se com as proprias coisas. Ora conformar-se com as coisas é conformar-se com as ideas archetypas da intelligencia divina, de que as coisas dependem essencialmente e que as coisas exprimem na sua realidade. Logo as ideas da nossa intelligencia emtanto \pg{465}  são dotadas de verdade \emph{logica}, emquanto imitam as idéas archetypas divinas, isto é, a verdade \emph{logica} de \textsc{Deus}. Logo a verdade \emph{logica} de \textsc{Deus} é a causa \emph{exemplar} da verdade \emph{logica} da intelligencia humana \footnote{« A verdade divina, diz S. Thomaz, é a medida de toda a verdade. Por quanto, a verdade da nossa intelligencia depende e é medida pela coisa, que está fóra da alma; pois dizemos que a nossa intelligencia é verdadeira, emquanto se conforma com a coisa. Mas a verdade da coisa é medida pela intelligencia divina, que é a causa de todos os seres creados, como a verdade das coisas artificiaes é medida pela arte do artista. Sendo, pois, \textsc{Deus} a primeira intelligencia e o primeiro intelligivel, é forçoso concluir que a sua verdade é a medida da verdade da intelligencia creada. Por isso a verdade divina é a \emph{primeira} verdade » (\emph{C}. \emph{Gentes}, l. 1, c. 62). --- Portanto a verdade \emph{logica}, propria da nossa intelligencia, pode dizer-se, ao mesmo tempo, \emph{unica} e \emph{multiplice}. E' \emph{unica} com relação a \textsc{Deus}, emquanto as nossas ideas são a expressão ou a semelhança, ainda que imperfeita, da unica e substancial verdade da intelligencia divina. E' \emph{multiplice} com relação ás \emph{formas} da intelligencia humana, que existem em muitas e distinctas intelligencias e representam muitos e diversos objectos.}.

45. A \textbf{verdade} divina é a \textbf{causa} final de toda a \textbf{verdade creada}. --- \textsc{Deus} não só é o principio, mas tambem é o fim de toda a creação; de modo que toda a perfeição creada não existe senão para imitar e revelar, embora de um modo inadequado, a correspondente perfeição divina. Ora a verdade creada é um effeito, e a sua existencia não pode ter e não tem outro fim, senão o de imitar e manifestar a verdade divina. Logo a verdade divina é a causa \emph{final} de toda a verdade creada \footnote{A verdade creada é, como dissemos, \emph{ontologica} e \emph{logica}, conforme se considera no \emph{objecto} ou na \emph{intelligencia}. Ora --- a verdade \emph{ontologica}, que é o proprio \emph{objecto} emquanto se conforma com a intelligencia, imita e manifesta a verdade \emph{ontologica} de \textsc{Deus}, porque todo o ente participado revela o ente por essencia, do qual deriva; --- e a verdade \emph{logica}, que é o verbo da intelligencia emquanto se conforma com o objecto, revela a verdade da intelligencia divina. Por isso a verdade creada, emquanto imita e revela a verdade divina, tende para esta, como para o seu fim natural; e d'este modo a verdade divina é a causa \emph{final} da verdade creada.}.

46. \textbf{Bondade} de \textbf{Deus}. --- \emph{Bondade} de \textsc{Deus} é o attributo, que consiste na \emph{conveniencia} da infinita perfeição de \textsc{Deus} com a sua santissima vontade. --- A \emph{bondade}, pois, é a propria entidade, a propria perfeição do ente, emquanto \emph{convem} á vontade, e attrahe-a ou pode attrahil-a \footnote{S. Thomaz trata da \emph{bondade} de \textsc{Deus} --- na \emph{Sum}. \emph{Th}., p. I, q. 6 --- c. \emph{Gentes}, l. I, c. 37 e seg. --- \emph{De verit}., q. 21.}.

47. \textbf{Deus} é bom. --- Um ente é e diz-se \emph{bom}, emquanto é dotado de perfeição, que tem a força de attrahir a vontade.

\pg{466} Ora Deus contém em si toda a perfeição, pela qual é objecto de um amor sem limites da parte da sua vontade; pois todo o ente não pode deixar de amar a propria perfeição. Logo \textsc{Deus} é \emph{bom}. --- \textsc{Deus} é \emph{summa} e \emph{infinitamente bom}, porque \emph{summa} e \emph{infinita} é a sua entidade ou perfeição \footnote{Deus é \emph{bom}, e \emph{summa} e \emph{infinitamente} bom, porque é \emph{perfeito}, e \emph{summa} e \emph{infinitamente} perfeito. A bondade está na razão \emph{directa} da perfeição; porque a bondade não é senão a perfeição, emquanto se refere á vontade. Diz o P. Lessio; « Imagina uma infinita luz, um infinito poder, uma infinita sabedoria, uma infinita belleza, uma infinita doçura, uma infinita alegria, uma infinita riqueza, uma infinita gloria, uma infinita majestade, uma infinita santitade. Pois bem, a Bondade divina contém todas estas perfeições como um oceano infinito de todos os bens. E não só contem todo o bem, mas contém-no de um modo infinitamente elevado; porque não o contem por muitas e diversas formas ou qualidades, mas por uma simplicissima e elevadissima forma, que é a sua essencia. Este modo de conter as perfeições, torna estas infinitamente mais excellentes, do que se existissem distinctas uma de outra. (\emph{De perfectione divina}, l. 7, e. 1, n. 2).}.

48. \textbf{Deus é a bondade.} --- A \emph{bondade} de \textsc{Deus} é a sua perfeição; e por isso tal é a natureza da sua bondade, qual é a natureza da sua perfeição. Ora a perfeição, em \textsc{Deus}, não é uma coisa accidental, addicionada á essencia, mas é uma coisa substancial, é a divina essencia; porque toda a perfeição é um ente, e a essencia divina é o proprio ser. Logo \textsc{Deus}, assim como é a sua essencia, assim tambem é a sua \emph{bondade} \footnote{Se \textsc{Deus} não fosse a sua bondade, isto é, se não fosse bom pela sua \emph{essencia}, seria bom \emph{por participação}. Mas então, como o ente, que é bom \emph{por participação}, presuppõe outro ente, anterior a si, do qual recebeu a bondade, Deus não seria a Causa primeira, o Motor immovel, o Ente absoluto e infinito: o que é absurdo. Logo Deus é bom pela sua essencia, isto é, é a propria bondade. --- Se Deus é a propria Bondade, é claro que exclue absolutamente todo e qualquer \emph{mal}, que é opposto ao \emph{bem}. Alem d'isso, se o mal podesse existir em Deus, existiria como uma coisa \emph{essencial}, seria a propria \emph{essencia} de Deus, porque tudo o que está em Deus é o proprio Deus: o que é absurdo, porque o mal, sendo privação do bem, é nada, e o nada não pode constituir nenhuma essencia, muito menos a perfeitissima essencia divina. (Cf. c. \emph{Gentes}, l. I, c. 38 e 39).}.

49. \textbf{A bondade divina é a causa efficiente de toda a bondade creada.} --- A bondade divina será a causa \emph{efficiente} de toda a bondade creada, se as perfeições, de que são dotadas as creaturas, derivarem da unica e semplicissima bondade de \textsc{Deus}. Ora é certo que todas as perfeições, de que são dotadas as creaturas, derivam da unica e semplicissima bondade de \textsc{Deus}; porque as perfeições das creaturas são finitas, contingentes, mutaveis, e por isso derivam da perfeição divina, que é infinita \pg{467} , necessaria, immutavel. Logo a bondade divina é a causa \emph{efficiente} de toda a bondade creada \footnote{Toda a bondade das creaturas deriva, como da causa \emph{efficiente}, da bondade de Deus. Dizem os escholasticos: \emph{quod est maximum in unoquoque genere est causa aliorum quae sunt in illo genere}. --- Mas que differença entre a bondade de Deus e a das creaturas! A bondade de Deus é imparticipada e fonte inexhaurivel de todo o bem; a bondade das creaturas é participada e deriva de Deus. --- A bondade de Deus é illimitada, porque não ha quem a possa limitar; a bondade das creaturas é limitada, porque é recebida numa essencia finita. --- A bondade de Deus é substancial e identifica-se com a sua essencia; a bondade das creaturas é uma forma accidental e addicionada á essencia. --- A bondade de Deus é immutavel, não pode estar sujeita ao augmento nem á diminuição; a bondade das creaturas é capaz de augmento até ao infinito e de diminuição até ao nada. Por causa d'esta excellencia da bondade divina sobre a creada disse \textsc{JESUS CHRISTO}: « Só Deus é bom: \emph{Nemo bonus nisi solus Deus} » (Luc., XVIII, 19).}.

50. \textbf{A bondade divina é a causa exemplar da bondade creada.} --- O ente, que possue uma determinada qualidade por participação, não se diz dotado d'essa qualidade, se não emquanto manifesta uma certa \emph{semelhança} com o ente, que possue a mesma qualidade por essencia; assim como o ferro não se diz igneo, senão emquanto tem uma certa semelhança com o fogo. Ora \textsc{Deus} é bom por essencia, e todos os seres creados são bons por participação. Logo todo o ente creado diz-se \emph{bom}, emquanto tem uma certa semelhança com a bondade de Deus. Logo a bondade divina é a causa \emph{exemplar} da bondade creada \footnote{Cf. \emph{c}. \emph{Gentes}, l. 1, c. 40. --- A bondade divina é a causa \emph{exemplar}, mas não é a causa \emph{formal} da bondade creada, porque a bondade das creaturas não é uma e a mesma coisa com a bondade divina. A bondade creada é \emph{distincta} da bondade divina. --- Por isso a bondade creada é \emph{unica} e \emph{multiplice}. E' \emph{unica}, emquanto se refere a Deus, como ao seu primeiro principio \emph{efficiente} e \emph{exemplar}. E' \emph{multiplice}, emquanto a bondade ou perfeição creada se acha multiplicada tantas vezes, quantas são as creaturas.\par Advertimos que as creaturas não participam da \emph{bondade} de \textsc{Deus} pelo mesmo modo, por que participam da sua \emph{verdade}. Com effeito, as creaturas participam da \emph{bondade} de \textsc{Deus} --- não só \emph{extrinsecamente}, emquanto derivam de \textsc{Deus}, --- mas tambem \emph{intrinsecamente}, emquanto possuem em si mesmas a forma de bondade, que é uma semelhança, embora infinitamente remota e deficiente, da propria bondade de \textsc{Deus}; ao passo que ellas se dizem \emph{verdadeiras}, --- não \emph{intrinsecamente}, pois a verdade reside, como no seu proprio sujeito, na \emph{intelligencia}, e não no \emph{objecto}, --- mas apenas \emph{extrinsecamente}, emquanto os seres creados, por isso mesmo que são creados, se referem, e de um modo \emph{necessario}, á intelligencia divina. (Cf. Sum. \emph{Th}., p. I, q. 6, a. 4).}.

51. \textbf{A bondade divina é a causa final da bondade creada.} --- A bondade divina será a causa \emph{final} da bondade creada, se \textsc{Deus}, infundindo a bondade nas creaturas, tiver tido o fim de

\pg{468} communicar e manifestar a sua bondade infinita. Ora tal foi, effectivamente, o fim do Creador. Logo a bondade divina é a causa \emph{final} da bondade creada \footnote{Não pode duvidar-se de que Deus creou o mundo com o fim de communicar a outros seres a propria bondade. Porquanto, todo o agente, quando opéra, tem em vista a propria bondade ou perfeição --- para a adquirir, se não a possue, como o agente creado, --- e para a communicar, se a possue, como o agente increado.}.

52. \textbf{Belleza} de \textbf{Deus}. --- \emph{Belleza} de \textsc{Deus} é o attributo, que consiste na infinidade das perfeições divinas, reduzidas a uma unica e simplicissima unidade, e sobremaneira agradaveis a quem as contempla \footnote{Como dissemos na \emph{Ontologia}, a \emph{belleza é a conveniente unidade na variedade}, \emph{que agrada a quem a vê: pulchritudo est multitudinis debitae debita unitas}, \emph{quae visa placet}. Na belleza devemos distinguir dois elementos --- o elemento \emph{objectivo} e o \emph{subjectivo}. O elemento \emph{objectivo}, que constitue propriamente a essencia da belleza, é a \emph{conveniente unidade na variedade}, isto é, a devida multiplicidade reduzida á devida unidade; o elemento \emph{subjectivo}, que não é propriamente \emph{constitutivo} mas \emph{effeito} da belleza, é o \emph{agrado}, que a propria belleza produz em quem a contempla. --- O elemento \emph{objectivo é composto}. \emph{A multiplicidade} ou \emph{variedade} é a \emph{materia} da belleza; a \emph{unidade} é a \emph{forma}, como diz S. Agostinho: « Omnis pulchritudinis \emph{forma} est \emph{unitas} » (\emph{Ep}. 18, al. 63, ad Coelest. n. 2). --- O elemento \emph{subjectivo} indica que a belleza, como a verdade e a bondade, é uma coisa \emph{relativa}. A verdade refere-se á faculdade cognoscitiva, a bondade á appetitiva, a belleza refere-se á cognoscitiva e á appetitiva; porque todos dizem que \emph{pulchrum est quod visum placet}, ora uma coisa \emph{vé-se} pela faculdade cognoscitiva, e \emph{agrada} á appetitiva. --- A belleza é um attributo \emph{transcendental}; porque é proprio de todo o ente, emquanto \emph{tal}. Todo o ente é \emph{multiplo} (real ou logicamente) e \emph{uno}, e a sua belleza é maior ou menor, conforme o maior ou menor numero de partes e conforme a maior ou menor harmonia e união, que existe entre ellas.}.

53. \textbf{Deus} é \textbf{bello}. --- Consistindo a belleza na multiplicidade das perfeições devidas, reduzida á conveniente unidade, \textsc{Deus} será \emph{bello}, se possuir todas as perfeições, que são proprias do Ente increado, e se todas essas perfeições forem reduzidas á devida unidade; de modo que essa unidade em tanta variedade possa agradar a quem a contempla. Ora \textsc{Deus} possue perfeições infinitas, e todas essas perfeições constituem a unica e simplicissima Essencia divina, e tal unidade numa variedade infinita não pode deixar de agradar immensamente a quem a contempla. Logo \textsc{Deus} é \emph{bello} \footnote{A conclusão é evidente, e prova-se tambem com argumentos \emph{a-posteriori}. O mundo, creado por Deus, é dotado de belleza, e é por isso que foi chamado \emph{cosmos} pelos gregos, e \emph{mundo} pelos latinos. Ora toda a perfeição do effeito deve preexistir na sua Causa efficiente. --- Além d'isso,\pgn{469}a belleza, que se admira nas creaturas, é mutavel e imperfeita. Ora tudo o que é mutavel e imperfeito suppõe e exige uma causa absolutamente immutavel e perfeita. Logo deve existir um Ente immutavel e perfeitamente bello, de que deriva toda a belleza, espalhada nas creaturas.}.

\pg{469} 54. \textbf{Deus} é a \textbf{sua belleza}. --- Todas as perfeições divinas são uma e a mesma coisa com a essencia divina; porque \textsc{Deus} é perfeito pela sua essencia. Ora a belleza é uma perfeição. Logo a belleza é uma e a mesma coisa com a essencia divina; e por isso \textsc{Deus}, assim como é a sua essencia, assim tambem é a sua \emph{belleza}. --- Além d'isso, consistindo a belleza divina numa infinita variedade de perfeições, reduzidas a uma simplicissima unidade, é evidente que \textsc{Deus} será a sua belleza, se fôr a sua perfeição e a sua unidade. Ora \textsc{Deus} é a sua perfeição e a sua unidade; porque tudo o que está em \textsc{Deus}, é \textsc{Deus}. Logo \textsc{Deus} é a sua \emph{belleza} \footnote{\textsc{Deus} é a propria belleza! Quem poderá comprehender e exprimir todo o encanto, toda a suavidade, que esta expressão encerra?... A belleza constitue a propria Essencia de \textsc{Deus}!... Os escriptores dizem que a belleza de Deus é \emph{summa}, \emph{infinita}, \emph{incorruptivel}, \emph{immutavel}, \emph{amabilissima}; mas essas palavras levam-nos ao perfeito conhecimento da belleza divina? Não. Só \textsc{Deus} contempla e comprehende a sua belleza infinita e eterna; só \textsc{Deus} ama a sua belleza com amor infinito e eterno; só \textsc{Deus} é infinitamente bemaventurado na contemplação e no amor da sua belleza... A creatura racional vê um pallido reflexo d'esta belleza infinita nas coisas visiveis, e d'esse reflexo, depurado das imperfeições, a intelligencia forma um ideal de belleza, que é superior a toda a belleza creada e representa, ainda que de um modo muito inadequado e deficiente, a ineffavel belleza divina. Esse ideal é sufficiente para alma se destacar de todos os encantos creados e se unir intimamente, pela fé e pelo amor, ao Summo Bem... Isto é muito, mas não é tudo. Deus quiz, na sua infinita misericordia, elevar o homem a uma ordem sobrenatural, que tem por fim a visão \emph{clara} e \emph{intuitiva} da propria essencia de Deus. \emph{Videbimus Deum sicuti est}! A nossa intelligencia verá, pela graça de Deus, a propria essencia divina como é em si mesma, e por isso infinitamente bella, a propria belleza! Não será uma visão \emph{comprehensiva} (porque o finito não pode comprehender o infinito), mas será verdadeira \emph{visão}. Esta visão transformará a alma e a tornará semelhante a Deus, e será principio de um amor e de uma felicidade sem limites! Felizes os que, conscientes da sua excelsa dignidade e do seu eterno destino, desprezam as bellezas caducas e miseraveis das creaturas, e collocam as suas esperanças e os seus amores na unica e immutavel Belleza divina!}.

55. A \textbf{belleza} divina é a \textbf{causa efficiente} de \textbf{toda} a \textbf{belleza creada}. --- A belleza divina será a causa \emph{efficiente} de toda a belleza creada, se esta derivar de \textsc{Deus}, emquanto e porque Elle é dotado de belleza infinita. Ora é certo que a belleza creada deriva de \textsc{Deus}, emquanto e porque Elle é dotado de belleza infinita; porque ninguem pode communicar a outros uma \pg{470}  perfeição, senão emquanto e porque a possue. Logo a belleza divina é a causa \emph{efficiente} de toda a belleza creada \footnote{De \textsc{Deus}, como da causa \emph{efficiente}, deriva toda a belleza, como toda e qualquer qualidade das creaturas. --- A belleza creada é multipla e diversa, porque muitos e diversos são os entes creados. São bellos os seres \emph{inorganicos}, sobretudo os crystaes e as pedras preciosas; principalmente a \emph{luz}, que é esplendida em si e torna esplendidos os outros corpos, e por isso é imagem e symbolo da Belleza divina. --- Mais bellos são os seres \emph{vivos}; sobretudo o homem, em cuja alma espiritual \textsc{Deus} gravou a propria imagem e semelhança. --- Immensamente, ineffavelmente bella é a SS. Humanidade de N. S. Jesus \textsc{Christo}, pela integridade, regularidade e harmonia das suas feições immaculadas, por aquelle esplendor de magestade e de graça, com que a Divindade, que nella habita substancialmente, a circumda e glorifica. --- Mas de todas essas bellezas a fonte unica é a Belleza divina.}.

56. \textbf{A belleza divina é a causa exemplar de toda a belleza creada.} --- Todo o agente tende a imprimir e imprime, de algum modo, no effeito a semelhança da propria forma, por que opera. Ora Deus opera pela sua essencia, que é bella e é a propria belleza. Logo Deus, produzindo as creaturas, imprime nellas a semelhança da sua belleza. Logo a belleza divina é a causa \emph{exemplar} de toda a belleza creada \footnote{Deve haver uma \emph{analogia}, uma \emph{semelhança} entre a belleza divina e a belleza creada, como ha sempre analogia ou semelhança entre a causa e o effeito. --- Mas a \emph{analogia}, a \emph{semelhança} não é \emph{egualdade}. \textsc{Deus} não é causa \emph{univoca} das creaturas; por isso a forma das creaturas não pode ser \emph{especificamente egual} á forma do Creador. Portanto a belleza divina é infinitamente mais pura, mais elevada, mais attrahente. Ora, se a belleza das creaturas, embora imperfeita e caduca, é digna de estima e tanto prende o coração, de quanta estima não é digna e quanto não deve ser amada a infinita e purissima belleza de \textsc{Deus}?!}.

\textbf{57. A belleza divina é a causa final da belleza creada.} --- A belleza divina será a causa \emph{final} da belleza creada, se \textsc{Deus} dotou de belleza as creaturas, para que estas representassem a sua belleza infinita. Ora a intenção de \textsc{Deus}, na producção do mundo, foi effectivamente a de revelar nas visiveis perfeições das creaturas as proprias perfeições invisiveis. Logo a belleza divina é a causa \emph{final} da belleza creada \footnote{Se a belleza divina é a causa \emph{final} da belleza creada, esta belleza creada não pode nem deve considerar-se como o \emph{termo} e o \emph{fim ultimo} das nossas tendencias e dos nossos desejos, mas unicamente como o \emph{meio} e o \emph{degrau} para nos elevarmos á contemplação e ao amor da belleza divina. O espirito, que assim considera a belleza e, em geral, as perfeições das creaturas, não pode deixar de se unir intimamente ao infinito Bem.\par Concluimos este artigo e confirmamos as nossas considerações com as palavras do nosso Angelico Doutor: « Consideratio creaturarum animos hominum in amorem divinae bonitatis accendit. Quidquid enim bonitatis et perfectionis\pgn{471}in diversis creaturis particulariter distributum est, totum in Deo universaliter est adunatum, sicut in fonte totius bonitatis. Si igitur creaturarum bonitas, pulchritudo et suavitas sic animos hominum allicit; ipsius Dei fons bonitatis, rivulis bonitatis in singulis creaturis repertis diligenter comparatus, animos hominum inflammatos totaliter ad se trahet » (C. \emph{Gentes}, l. II, c. 2).}.

\pg{471} CAPITULO TERCEIRO

FACULDADES E \textbf{OPERAÇÕES} IMMANENTES \textbf{DE DEUS}

ARTIGO I.

\textbf{Intelligencia de Deus, sua natureza e objecto.}

58. \textbf{Faculdades} de \textbf{Deus}. --- \emph{Faculdades} de \textsc{Deus} são as perfeições ou os attributos divinos, que se referem á \emph{operação}. --- Estas faculdades são duas: a \emph{intelligencia} e a \emph{vontade} \footnote{Depois de termos tratado dos attributos \emph{quiescentes} de \textsc{Deus}, devemos agora tratar dos seus attributos \emph{operativos}, e, por isso, das \emph{faculdades} e \emph{operações}. --- Sendo a operação uma consequencia e uma expressão da natureza, nunca alcançaremos um pleno conhecimento de um ente, se não conhecermos convenientemente a sua operação. Por quanto, pelo modo e pela qualidade da operação conhecemos a natureza e a qualidade da virtude ou faculdade operativa, e essa virtude ou faculdade manifesta a essencia do ente, visto que um ente opera de um certo e determinado modo, pelo facto de possuir uma certa e determinada essencia. --- Ora, ha duas. especies de operações; uma que permanece no agente, e é perfeição d'elle; outra que tem o termo fóra do agente, e é perfeição do effeito externo. A primeira diz-se \emph{immanente}, porque tem o principio e o termo no mesmo agente; a segunda chama-se \emph{transeunte}, porque tem o termo fóra do agente. --- A primeira precede a segunda. Por isso, tratando das operações divinas, primeiramente fallaremos das \emph{immanentes}, e depois das \emph{transeuntes}.}.

59. Intelligencia de \textbf{Deus}. --- \emph{Intelligencia} de \textsc{Deus} é a faculdade, pela qual Elle conhece a sua essencia e todas as coisas. --- Da intelligencia procede a \emph{sciencia}. --- Diremos da \emph{existencia} d'esta faculdade, da sua \emph{natureza} e \emph{objecto} \footnote{S. Thomaz trata d'este assumpto sobretudo --- na \emph{Sum. Th.}, p. I, q. 14 --- c. \emph{Gentes}, l. I, c. 44 e seg. --- \emph{De Verit.}, q. 2 e 3, --- I \emph{Sent.}, Dist. 35.}.

60. \textbf{Deus é dotado de intelligencia.}

a) Toda a substancia \emph{espiritual,} ou independente da \emph{materia}, é dotada de \emph{intelligencia}. Ora \textsc{Deus} é \emph{purissimo espirito}. Logo \textsc{Deus} é dotado de \emph{intelligencia} \footnote{O grau do conhecimento de um ser é sempre proporcionado ao grau da sua immaterialidade; de modo que, quanto mais um ente se afastar\pgn{472}da materia, tanto mais elevado será o seu conhecimento. Assim a alma das plantas, por estar inteiramente ligada á materia, é destituida de todo e qualquer conhecimento. A alma dos animaes, embora não possa operar sem o organismo, todavia recebe as especies \emph{sensiveis} dos objectos, e por isso os animaes são dotados de conhecimento, embora no infimo grau. A alma humana, por ser \emph{espiritual}, isto é, independente da materia, é dotada de conhecimento \emph{intellectual}, ou de intelligencia. \textsc{Deus}, sendo um espirito \emph{purissimo}, não pode deixar de ser summamente intelligente. Mas escutemos o Angelico Doutor: « Cognoscentia a non cognoscentibus in hoc distinguuntur, quia non cognoscentia nihil habent nisi formam suam tantum; sed cognoscens natum est habere formam (\emph{intentionalem}) etiam rei alterius: nam species cogniti est in cognoscente. Unde manifestum est quod natura rei non cognoscentis est magis coarctata et limitata. Natura autem rerum cognoscentium habet maiorem amplitudinem et extensionem. Propter quod dicit Philosophus (\emph{De anima}, l. 3, t. 37) quod anima est quodammodo omnia (\emph{non realiter}, \emph{sed intentionaliter}). --- Coarctatio autem formae est per materiam. Unde et supra diximus quod formae, secundum quod sunt magis immateriales, secundum hoc magis accedunt ad quamdam infinitatem. Patet igitur quod \emph{immaterialitas} alicujus rei est \emph{ratio} quod sit \emph{cognoscitiva}; et secundum modum immaterialitatis est modus cognitionis. Unde dicitur (\emph{De anima}, l. II, t. 124) quod plantae non cognoscunt propter suam materialitatem. Sensus autem cognoscitivus est, quia receptivus est specierum sine materia; et intellectus adhuc magis cognoscitivus, quia magis separatus a materia, et immixtus; ut dicitur (\emph{De anima}, l. III. t. 4). Unde cum Deus sit in summo immaterialitatis, sequitur quod ipse sit in summo cognitionis » (\emph{Sum. Th.}; p. I, q. 14, a. 1).}.

\pg{472} b) Deus é infinitamente perfeito, e por isso não pode faltar-lhe alguma perfeição \emph{absoluta}. Ora a intelligencia é uma perfeição \emph{absoluta}. Logo \textsc{Deus} é dotado de \emph{intelligencia} \footnote{A \emph{intelligencia}, como tambem a \emph{sciencia}, que é o fructo do exercicio d'esta faculdade, é uma perfeição \emph{absoluta}, e por isso convem e attribue-se a Deus em sentido \emph{positivo} e \emph{proprio}, e não num sentido \emph{metaphorico}, como lhe se attribue a \emph{ira} ou a \emph{misericordia}. De facto, a \emph{intelligencia}, e por isso a \emph{sciencia}, não encerra, no seu conceito formal, imperfeição alguma e não se oppõe a uma perfeição mais elevada, e portanto deve attribuir-se a Deus no sentido \emph{proprio}, no mesmo sentido em que se attribue ao homem e ao Anjo; embora em Deus se encontre de um modo infinitamente superior ao modo por que se encontra nas creaturas, pois em Deus é uma perfeição substancial, que se identifica com a essencia, emquanto nas creaturas é uma perfeição accidental, que se distingue da propria essencia. Logo a \emph{intelligencia}; como tambem a \emph{sciencia}, deve attribuir-se a \textsc{Deus}, ainda que se lhe não possa attribuir a imperfeição, de que se acha cercada nas creaturas. Diz S. Thomaz: « Quia perfectiones procedentes a Deo in creaturas altiori modo sunt in Deo, oportet quod quandocumque nomen sumptum a quaecumque perfectione creaturae Deo attribuitur, secludatur ab ejus significatione omne illud quod pertinet ad imperfectum modum, qui competit creaturae. Unde scientia non est qualitas in Deo, vel habitus, sed substantia et actus purus » (\emph{Sum. Th.}, p. I, q. 14, a. 1 ad 1).}.

\pg{473} c) \textsc{Deus} estabeleceu a ordem, que resplandece no mundo. Ora esta ordem é effeito de uma causa infinitamente intelligente. Logo \textsc{Deus} é dotado de \emph{intelligencia} \footnote{Não adduzimos mais argumentos em defeza d'esta proposição, porque é uma consequencia do que dissemos ácerca da \emph{existencia} e da \emph{essencia} de \textsc{Deus}. Todos os povos admittem que o Ente supremo recebe as homenagens dos mortaes, ouve as supplicas dos que recorrem a Elle, julga as acções humanas, e recompensa ou condemna conforme a qualidade d'essas acções, e por isso admittem que \textsc{Deus} é dotado de \emph{intelligencia}.}.

\textbf{61. A intelligencia de Deus identifica-se com a sua essencia, e é acto purissimo e infinito.}

a) Identifica-se com a sua essencia. --- \textsc{Deus} é \emph{simplicissimo}. Ora, se a sua intelligencia não se identificasse com a sua essencia, \textsc{Deus} seria \emph{composto} de \emph{essencia} e de \emph{attributos}: o que é falso. Logo a intelligencia de \textsc{Deus} identifica-se com a sua essencia. --- Por isso \textsc{Deus}, assim como é a sua \emph{essencia}, assim tambem é a sua \emph{intelligencia}.

b) E' acto purissimo. --- Se a intelligencia de \textsc{Deus} não fosse \emph{acto purissimo}, mas fosse \emph{potencial}, isto é, passasse da \emph{potencia} para o \emph{acto}, ella, como a nossa, seria \emph{mutavel} e \emph{composta}. Ora a intelligencia de \textsc{Deus}, identificando-se com a sua essencia, é \emph{immutavel} e \emph{simplicissima}. Logo é \emph{acto purissimo}.

c) E' acto infinito. --- A intelligencia de \textsc{Deus} identifica-se com a essencia e é acto purissimo. Ora a essencia de \textsc{Deus} é infinita, porque \textsc{Deus} é o proprio ser; e um acto purissimo é tambem infinito, porque é sempre a potencia que limita o acto. Logo a intelligencia de \textsc{Deus} é acto \emph{infinito} \footnote{A conclusão é evidente. Deus é simplicissimo; e tudo o que está em \textsc{Deus} identifica-se com \textsc{Deus}, é o proprio \textsc{Deus}. Por isso a intelligencia, como tambem o acto de entender, são a propria essencia divina; porque, se fossem, como em nós, uma perfeição distincta, haveria em \textsc{Deus} composição de \emph{potencia} e de \emph{acto}: o que repugna. « Em \textsc{Deus}, diz S. Thomaz, a intelligencia, o intelligente, o objecto entendido, a especie intelligivel, e o acto de entender são uma e a mesma coisa. D'ahi se segue que, embora seja \textsc{Deus} infinitamente intelligente, nem por isso se admitte qualquer multiplicidade na sua substancia » (\emph{Sum}. \emph{Th.}, p. I, q. 14, a. 4).}.

\textbf{62. Deus conhece tudo pela sua essencia.}

a) Se \textsc{Deus} conhecesse as coisas por uma especie intelligivel, distincta da sua essencia, essa especie seria o \emph{principio formal} do conhecimento divino, e \textsc{Deus} seria \emph{dependente} de um outro ente na operação e por isso no ser; porque a operação e o ser, em \textsc{Deus}, são uma e a mesma coisa. Ora \textsc{Deus} é a Causa primeira, absolutamente independente de todo e qualquer ente. Logo \textsc{Deus} conhece tudo pela sua \emph{essencia}.

\pg{474} b) \textsc{Deus} conhecerá tudo pela sua \emph{essencia}, se esta representar \emph{adequadamente} tudo o que Elle entende --- o seu proprio ser e os entes finitos, existentes e possiveis; porque, então, não ha razão sufficiente para procurar outra especie intelligivel. Ora a essencia divina representa \emph{adequadamente} --- o ser divino, porque, em \textsc{Deus}, a essencia e o ser são uma e a mesma coisa, --- e todo o ente finito, existente e possivel, porque a essencia divina é a causa efficiente e exemplar de todas as coisas. Logo \textsc{Deus} conhece tudo pela sua \emph{essencia}.

c) A especie intelligivel, principio formal do conhecimento, refere-se á intelligencia, como o acto á potencia. Portanto essa especie, se fosse, em \textsc{Deus}, distincta da essencia, seria \emph{acto} da intelligencia divina, a qual, por consequencia, seria \emph{potencia}. Ora não pode haver, em Deus, composição de acto e de potencia. Logo Deus conhece tudo pela sua \emph{essencia} \footnote{Cf. c. \emph{Gentes}, l. I, c. 47. --- Para a intelligencia do assumpto, lembramos certos principios. O conhecimento consiste nisto que o cognoscente se torna semelhante ao conhecido. Este une-se áquelle por uma forma ou imagem intencional, que se chama \emph{especie intelligivel}, e que actua na intelligencia e lhe representa o objecto. Por isso a perfeição do conhecimento depende inteiramente da perfeição d'esta forma ou especie, a qual, pela sua vez, deve ser proporcionada á virtude ou energia intellectual, que ella determina ao acto. D'ahi se segue --- que, quanto mais perfeita fôr a intelligencia, tanto mais excellente deve ser a especie intelligivel, --- e que, quanto mais excellente fôr a especie, tanto mais perfeita será a sua maneira de representar. Portanto, todo o conhecimento intellectual exige e suppõe tres coisas: a intelligencia cognoscente, o objecto conhecido, e a especie intelligivel. O objecto conhecido chama-se objecto \emph{material} ou \emph{quasi-material} do conhecimento. A razão, pela qual o objecto \emph{material} determina a intelligencia e que é representada ou expressa pela especie intelligivel, diz-se objecto \emph{formal}. Sendo assim, é claro que o objecto \emph{formal}, quando se trata das faculdades \emph{cognoscitivas}, coincide de algum modo com a \emph{especie intelligivel}; porque o objecto \emph{material} não actua na intelligencia senão pela especie e pelo que está representado na especie. Por isso, quando se trata do objecto \emph{formal} de uma intelligencia, procuramos e indicamos a \emph{especie intelligivel}, que actua naquella intelligencia e lhe torna as coisas presentes sob um certo e determinado aspecto. --- O objecto \emph{material} de toda e qualquer intelligencia divide-se em \emph{primario} e \emph{secundario}. \emph{Primario} é o que é proporcionado á intelligencia, é conhecido por si mesmo e por uma especie adequada, e esse conhecimento é a razão pela qual são conhecidas todas as outras coisas; assim o objecto \emph{primario} da nossa intelligencia, no estado presente, é a essencia das coisas materiaes. \emph{Secundario} é o que não é proporcionado á intelligencia, não é conhecido por si mesmo, mas por outra coisa e por uma especie inadequada; assim o objecto \emph{secundario} da nossa intelligencia é constituido por tudo o que podemos conhecer por meio das especies, derivadas das coisas sensiveis ou materiaes.\par Applicando estas noções á intelligencia divina, é evidente que o seu\pgn{475}objecto \emph{formal} não pode ser uma forma ou especie intelligivel, distincta da essencia divina; porque, como dissemos no texto, a intelligencia divina, sendo acto purissimo e infinitamente perfeito, não pode ser aperfeiçoada e completada, em ordem ao acto, por uma especie, distincta d'ella e da essencia. --- Emquanto ao objecto \emph{material}, a intelligencia divina, como vamos provar, comprehende o proprio ser divino e todas as outras coisas, que podem ser de algum modo conhecidas; porque a essencia divina, sendo forma perfeitissima e de infinito valor, representa tudo á intelligencia divina, com a qual é uma e a mesma coisa.\par Portanto, se nos perguntarem se para o acto da intelligencia divina é necessaria a \emph{especie intelligivel}, respondemos que, se a especie \emph{intelligivel} se toma no sentido rigoroso emquanto significa «a semelhança de uma coisa, a qual, á maneira de um principio subjectivo, realmente distincto da intelligencia, completa e determina esta faculdade, por si mesma incompleta e indifferente, ao conhecimento de um determinado objecto», então devemos dizer que ella não é necessaria nem pode existir \emph{formalmente} em Deus, porque não pode admittir-se que a intelligencia divina seja por si incompleta e deficiente com relação ao seu acto; --- mas, se a \emph{especie intelligivel} se toma num sentido \emph{lato}, emquanto significa «a razão sufficiente pela qual a intelligencia conhece um objecto, isto é, emquanto é a semelhança da coisa, que torna a intelligencia capaz de conhecer a mesma coisa», ella é necessaria para o acto da intelligencia divina e existe \emph{formalmente} um \textsc{Deus}, porque, se \textsc{Deus} conhece, deve existir nelle a razão sufficiente do seu conhecimento. --- Esta especie, tomada no sentido \emph{lato} e existente em \textsc{Deus}, não é uma coisa \emph{realmente} distincta da essencia de Deus, nem da sua intelligencia, nem do seu acto intellectual, mas identifica-se com cada uma d'estas coisas, como ensina a cada passo S. Thomaz. --- Dissemos que a especie, em \textsc{Deus}, não se distingue \emph{realmente} da \emph{essencia}, nem da \emph{intelligencia}, nem do \emph{acto intellectual}; porque, \emph{logicamente}, distingue-se da \emph{intelligencia} e do \emph{acto intellectual}. Porquanto, a \emph{especie} considera-se como um principio, que \emph{torna} a \emph{intelligencia capaz} de perceber um certo objecto, e por isso \emph{logicamente} distingue-se da intelligencia e do seu acto. --- Mas nem \emph{logicamente} se distingue da \emph{essencia divina}. Por quanto, a especie deve \emph{conter em si} o \emph{objecto cognoscivel}, para que possa determinar a intelligencia á percepção d'elle; ora a essencia divina, contendo em si todas as coisas, pode e deve exercer o munus da dita especie. Diz S. Thomaz: «Species intelligibilis in Deo nihil aliud est quam \emph{essentia sua}, in qua omnes species rerum comprehenduntur» (\emph{Sum}. \emph{Th}., p. I, q. 14, a. 5 ad 3).\par Da doutrina exposta resultam os seguintes corollarios: --- 1.º) \emph{Unica} é a especie intelligivel em \textsc{Deus}, porque unica é a sua essencia, capaz de determinar por si só á intelligencia divina ao conhecimento de todas as coisas; --- 2.º) Pela mesma unica especie, que é a sua essencia, \textsc{Deus} conhece o proprio ser e todas as coisas, --- o proprio ser pela essencia, que é a mesma coisa com o mesmo ser, --- todas as outras coisas pela essencia, que é uma\pgn{476}semelhança d'essas coisas; --- 3.º) o conhecimento divino não é \emph{especificado} pela especie propria das outras coisas, mas unicamente pela divina essencia, e por isso é sempre infinito (\emph{Sum}. \emph{Th}., p. I, q. 14, a. 5 ad 3).}.

\pg{475} 63. \textbf{Deus conhece perfeitamente} o \textbf{proprio ser}.

a) \textbf{Conhece} o proprio ser. --- A intelligencia de \textsc{Deus}, sendo um acto purissimo e de infinita energia, deve conhecer tudo o que é intelligivel e lhe está presente. Ora o ser divino, não só é intelligivel, porque é immaterial, mas tambem está intimamente \pg{476}  presente á intelligencia, porque é uma e a mesma coisa com ella. Logo \textsc{Deus} conhece o proprio ser \footnote{A conclusão é evidente, e é um corollario, que deriva da natureza do conhecimento intellectual. Este, por ser immaterial e inorganico, é essencialmente capaz de perfeita reflexão, pela qual se conhece a existencia e a natureza do cognoscente. Deveriamos negar á intelligencia divina uma perfeição, que attribuimos á intelligencia humana? Não, com certeza.}.

\textbf{b) Conhece perfeitamente o proprio ser.} --- A perfeição do conhecimento, que consiste na completa adequação entre a cognoscibilidade do objecto e o proprio acto cognoscitivo, depende de duas condições: que a especie intelligivel --- reproduza fielmente o objecto, --- e se una intimamente com a faculdade. Ora a essencia divina, que faz as vezes da especie intelligivel no conhecimento de \textsc{Deus}, não só exprime fielmente o ser divino e se une intimamente com a intelligencia divina, mas identifica-se com o mesmo ser e com a mesma intelligencia. Logo \textsc{Deus} conhece \emph{perfeitamente} o proprio ser \footnote{Podemos expor o argumento numa outra forma. Para a percepção intellectual são necessarios tres elementos: a faculdade intellectual, o objecto intelligivel, a união do objecto com a faculdade; e quanto mais perfeitos forem estes elementos, tanto mais perfeita será a percepção. Ora a intelligencia divina é infinitamente perspicaz, --- o ser divino é infinitamente intelligivel, --- a união do ser com a intelligencia não poderia ser mais intima, porque são uma e a mesma coisa. Logo a percepção que \textsc{Deus} tem do proprio ser é infinitamente perfeita. --- Esta percepção chama-se tambem \emph{comprehensão}; porque \emph{comprehender} um objecto é conhecel-o tanto, quanto elle é cognoscivel, de modo que á intelligencia não escape nenhum elemento, nenhuma propriedade, nenhuma efficacia, de que elle é dotado. Ora o conhecimento de Deus é proporcionado á sua cognoscibilidade, isto é, \textsc{Deus} se conhece a si mesmo tanto, quanto Elle é cognoscivel; porque esse conhecimento e essa cognoscibilidade fundam-se na mesma purissima actualidade do ser divino. Logo \textsc{Deus} \emph{comprehende} o proprio ser. Diz S. Thomaz: «Tunc dicitur aliquid \emph{comprehendi}, quando res cognoscitur ita perfecte, sicut cognoscibilis est... Manifestum est autem quod \textsc{Deus} ita perfecte cognoscit seipsum, sicut perfecte cognoscibilis est... Tanta est autem virtus Dei in cognoscendo, quanta est actualitas ejus in existendo; quia per hoc quod actu est et ab omni materia et potentia separatus, \textsc{Deus} cognoscitivus est. Unde manifestum est quod tantum seipsum cognoscit, quantum cognoscibilis est. Et propter hoc seipsum perfecte \emph{comprehendit}» (\emph{Sum}. \emph{Th}., p. I, q. 14, a. 8).\par Que \textsc{Deus} conheça o proprio ser, como objecto \emph{primario} da sua intelligencia, é fóra de toda a duvida. Objecto \emph{primario} da intelligencia é, como dissemos, o que é proporcionado á intelligencia, é conhecido por si mesmo e por uma especie adequada, e esse conhecimento é a razão pela qual são conhecidas as outras coisas. Ora o ser divino --- é o unico objecto proporcionado\pgn{477} á intelligencia divina, porque só esse ser é infinito como infinita é essa intelligencia, --- só elle é conhecido por si mesmo e por uma especie \emph{adequada}, porque só elle está intimamente unido com a intelligencia divina e só elle é representado por uma especie \emph{adequada}, que é a divina essencia, --- e o conhecimento do ser divino é a razão pela qual são conhecidas todas as coisas, porque \textsc{Deus} conhece tudo na sua essencia e pela sua essencia. Logo o ser divino é o objecto \emph{primario} da intelligencia divina.}.

%%>>> cap-7

\chapter*{CAPITULO III}
\addcontentsline{toc}{chapter}{CAPITULO III}

\pg{477} 64. \textbf{Deus} conhece \textbf{todas as} outras \textbf{coisas}, \textbf{existentes} e \textbf{possiveis}.

a) A intelligencia de Deus, identificando-se com o proprio ser d'Elle, é infinitamente perfeita. Ora uma intelligencia infinitamente perfeita não pode deixar de conhecer tudo o que é cognoscivel, isto é, tudo o que tem alguma razão de entidade, quer exista quer não exista; porque, se uma coisa só, existente ou possivel, lhe escapasse, ella seria necessariamente limitada e imperfeita. Logo \textsc{Deus} conhece todas as outras coisas, existentes e possiveis.

b) \textsc{Deus} conhece \emph{perfeitamente} o proprio ser, isto é, conhece-o tanto, quanto é cognoscivel. Ora \textsc{Deus} não conheceria perfeitamente o proprio ser, se não conhecesse tambem todas as outras coisas, existentes e possiveis. Por quanto, um ente não conhece perfeitamente o proprio ser, se não conhece perfeitamente a efficacia e a extensão do seu poder, --- e não conhece perfeitamente a efficacia e a extensão do seu poder, se não conhece todos os effeitos, que esse poder produz ou pode produzir. Mas o poder de \textsc{Deus} extende-se a todas as coisas, existentes e possiveis; porque Elle é a primeira causa efficiente de tudo o que existe ou pode existir. Logo \textsc{Deus} conhece todas as outras coisas, existentes e possiveis\footnote{E' o argumento de S. Thomaz, que diz: « Respondeo dicendum quod necesse est Deum cognoscere alia a se. Manifestum est enim quod seipsum perfecte intelligit; alioquin suum esse non esset perfectum, cum suum esse sit suum intelligere. Si autem perfecte aliquid cognoscitur, necesse est quod virtus ejus perfecte cognoscatur. Virtus autem alicujus rei perfecte cognosci non potest, nisi cognoscantur ea, ad quae virtus se extendit. Unde, cum virtus divina se extendat ad alia, eo quod ipsa est prima Causa effectiva omnium entium, necesse est quod Deus alia a se cognoscat » (\emph{Sum}. \emph{Th}., p. I, q. 14, a. 5).\par Do argumento conclue-se --- que \textsc{Deus} conhece todas as outras coisas, existentes e possiveis, --- \emph{não immediatamente}, em \emph{si mesmas}, --- mas em \emph{si mesmo}, isto é, na propria \emph{essencia} como num \emph{meio conhecido}, --- e que por isso todas as coisas constituem o objecto \emph{secundario} da intelligencia divina. Na verdade, o que se conhece immediatamente em si mesmo, e por uma especie adequada ao proprio cognoscivel, é o objecto \emph{primario}; emquanto o que se conhece mediatamente num outro ente, e por uma especie propria d'esse outro ente, é o objecto \emph{secundario}. Ora \textsc{Deus} conhece todas as outras coisas, existentes e possiveis, pela sua essencia, isto é, por uma\pgn{478}especie, que é adequada unicamente ao seu ser mas excede immensamente todas as outras coisas, nella \emph{virtualmente contidas}, e por isso conhece o proprio ser como objecto \emph{primario}, e todas as outras coisas como objecto \emph{secundario}. --- E não pode ser d'outro modo. A intelligencia divina não entende as coisas, senão pela sua essencia, que exerce o munus de especie intelligivel. Ora a essencia determina necessariamente a intelligencia a conhecer \emph{immediata} e \emph{directamente} a mesma essencia, e \emph{mediata} e \emph{indirectamente} as outras coisas, que estão nella virtualmente contidas. Por quanto, \textsc{Deus} só conhece as outras coisas, porque conhece a sua essencia, --- não emquanto é uma essencia, --- mas emquanto pode ser e é diversamente participada e imitada por muitas coisas. Diz S. Thomaz (l. c.): « Dupliciter aliquid cognoscitur: uno modo \emph{in seipso}, alio modo \emph{in altero}. In seipso quidem cognoscitur aliquid, quando cognoscitur per \emph{speciem adaequatam ipsi cognoscibili}; sicut cum oculus videt hominem per speciem hominis. In alio autem videtur id quod videtur \emph{per speciem continentis}; sicut cum pars videtur in toto per speciem totius, vel cum homo videtur in speculo per speciem speculi, vel quocumque alio modo contingat aliquid in alio videri. Sic igitur dicendum est, quod Deus seipsum videt in seipso, quia seipsum videt per essentiam suam. Alia autem a se videt non in seipsis, sed in se ipso, inquantum essentia sua \emph{continet} similitudinem aliorum ab ipso ».\par Como apparece das noções expostas, o que é objecto \emph{primario} para a intelligencia humana é \emph{secundario} para a intelligencia divina, e o que é \emph{secundario} para nós é \emph{primario} para \textsc{Deus}. Na verdade, o objecto \emph{primario} da nossa intelligencia, no estado presente, é a essencia abstracta das coisas materiaes; porque esta essencia --- é proporcionada á natureza e ao modo de ser da nossa intelligencia, --- é conhecida por uma especie adequada, --- e é a razão, pela qual conhecemos todas as outras coisas, e o proprio Deus, que, por isso, é objecto \emph{secundario} da nossa intelligencia. Ma o objecto \emph{primario} da intelligencia divina é a propria essencia divina; porque só ella --- é proporcionada á natureza e ao modo de ser da intelligencia divina, --- é conhecida por uma especie adequada, --- e é a razão pela qual \textsc{Deus} conhece todas as outras coisas, que, por isso, constituem o objecto \emph{secundario} da intelligencia divina. Nós não podemos conhecer a \textsc{Deus}, senão depois de conhecermos as creaturas; \textsc{Deus} não pode conhecer as creaturas, senão \emph{depois} de se ter conhecido a si mesmo. --- Por isso a essencia divina é o meio \emph{in quo} \textsc{Deus} percebe todas as coisas; porque, como dissemos, esse meio \emph{primeira} ou \emph{directamente} é percebido em si mesmo e \emph{depois} leva á percepção de outro objecto, e portanto é, ao mesmo tempo, \emph{objecto} e \emph{meio}. --- E' claro que, no caso presente, não se trata de prioridade de \emph{tempo}, mas de \emph{natureza}; porque \textsc{Deus}, com um só acto, que é a sua substancia, vê tudo numa unica coisa, que é a sua essencia (\emph{Sum}. \emph{Th}., p. I, q. 14, a. 7).\par Do que temos dito resulta que a sciencia divina, emquanto aos objectos entendidos, pode e deve dividir-se em sciencia de \emph{visão} e de \emph{simples intelligencia}. --- Sciencia \emph{de visão} é aquella, pela qual \textsc{Deus} conhece o proprio ser e todas as creaturas \emph{existentes} numa determinada differença de tempo, --- e por isso passadas, presentes e futuras. Sciencia de \emph{simples intelligencia}\pgn{479}é aquella, pela qual \textsc{Deus} conhece as coisas meramente \emph{possiveis}, que não existem e nunca hão de existir. --- E' a divisão dada por S. Thomaz (\emph{Sum}. \emph{Th}., p. I, q. 14, a. 9), e admittida pela eschola thomista, que a proclama \emph{adequada}, porque entre a sciencia de \emph{visão} e a de \emph{simples intelligencia} não pode conceber-se outra, sendo \emph{adequada} a divisão do objecto. --- Outros theologos, como veremos, admittem uma terceira especie de sciencia, que chamam \emph{media} e que teria por objecto os \emph{futuros condicionados livres}.\par Para a solução das difficuldades contra a doutrina exposta, indicamos os seguintes principios. --- 1º) O objecto conhecido, representado na especie intelligivel, aperfeiçõa o cognoscente --- não emquanto ao ser \emph{real} d'elle, --- mas emquanto ao ser \emph{intencional}; porque a especie intelligivel é acto ou principio formal da intelligencia, e todo o acto ou toda a forma aperfeiçõa a potencia. Ora a especie, que se une com a intelligencia divina, não é uma coisa distincta de \textsc{Deus}, mas é a sua propria essencia; e por isso só \textsc{Deus} é a perfeição de si mesmo. (\emph{Sum}. \emph{Th}., p. I, q. 14, a. 5 ad 4) --- 2º) O meio, pelo qual se conhece uma coisa, deve ser \emph{proporcionado} á propria coisa --- não emquanto ao \emph{ser}, como se o meio e a coisa devessem ter o mesmo ser real, --- mas emquanto á \emph{representação}, de modo que o meio represente a coisa. Ora a essencia divina, embora exceda infinitamente o ser das creaturas, representa-as perfeitamente, e por isso é um meio proporcionado para o conhecimento d'ellas. (\emph{De} \emph{Ver}., q. 2, a. 3 ad 4). --- 3º) O conhecimento de \textsc{Deus} pode chamar-se \emph{multiplice} e \emph{diverso} --- emquanto ao \emph{objecto conhecido}, que é multiplice e diverso, --- mas não emquanto ao \emph{principio} ou ao \emph{meio} de conhecimento, que é unico, porque é a propria essencia divina. A \emph{unidade} do conhecimento, se se considera no sujeito cognoscente, depende da \emph{unidade} da forma intelligivel, que, em \textsc{Deus}, é a propria essencia divina. Esta, pela sua perfeição infinita, representa tudo o que tem razão de \emph{ente}, e d'ella deriva a multiplicidade e diversidade dos objectos entendidos. --- 4º) A sciencia de \textsc{Deus} é, como veremos, a \emph{causa} das coisas, emquanto \textsc{Deus} produz as creaturas em conformidade com as ideas exemplares da sua intelligencia; mas é causa juntamente com a vontade e com o poder. Por isso, embora \textsc{Deus} conheça as coisas desde a eternidade, não se segue d'ahi que desde a eternidade existam as coisas. O acto da vontade divina com relação ás creaturas é necessariamente livre, e por isso \textsc{Deus} pode querer e deixar de querer a existencia das creaturas, e pode querel-a nesta ou naquella differença de tempo. (\emph{Sum Th}., p. I, q. 14, a. 8) --- 5º) A infinita perfeição de \textsc{Deus} exige que Elle --- não só conheça tudo o que, de algum modo, é, --- mas tambem que o conheça de um modo perfeitissimo, proporcionado á excellencia do ser divino. Ora é mais perfeito que \textsc{Deus} conheça as coisas na divina essencia, do que na essencia d'ellas. --- Nem por isso este conhecimento é indeterminado ou imperfeito; porque, como diremos, a essencia divina exprime \emph{distincta} e \emph{adequadamente} todas as creaturas, e por isso representa não só os elementos mais \emph{communs} d'ellas, mas tambem os mais \emph{proprios}, até ás ultimas differenças (\emph{Sum}. \emph{Th}., p. I, q. 13, a. 6).\par Tendo provado que \textsc{Deus}, cuja intelligencia e sciencia são uma e a mesma coisa com a substancia, conhece na sua essencia tudo o que é de\pgn{480}algum modo cognoscivel, não seria absolutamente necessaria a questão relativa a alguns entes particulares --- se e como \textsc{Deus} os conhece. Todavia, como a nossa intelligencia limitada encontra muita difficuldade em conceber o modo por que a intelligencia divina conhece este ou aquelle ente, julgamos util tratar este assumpto. --- Os objectos, que são o termo do conhecimento divino, podem considerar-se em ordem á \emph{essencia}, ou em ordem á \emph{existencia}. Se se consideram em ordem á \emph{essencia}, podemos perguntar: se e como \textsc{Deus} conhece os entes \emph{singulares} na sua \emph{materialidade}, e os \emph{males}. Se se consideram em ordem á \emph{existencia}, podemos perguntar: se e como \textsc{Deus} conhece os \emph{futuros necessarios}, os \emph{futuros livres}, quer \emph{absolutos} quer \emph{condicionaes}, e todos os \emph{possiveis}. Vamos responder a cada uma d'estas perguntas.\par a) \textsc{Deus} \emph{conhece os entes singulares na sua materialidade}. \textsc{Deus} deve conhecer tudo o que Elle produz pela sua arte. Ora a arte divina não \emph{só} produz a \emph{forma}, pela qual o ente é collocado numa certa \emph{especie}, mas tambem produz a \emph{materia}, que é principio de \emph{individuação}. Logo \textsc{Deus} conhece todos os entes \emph{materiaes}, tanto com relação á \emph{materia}, quanto com relação á \emph{forma}, e por isso conhece os entes \emph{singulares} na sua \emph{materialidade}. (Cf. \emph{De} \emph{Ver}., q. 2, a. 5). --- Alem d'isso, é uma perfeição para a nossa intelligencia o conhecimento dos entes \emph{singulares} e \emph{materiaes}. Ora toda a perfeição da creatura deve encontrar-se, num grau mais elevado, no Creador. Logo \textsc{Deus} conhece os entes \emph{singulares} e \emph{materiaes} de um modo mais perfeito, do que nós. Nós conhecemos, por meio da \emph{intelligencia}, o \emph{universal} e o \emph{immaterial}, e, por meio dos \emph{sentidos}, o \emph{singular} e o \emph{material}; \textsc{Deus} conhece o \emph{universal} e o \emph{singular}, o \emph{immaterial} e o \emph{material} pela mesma faculdade intellectual. --- Diz S. Thomaz: « Cum Deus sit causa rerum per suam scientiam, intantum se extendit scientia Dei, inquantum se extendit ejus causalitas. Unde, cum virtus activa Dei se extendat non solum ad formas, a quibus accipitur ratio universalis, sed etiam usque ad materiam, necesse est quod scientia Dei usque ad singularia se extendat, quae per materiam individuantur. Cum enim sciat alia a se per essentiam suam, inquantum est similitudo rerum; velut principium activum earum, necesse est quod essentia sua sit principium sufficiens cognoscendi omnia, quae per ipsum fiunt, non solum in universali, sed etiam in singulari. Et esset simile de scientia artificis, si esset productiva totius rei, et non formae tantum » (\emph{Sum}. \emph{Th}., p. I, q. 14, a. 11).\par Para a solução das dificuldades, notamos com S. Thomaz (\emph{art}. \emph{cit}.): --- 1º) A nossa intelligencia abstrahe dos principios individuantes a especie intelligivel. Por isso a especie intelligivel da nossa intelligencia não pode representar os principios individuantes, e não podemos ter um conhecimento \emph{directo} dos entes singulares e materiaes. Mas a especie intelligivel da intelligencia divina, que é a propria essencia de \textsc{Deus}, não é immaterial por abstracção, mas por si mesma, e é o principio de todos os elementos, quer especificos, quer individuaes, que entram na composição da coisa. Por isso \textsc{Deus} conhece, pela sua essencia, não só o universal, mas tambem o singular. --- 2º) Ainda que a especie da intelligencia divina não contenha as condições materiaes, como a especie recebida nos nossos sentidos, comtudo pode representar e representa o immaterial e o material.\pgn{481}Para a perfeição do conhecimento não se requer que a coisa se encontre no cognoscente no seu modo de ser; mas basta que a coisa se conheça como ella é. --- 3º) A materia, embora, por ser pura potencia, seja muito dessemelhante de \textsc{Deus}, que é puro acto, comtudo, emquanto é, tem uma certa semelhança com o \emph{ser} divino.\par b) \textsc{Deus} \emph{conhece os males}. --- Quem entende perfeitamente uma coisa, deve tambem entender tudo o que a essa coisa pode acontecer. Ora ha \emph{bens}, que podem ser corrumpidos pelos \emph{males}. Logo \textsc{Deus} não entenderia perfeitamente os \emph{bens}, se não entendesse tambem os \emph{males}... Mas uma coisa emtanto é cognoscivel, emquanto é. Ora, consistindo o mal na \emph{privação} do bem, e a \emph{privação} sendo \emph{nada}, o mal não se conhece em si mesmo, mas no bem, de que é privação; e por isso \textsc{Deus}, pelo facto de conhecer os bens, conhece tambem os males, como pela luz se conhecem as trevas. --- Notamos --- 1º) a sciencia de \textsc{Deus} não é causa do mal, mas é causa do bem, pelo qual conhecemos o mal; --- 2º) o conhecimento de uma coisa por meio de outra é imperfeito, quando aquella coisa é cognoscivel por si. Ora o mal não é cognoscivel por si, porque consiste essencialmente na privação do bem, e por isso o mal não pode ser definido, nem conhecido, senão por meio do bem. Diga-se o mesmo com relação ao \emph{peccado}, e, em geral, a todas as \emph{privações} (\emph{Sum}. \emph{Th}., p. I, q. 14, a. 10).\par c) \textsc{Deus} \emph{conhece os futuros contingentes absolutos}. --- \emph{Futuro} é o ente que actualmente não existe, mas que existirá com certeza numa certa differença de tempo. Distingue-se do \emph{possivel}, que nunca existirá. --- Divide-se o \emph{futuro} em \emph{necessario} e \emph{contingente}. \emph{Necessario} é o que deriva de uma causa \emph{necessaria}, isto é, que não pode deixar de produzir um certo effeito. \emph{Contingente} é o que deriva de uma causa \emph{contingente}, isto é, que pode produzir ou deixar de produzir um certo effeito. --- Se a causa \emph{contingente} é \emph{livre}, o effeito futuro chama-se \emph{contingente livre}. --- O futuro \emph{contingente} subdivide-se em \emph{absoluto} e \emph{condicionado}. O contingente \emph{absoluto} é o que será certamente produzido pela causa contingente. O contingente \emph{condicionado} é o que nunca será produzido pela causa contingente, mas seria por ella produzido, se se verificasse esta ou aquella condição. --- O \emph{futuro contingente condicionado} costuma chamar-se \emph{futurivel}. --- Feitas estas distincções, provamos a nossa proposição com as palavras de S. Thomaz (\emph{Sum}. \emph{Th}., p. I, q. 14, a. 13).\par O \emph{futuro contingente absoluto} pode considerar-se sob dois aspectos --- em \emph{si mesmo} e \emph{na sua causa}. Considerado \emph{em si mesmo}, este contingente não é \emph{futuro}, mas é \emph{presente}, --- nem é \emph{indifferente} para a existencia e para a não existencia, mas é \emph{determinado} para a existencia, e, sob este aspecto, pode ser objecto de um conhecimento certo e infallivel. Considerado \emph{na sua causa}, este contingente é \emph{futuro} e é \emph{indifferente} para a existencia e para a não existencia, porque a \emph{causa contingente} é indifferente com relação aos dois termos oppostos; e por isso o mesmo contingente não pode ser conhecido com certeza. Portanto quem conhece o effeito contingente só na respectiva causa, só pode ter d'elle um conhecimento conjectural. \textsc{Deus}, porém, conhece todos os contingentes, não só emquanto estão contidos nas proprias causas, mas tambem emquanto subsistem actualmente em si mesmos.\par \pgn{482}E, ainda que os contingentes recebam a existencia successivamente, todavia \textsc{Deus} não os conhece successivamente, como nós, mas conhece-os todos \emph{ao mesmo tempo}; porque o seu conhecimento é eterno, como eterno é o seu ser, e a eternidade contem todo o tempo e corresponde a todas as differenças do tempo. Por isso, todas as coisas, que existem no tempo, estão presentes, desde a eternidade, a \textsc{Deus}; não só porque \textsc{Deus} possue as idéas exemplares de todas as coisas, mas tambem porque o seu olhar vê, desde a eternidade, todas as coisas, como ellas existem em si mesmas. E', pois, manifesto que os contingentes são infallivelmente conhecidos por \textsc{Deus}, emquanto estão presentes aos olhos divinos, e comtudo são futuros contingentes, se se compararem com as suas causas proximas.\par \emph{A certeza} da sciencia divina com relação aos futuros contingentes \emph{absolutos} não destroe a \emph{liberdade} dos nossos actos. Na verdade, o conhecimento de \textsc{Deus} é verdadeiro. Ora o conhecimento, para que seja verdadeiro, deve conformar-se com o seu objecto; de modo que se limite a descobrir a natureza do objecto, sem alteral-a. Ora a natureza do futuro \emph{contingente} consiste em que elle possa existir ou deixar de existir. Logo esta indifferença não pode ser destruida ou alterada pela sciencia divina. --- Além d'isso, a sciencia divina refere-se aos futuros contingentes, como a nossa vista se refere aos objectos presentes. Ora é claro que os objectos, que nos estão presentes, não perdem a contingencia, e não se tornam necessarios pelo facto de serem vistos por nós, mesmo de um modo infallivel. Logo nem os futuros contingentes deixam de ser contingentes pelo facto de serem conhecidos infallivelmente por \textsc{Deus}. Diz S. Thomaz: « Sicut.... noster visus non fallitur unquam, videns contingentia ut sunt praesentia, et tamen ex hoc non removetur, quin illa \emph{contingenter} eveniant; ita \textsc{Deus} infallibiliter videt omnia contingentia, sive quae nobis sunt praesentia, sive quae praeterita, sive quae futura, quia sibi non sunt futura, sed ea inspicit esse tunc quando sunt; unde ex hoc non removetur quin \emph{contingenter} eveniant » (\emph{De} \emph{Ver}., q. 2, a. 12). --- Em poucas palavras: o acto \emph{contingente}, que \textsc{Deus} prevê, realizar-se-ha \emph{infallivelmente}, mas não \emph{necessariamente}. Se a \emph{necessidade} tira a liberdade, não a tira a \emph{infallibilidade}.\par Objectam: Tudo o que \textsc{Deus} conhece desde a eternidade, como existente, é um ente \emph{necessario}. Ora nenhum futuro \emph{contingente} é \emph{necessario}. Logo \textsc{Deus} não conhece os futuros contingentes.\par Resposta. Tudo o que \textsc{Deus} conhece desde a eternidade, como existente, quer seja ente \emph{necessario} com relação ás suas causas proximas, quer seja \emph{contingente}, existe \emph{necessariamente}, porque é conhecido por \textsc{Deus} como presente, e não pode deixar de existir o que \textsc{Deus} conhece como existente; aliás haveria engano na sciencia divina. Mas a \emph{necessidade} de existencia, que se attribue a um effeito, é duplice: uma \emph{absoluta}, quando o effeito deriva de causas necessarias; --- outra \emph{condicionada}, quando o effeito deriva de causas contingentes, porque tudo o que é, emquanto é, não pode deixar de ser. Ora esta necessidade \emph{condicionada} deve attribuir-se aos effeitos \emph{contingentes}, cuja existencia é vista por \textsc{Deus}; assim como é necessario que Socrates\pgn{483} escreva, quando escreve, e nós vemos que escreve. Portanto, se com o futuro contingente não é compativel a necessidade \emph{absoluta}, é certamente compativel a necessidade \emph{condicionada}. (\emph{Cf}. \emph{S}. \emph{Th}., p. I, q. 14, a. 3 ad 3; \emph{De} \emph{Ver}., q. 2, a. 12; c. \emph{Gentes}, l. I, c. 67).\par d) \textsc{Deus} \emph{conhece} os \emph{futuros contingentes condicionados}. --- A questão limita-se aos \emph{contingentes condicionados livres}. --- Que \textsc{Deus} conheça estes \emph{futuros condicionados livres}, todos os philosophos espiritualistas o admittem. Esses futuros encerram ou constituem uma verdade, uma entidade; e tudo o que é verdade e entidade não pode escapar á sciencia divina. --- A difficuldade não está no \emph{facto}, está no \emph{modo} de explicar o facto. Com effeito, não podemos dizer que \textsc{Deus} vê esses futuros \emph{condicionados livres}, como vê os futuros \emph{contingentes} ou os \emph{livres absolutos}; porque estes estão \emph{presentes} aos olhos de \textsc{Deus}, e são vistos emquanto \emph{presentes}, porque existem numa differença de tempo, e a duração simultanea da eternidade divina abrange todas as differenças do tempo, mas os futuros \emph{condicionados livres} não são e não serão nunca \emph{presentes}. --- Como, pois, vê \textsc{Deus} os futuros \emph{condicionados livres}? --- Ha duas respostas, dadas por duas escholas, egualmente celebres, --- a \emph{thomista} e a \emph{molinista}. Vamos expôr estas respostas com toda a lealdade e com a clareza possivel em materia tão intricada.\par a) Resposta da eschola \emph{thomista}. --- \textsc{Deus} não vê os futuros \emph{condicionados livres} em si mesmos, na propria actualidade, porque não existem e não existirão em nenhuma differença de tempo; por isso deve vel-os na causa, que os contém e que teria podido dar-lhes a existencia. Esta causa não pode deixar de ser a propria Causa primeira, principio adequado de tudo o que existe ou que teria podido existir. \textsc{Deus} vê tudo isso nos \emph{seus decretos}.\par A essencia divina, --- emquanto é imitavel ao infinito, é, como veremos, o fundamento de todas as possibilidades necessarias e eternas, --- mas, emquanto a realisação de algumas d'estas possibilidades é decretada pela vontade divina, é o fundamento de todas as existencias contingentes, --- presentes, passadas ou futuras. Por quanto, as coisas --- emtanto são possiveis, emquanto \textsc{Deus} é, --- e, emtanto são existentes, emquanto \textsc{Deus} as quer ou permitte. O conhecimento das primeiras é, como dissemos, o objecto da sciencia de \emph{simples intelligencia}; o conhecimento das segundas é o objecto da sciencia de \emph{visão}. --- A sciencia de \emph{simples intelligencia} extende-se á \emph{essencia} e ás \emph{necessarias consequencias} das coisas. A sciencia de \emph{visão} abrange todas as realidades contingentes, queridas ou permittidas nos decretos eternos, não exceptuados os \emph{futuros livres}, que permanecem \emph{livres}, não obstante a infallivel efficacia da Causa primeira. Ora estes decretos eternos, infinitamente perfeitos, devem extender-se a todos os casos, não só aos existentes mas tambem aos meramente possiveis, e nelles \textsc{Deus} lê não só tudo o que acontece, mas tambem tudo o que aconteceria, se se verificassem certas condições, que nunca se hão de verificar. Tal é a sciencia dos \emph{futuriveis}.\par Por isso, para que \textsc{Deus} saiba o que fará Pedro, se se achar nestas ou naquellas circumstancias, deve apenas consultar o decreto geral da sua Pro\pgn{484}videncia, que regula o nosso livre arbitrio, e ver as soluções, que nelle estão virtualmente contidas para cada um dos casos.\par O \emph{Cardeal Zigliara} escreve as seguintes palavras: « \textsc{Deus} nihil potest cognoscere per aliquid sibi extrinsecum et a se independens, a quo cognitionem accipiat et in cognoscendo dependeat, sed omnia cognoscit in sua essentia ut in medio \emph{in quo} cognito. Sed essentia divina non est causa futuri, sive absoluti sive conditionati, praecise ut essentia est; secus enim, cum ipsa sit eodem modo exemplar seu idea omnium possibilium, aut haec omnia dicenda essent futura, aut non esset ratio cur aliqua prae aliis possibilibus ad existentiam pervenirent. Ergo essentia divina est medium in quo Deus videt quaecumque futura ut est determinata per liberum decretum divinae voluntatis, determinantis alia possibilia aliquando futura, alia vero non; scilicet \textsc{Deus} videt in sua essentia omnia futura, quia libere voluit esse futura: quod est videre futura in suo libero decreto » (\emph{Sum}. \emph{Phil}., \emph{Theol}., l. III, c. 2, a. 3).\par b) Resposta da eschola \emph{molinista}. --- A sciencia de \emph{visão} nos decretos divinos e a sciencia de \emph{simples intelligencia} não explicam a presciencia dos futuros \emph{condicionados livres}. E' necessaria admittir em \textsc{Deus} uma terceira sciencia --- a \emph{sciencia media}. --- E' claro que estas divisões da sciencia divina, unica e simplicissima em si mesma, se referem á variedade dos objectos, e significam a nossa maneira de conceber as coisas divinas.\par O que é a \emph{sciencia media}? Ella não é só a sciencia dos futuros condicionados livres. Se assim fosse, não haveria divergencia de opiniões; porque todos admittem que \textsc{Deus} tem o conhecimento d'esses objectos. A \emph{sciencia media} é uma explicação systematica do \emph{meio}, em que \textsc{Deus} vê os futuros condicionados livres, --- é a sciencia dos \emph{futuriveis}, independente dos decretos divinos e anterior a elles. Diz o theologo \emph{Tournely}: « scientia media est scientia conditionatorum, independens ab omni decreto absoluto et efficaci, eoque anterior » (\emph{De scientia Dei}, a. 5). Portanto, os effeitos do nosso livre arbitrio são independentes da vontade divina; e, se \textsc{Deus} os conhece, deve conhecel-os anteriormente a todo e qualquer decreto.\par Mas se \textsc{Deus} não conhece os \emph{futuriveis} nos decretos da sua vontade, em que \emph{meio} os conhece? --- \emph{Molina} responde que Deus os conhece nas \emph{causas livres}, que são capazes de produzir esses futuros em certas circumstancias e das quaes Elle tem uma comprehensão perfeita, uma \emph{supercomprehensão}. Diz assim: « \textsc{Deus} tem uma profunda e insondavel comprehensão de cada livre arbitrio; Elle vê claramente o que cada uma das causas livres faria, de modo espontaneo, nestas e naquellas circumstancias, e mesmo na infinidade das circumstancias. Esta vista de \textsc{Deus} é o que nós chamamos \emph{sciencia media} » (\emph{Concord}., q. 14, a. 13, disp. 52).\par Todavia este \emph{meio} é rejeitado hoje por todos os philosophos, incluidos os neo-molinistas. Diz o \emph{Card}. \emph{Mazzella}: « \textsc{Deus} não pode conhecer os \emph{futuriveis} nas suas causas proximas. Porquanto os actos \emph{livres} não têm um nexo \emph{necessario} com as proprias causas, embora dispostas completa e proximamente para o acto; pois, se a causa livre não fosse, mesmo nesse momento, indifferente para operar e deixar de operar, já não seria livre. E este argumento é invincível; porque aqui não se trata de um conhecimento \emph{conjectural}, mas de um conhecimento \emph{infallivel} dos futuriveis, com o qual é \emph{metaphysicamente} incompativel o erro, qual é o conhecimento de \textsc{Deus} » (\emph{De gratia}, disp. 3, a. 7).\par Nem vale recorrer á infinita energia da vista divina e á sua \emph{supercomprehensão}. Quanto mais profundo e penetrante é o olhar de quem contempla a natureza de uma causa livre, tanto mais clara e distinctamente conhece que os effeitos, que ella contém, são ainda indeterminados e meramente possiveis. Logo \textsc{Deus} não pode ver, nas causas livres, as determinações, que ainda não existem; a não ser que se diga que estas determinações existem: o que seria a negação da liberdade. Por isso, esta explicação leva logicamente ou a negar a liberdade humana, ou a reduzir a sciencia divina á condição de sciencia \emph{conjectural}.\par O \emph{meio}, proposto por \emph{Molina} e rejeitado pelos proprios sequazes, ainda não foi substituido. Geralmente, os \emph{neo-molinistas} affirmam que \textsc{Deus} conhece os futuriveis exclusivamente na contemplação das possibilidades eternas ou das essencias das coisas, independentemente de todo e qualquer decreto. --- Esta explicação não serve para o escopo. Comprehende-se facilmente que a Essencia de \textsc{Deus} possa representar e represente á sua intelligencia todas as possibilidades eternas e necessarias, isto é, todas as essencias das coisas contingentes e todos os effeitos \emph{necessarios}, absolutos ou condicionados, que estão nellas contidos. Mas, se os futuros \emph{livres} estivessem tambem contidos nas essencias das coisas, a futurição d'elles, longe de ser livre, seria absolutamente necessaria, e tudo seria regulado pela fatalidade. Logo \textsc{Deus} não pode ver na essencia das coisas os futuros \emph{livres}, que ahi se não encontram.\par Outros, então, appellam para o \emph{meio das contradictorias}. Todos os futuriveis, dizem elles, podem ser representados sob a forma de proposições \emph{contradictorias}; por exemplo: se Pedro é tentado nestas circumstancias, peccará ou deixará de peccar. Não ha meio termo. E, como estas proposições não podem ser ambas verdadeiras ao mesmo tempo, mas é necessario que uma seja verdadeira, e outra falsa, pela essencia mesma das coisas e anteriormente a todo o decreto divino, é claro que a intelligencia divina, a qual conhece a verdade toda, deve conhecer com certeza absoluta qual d'estas proposições é a verdadeira, e qual é a falsa. --- Mas, se o principio de contradicção é sufficiente para mostrar que as duas contradictorias não se podem realisar simultaneamente, não o é para mostrar qual das duas se realisará em certas e determinadas circumstancias; pois o livre arbitrio torna-as egualmente possiveis. Se uma das duas fosse determinada para a existencia pelo principio de contradicção, ella seria necessaria e não livre. E, se ambas são \emph{possiveis}, qual é a razão sufficiente, pela qual \textsc{Deus} vê que esta será realisada, e aquella não? Para responder a esta pergunta, não ha outro remedio senão recorrer á \emph{supercomprehensão} de \emph{Molina}.\par Em vista d'isso, não é para admirar que o \emph{P}. \emph{de Regnon}, excellente philosopho e theologo, conclua a exposição da controversia com estas pa\pgn{486}lavras: « E' um mysterio, um mysterio insondavel (o da sciencia divina dos \emph{futuriveis}); nenhuma das varias explicações apresentadas satisfaz completamente. Devemos renunciar á explicação do \emph{modo} por que \textsc{Deus} vê os futuriveis. Explicar esta sciencia é obra de \emph{dilettantismo} philosophico! ». (\emph{Bannez} et \emph{Molina} p. 113-115).\par Expuzemos, em resumo, a controversia ácerca da sciencia dos \emph{futuriveis} com a palavras de um philosopho moderno (\emph{Farges}, l'Idée de Dieu, pag. 370-377). Não temos competencia para emittir um juizo definitivo numa questão, que tem cançado as mais profundas intelligencias. Só dizemos, com o citado escriptor, que, devendo escolher entre o mysterio da \emph{sciencia media} e o mysterio dos \emph{decretos divinos}, preferimos este áquelle. A efficacia da Causa primeira, muito diversa da efficacia das causas segundas, pode mover a nossa natureza, sem fazer-lhe violencia, e produzir tudo em nós, mesmo a nossa liberdade, que não pode; sem sacrilegio, subtrahir-se á universal causalidade de \textsc{Deus}, segundo a expressão de S. Agostinho. Este padre da Egreja, criticando \emph{Cicero} e outros pagãos, que não admittiam em \textsc{Deus} o conhecimento dos \emph{futuriveis}, diz que a theoria d'elles « \emph{dum vult homines facere liberos, facit sacrilegos} » (\emph{De civ}. \emph{Dei}, l. 5, c. 9).\par e) \textsc{Deus} \emph{conhece todos os possiveis}. --- Tratamos dos \emph{meramente possiveis}, isto é, dos entes, que, embora não lhes repugne a existencia, comtudo não a têm e não a terão jamais. --- Estes \emph{possiveis}, como dissemos na \emph{Ontologia}, têm o seu fundamento na essencia divina. --- Que \textsc{Deus} os conheça é fóra de duvida. A intelligencia divina, sendo infinita, conhece tudo o que é cognoscivel. Ora todos os \emph{possiveis} são cognosciveis nos seus elementos \emph{communs} e \emph{individuaes}. Logo \textsc{Deus} conhece-os. --- Além d'isso, \textsc{Deus} conhece e comprehende a sua essencia e potencia infinita. Comprehendendo a sua essencia infinita, \textsc{Deus} deve conhecer todos os modos por que ella pode ser imitada e representada, --- e, comprehendendo a sua potencia infinita, deve conhecer todos os effeitos, que ella pode produzir. --- Finalmente, \textsc{Deus} pode produzir tudo o que é \emph{possivel}, e produz tudo pela sua intelligencia e vontade, isto é, emquanto conhece e quer as coisas. Ora quem por tal forma produz ou pode produzir as coisas, não pode deixar de as conhecer. Logo \textsc{Deus} conhece todos os possiveis.\par E, como a multidão dos possiveis é infinita, podemos dizer que \textsc{Deus} conhece objectos \emph{infinitos}. Diz S. Thomaz: « \textsc{Deus} conhece perfeitamente o seu poder. Ora o poder não pode ser conhecido perfeitamente, se não forem conhecidas todas as coisas, ás quaes pode extender-se; porque é nellas que deve fundar-se o criterio e a quantidade da sua efficacia. Ora o poder de \textsc{Deus}, sendo infinito, extende-se a objectos infinitos. Logo \textsc{Deus} conhece objectos infinitos. (\emph{C}. \emph{Gentes}, l. I, c. 69). --- Todavia, como adverte o S. Doutor (l. c.), \textsc{Deus} não conhece objectos infinitos pela sciencia \emph{de visão}, porque objectos \emph{infinitos} não existem \emph{actualmente}, nem existiram, nem existirão; mas conhece objectos infinitos pela sciencia de \emph{simples intelligencia}, porque \textsc{Deus} conhece infinitas coisas, que não são, nem serão, nem foram, e que todavia existem na potencia da creatura, como tambem conhece infinitas coisas, que existem na sua potencia (que Elle pode produzir), mas que não são, nem serão, nem foram. --- Este conhecimento não é confuso, nem\pgn{487}indeterminado; mas é infinitamente distincto e proprio, porque a essencia divina, que faz o officio da especie intelligivel, é uma \emph{semelhança sufficiente} de todas as coisas, que existem ou podem existir, não só emquanto aos elementos \emph{communs}, mas tambem emquanto aos elementos \emph{proprios} de cada ente. (\emph{Sum}. \emph{Th}., p. I, q. 14, a. 13).\par Objectam: Se \textsc{Deus} conhecesse infinitas coisas, conheceria um numero infinito. Ora um numero infinito não pode ser conhecido, porque é uma contradicção, um nada. Logo \textsc{Deus} não conhece infinitas coisas. --- Respondemos, \emph{negando} a proposição \emph{maior}. A \emph{multidão}, absolutamente considerada, não pode nem deve confundir-se com o \emph{numero}. O \emph{numero} é, sem duvida, uma multidão, mas é uma multidão, medida pela \emph{unidade}, (porque é formado pela addição de unidades), e por isso não pode ser infinito. Mas \textsc{Deus}, conhecendo objectos infinitos, --- não os conhece como se conhece um numero, addicionando partes a partes e unidades a unidades, --- mas conhece-os todos e ao mesmo tempo, e não os conhece em si mesmos mas sua essencia divina, emquanto vê que esta pode ser imitada e participada por uma illimitada multidão de creaturas. Diz S. Thomaz: « \textsc{Deus} non sic cognoscit infinitum vel infinita, quasi enumerando partem, cum cognoscat omnia simul, non successive » (\emph{Sum}. \emph{Th}. p. I, q. 14, a. 12 ad 1; Cf. c. \emph{Gentes}, l. I. c. 69). E' falso, pois, que \textsc{Deus}, se conhecesse infinitas coisas, conheceria um numero infinito.}.

\pg{478} 65. \textbf{Deus} tem um conhecimento proprio de todas as coisas.

a) Um ser intelligente tem um conhecimento \emph{proprio} de uma coisa, quando a conhece --- não só nos seus elementos \pg{479}  \emph{communs} e \emph{universaes}, pelos quaes convem com as outras coisas do mesmo genero ou da mesma especie, --- mas tambem nos \pg{480}  seus elementos \emph{proprios} e \emph{individuaes}, pelos quaes se distingue de todas as outras coisas do mesmo genero e da mesma \pg{481}  especie. Ora \textsc{Deus} não pode deixar de conhecer todas as coisas tambem nos seus elementos \emph{proprios} e \emph{individuaes}; porque, se as conhecesse somente nos elementos \emph{communs} e \emph{universaes}, \pg{482}  o seu conhecimento seria imperfeito: o que se não pode admittir em \textsc{Deus}, cujo entender é uma e a mesma coisa com a sua perfeitissima essencia. Logo \textsc{Deus} tem um conhecimento \emph{proprio} de todas as coisas.

b) \textsc{Deus} é causa de todas as coisas pela sua intelligencia, e per isso deve conhecel-as como são em si mesmas. Ora todas as coisas são em si mesmas o que são --- não só pela natureza, que lhes é commum com outras da mesma especie, --- \pg{483} \pg{484}  mas tambem pelas ultimas differenças e propriedades individuaes, que as distinguem de todas as outras, embora pertencentes á mesma especie. Logo \textsc{Deus} tem um conhecimento \emph{proprio} de todas as coisas.

c) \textsc{Deus} comprehende perfeitamente o proprio ser, e por isso todas as coisas, que estão contidas na sua potencia \emph{activa}, \pg{485} \pg{486}  isto é, a que a sua potencia \emph{activa} pode extender-se. Ora todas as coisas estão contidas na potencia \emph{activa} de \textsc{Deus} emquanto \pg{487} quanto ás formas \emph{proprias} e ás differenças \emph{ultimas}, porque \textsc{Deus} é o principio de todo o ente. Logo \textsc{Deus} tem um conhecimento \emph{proprio} de todas as coisas \footnote{Para evitar equivocos, advertimos que o conhecimento se pode chamar \emph{proprio} em dois sentidos, muito diversos, --- ou emquanto a especie, que informa a intelligencia, representa uma coisa com aquellas determinadas perfeições, e não outra, --- ou emquanto a especie é tão perfeita, que representa todas e cada uma das coisas com as perfeições, que lhe são proprias. No primeiro caso, a especie é imperfeitissima; no segundo, é infinitamente perfeita. Neste segundo sentido dizemos que \textsc{Deus} conhece todas as coisas com um conhecimento \emph{proprio}, emquanto as conhece pela sua Essencia infinita, representativa de todas as coisas com as propriedades e perfeições, proprias de cada uma. --- Do mesmo modo, a especie \emph{propria} --- pode representar o \emph{universal}, como são todas as especies abstractas das coisas materiaes, que informam a nossa intelligencia no estado de união e que lhe apresentam o objecto \emph{proprio}, e, então, o \emph{proprio} não se oppõe ao \emph{commum}, mas ao \emph{improprio}, ao \emph{indirecto}, --- ou pode representar o \emph{singular}, o \emph{individuo}, e, então, o \emph{proprio} oppõe-se ao \emph{commum}. Ora a essencia divina não se diz especie \emph{propria} no \emph{primeiro} sentido, mas no \emph{segundo}, emquanto entendemos dizer que a essencia divina, pela sua infinita perfeição, representa as coisas como são em si mesmas, isto é, na sua singularidade.\par O 1.º argumento, adduzido no texto, é convincente. Conhecer uma coisa nos seus elementos \emph{communs}, p. e., emquanto é \emph{ente}, é conhecel-a \emph{imperfeitamente}, porque não se conhecem os elementos \emph{proprios}, isto é, as ultimas perfeições, que completam uma coisa e a distinguem das outras. Diz S. Thomaz: « Intelligere aliquid in \emph{communi} et non \emph{in speciali}, est \emph{imperfecte} aliquid cognoscere. Unde intellectus noster, dum de potentia in actum reducitur, pertingit prius ad cognitionem \emph{universalem} et \emph{confusam} de rebus,\pgn{488}quam ad \emph{propriam} rerum cognitionem, sicut de imperfecto ad perfectum procedens. Si igitur cognitio Dei de rebus aliis a se esset \emph{in universali} tantum, et non \emph{in speciali}, sequeretur quod \emph{ejus intelligere non esset omnibus modis perfectum}, et per consequens nec ejus \emph{esse} » (\emph{Sum}. \emph{Th.}, p. I, q. 14, a. 6; cf. \emph{c. Gentes}, l. I, c. 5).\par A mesma doutrina é exposta por S. Thomaz no magnifico trecho seguinte. Diz assim: « Toda a perfeição, que se encontra na creatura, deve preexistir, num grau infinitamente mais elevado, no proprio Creador. Ora perfeição das creaturas não é apenas o que lhes é \emph{commum}, isto é, o ser, mas é tambem o que lhes é \emph{proprio}. E por isso todas as coisas creadas preexistem em \textsc{Deus} não só emquanto aos elementos, pelos quaes convêm com outras, mas tambem emquanto aos elementos, pelos quaes se distinguem entre si. Contendo todas as perfeições, a essencia divina refere-se a todas as essencias das coisas, \emph{não como} o \emph{commum ao proprio}, \emph{mas como} um \emph{acto perfeito aos imperfeitos}. Ora é evidente que, por meio de um acto perfeito, podemos conhecer os actos imperfeitos não só com um conhecimento \emph{commum}, mas tambem com um conhecimento \emph{proprio}; assim quem conhece o \emph{homem} tem um conhecimento \emph{proprio} do \emph{animal}. Portanto, contendo eminentemente a Essencia divina todas as perfeições, de que são dotadas as creaturas, \textsc{DEUS} pode conhecer em si \emph{mesmo} todas as coisas com um conhecimento \emph{proprio}» (\emph{Sum}. \emph{Th.}, l. c.).\par Para a solução das varias difficuldades contra a verdade exposta e defendida no texto, fazemos as observações seguintes: 1.º) O conhecimento exige a união do objecto conhecido com o cognoscente por meio de uma especie ou forma intencional, que actue como principio formal e represente o proprio objecto. E, para que o conhecimento seja \emph{proprio}, basta que a especie represente o objecto tambem emquanto aos elementos \emph{proprios} e \emph{individuaes}, e não é necessario que ella seja por tal modo proporcionada ao objecto, que represente exclusivamente este, e não um outro. Ora a Essencia divina, representando todas as coisas emquanto aos elementos \emph{communs} e \emph{proprios}, é principio de um conhecimento \emph{proprio}, embora não seja limitada aos objectos, que representa. --- 2.º) Pelo facto de a Essencia divina ser um \emph{unico} meio de conhecimento, não se segue que \textsc{Deus} não possa ter um conhecimento \emph{proprio} e \emph{distincto} de \emph{muitos} objectos; porque nenhum d'estes objectos imita \emph{completamente} e reproduz \emph{todas} as perfeições, que ella encerra, mas cada um reproduz, e de um modo deficientissimo, uma ou outra perfeição. Por isso a unica Essencia divina, excedendo infinitamente todos os objectos representados, pode representar objectos infinitos, sem nunca se exgottar. Se houvesse uma imagem de \textsc{Deus}, que imitasse completamente e exprimisse tudo o que está em \textsc{Deus}, essa imagem não poderia deixar de ser \emph{unica}. (Cf. \emph{De} Ver., q. 2, a. 4 ad 2). --- 3.º) \textsc{Deus} conhece todas as coisas pela sua Essencia, que é causa \emph{efficiente} e \emph{extrinseca} de todas as coisas, e por isso mesmo o conhecimento, que \textsc{Deus} tem de todas as coisas, é eminentemente \emph{proprio} e \emph{distincto}, visto ser a Essencia divina o principio de todo o \emph{ser}, de que as coisas são constituidas.}.

\pg{488} 66. A \textbf{sciencia divina} é \textbf{causa das coisas creadas}. --- A sciencia de \textsc{Deus} está para as coisas creadas, como a sciencia do artifice está para as coisas artificiaes. Ora a sciencia do artifice é causa das coisas artificiaes; porque o artifice opera pela sua

\pg{489} intelligencia. Logo tambem a sciencia de \textsc{Deus} é causa das coisas creadas. --- A sciencia divina é causa das coisas creadas, emquanto está unida, para o effeito, com o acto da vontade \footnote{A conclusão é certissima. A verdade e a certeza da sciencia das coisas, que a mesma sciencia não produz, mas suppõe já produzidas, depende das proprias coisas e é por ellas medida. Por isso, se a sciencia de \textsc{Deus} não fosse causa productiva ou efficiente das coisas, a sua verdade e certeza dependeria das coisas creadas e seria medida por ellas. Ora isto é falso. Deus não conhece as coisas, porque são: mas as coisas são, porque \textsc{Deus} as conhece. (\emph{Sum}. \emph{Th}., p. I, q. 14, a. 8). --- D'ahi a differença entre a nossa sciencia e a de \textsc{Deus}. As coisas naturaes são anteriores á nossa sciencia, e são a medida d'esta; porque a nossa sciencia é verdadeira, emquanto se adapta e corresponde aos objectos percebidos; --- ao passo que a sciencia divina é anterior ás coisas naturaes, e é a medida d'estas; porque as coisas são verdadeiras, emquanto se adaptam e correspondem aos archetypos divinos. Por isso as coisas creadas são intermedias entre a sciencia divina e a humana (\emph{Sum}. \emph{Th}., l. c. ad 3). --- Mas a sciencia divina é causa das coisas, não emquanto é um simples conhecimento das coisas, mas emquanto está unida com o acto da vontade, isto é, \textsc{Deus} produz as coisas emquanto as conhece e as quer. Com effeito, se a sciencia, só por si, fosse a causa efficiente e sufficiente das coisas, tudo o que \textsc{Deus} conhece existiria: o que é falso. Por isso a sciencia divina, que é causa das coisas e que está unida com a vontade, chama-se propriamente \emph{sciencia de approvação.} D'ahi se segue: 1.º) que \textsc{Deus} não é causa do \emph{peccado;} porque \textsc{Deus} não \emph{quer} e não \emph{approva} a deformidade do peccado, mas só \emph{permitte;} --- 2.º) que, embora a sciencia divina seja eterna, comtudo as creaturas não são eternas; porque, repetimos, as coisas derivam da sciencia, por meio da vontade, e por isso não é necessario que ellas existam desde que são conhecidas, mas devem existir quando o tiver determinado a vontade. (Cf. \emph{De Ver.}, q. 2, a. 14 ad 2).\par Extendendo-se a sciencia divina a todas as coisas, existentes e possiveis, e sendo ella productiva dos seres creados, é evidente que na intelligencia divina ha muitas \emph{idéas,} quer a \emph{idéa} se tome em ordem ao \emph{conhecimento} e signifique a semelhança intencional dos objectos percebidos, quer se tome em ordem á \emph{operação} e denote a forma exemplar ou o modelo da coisa, que se deve fazer. Com effeito, \textsc{Deus} é dotado de sciencia perfeitissima, identificada com a sua essencia, e é causa de todas as coisas pela sua intelligencia. Ora seria impossivel tanto a sciencia de coisas diversas, como a producção das mesmas pela intelligencia, se em \textsc{Deus} não existisse uma pluralidade de \emph{idéas,} quer a \emph{idéa} signifique a semelhança intencional do objecto, que foi percebido, quer signifique o modelo, segundo o qual se ha de fazer uma coisa. --- Que em \textsc{Deus} existam muitas \emph{idéas,} emquanto estas significam as semelhanças intencionaes dos objectos percebidos, entende-se facilmente. Que em \textsc{Deus} existam muitas \emph{idéas,} emquanto estas significam as formas exemplares, os modelos das coisas creadas, deduz-se da soberana independencia de \textsc{Deus} na operação, como no ser. \textsc{Deus} não é uma causa \emph{necessaria,} que opéra em virtude da sua forma natural, e que por isso só pode produzir um effeito determinado; mas é a causa infinitamente \emph{livre,} que\pgn{490}opera pela sua intelligencia, e que por isso deve ter e tem as formas exemplares á semelhança das quaes produz todas as coisas. --- Logo todas as coisas, feitas por \textsc{Deus}, têm um duplice ser: um em si mesmas, emquanto são realmente distinctas de \textsc{Deus}, mutaveis, compostas, finitas; outro na arte divina, isto é, nos divinos archetypos ou exemplares, e então não se distinguem realmente de \textsc{Deus}, e constituem a propria vida de \textsc{Deus}.\par Esta \emph{multiplicidade} de idéas não importa \emph{composição}, e por isso não é contraria á \emph{simplicidade} de \textsc{Deus}. Dá-se a \emph{composição}, opposta á \emph{simplicidade}, numa intelligencia, que percebe muitas e diversas coisas, quando o principio determinativo da intelligencia, que é a especie intelligivel, ou é \emph{realmente distincto} da intelligencia, ou é \emph{multiplice;} porque, em ambos os casos, ha composição de \emph{potencia} e \emph{acto}. Ora nada d'isso se verifica com relação á intelligencia divina. Porquanto, a essencia divina, que faz o officio de especie intelligivel, identifica-se com a intelligencia divina e é \emph{simplicissima}, embora possa ser imitada pelas creaturas por modos infinitos. Por isso, entendendo a intelligencia divina todas as coisas por uma unica especie, --- pela essencia divina, que é semelhança de si mesma e de todas as coisas, a multidão das coisas entendidas não importa composição na propria intelligencia. --- A \emph{variedade} e \emph{multiplicidade} de ideas não diz respeito á divina essencia, mas aos termos ou ás coisas, que a podem imitar. (Cf. \emph{Sum}. \emph{Th}., p. I, q. 15; \emph{c. Gentes}, l. I, c. 51 e seg.; \emph{De Ver.}, q. 3, a. 2).\par Para a solução das difficuldades, tenhamos presentes os seguintes principios: 1.º) Ha uma grande differença entre a semelhança, que as creaturas têm com a natureza divina, e a semelhança, que as mesmas creaturas têm com os divinos archetypos. A primeira semelhança é imperfeitissima, porque as creaturas imitam a natureza divina por um modo infinitamente deficiente e remoto. A segunda é perfeitissima, porque as creaturas representam adequadamente, na sua entidade, os divinos archetypos, em conformidade com os quaes foram produzidas. (Cf. \emph{Sum}. \emph{Th}., p. I, q. 44, a. 3 ad 1). --- 2.º) Se o conhecimento divino se refere \emph{ao cognoscente}, \textsc{Deus} conhece o proprio ser e as outras coisa pelo \emph{mesmo modo}, isto é, pela sua \emph{essencia}. Se o conhecimento se refere \emph{ás coisas conhecidas}, \textsc{Deus} conhece o proprio ser e as outras coisas por um modo \emph{differente;} porque conhece o proprio ser pela sua \emph{essencia}, emquanto esta o representa por \emph{identidade}, e conhece as outras coisas pela sua \emph{essencia}, emquanto esta as representa por \emph{semelhança}. Por isso é necessario admittir em \textsc{Deus} muitas e diversas idéas, visto ser muitos e diversos os modos, por que pode ser imitada pelas creaturas a unica e simplicissima essencia divina. --- 3.º) As coisas imitam a essencia divina pelo facto de \textsc{Deus} a conhecer como imitavel; e não vice-versa, isto é, não é pelo facto de as coisas imitarem a essencia divina que \textsc{Deus} a conhece como imitavel. Por isso, na divina intelligencia, que conhece a sua essencia emquanto imitavel, subsistem todos os archetypos, á semelhança dos quaes são feitas todas as coisas. --- 4.º) As idéas, ou os archetypos, emquanto subsistem na divina intelligencia, são uma e a mesma coisa com ella; mas, emquanto se consideram fóra de \textsc{Deus}, são nada; porque, fóra de \textsc{Deus}, só existem entes \emph{reaes}.\par Do que deixamos exposto vê-se quão absurdas são as theorias dos \emph{pantheistas} e de alguns \emph{theistas} ácerca da \emph{sciencia divina}. --- \emph{Spinoza}, ensinando\pgn{491}que \textsc{Deus} é \emph{impessoal}, conclue que Elle não é intelligente senão no homem. --- \emph{Fichte}, \emph{Schelling}, \emph{Hegel}, e os sequazes do \emph{progresso indefinido}, affirmam que \textsc{Deus} conhece todas as coisas não \emph{simultanea} mas \emph{successivamente}, em virtude das suas evoluções. \emph{Renan} diz que \textsc{Deus} não chegou ainda a conhecer, de um modo perfeito, o mundo, mas que, estando Elle \emph{in fieri}, talvez chegue a possuir uma sciencia infinita. --- \emph{Cicero}, os \emph{Sociniamos}, e, nos tempos modernos, \emph{Ramiron}, disseram que \textsc{Deus} não tem conhecimento, pelo menos \emph{certo}, dos \emph{futuros livres;} pois, na sua opinião, a presciencia divina não pode conciliar-se com a liberdade humana.}.

\pg{490} 67. \textbf{Propriedades} da \textbf{sciencia divina}. --- As principaes são \emph{quatro}: \emph{é substancial}, \emph{infinita}, \emph{intuitiva}, \emph{immutavel}.

\pg{491} a) E' \emph{substancial} emquanto ao \emph{sujeito}. A sciencia, em \textsc{Deus}, não é uma coisa accidental, como é em nós, mas é uma e a mesma coisa com a sua substancia. --- Identificando-se com a substancia divina, a sciencia divina é \emph{eterna} \footnote{Se a sciencia de \textsc{Deus} não se identificasse com a sua substancia, a Essencia de \textsc{Deus} referir-se-hia ao conhecimento como a potencia ao acto, e \textsc{Deus} receberia de uma coisa distincta d'Elle a perfeição. Ora isto é absurdo, porque \textsc{Deus} é acto purissimo e é perfeição de si mesmo. Por isso diz S. Thomaz que « in Deo intellectus, intelligens, et id quod intelligitur, et species intelligibilis, et ipsum \emph{intelligere}, sunt omnino unum et idem » (\emph{Sum}. \emph{Th}., p. I, q. 14, a. 4, 15).}.

b) E' \emph{infinita} emquanto ao \emph{objecto}. Com effeito, \textsc{Deus} conhece tudo o que é, de algum modo, cognoscivel; e por isso a sua sciencia extende-se a todas as coisas, --- existentes e possiveis, --- passadas, presentes e futuras, --- necessarias e livres, --- absolutas e condicionadas \footnote{« A extensão da sciencia ou do conhecimento de um cognoscente é proporcionado á forma, que é o principio do proprio conhecimento. Assim, a especie \emph{sensivel}, que está nos sentidos, é a semelhança de um só individuo, e por isso não faz conhecer senão aquelle individuo. A especie \emph{intelligivel} da nossa intelligencia é a semelhança do objecto emquanto á \emph{natureza} da \emph{especie}, que pode ser participada por infinitos individuos; e por isso a nossa intelligencia, por meio da especie intelligivel de homem, conhece, de algum modo, uma multidão infinita de homens, --- não emquanto aos elementos \emph{individuaes}, mas emquanto aos elementos \emph{especificos}. A essencia divina, pela qual a intelligencia divina conhece, representa uma multidão infinita de coisas, existentes e possiveis; e por isso a sciencia de \textsc{Deus} é infinita ». (\emph{Sum}. \emph{Th}., p. I, q. 14, a. 12).}.

c) E' \emph{intuitiva} emquanto ao \emph{meio}. \textsc{Deus}, conhecendo a sua essencia, conhece nella todas as coisas. Ora \textsc{Deus} conhece a sua essencia com um unico acto \emph{intuitivo}, de infinita energia e extensão. Logo a sciencia divina é \emph{intuitiva} \footnote{No conhecimento \emph{discursivo}, que é opposto ao \emph{intuitivo} e que é sempre acompanhado por um \emph{movimento}, a intelligencia passa da \emph{potencial} para a \emph{actual} percepção da conclusão, e esta não a conhece pelo mesmo acto, por que conhece as premissas, mas por um novo acto, que succede ao primeiro. Por isso o conhecimento \emph{discursivo} suppõe sempre uma composição e\pgn{492}distincção de potencia e de acto, e um conjuncto de sciencia e de ignorancia. Ora, como nada d'isso pode admittir-se em \textsc{Deus}, concluimos que a sciencia de \textsc{Deus} não é \emph{discursiva}, mas \emph{intuitiva}. (\emph{Sum}. \emph{Th}., p. I, q. 14, a. 7).}.

\pg{492} d) E' \emph{immutavel} emquanto á \emph{natureza}. --- A sciencia divina identifica-se, como dissemos, com a substancia divina, e por isso é dotata dos mesmos caracteres. Ora a substancia divina é \emph{immutavel}. Logo tambem a sciencia divina é \emph{immutavel} \footnote{A sciencia divina refere-se ás creaturas, não emquanto estas existem em si mesmas, mas emquanto existem em \textsc{Deus}. E, como em \textsc{Deus} existem por um modo \emph{invariavel}, assim a sciencia, que \textsc{Deus} tem das creaturas, é invariavel. --- \textsc{Deus} conhece a \emph{variabilidade} ou \emph{mutabilidade} das coisas; mas a sua sciencia é invariavel ou immutavel. --- A sciencia de \textsc{Deus} seria variavel só no caso de Elle chegar a conhecer uma coisa, que antes ignorava. Ora isto é impossivel. \textsc{Deus}, na sua eternidade, vê tudo o que existe no tempo e conhece tudo o que é cognoscivel. (\emph{Sum}. \emph{Th}., p. I, q. 14, a. 15).}.

68. \textbf{Termo immanente da sciencia diviva.} --- O acto da intelligencia divina, pelo qual \textsc{Deus} conhece o proprio ser e todas as coisas, existentes e possiveis, deve ter o seu termo. E, como o acto de entender é formalmente \emph{immanente}, o termo tambem deve ser \emph{immanente}, isto é, deve permanecer no principio, de que deriva. Este termo \emph{immanente}, que é perfeitissima expressão de tudo o que \textsc{Deus} conhece e por isso é imagem semelhantissima do proprio ser divino e de todas as coisas, é e diz-se Verbo. --- O Verbo, em \textsc{Deus}, é unico, infinito, substancial; como unico, infinito, substancial é o acto da intelligencia divina, de que procede \footnote{O \emph{verbo mental} encontra-se em toda a natureza intellectual, e por isso, como no homem, tambem em \textsc{Deus}. Já tivemos occasião de indicar as \emph{analogias} entre o verbo humano e o verbo divino. Não insistimos no assumpto. A razão humana chega a entender que, emquanto o verbo humano não é a propria substancia da alma, mas é uma coisa accidental, o verbo divino, ao contrario, não é uma coisa accidental, mas é o proprio \textsc{Deus}. Porquanto, em \textsc{Deus}, a intelligencia, o acto intellectual, o principio de que o verbo deriva, e o proprio verbo não se distinguem da essencia, mas são uma e a mesma coisa com a essencia. Tal é a conclusão da razão. (\emph{Cf}. \emph{Sum}. \emph{Th}., p. I, q. 27, 34; \emph{c}. \emph{Gent}., 1. IV, c. 16; \emph{De Verit}., q. 4; \emph{De dif}. \emph{Verb}. \emph{div}. \emph{et hum}.).\par Mas a Fé descobre horisontes invisiveis á razão humana, e diz-nos que o Verbo, em \textsc{Deus}, é uma verdadeira \emph{Pessoa}, subsistente na divina natureza e distincta, pela relação da origem, do Principio, de que deriva e que é tambem \emph{Pessoa}, porque tambem Elle subsiste na mesma natureza divina. Na verdade, chama-se \emph{pessoa} quem subsiste numa natureza intellectual e distingue-se de todos os outros, que subsistem na mesma natureza, tanto quando a natureza é uma e a mesma \emph{especificamente} em todos, o que se verifica nas creaturas, como quando a natureza é uma e a mesma \emph{numericamente}, o que se verifica em \textsc{Deus}.\par O Verbo divino procede do seu Principio por verdadeira \emph{geração}, e por\pgn{493}isso é verdadeiro \emph{Filho} de \textsc{Deus}. A \emph{geração} é a origem de um ser vivo de outro ser vivo, que communica ao primeiro a propria natureza em virtude da sua intrinseca actividade e a propria semelhança em consequencia da indole da origem. Ora todos estes elementos encontram-se na processão do Verbo divino. Na verdade, o Verbo --- procede do seu Principio, como \textsc{Deus} vivo de \textsc{Deus} vivo, --- recebe do seu Principio não uma natureza da mesma especie, mas a propria identica natureza, --- procede em virtude da intrinseca actividade e fecundidade do seu Principio, porque a intelligencia divina possue sempre a energia productora do seu Verbo, --- procede como imagem do seu Principio, por isso mesmo que procede como Verbo; porque o verbo é essencialmente imagem do objecto entendido pela intelligencia, e é imagem da propria intelligencia, quando esta se entende a si mesma, e, como a intelligencia divina produz o seu Verbo entendendo-se e comprehendendo-se a si mesma, por isso o Verbo é essencialmente imagem da mesma intelligencia, de que procede. Logo o Verbo procede de \textsc{Deus} por verdadeira geração, e por geração infinitamente perfeita e superior a toda a geração das creaturas; e assim o Verbo é \emph{Filho} de \textsc{Deus}.\par O Verbo é \emph{Filho} de \textsc{Deus}, e é \emph{todo} Filho de \textsc{Deus}, porque a \textsc{Deus} se attribue tudo o que concorre para a geração do Verbo. Para, entendermos isto, consideremos o que se verifica na nossa intelligencia. Quando esta pensa num objecto, concebe o seu verbo, que é uma imagem ideal do proprio objecto. Este verbo mental é filho --- tanto da intelligencia, que o produz, --- como do objecto, que, pela sua especie, fecunda a intelligencia e a determina a produzir o verbo, semelhante a si mesmo. Por isso, quando a intelligencia pensa num objecto, distincto de si mesma, o objecto entendido é como o pae do verbo, concebido pela intelligencia, porque é principio \emph{activo}, emquanto actua na faculdade, e a intelligencia é como a mãe do verbo, porque é principio \emph{passivo}, emquanto concebe e traz no seu seio o verbo. Mas, quando a intelligencia pensa em si mesma, então ella é, ao mesmo tempo, pae e mãe do verbo, e a ella deve attribuir-se tudo o que concorre para a geração do mesmo verbo. E' o que se verifica em \textsc{Deus}. A intelligencia divina produz o seu Verbo pela comprehensão de si mesma, e por isso ella é o unico principio gerador do Verbo.\par O Verbo é todo Filho de \textsc{Deus}, e é Filho \emph{unigenito}. O Filho de \textsc{Deus} não pode deixar de ser \emph{unico}, tanto se se considera o \emph{acto}, que O gera, como se se considera o termo d'esse acto. --- Se se considera o \emph{acto}, este é \emph{unico}, simplicissimo, purissimo, infinito, e infinitamente intelligente, e por isso comprehende a essencia divina e nella tudo o que é cognoscivel, e por isso \textsc{Deus} exprime pelo Verbo tudo o que Elle comprehende. --- Se se considera o \emph{termo} desse acto, que é o \emph{Verbo}, este é infinitamente perfeito, e por isso é uma expressão adequada da intelligencia divina, de que procede,\pgn{494}Ora se o \emph{acto} do entender divino é \emph{unico}, e se o Verbo divino exprime \emph{tudo} o que \textsc{Deus} entende, o Verbo deve ser \emph{unico}. A repetição ou multiplicação do Verbo, em \textsc{Deus}, significaria impotencia, imperfeição. Poderá haver, na creação, innumeraveis expressões d'esse Verbo, mais ou menos inadequadas; mas estas, embora multiplicadas ao infinito, não exprimirão nunca tudo o que o unico Verbo de \textsc{Deus} exprime. --- Este Verbo, tão glorioso na sua Divindade e depois tão humilhado na sua Humanidade, é o nosso querido \textsc{Jesus}... Outras noções relativas a este assumpto encontram-se nos tratados de Theologia revelada. (\emph{Cf}. \emph{Sum}. \emph{Th}., p. I, q. 27; p. III, q. 23; \emph{c}. \emph{Gent}., 1. IV, c. 11, 13).}.

%%>>> art-42

\section*{Artigo II.\\ Faculdade e operação volitiva immanente.}
\addcontentsline{toc}{section}{Artigo II. --- Faculdade e operação volitiva immanente.}

\pg{493} 69. \textbf{Vontade de Deus}. --- \emph{Vontade} de \textsc{Deus} é a faculdade, pela qual Elle ama o bem, percebido pela sua intelligencia.

\pg{494} --- Acto da vontade é o \emph{amor}. --- Diremos da \emph{existencia} d'esta faculdade, da sua \emph{indole} e \emph{objecto} \footnote{\emph{S. Thomaz} trata da \emph{vontade} de \textsc{Deus} na --- \emph{Sum}. \emph{Th}., p. I, q. 19; --- \emph{c}. \emph{Gent}., 1. I, c. 62 e seg.; --- \emph{De Ver}., q. 25; --- I \emph{Sent}., Dist. 45.}.

70. \textbf{Deus é dotado de vontade.}

a) Toda a substancia espiritual, que é dotada de \emph{intelligencia}, é tambem dotada de \emph{vontade}; porque essa substancia, percebendo pela intelligencia o bem conveniente, não pode deixar de ter uma inclinação natural para esse bem, de modo que tenda para elle, quando o não possue, e descance nelle, quando o possue; e tal inclinação é a \emph{vontade}. Ora \textsc{Deus} é dotado de \emph{intelligencia}. Logo é tambem dotado de \emph{vontade}.

b) A forma, considerada pela intelligencia, não determina o agente ao acto externo, se não por meio da vontade, cujo objecto é o fim ou o bem; como pode verificar-se no artifice, que concebeu na intelligencia a forma da casa. Ora \textsc{Deus}, artifice sapientissimo, produz as suas operações \emph{transeuntes} na creação conforme a idéa ou forma exemplar, concebida na sua intelligencia. Logo deve admittir-se em \textsc{Deus} a vontade, que o move a operar. --- Por isso em \textsc{Deus}, que todas as creaturas produziu á imitação da sua essencia, a intelligencia \emph{dirige}, a vontade \emph{impera}, o poder \emph{executa}.

c) \textsc{Deus} é infinitamente perfeito. Ora não o seria, se não fosse dotado de \emph{vontade}; não só porque a vontade é uma perfeição \emph{absoluta}, mas tambem porque, se não possuisse a vontade, nem poderia possuir a \emph{santidade}, a \emph{justiça}, a \emph{bondade}, a \emph{felicidade}, e todos os attributos, que residem nessa faculdade. Logo \textsc{Deus} é dotado de \emph{vontade} \footnote{A vontade, em nós, tem um duplice acto --- \emph{appetecer} o bem não possuido, e \emph{amar} o bem possuido e \emph{gosar} d'elle. O primeiro acto indica \emph{imperfeição}, os outros não. Por isso quando o nome de \emph{vontade} se transfere das creaturas a \textsc{Deus}, não se applica a \textsc{Deus} senão emquanto significa perfeição, isto é, emquanto \textsc{Deus} ama o que possue, porque possue sempre o bem, que é seu objecto. D'onde se segue que a vontade divina impropria e falsamante se chama \emph{appetite}, se esta palavra se toma no sentido rigoroso,\pgn{495}emquanto denota a tendencia para um bem não possuido. --- O objecto da vontade de \textsc{Deus} é a sua propria bondade, que se identifica com a sua essencia. Por isso a vontade de \textsc{Deus} não é movida por um bem extranho, mas por Elle só, e diz-se \emph{movida}, emquanto se chama \emph{movimento} tambem o acto da intelligencia e da vontade.}.

\pg{495} \textbf{71. A vontade de Deus identifica-se com a sua essencia, e é acto purissimo, e infinitamente perfeita.}

a) Identifica-se com a sua essencia. --- Se a \emph{vontade} de \textsc{Deus} não se identificasse com a sua \emph{essencia}, Elle seria \emph{composto} de \emph{essencia} e de \emph{attributos}. Ora \textsc{Deus} não é composto de \emph{essencia} e de \emph{attributos}. Logo a \emph{vontade} de \textsc{Deus} identifica-se com a sua \emph{essencia}. --- Por isso \textsc{Deus} \emph{quer} pela sua essencia, como \emph{conhece} pela sua essencia.

b) E' acto purissimo. --- A \emph{vontade} de \textsc{Deus} identifica-se com a sua essencia. Ora a essencia de \textsc{Deus} é \emph{acto purissimo}, excluindo toda a sombra de potencia e de imperfeição. Logo a \emph{vontade} de \textsc{Deus} é \emph{acto purissimo}.

c) E' infinitamente perfeita. --- Uma vontade, que ama uma bondade infinita tanto, quanto esta é amavel, é infinitamente perfeita. Ora a vontade divina, por ser adequada á divina bondade, ama esta bondade tanto, quanto é amavel, isto é, com um amor simplesmente infinito. Logo a divina vontade é \emph{infinitamente perfeita} \footnote{Todas as faculdades de um ente, como tambem os seus actos, seguem a natureza do mesmo ente. --- Em \textsc{Deus}, ente \emph{simplicissimo}, a vontade, o acto volitivo, a substancia, são uma e a mesma coisa na \emph{realidade}, havendo apenas entre ellas uma distincção de \emph{razão}. No homem, como nas outras creaturas intellectuaes, a vontade, o acto volitivo, e substancia, são coisas \emph{realmente} distinctas. --- Além d'isso, a vontade de \textsc{Deus} está \emph{sempre em acto}, isto é, \emph{quer sempre}, e permanece sempre immovel e invariavel, embora queira coisas diversas, e livremente se refira aos effeitos, que imitam a sua bondade. A creatura ora quer, ora não quer, passa da potencia para o acto, e está sujeita a mudanças. --- O acto da vontade divina produz a bondade e a communica ás coisas, que ella quer. O acto da vontade creada é produzido e medido pelo bem, que se encontra nas coisas.\par *\par A vontade divina, embora una e simplicissima, divide-se --- emquanto ao \emph{modo}, por que se manifesta, em vontade de \emph{beneplacito} e de \emph{signal}, --- emquanto ao \emph{effeito}, que produz, em \emph{antecedente} e \emph{consequente}, --- emquanto ao \emph{objecto}, que attinge, em \emph{necessaria} e \emph{livre}.\par \emph{a}) Vontade de \emph{beneplacito} e de \emph{signal}. A primeira é a vontade \emph{propriamente} dita; a segunda é a vontade tomada em sentido \emph{metaphorico}, e consiste num \emph{signal} externo, com que \textsc{Deus} manifesta a sua vontade, assim o \emph{castigo}, que \textsc{Deus} inflige ao peccador, é um \emph{signal} de \emph{ira} na sua vontade.\par \emph{b}) Vontade \emph{antecedente} e \emph{consequente}. Pela primeira, \textsc{Deus}, movido pela\pgn{496}sua benevolencia, quer um effeito, independentemente de toda e qualquer condição; assim \textsc{Deus} quer a salvação da todos os homens. Pela segunda, \textsc{Deus}, suppostas certas circumstancias, quer um effeito ainda que extranho á sua bondade; assim \textsc{Deus} pune o homem, se elle morrer em peccado.\par \emph{c}) Vontade \emph{necessaria} e \emph{livre}. E' \emph{necessaria}, quando não pode deixar de produzir o seu acto; assim \textsc{Deus} ama a sua bondade necessariamente. E' \emph{livre}, quando pode deixar de produzir o seu acto; assim \textsc{Deus} ama as coisas creadas livremente. --- Na \emph{supposição} que \textsc{Deus} queira uma coisa, não pode deixar de a querer; porque a vontade de \textsc{Deus} é immutavel. Mas, então, é uma necessidade \emph{hypothetica}, não \emph{absoluta}.}.

\pg{496} \textbf{72. Deus quer ou ama tudo pela sua bondade.}

a) Uma volição perfeitissima, como é toda a volição divina, deve ter um perfeitissimo \emph{objecto formal}, ou \emph{motivo}, e um perfeitissimo modo de querer. Ora o perfeitissimo \emph{objecto formal}, ou \emph{motivo}, da volição divina não pode deixar de ser a \emph{Bondade divina}, e o perfeitissimo modo de querer é querer ou amar essa Bondade por si mesma e sobre todas as coisas e em todas as coisas, e querer ou amar as outras coisas por ella e nella. Logo \textsc{Deus} quer ou ama tudo pela sua \emph{bondade}.

b) O objecto \emph{formal} dá a cada volição a perfeição especifica, ou quasi-especifica; porque é pelos objectos \emph{formaes} que se especificam os actos. Ora é impossivel que um acto de \textsc{Deus} seja especificado por uma coisa, distincta da sua essencia, e por isso da sua bondade, que é uma e a mesma coisa com a essencia. Logo é impossivel que o objecto \emph{formal} da volição de \textsc{Deus} seja uma coisa distincta da sua bondade. Logo \textsc{Deus} quer ou ama tudo pela sua \emph{bondade}.

c) O objecto \emph{formal}, como toda a causa \emph{final}, refere-se á vontade, como o motor ao movido. Logo se o objecto \emph{formal}, ou o \emph{motivo}, da volição de \textsc{Deus} fosse uma coisa distincta da sua bondade, \textsc{Deus} seria movido e determinado ao acto por um outro motor. Ora repugna que \textsc{Deus}, primeiro motor e infinitamente independente, receba de uma coisa extranha o motivo e o impulso da sua volição. Logo \textsc{Deus} quer ou ama tudo pela sua \emph{bondade} \footnote{Em todas as faculdades devemos distinguir o objecto \emph{formal} e o objecto \emph{material}. O objecto \emph{material} é o que a faculdade attinge; o \emph{formal} é o \emph{motivo}, ou \emph{aspecto}, pelo qual, ou sob o qual, a faculdade attinge o objecto \emph{material}. Dá-se o mesmo com relação á vontade. O objecto \emph{formal}, ou o \emph{motivo}, exerce uma verdadeira \emph{causalidade}, actua na vontade, e a leva a praticar o acto. --- Quando se trata da vontade de \textsc{Deus}, devemos tambem admittir um objecto \emph{formal}, um \emph{motivo}, uma \emph{razão sufficiente}, pela qual essa faculdade produz o proprio acto; aliás o acto não seria racional. Mas esse motivo, essa \emph{razão sufficiente}, não exerce na vontade divina uma\pgn{497}verdadeira \emph{causalidade}, porque a idéa de \emph{causalidade} desperta a idéa de um \emph{effeito}, de um \emph{producto}, que começa a existir, quando antes não existia. --- Ora o motivo, \emph{sufficiente} e \emph{supremo}, de todas as volições divinas é a \emph{bondade divina}. \textsc{Deus}, assim como vê tudo com um só acto na sua essencia e pela sua essencia, assim tambem com um só acto quer tudo na sua bondade e pela sua bondade. Só este motivo é adequado á perfeição infinita da sua vontade. (\emph{Cf}. \emph{Sum}. \emph{Th}., p. I, q. 19, a. 5). --- \textsc{Deus}, nas suas obras, pode ter em vista qualquer outro motivo, como seria o bem, a felicidade do homem; mas o motivo \emph{supremo}, que se encerra nos outros motivos e a todos confere a \emph{appetibilidade}, é sempre a infinita bondade de \textsc{Deus}, que deseja communicar-se a outros entes.\par Advirta-se que \textsc{Deus}, embora tudo queira e opere pela sua bondade e gloria, comtudo ama as creaturas com perfeitissimo amor de \emph{amizade} e de \emph{benevolencia} (que procura o bem da coisa amada); porque a maxima gloria extrinseca, que \textsc{Deus} pode receber das creaturas, está inseparavel e essencialmente unida ao maximo bem, que \textsc{Deus} pode dar á creatura e que consiste na visão da essencia divina e no correspondente amor beatifico.}.

\pg{497} 73. \textbf{Deus} quer ou ama o proprio ser como fim.

a) Um ente quer ou ama um objecto como \emph{fim}, quando esse objecto é amavel \emph{por si mesmo}, e communica a sua amabilidade a todos os outros objectos, que se tornam, por isso, termos da vontade. Ora \textsc{Deus}, sendo a propria bondade infinita, é infinitamente amavel \emph{por si mesmo}, e todas as outras coisas não são amaveis, senão emquanto participam da bondade divina. Logo \textsc{Deus} quer ou ama o proprio ser como \emph{fim}.

b) Um objecto ama-se como \emph{fim}, quando é proporcionado e adequado á natureza e excellencia da faculdade appetitiva; porque essa proporção e adequação termina e satisfaz por completo a tendencia da faculdade. Ora o ser divino, por ser uma bondade infinita, é o \emph{unico} objecto proporcionado e adequado á natureza e excellencia da vontade divina; porque esta é uma e a mesma coisa com aquelle. Logo \textsc{Deus} quer ou ama o proprio ser como \emph{fim}\footnote{Esta e as seguintes conclusões são simples corollarios da conclusão precedente --- que \textsc{Deus} quer ou ama tudo pela sua bondade.}.

74. \textbf{Deus} quer ou ama as outras coisas como meios.

a) Deus quer ou ama as outras coisas. --- \textsc{Deus}, amando o proprio ser \emph{por si mesmo} e como \emph{fim}, porque esse ser é a bondade infinita, não pode deixar de amar tudo o que é participação ou semelhança d'essa bondade; porque não se ama a si mesmo quem não ama a propria semelhança. Ora as outras coisas são participações ou semelhanças da bondade divina. Logo \textsc{Deus} quer ou ama as outras coisas.

b) Deus quer ou ama as outras coisas como meios. --- \textsc{Deus}, pelo \pg{498}  facto de amar, como \emph{fim}, o proprio ser, que é bondade \emph{infinita}, não póde deixar de amar tudo o que, de algum modo, concorre para a perfeição da mesma bondade. Ora as outras coisas, embora, por serem finitas, não possam concorrer para o augmento da bondade infinita, comtudo, por serem semelhanças d'essa bondade, concorrem para a sua manifestação e glorificação; e por isso têm razão de \emph{meios}. Logo \textsc{Deus} quer ou ama as outras coisas como \emph{meios} \footnote{Se \textsc{Deus} quer e ama o proprio ser como \emph{fim} e as outras coisas como \emph{meios}, é claro que o ser divino é o objecto \emph{primario} da vontade divina, e as outras coisas constituem o objecto \emph{secundario}. De facto, objecto \emph{primario} de uma faculdade é o que é attingido \emph{por si}; o objecto \emph{secundario} é o que não é attingido \emph{por si mesmo}, mas por meio de \emph{outro} objecto. Ora a vontade divina ama o ser divino \emph{por si mesmo}, porque é bondade infinita, e não ama as outras coisas senão porque se ama a si mesmo. --- D'onde se vê que a tendencia para o bem e a propria posse do bem não excluem o amor de outras coisas, quando estas estão para esse bem, como meios para o fim.\par Como o bem se encontra propriamente no ente e a verdade na intelligencia, o objecto \emph{primario} da \emph{intelligencia} e o objecto \emph{primario} da vontade --- nalguma coisa convém entre si, --- e noutra differem. Convém em que, tanto para uma faculdade como para outra, objecto \emph{primario} é o que é attingido \emph{por si}, e não por outro ou noutro. --- Differem em que, emquanto o objecto \emph{primario} da intelligencia é constituido por uma coisa, diversa segundo a diversidade da intelligencia, o objecto \emph{primario} da vontade não é constituido por uma coisa, diversa segundo a diversidade da vontade. Com effeito, o objecto \emph{primario} da intelligencia divina é o \emph{ser divino}, e o objecto \emph{primario} da intelligencia humana, no estado da união com o corpo, é a \emph{essencia abstracta das coisas materiaes}; emquanto o objecto \emph{primario} tanto da vontade divina como da humana é a \emph{felicidade}, que consiste na posse de \textsc{Deus}, ainda que a nossa vontade tenda para \textsc{Deus} emquanto tende para a felicidade e não emquanto Deus é o verdadeiro objecto da felicidade, porque tendemos para o bem como o concebemos.\par Embora \textsc{Deus} e a creatura racional queiram alguns bens, como meios para o fim, comtudo a tendencia da vontade creada para esses objectos \emph{secundarios} differe da tendencia da vontade divina para os mesmos objectos. Com effeito, a vontade creada quer os meios, para que a levem á consecução da felicidade, que não possue mas que deseja possuir; e por isso os taes meios tem a indole ou o caracter de \emph{bens uteis}. \textsc{Deus}, porém, não quer os meios, isto é, os objectos \emph{secundarios}, para que o levem ao alcance da felicidade, porque a felicidade de \textsc{Deus} é a sua essencia, que ama e em que se deleita, mas quer esses objectos para que lhes possa communicar a propria bondade, e d'este modo manifeste exteriormente as riquezas d'essa bondade. Portanto, emquanto a creatura se serve das coisas, que são ordenadas ao fim, para vantagem propria, \textsc{Deus} serve-se de nós para a sua bondade e para a vantagem nossa. Escutemos S. Thomaz: « Sicut Deus cognoscendo essentiam suam, cognoscit omnia, quae sunt ab eo, in quantum sunt similitudo quaedam veritatis ejus; ita etiam volendo vel amando essentiam suam, vult omnia, quae sunt ab eo, in quantum habent similitudinem bonitatis eius.\pgn{499}Unde id quod est volitum \emph{primo} ab eo, est bonitas sua tantum. Alia vero vult in ordine ad bonitatem suam; non autem hoc modo ut per ea aliquid bonitatis acquirat, sicut nos facimus circa alios bene operando, sed ita quod eis de bonitate sua aliquid largiatur; et ideo liberalitas est quasi proprium ipsius, secundum \emph{Avicennam}; quia ex operatione sua non intendit aliquod sibi commodum provenire, sed vult bonitatem suam in alios diffundere; et ideo \emph{Augustinus} dicit quod ipse utitur nobis ad bonitatem suam et utilitatem nostram » (I \emph{Sent}. \emph{Dist}. 45, q. 1, a. 2).\par Contra a ultima proposição --- que \textsc{Deus} quer ou ama outras coisas --- levantem-se muitas difficuldades. Apresentamos as principaes.\par Obj. 1.ª --- Quem possue o bem, que é objecto \emph{primario} da vontade, não pode desejar outra coisa; porque esse bem constitue o fim ultimo, a felicidade. Ora \textsc{Deus} possue a sua bondade infinita, que é objecto \emph{primario} da sua vontade. Logo não pode amar outras coisas.\par Resposta. O bem, que constitue o objecto da felicidade, não pode ser senão o objecto \emph{primario} da vontade, porque este e só este tem a indole, o caracter de \emph{fim ultimo}. E, como o fim ultimo é um só, é evidente que a vontade, a qual possue aquelle bem, não pode querer outro objecto \emph{primario}. Todavia isso não obsta a que a mesma vontade possa querer outros objectos ou bens \emph{secundarios} por amor do bem \emph{primario}, para o qual estão ordenados como meios para o fim; e, neste caso, os bens \emph{secundarios} constituem o objecto \emph{material} da vontade, porque o objecto \emph{formal} é constituido pelo \emph{fim}, de que deriva toda a razão de \emph{bem}. Por isso \textsc{Deus}, a quem basta a sua bondade, pode querer, por amor da sua mesma bondade, outras coisas, que a Elle se refiram. Diz S. Thomaz: « Ex hoc quod voluntati divinae sufficit sua bonitas, non sequitur quod nihil aliud velit, sed quod nihil aliud vult nisi ratione suae bonitatis. Sicut etiam intellectus divinus, licet sit perfectus ex hoc ipso quod essentiam divinam cognoscit, tamen in ea cognoscit alia ». (\emph{Sum}. \emph{Th}., p. I, q. 19, a. 2 ad 3).\par Obj. 2.ª --- \textsc{Deus}, desde a eternidade, possua o seu fim, que é a sua bondade, e por isso não pode querer outras coisas; porque quem quer uma coisa pelo fim, deixa de querel-a, desde que seja alcançado o fim.\par Resposta. \textsc{Deus}, embora desde a eternidade possua o seu fim, que é a sua propria bondade, comtudo pode querer e amar outras coisas, como meios ordenados para o fim. Por quanto, como advertimos, a vontade divina não dispõe e ordena os meios para o fim pelo mesmo motivo, por que os dispõe e ordena a vontade humana. Esta dispõe e ordena os meios para o fim, com o intuito de alcançar um fim ou um bem, que não possue, e por isso, alcançado o fim, cessa o uso e o amor dos meios. A vontade divina, pelo contrario, quer as outras coisas pela sua bondade, emquanto, por amor da sua bondade, confere a outras coisas uma certa bondade, que seja semelhança e participação do bem absoluto, e por isso o amor dos \emph{meios} não só não cessa, quando se possue o \emph{fim}, mas funda-se propriamente na posse do mesmo \emph{fim}, ou no \emph{fim} emquanto possuido. Por isso mesmo que é a bondade infinita, \textsc{Deus} a derrama largamente sobre as outras coisas; e, amando a sua bondade, produz\pgn{500} a bondade nos outros. Diz S. Thomaz: « \textsc{Deus} non ordinat creaturas in finem bonitatis suae, quasi per eas bonitatem suam assequatur, sed ut ipsae creaturae divina operatione similitudinem aliquam divinae bonitatis acquirant, quod esse non posset, nisi eo volente et faciente » (I \emph{Sent}. Dist. 45, q. 1, a. 2 ad 4).\par Obj. 3.ª --- Quando são varios e diversos os \emph{objectos} da vontade, varios tambem e diversos são os \emph{actos} da mesma faculdade. Ora em \textsc{Deus} não pode haver senão um unico acto, como não ha senão um unico ser. Logo \textsc{Deus} não quer, não pode querer as outras coisas.\par Resposta. A unidade ou multiplicidade dos actos de uma faculdade depende, como dissemos, da unidade ou multiplicidade do objecto \emph{formal}, e não do objecto \emph{material}; porque é pelo objecto \emph{formal} que se especificam e distinguem os actos das faculdades. D'onde se segue que a unidade do acto é compativel com a multiplicidade dos objectos \emph{materiaes}, quando muitos objectos materiaes são attingidos sob um só aspecto \emph{formal}, segundo o qual movem e determinam a faculdade. Por isso, assim como a intelligencia divina percebe todas as coisas por um só acto, porque é determinada pela sua essencia, que é a forma intelligivel, representativa de todas as coisas, assim tambem a divina vontade quer todas as coisas com um só acto, porque é levada a querer pela sua bondade, a qual é a unica razão, por que \textsc{Deus} ama primeiramente o seu ser, e segundariamente as outras coisas. Diz S. Thomaz: « Sicut intelligere divinum est unum, quia multa non videt nisi in uno; ita velle divinum est unum et simplex, quia multa non vult nisi per unum, quod est bonitas sua » (\emph{Sum}. \emph{Th}., p. I, q. 19, a. 2 ad 4).\par Obj. 4.ª --- Não parece \emph{conveniente} que a infinita Magestade de \textsc{Deus} se abaixe até ás creaturas, amando-as e communicando-lhes a sua bondade.\par Resposta. Negamos que não \emph{convenha} á bondade divina communicar aos outros os proprios thesouros. Se as creaturas, tão pobres em bens e em generosidade, communicam aos outros os proprios bens, como não deverá admittir-se esse tendencia em \textsc{Deus}, tão rico e tão generoso? Nem devemos receiar pela Magestade de \textsc{Deus}. Elle encontra, na sua infinita sabedoria, o modo de se communicar, conveniente á sua grandeza e dignidade, e do acto da sua generosidade tirará sempre a sua maior gloria. Mas escutemos um magnifico trecho de S. Thomaz: « Res naturalis non solum habet naturalem inclinationem respectu proprii boni, ut acquirat ipsum cum non habet, vel ut quiescat in illo cum habet; sed etiam ut proprium bonum in alia diffundat, secundum quod possibile est. Unde videmus quod omne agens, inquantum est actu et perfectum, facit sibi simile. Unde et hoc pertinet ad rationem bonitatis, ut bonum quod quis habet aliis communicet, secundum quod possibile est. Et hoc praecipue pertinet ad bonitatem divinam, a qua per quamdam similitudinem derivatur omnis perfectio. Unde si res naturales, inquantum perfectae sunt, suum bonum aliis communicant, multo magis pertinet ad voluntatem divinam ut bonum suum aliis per similitudinem communicet, secundum quod possibile est (scilicet per similitudinem \emph{imperfectam} et \emph{analogam}). Sic igitur vult et se et alia: sed se ut finem, alia vero ut ad finem, inquantum condecet divinam bonitatem etiam alia ipsam participare » (\emph{Sum}. \emph{Th}., p. I, q. 19, a. 2).}.

\pg{499} 75. \textbf{Deus} quer ou \textbf{ama necessariamente} o proprio ser e \textbf{livremente as outras coisas}.

a) \textbf{Deus} quer ou ama necessariamente o proprio ser. --- Nenhuma faculdade pode ser indifferente com relação ao objecto, que \pg{500}  lhe é adequado e que por isso lhe satisfaz todas as tendencias; porque, n'esse caso, a indifferença seria uma desordem.

\pg{501} Ora o ser divino, sendo a propria bondade ou perfeição, é o unico objecto adequado á vontade divina e satisfaz plenamente todas as tendencias d'ella; porque em tal objecto a vontade encontra tudo o que pode desejar. Logo \textsc{Deus} quer ou ama \emph{necessariamente} o proprio ser \footnote{Este argumento pode tambem propor-se d'este modo. O ser divino, por ser a propria bondade infinita, constitue o objecto da felicidade de \textsc{Deus}. Ora a vontade não pode deixar de amar o objecto, que constitue a propria felicidade e que é conhecido como tal. Logo \textsc{Deus} quer ou ama necessariamente o proprio ser. --- Nem podia ser diversamente. O ser divino, sendo infinitamente perfeito, é infinitamente amavel. E poderia a vontade divina negar-lhe o seu amor? Não. E' certo que a vontade não quer ou não ama sem razão sufficiente. Basta reflectir um instante nas imperfeições dos bens, que nos cercam, para dar á vontade motivos sufficientes para os rejeitar. Mas a intelligencia divina, que penetra com o seu olhar no oceano infinito das amabilidades do ser divino, não encontra o minimo defeito, a mais leve imperfeição, que possa deixar indifferente a vontade. Esta, portanto, não pode deixar de amar o bem infinito, o unico digno d'ella, e de o amar na medida da sua bondade, e por isso infinitamente e sem medida. Em \textsc{Deus} ha uma equação perfeita entre a intelligencia, o amor e a plenitude do ser divino, objecto d'este conhecimento e d'este amor.}.

b) \textsc{Deus} quer ou ama livremente as outras coisas. --- Todo o bem, que --- ou não é \emph{adequado} á vontade, isto é, não enche plenamente a capacidade d'essa faculdade, --- ou não tem um nexo \emph{necessario} com esse objecto adequado, não attrahe \emph{necessariamente} a mesma vontade, mas deixa-a na sua indifferença ou liberdade. Ora nenhum bem \emph{finito}, --- por isso mesmo que é finito, é objecto \emph{adequado} á divina vontade, que é infinita, --- ou tem um nexo \emph{necessario} com o objecto \emph{adequado} da vontade divina, que é constituido pela propria bondade infinita, pois o infinito é independente do finito. Logo \textsc{Deus} quer ou ama \emph{livremente} as outras coisas \footnote{Alem d'isso, pertence á perfeição da vontade, que ella tenda para o objecto em conformidade com a indole e o merito do mesmo objecto; de modo que deve tender \emph{necessariamente} para o objecto \emph{necessario}, e \emph{livremente} para o objecto \emph{livre} ou \emph{contingente}. Ora o ser divino, por ser a propria bondade infinita, apresenta-se como o objecto summamente bom e necessario, e as outras coisas, que existem fóra de \textsc{Deus}, apresentam-se como bens finitos, e contingente ou livremente amaveis. Logo a vontade divina, sendo infinitamente perfeita, deve amar \emph{necessariamente} o proprio ser divino e \emph{livremente} as creaturas. --- Ha mais. Tudo o que \textsc{Deus} ama, ama-o emquanto se ama a si mesmo; porque, sendo ordenado o seu amor, não pode amar coisa alguma senão emquanto está ligada com o ultimo fim, de que deriva toda a amabilidade para os outros bens. Ora \textsc{Deus} conhece perfeitamente que nenhum amor de coisa finita está \emph{necessariamente} ligado com o amor do proprio ser; porque a sua bondade é absoluta e perfeitissima, e não precisa\pgn{502} de se associar com qualquer outro bem. Logo é impossivel que \textsc{Deus} queira ou ame \emph{necessariamente} outros bens, distinctos do proprio bem. --- Finalmente, a propria tendencia necessaria e irresistivel da vontade divina para o bem infinito é a razão da contingencia ou liberdade da sua tendencia para os bens finitos e imperfeitos. A infinita intelligencia de \textsc{Deus} descobre nestes bens uma multidão de motivos, que, se o podem levar a querel-os, podem-no tambem levar a não querel-os, ficando sempre livre a respeito d'elles. \textsc{Deus} não precisa de nenhum bem, de nenhuma coisa finita, para a sua perfeita felicidade; e por isso Elle é tao independente no seu amor, como o é no seu ser. Emquanto nós amamos as coisas \emph{porque} são \emph{bôas}; \textsc{Deus} ama as coisas \emph{para} as \emph{fazer bôas}, e dôar-lhes com o ser essa bondade, que nunca teriam tido sem Elle. Emquanto o nosso amor vai procurando o bem para o abraçar e encontrar na posse d'elle a felicidade, o amor divino \emph{produz} o bem e derrama-o sobre todas as coisas. Se, pois, \textsc{Deus} não depende, de modo algum, das creaturas, mas estas dependem, em tudo e por tudo, de \textsc{Deus}, é claro que o amor de \textsc{Deus} para com as creaturas é plenamente \emph{livre}.\par Mas escutemos S. Thomaz, que prova ambas as partes da nossa proposição. Diz assim: « Voluntas divina \emph{necessariam} habitudinem habet ad bonitatem suam, quae est \emph{proprium} ejus objectum. Unde bonitatem suam Deus \emph{ex necessitate} vult; sicut et voluntas nostra ex necessitate vult beatitudinem, sicut et quaelibet alia potentia habitudinem habet ad proprium et principale objectum, ut visus ad colorem, quia de sui ratione est ut ad illud tendat. Alia autem a se Deus vult in quantum ordinantur ad suam bonitatem, ut in finem. Ea autem quae sunt ad finem, non ex necessitate volumus, volentes finem, nisi sint talia, sine quibus finis esse non potest; sicut volumus cibum, volentes conservationem vitae, et navem, volentes transfretare. Non sic autem ex necessitate volumus ea, sine quibus esse potest, sicut equum ad ambulandum, quia sine hoc possumus ire; et eadem ratio est in aliis. Unde cum bonitas Dei sit perfecta, et esse possit sine aliis, cum nihil ei perfectionis ex aliis accrescat, sequitur quod alia a se eum velle non sit necessarium \emph{absolute}. Et tamen est necessarium \emph{ex suppositione}; supposito enim quod vult, non potest non velle, quia non potest voluntas eius mutari » (\emph{Sum}. \emph{Th}., p. I, q. 19, a. 3).\par Advirta-se o nexo d'esta proposição com a precedente. Na precedente, demonstrámos que \textsc{Deus}, com o mesmo simplicissimo acto da sua vontade, quer o proprio ser como \emph{fim} e as outras coisas como \emph{meios}. Nesta proposição, provamos que o mesmo simplicissimo acto, --- que, considerado em relação ao seu \emph{principio}, \emph{é necessario}, porque é uma e a mesma coisa com a essencia, a qual é necessaria, --- considerado em relação ao seu \emph{termo}, \emph{é necessario}, emquanto se refere á divina bondade, que é o seu objecto \emph{primario} e a razão pela qual ama todas as coisas, mas é \emph{contingente} ou \emph{livre}, emquanto se refere ás outras coisas, que são o seu objecto \emph{secundario}. --- Por isso esta questão attinge propriamente a \emph{liberdade divina}. Por quanto, se se admitte que o acto da vontade divina tem uma ordem ou relação \emph{necessaria} a \emph{tudo} o que \textsc{Deus} quer, \textsc{Deus} não pode chamar-se \emph{livre}. Se, pelo contrario, o mesmo acto da vontade divina tem uma relação \emph{necessaria} á divina bondade, de \pgn{503}modo que não pode deixar de a amar, e tem uma relação não \emph{necessaria} ás outras coisas, de modo que pode deixar de as amar, então deve admittir-se em \textsc{Deus} uma perfeitissima \emph{liberdade}, que não importa \emph{contingencia} e \emph{potencialidade} da parte do acto da vontade, mas da parte das coisas, queridas por \textsc{Deus}. --- \emph{Que seja assim}, isto é, que o mesmo acto da divina vontade seja \emph{necessario} em ordem á divina bondade, e seja \emph{livre} em ordem ás outras coisas, entende-se facilmente. \emph{Como possa ser assim}, isto é, como é que um e o mesmo acto acto é \emph{necessario} e \emph{livre} em ordem a termos diversos, não se comprehende pela limitação da nossa intelligencia e pela transcendencia das operações divinas. Alguma coisa dissemos a este respeito, tratando da \emph{immutabilidade} de \textsc{Deus}. Accrescentamos uma outra consideração. O acto da vontade divina não é \emph{causado} nem \emph{especificado} pelas coisas queridas, que a elle se referem de um modo contingente; porque estas não são \emph{necessariamente} ligadas com o fim, isto é, com a bondade divina, que pode subsistir sem ellas. Pelo contrario, o acto livre da vontade creada é \emph{causado} e \emph{especificado} pelo seu \emph{objecto}, porque depende d'este, e por isso a vontade creada não pode, com o mesmo acto, querer e deixar de querer, querer necessariamente e querer livremente. Sendo, pois, o acto das creaturas distincto da vontade, e referindo-se ao seu objecto como o effeito á causa, e multiplicando-se, por isso, conforme a multidão dos objectos e a diversa relação que tem com elles, não admira que o acto da vontade divina, --- que é uma e a mesma coisa com a propria vontade e que de maneira nenhuma é causado pelos objectos, que quer livremente, --- possa ser um e o mesmo, seja qual fôr a variedade dos objectos queridos, e a diversa relação que estes têm com elle. Concluimos que a unica e simplicissima volição divina, por ser infinita, pode corresponder e corresponde a muitas e diversas volições, e é \emph{necessaria} ou \emph{livre}, conforme o seu objecto é o bem infinito ou o bem limitado.\par Contra a segunda parte da proposição --- que \textsc{Deus} ama \emph{livremente} as outras coisas --- levantam muitas objecções. Expomos as principaes.\par Dizem: \textsc{Deus} \emph{necessariamente} conhece tudo o que conhece, porque a sua sciencia se identifica com a essencia. Ora tambem a vontade de \textsc{Deus} identifica-se com a sua essencia. Logo \textsc{Deus} tambem quer \emph{necessariamente} tudo o que quer.\par Resposta. Ha uma grande differença entre o acto da intelligencia e o da vontade, e esta differença deriva da diversa natureza do conhecimento e do appetite. Com effeito, o conhecimento faz-se emquanto o conhecido se encontra no cognoscente; por isso \textsc{Deus} conhece as coisas emquanto existem na sua essencia; e, como as coisas existem \emph{necessariamente} na essencia divina, assim \emph{necessariamente} \textsc{Deus} as conhece. Pelo contrario, a vontade refere-se ás coisas, emquanto e como estas existem em si mesmas, e, como as coisas não são em si necessarias, \textsc{Deus} não as quer \emph{necessariamente}. Por isso, devemos distinguir a proposição \emph{maior}, e dizer que \textsc{Deus} entende \emph{necessariamente} as coisas, não porque a sua sciencia se identifica com a essencia, mas porque a sciencia refere-se ás coisas conhecidas, emquanto são na essencia, e, como nesta existem \emph{necessariamente}, \emph{necessariamente} \textsc{Deus} as conhece. Como tambem devemos distinguir a proposição \emph{menor} e dizer que a \pgn{504}vontade de \textsc{Deus} se identifica com a essencia, mas nem sempre se refere ao que está na mesma essencia, porque pode referir-se a coisas, que estão fóra da essencia, e que por isso não são queridas \emph{necessariamente} por \textsc{Deus}. --- Devemos advertir que o acto de entender importa uma perfeição no ser intelligente, e o acto de querer importa uma perfeição na coisa querida, e que, por isso, se \textsc{Deus} não entendesse tudo, não seria absolutamente perfeito, mas, se não quizesse as outras coisas, não deixaria de ser perfeito. --- S. Thomaz resolve a difficuldade proposta com estas palavras: « Sicut divinum esse in se est necessarium, ita et divinum velle et divinum scire; sed divinum scire habet necessariam habitudinem ad scita, non autem divinum velle ad volita. Quod ideo est, quia scientia habetur de rebus secundum quod sunt in sciente; voluntas autem comparatur ad res secundum quod sunt in se ipsis. Quia igitur omnia alia habent necessarium esse secundum quod sunt in Deo; non autem secundum quod sunt in seipsis habent necessitatem absolutam, ita quod sint per se ipsa necessaria; propter hoc Deus quaecumque scit, ex necessitate scit; non autem quaecumque vult, ex necessitate vult » (\emph{Sum}. \emph{Th}., p. I, q. 19, a. 3 ad 6).\par Continuam: Se \textsc{Deus} não quer \emph{necessariamente} as creaturas, pode deixar de as querer. Ora este consequente é falso, porque quem pode deixar de querer o que quer, tem uma vontade \emph{mutavel}. Logo \textsc{Deus} quer \emph{necessariamente} as outras coisas, isto é, as creaturas.\par Resposta. Distinguimos a proposição \emph{maior}, dizendo --- ser verdade que, se \textsc{Deus} quer \emph{livremente} as creaturas, pode deixar de as querer em sentido \emph{absoluto}, --- mas ser falso que, se \textsc{Deus} quer \emph{livremente} as creaturas, pode deixar de as querer mesmo na \emph{supposição} que as tenha querido. Porquanto, tudo o que \textsc{Deus} quer, tanto \emph{necessaria} como \emph{livremente}, é desde a eternidade que o quer; porque o acto da vontade divina é medido, como o seu ser, pela eternidade. Mas ha esta differença entre a vontade necessaria e a livre: a vontade necessaria refere-se por tal modo ao seu objecto, que não pode absolutamente deixar de o querer, emquanto a vontade livre refere-se por tal modo ao seu objecto, que absolutamente podia não tel-o querido, embora, se por \emph{supposição} o quiz, não possa deixar de o ter querido ou de o querer. E esta necessidade \emph{hypothetica} ou de \emph{supposição} chama-se necessidade de \emph{immutabilidade}, porque é propriamente a \emph{immutabilidade} da vontade divina que exige que \textsc{Deus} não possa deixar de querer o que desde a eternidade quiz. Portanto a necessidade \emph{hypothetica} ou de \emph{supposição} não tem o mesmo caracter em \textsc{Deus} e nas creaturas; porque nas creaturas suppõe successão e mudança; o que não suppõe em \textsc{Deus}.\par Insistem: \emph{Convém} á divina bondade communicar aos outros a propria perfeição. Ora \textsc{Deus} não pode deixar de fazer o que \emph{convém} á sua bondade. Logo \textsc{Deus} quer \emph{necessariamente} as outras coisas.\par Resposta. São sobretudo os sequazes de \emph{Rosmini} que propugnam um nexo \emph{necessario}, não \emph{physico} mas \emph{moral}, entre a divina bondade e a existencia das creaturas (Cf. Rosmini, \emph{Teosofia}, v. I). Mas o erro é manifesto. Ha duas especies de \emph{conveniencia} ou \emph{decencia}: uma, quando o contrario \emph{não convém}; outra, quando tambem o contrario \emph{convém}. Ora não é a \emph{primeira} especie que se verifica na creação, mas é a \emph{segunda}. Com effeito, é \emph{conveniente}\pgn{505} que a bondade divina, sendo infinita, communique aos outros os seus thesouros; mas, se \textsc{Deus} os não quizer communicar, não ha nisso nenhum \emph{inconveniente} para a sua bondade, sendo esta perfeitissima e absolutamente independente de qualquer creatura. Inconveniente e absurdo seria o acto divino, se se admittisse nelle uma \emph{necessidade}, embora \emph{moral}, a respeito das creaturas; porque, em tal caso, \textsc{Deus} seria dependente das mesmas creaturas. --- Em poucas palavras: Assim como \emph{convém} que \textsc{Deus} queira outras coisas, porque Elle é a fonte inexhaurivel de todo o bem, assim tambem \emph{convém} que \textsc{Deus} não queira outras coisas, porque Elle é um bem infinito, sufficientissimo a si mesmo. A \emph{conveniencia} ou \emph{decencia} de \textsc{Deus} é compativel tanto com a hypothese da creação, como com a de \emph{não-creação}.\par *\par Das coisas expostas resultam os seguintes corollarios:\par a) \textsc{Deus} é infinitamente \emph{bemaventurado}. \emph{A bemaventurança}, segundo \emph{Boecio} (\emph{De cons}., l. IV) é « estado perfeito pela reunião de todos os bens --- \emph{status omnium bonorum aggregatione perfectus} ». --- A bemaventurança encerra --- a perfeição absoluta, --- o conhecimento d'esta perfeição, --- a complacencia na sua posse, --- a isenção de todos os males, presentes e futuros, --- e o dominio supremo sobre todas as coisas contingentes. Ora \textsc{Deus} --- contém na sua simplicissima essencia a \emph{plenitude} da perfeição, --- conhece e comprehende \emph{adequadamente} a propria essencia e nella todas as coisas, --- quer e ama \emph{necessariamente} a propria bondade ou perfeição, --- não pode, por ser \emph{immutavel}, perder algum bem ou adquirir algum mal, --- tem o dominio supremo e absoluto sobre todas as creaturas. Logo \textsc{Deus} é \emph{bemaventurado}. E é bemaventurado \emph{infinitamente}, porque \emph{infinita} é a sua \emph{perfeição}, \emph{infinito} é o \emph{conhecimento} da sua intelligencia, \emph{infinito} é o \emph{amor} da sua vontade, e por isso é a propria \emph{bemaventurança}. --- \textsc{Deus} não só é a bemaventurança de si mesmo, mas é tambem a bemaventurança de todas as creaturas intelligentes. Por quanto, a bemaventurança das creaturas não pode deixar de ser uma participação finita da infinita bemaventurança de \textsc{Deus}, summo bem e fonte de todo o bem. --- \textsc{Deus} é a bemaventurança \emph{objectiva} das creaturas, emquanto é conhecido, amado e possuido (\emph{Sum}. \emph{Th}., p. I, q. 27).\par \emph{b}) \textsc{Deus} é dotado de \emph{liberdade}. E' uma consequencia do que dissemos acerca da vontade de \textsc{Deus} em relação ás creaturas. Diz S. Thomaz: « \emph{Liberum arbitrium habemus respectu eorum, quae non necessario volumus, vel naturali instinctu. Non enim ad liberum arbitrium pertinet quod volumus esse felices, sed ad naturalem instinctum. Unde et alia animalia, quae naturali instinctu moventur ad aliquid, non dicuntur libero arbitrio moveri. Cum igitur Deus ex necessitate suam bonitatem velit, alia vero non ex necessitate, respectu illorum, quae non ex necessitate vult, liberum arbitrium habet} » (\emph{Sum}. \emph{Th}., p. I, q. 19, a. 10). --- Esta liberdade é \emph{perfeitissima}. Com effeito, o maior ou menor grau de perfeição na liberdade consiste na maior ou menor independencia, de que a vontade é dotada com relação a um determinado bem. Ora \textsc{Deus} é infinitamente independente de todos os bens creados, porque de nada precisa para a sua felicidade. Logo a liberdade de \textsc{Deus} é \emph{perfeitissima}. O \emph{poder peccar} não constitue a \emph{essencia} da liberdade, mas, antes, é um seu \emph{defeito}, e por isso não pode encontrar-se em \textsc{Deus}. Diz S. Thomaz: « Quod liberum arbitrium diversa eligere possit, servato ordine finis, hoc pertinet ad \emph{perfectionem} libertatis ejus; sed quod eligat aliquid, divertendo ab ordine finis, hoc pertinet ad \emph{defectum} libertatis » (\emph{Sum}. \emph{Th}., p. I, q. 62, a. 8 ad 3).\par c) \textsc{Deus} não quer o mal \emph{physico} de um modo \emph{directo}. Porquanto, a vontade divina não tem e não pode ter outro objecto a não ser a sua infinita bondade, e, se ama outras coisas, ama-as emquanto participam da mesma bondade. Ora todo o mal, quer \emph{physico} quer \emph{moral}, é \emph{privação} de algum bem, e por isso, não tendo em si uma razão de appetibilidade, não pode ser objecto \emph{directo} da vontade divina. --- Todavia, se \textsc{Deus} não pode querer \emph{directamente} o mal \emph{physico}, pode querel-o \emph{indirectamente}, emquanto esse mal está \emph{accidentalmente} unido a um \emph{bem commum} ou a um bem de uma \emph{ordem mais elevada}; assim não repugna que \textsc{Deus}, querendo a \emph{justiça}, que é um bem, queira a \emph{pena} do peccador, que é um mal \emph{physico}, e, querendo a ordem do mundo, queira a natural corrupção de muitas coisas. --- Mas o mal \emph{moral}, que é o \emph{peccado}, \textsc{Deus} não o quer \emph{de modo algum}. O mal \emph{moral} contem a aversão da creatura do seu Creador, para quem, emquanto é Summo Bem, todas as coisas devem estar naturalmente ordenadas. Ora \textsc{Deus} não pode \emph{querer}, nem \emph{indirectamente}, tudo o que se oppõe á essencial ordem das coisas. Logo \textsc{Deus} não \emph{quer}, nem \emph{indirectamente}, o mal \emph{moral}. Não o quer, mas \emph{permitte}, e esta \emph{permissão}, diz S. Thomaz, é um \emph{bem}: « Deus vult \emph{permittere} mala fieri: et hoc est \emph{bonum} » (\emph{Sum}. \emph{Th}., p. I, q. 19, a. 9).}.

\textbf{76. A vontade divina é causa das coisas creadas.} --- Se \textsc{Deus}, na creação, operasse --- não pela deliberação da vontade, --- mas pela necessidade da natureza, operaria sempre a mesma \pg{502} \pg{503}  coisa e sempre do mesmo modo, e produziria tudo o que pode produzir, a não ser que encontrasse impedimento; porque assim opera quem opera por necessidade de natureza. Ora \pg{504}  \textsc{Deus} produziu muitas e diversas coisas, e por muitos e diversos modos, nem produziu tudo o que pode produzir, embora \pg{505}  não possa ser e não seja contrariado por ninguem. Logo \textsc{Deus} não opera pela necessidade da natureza, mas pela deliberação \pg{506}  liberação da sua vontade. Logo a vontade divina é causa das coisas creadas \footnote{Vejam-se outras razões, que traz S. Thomaz sobre este assumpto (\emph{Sum}. \emph{Th}., p. I, q. 19, a. 4), especialmente a \emph{primeira}, com que prova que na ordem das causas efficientes, o agente pela intelligencia e pela vontade deve preceder o agente pela natureza, porque só o agente pela intelligencia e pela vontade pode determinar o fim ao agente pela natureza, e todo o agente pela natureza tende para o seu fim; e por isso \textsc{Deus}, que é o primeiro na ordem dos agentes, deve operar pela intelligencia e pela vontade. --- A verdade d'esta proposição torna-se ainda mais clara, se se considera que a potencia divina, por ser infinita, não é por si determinada para produzir esta ou aquella coisa, e que a sciencia divina, sendo por si indifferente com relação a esta ou aquella especie de seres, não pode determinar a potencia a crear este ou aquelle ser. Logo só a vontade pode determinar tal indifferença, decretando que, entre possiveis infinitos, se produza este ou aquelle ser. (Cf. I \emph{Sent}., Dist. 45, a. 3).\par Sendo necessario, para a producção das coisas, o intervento da intelligencia e da vontade de \textsc{Deus}, pode perguntar-se: Qual das duas faculdades é a causa \emph{mais proxima} do effeito? --- Respondemos que a causa mais proxima é a \emph{vontade}. A sciencia nada pode fazer por si, sem o imperio da vontade, que tudo determina e move. (Cf. \emph{Sum}. \emph{Th}., p. I, q. 19, a. 4 ad 4).}.

77. \textbf{Propriedades do amor divino.} --- O amor divino é \emph{substancial}, \emph{immutavel}, \emph{efficacissimo}, \emph{benevolo}.

a) E' \emph{substancial}. Em \textsc{Deus} o amor, como o acto intellectual, identifica-se com o ser. Por isso \textsc{Deus} é \emph{amor}.

\pg{507} b) E' \emph{immutavel}. Identificando-se com o proprio ser, o amor divino é immutavel, como immutavel é o ser.

c) E' \emph{efficacissimo}. O homem ama as coisas, porque são bôas; \textsc{Deus}, pelo contrario, ama as coisas, para as tornar bôas. Por isso, o nosso amor é \emph{effeito} da bondade das creaturas; mas o amor divino é \emph{causa} de toda a bondade, de que são dotados os seres finitos \footnote{Por isso \textsc{Deus} não ama as creaturas com amor \emph{affectivo}. O amor diz-se \emph{affectivo}, quando se ama um objecto pela perfeição, que elle encerra em si e por si; de modo que o amor é um \emph{effeito} da mesma perfeição. Ora \textsc{Deus} não ama nem pode amar as coisas creadas por alguma perfeição, que ellas possuam em si e por si; porque as creaturas, antes de serem objecto da vontade divina, não tinham em si perfeição alguma. Logo \textsc{Deus} não ama as coisas creadas com amor \emph{affectivo}. --- Se \textsc{Deus} não ama as creaturas com amor \emph{affectivo}, ama-as com amor \emph{effectivo}. O amor é \emph{effectivo}, quando emtanto se ama un objecto, emquanto lhe se confere gratuitamente alguma perfeição. Ora \textsc{Deus} ama as creaturas, emquanto lhes concedeu todas as perfeições, de que são dotadas. Logo \textsc{Deus} ama as coisas creadas com amor \emph{effectivo}. Diz S. Thomaz: « Cum amare nihil aliud sit quam velle bonum alicui, manifestum est quod Deus omnia quae sunt amat, non tamen eo modo sicut nos. Quia enim voluntas nostra non est causa bonitatis rerum, sed ab ea movetur sicut ab objecto; amor noster, quo bonum alicui volumus, non \emph{est causa bonitatis} ipsius; sed e converso bonitas ejus, vel vera vel aestimata, \emph{provocat amorem}, quo ei volumus et bonum conservari quod habet, et addi quod non habet, et ad hoc operamur. \textsc{SED} AMOR \textsc{DEI} EST INFUNDENS ET CREANS BONITATEM IN REBUS » (\emph{Sum}. \emph{Th}., p. I, q. 20, a. 2).\par Por isso --- 1.º) \textsc{Deus} ama \emph{todas} as coisas \emph{existentes}. Porquanto, estas, emquanto \emph{existem}, são bôas; porque todo o ente é bom, e toda a perfeição é um bem. Ora as coisas não existem senão pela vontade de \textsc{Deus}. Logo todo o bem, que as creaturas existentes possuem, deriva da vontade de \textsc{Deus}, e por isso a \emph{todas} as creaturas existentes \textsc{Deus} quer algum bem. Ora querer um bem é amar. Logo \textsc{Deus} ama \emph{todas} as creaturas \emph{existentes}. (\emph{Sum}. \emph{Th}. ib.). --- 2.º) \textsc{Deus} mais ama as creaturas, que são \emph{mais perfeitas}. Uma creatura é melhor que outra, quando é dotada de natureza mais nobre, ou de uma bondade mais perfeita e elevada. E, como toda a bondade ou perfeição das creaturas deriva da vontade de \textsc{Deus}, é claro que \textsc{Deus} mais ama as creaturas mais perfeitas, porque lhes quiz um bem maior, uma perfeição mais elevada. --- Note-se que o amor se diz mais ou menos intenso emquanto ao \emph{objecto}, isto é, emquanto ao maior ou menor bem que se quer ou confere; porque, da parte do seu \emph{principio}, que é de \textsc{Deus}, o amor é simplicissimo e não admitte o mais nem o menos, mas é sempre egual (\emph{Ib}., a. 3). --- 3.º) Se considerarmos a indole da \emph{natureza creada}, \textsc{Deus} mais ama os anjos que os homens, mais os homens que as creaturas inferiores; porque a natureza dos anjos é superior á dos homens, e a dos homens é superior á das outras creaturas inferiores (\emph{Ib}., a. 4). --- 4.º) Se considerarmos o estado \emph{moral} das creaturas \textsc{Deus} mais ama os justos do que os peccadores; porque os justos, além da bondade \emph{physica}, que lhes é commum com os peccadores, possuem a bondade \emph{moral}, de que os outros estão privados (\emph{Ib}., a. 4). --- 5.º) A' per\pgn{508}gunta: se \textsc{Deus} ame os meramente \emph{possiveis}, --- não emquanto ao \emph{ser}, que elles têm em \textsc{Deus} (porque esse ser, identificando-se com o ser divino, é objecto de um amor \emph{necessario} da vontade divina), --- mas emquanto ao \emph{ser}, que os \emph{possiveis} teriam em si mesmos, se fossem creados por \textsc{Deus}, --- os philosophos e os theologos não dão a mesma resposta. A opinião mais provavel é --- que \textsc{Deus} ama as creaturas meramente \emph{possiveis} (que nunca terão existencia) com \emph{amor inefficaz de simples complacencia}, porque tambem ellas têm uma \emph{certa} entidade (embora \emph{não-actual}), constituida por notas especificas e individuaes, e, se \textsc{Deus} quizesse, poderiam existir em si mesmas, --- e que esse amor é \emph{necessario}, porque \textsc{Deus} não pode deixar de amar tudo o que tem uma relação necessaria com a sua essencia e omnipotencia, como é o ente meramente \emph{possivel},}.

\pg{508} d) E' \emph{benevolo}. \textbf{Deus}, amando-nos desde a eternidade e sem \textbf{algum merito} nosso, deu-nos a existencia, conserva-nos e rege-nos com paterna e sabia providencia, e todas as coisas \textbf{dispõe} para o nosso bem, para a nossa eterna felicidade \footnote{Quaes são os nossos deveres em vista dos innumeraveis beneficios, que o amor divino tem espalhado e espalha continuamente sobre nós? Todos esses deveres, como os beneficios de \textsc{Deus}, estão resumidos e concentrados na celebre formula do Apostolo S. João: « \emph{Nos ergo diligamus} Deum, \emph{quoniam} \textsc{Deus} \emph{prior dilexit nos} » (I Ep., IV, 19).}.

78. \textbf{Termo immanente} do \textbf{amor divino}. --- O acto da vontade divina, pelo qual \textsc{Deus} ama o proprio ser e todas as coisas, que participam da sua bondade infinita, deve ter o seu termo, como todo e qualquer acto. E, como o acto da vontade é formalmente \emph{immanente}, tambem \emph{immanente} deve ser o seu termo. Este termo \emph{immanente} do amor divino costuma chamar-se \emph{espirito}. --- O \emph{espirito}, em \textsc{Deus}, é unico, infinito, substancial, como unico, infinito, substancial é o acto da vontade divina, de que procede \footnote{\textsc{Deus} é purissimo espirito. Portanto o seu primeiro acto é o \emph{pensamento}. O pensamento de \textsc{Deus} não é, como o nosso, multiplice, vario, imperfeito, mas é um acto simplicissimo, unico, infinitamente perfeito. Fructo e termo d'este simplicissimo acto é o \emph{verbo}. --- Mas o espirito não só pensa, ama tambem. \textsc{Deus}, conhecendo-se a si mesmo no seu verbo, e, vendo-se infinitamente perfeito, não pode deixar de se amar com um amor, proporcionado á sua amabilidade. O termo d'este amor é o \emph{espirito}. O \emph{espirito} é o termo \emph{necessario} da vontade divina; porque o acto d'esta faculdade não é \emph{livre}, mas é \emph{natural}, isto é, nasce da mesma faculdade considerada como \emph{natureza}. O \emph{espirito é unico}, como \emph{unico} é o \emph{verbo}. Uma vontade infinita não pode exhalar senão um unico espirito, como uma intelligencia infinita não pode gerar senão um unico verbo. Uma variedade de termos suppõe uma repetição de actos, e por isso uma limitação nas faculdades operativas.\par A Fé vai mais além, e ensina que, em \textsc{Deus}, o \emph{Espirito} é uma Pessoa, subsistente na unica essencia divina, e procedente do Pae e do Filho, como de unico principio. \textsc{Deus} pensa e ama com toda a sua substancia, e toda a substancia Elle transfunde nos termos do seu pensamento e do seu amor.\pgn{509}Como o Pae, entendendo-se a si mesmo, se transfunde todo no Verbo, assim Pae e Filho, amando-se, se transfundem no termo do seu amor, que é o \emph{Espirito Santo}. Assim, o Espirito Santo é o termo substancial e pessoal do mutuo amor do Pae e do Filho. Se a intelligencia do Pae é o principio do Filho, a vontade do Pae e do Filho é o principio do Espirito Santo. --- Procedendo da \emph{intelligencia} do Pae, o Verbo é verdadeiro Filho, porque, como dissemos, a intelligencia, quando entende, torna semelhante a si mesma o objecto entendido, --- semelhança indispensavel para que uma origem possa chamar-se \emph{geração}; --- procedendo da \emph{vontade} do Pae e do Filho, o Espirito Santo não procede por \emph{geração}, porque a vontade, longe de tornar semelhante a si mesma o objecto querido, torna-se semelhante a elle. --- Como o Verbo não dá origem a outro verbo, porque é Elle o Verbo, que procede; assim o Espirito Santo não dá origem a outro Espirito, porque é Elle o Espirito que procede. São tres, portanto, as Pessoas divinas; porque dois são os actos \emph{immanentes} de \textsc{Deus} --- o conhecimento e o amor, e dois são os termos d'estes actos. --- O Pae, que não procede, é a \emph{primeira} Pessoa; o Filho, que procede da intelligencia do Pae, é a \emph{segunda}; o Espirito Santo, que procede da vontade do Pae e do Filho, é a \emph{terceira} Pessoa, porque o acto da vontade é posterior, na origem, ao acto da intelligencia. --- A theologia \emph{dogmatica} apresenta noções mais minuciosas e profundas acerca d'este altissimo mysterio da SS. Trindade, á qual seja gloria, benção e amor por toda a eternidade!}.

%%>>> art-43

\section*{Artigo III.\\ Vida divina, sua existencia e natureza.}
\addcontentsline{toc}{section}{Artigo III. --- Vida divina, sua existencia e natureza.}

{\small\noindent\textsc{Summario}: \textsc{Summario:} --- Creação, sua natureza, principio e fim. --- Conservação, sua divisão e necessidade. --- Concurso, sua natureza e extensão. --- Providencia, sua existencia, extensão e propriedades.\par}

\pg{509} 79. \textbf{Vida divina}. --- \emph{Vida divina} é o attributo, pelo qual \textsc{Deus} produz operações \emph{immanentes}, que têm o \emph{principio} e o \emph{termo} no mesmo sujeito. --- Diremos da sua \emph{existencia} e \emph{natureza} \footnote{A palavra \emph{vida}, embora formada pelo facto de se perceberem, nalguns entes, umas \emph{operações immanentes}, que tinham nelles o seu principio e o seu termo, comtudo emprega-se para denotar a \emph{substancia}, que tem esse poder de se mover a si mesma, ou de produzir as operações \emph{immanentes}. Por isso a \emph{vida} toma-se em acto \emph{primeiro} e em acto \emph{segundo}. Tomada em acto \emph{primeiro}, significa a propria \emph{substancia} do ser vivo; d'ahi o adagio: \emph{Vivere viventibus est esse}. Tomada em acto \emph{segundo}, significa a \emph{operação vital}, isto é, a operação propria dos seres vivos. Diz S. Thomaz: « Vita dicitur dupliciter: uno modo ipsum \emph{esse} viventis...; alio modo ipsa \emph{operatio} viventis, secundum quam principium vitae in actu reducitur » (\emph{Sum}. \emph{Th}. 1-2, q. 3, a. 2 ad 1). Por isso a vida é um attributo \emph{substancial}. (Cf. \emph{Sum}. \emph{Th}., p. I, q. 18; --- c. \emph{Gentes}, I. c. 97 e seg.).}.

\textbf{80. Deus é dotado de vida.}

a) Um ser é dotado \emph{de vida}, quando produz ou é capaz de produzir operações \emph{immanentes}. Ora \textsc{Deus}, como provámos, produz operações \emph{immanentes}, que são as da sua \emph{intelligencia} e da sua \emph{vontade}. Logo \textsc{Deus} é dotado de \emph{vida}.

\pg{510} b) O ser divino deve conter, na sua simplicissima unidade,

toda a plenitude da perfeição, e por isso todas as perfeições

\emph{absolutas}. Ora a \emph{vida} é uma perfeição \emph{absoluta}, porque no seu

conceito não encerra imperfeição alguma e não se oppõe a

perfeições mais elevadas. Logo \textsc{Deus} é dotado de \emph{vida} \footnote{E' certo que tudo o que se encontra no effeito deve encontrar-se na causa --- ou de um modo \emph{eminente}, --- ou de um modo \emph{proprio} e \emph{formal}. Por isso, em \textsc{Deus}, causa primeira de todas as creaturas, devem encontrar-se todas as perfeições creadas --- ou de um modo \emph{eminente}, --- ou de um modo \emph{proprio} e \emph{formal}. Ora a \emph{vida}, considerada no seu conceito \emph{essencial}, não encerra senão \emph{perfeição}; porque \emph{viver} é poder produzir um acto, independentemente de um principio \emph{extrinseco}, e isto é perfeição. --- E' verdade que a \emph{vida} se faz consistir tambem no \emph{movimento immanente}, e que em \textsc{Deus} não ha nem pode haver \emph{movimento}, porque Elle é purissimo \emph{acto}, sem sombra de \emph{potencia}, e todo o \emph{movimento} importa uma composição de \emph{potencia} e de \emph{acto}. Mas, quando se diz que a \emph{vida} consiste num \emph{movimento immanente}, a palavra \emph{movimento} deve tomar-se em sentido \emph{latissimo}, emquanto significa toda e quelquer \emph{operação immanente} no sujeito, quer essa operação seja realmente distincta do sujeito, quer não. Neste sentido, essa definição de \emph{vida} pode applicar-se tambem á vida de \textsc{Deus}. Se, porém, a palavra \emph{movimento} se toma em sentido \emph{rigoroso}, emquanto importa uma passagem da potencia para o acto, e uma distincção \emph{real} outre o sujeito operante e a sua operação immanente, então essa definição de \emph{vida}, --- se convem á vida das \emph{creaturas}, que, nas operações, passam da potencia para o acto, e cujo acto é realmente distincto da substancia, --- não pode convir á vida do \emph{Creador}, que é acto purissimo e simplicissimo. --- Mas, repetimos, o \emph{essencial} para o conceito e natureza da \emph{vida} é que a operação seja \emph{immanente}, independente de um principio \emph{extrinseco}. Que essa operação \emph{immanente} se exerça com ou sem uma passagem da potencia para o acto, com ou sem composição de substancia e de acto, isto é \emph{accidental} para o conceito e natureza da \emph{vida}.}.

81. \textbf{Deus} é a \textbf{sua vida}.

a) A vida, nos entes vivos, é o proprio ser. Ora \textsc{Deus} é

o seu ser. Logo \textsc{Deus} é a sua vida.

b) Entender e querer é viver; porque os actos da intelligencia e da vontade são vitaes. Ora \textsc{Deus} é o seu entender e o seu querer. Logo é o seu viver e a sua vida \footnote{Alem d'isso, se \textsc{Deus} não fosse a sua vida, seria vivo, não por essencia, mas por participação, e por isso teria recebido de outro ente superior a vida, de que é dotado. Ora é absurdo dizer que Ser infinito tenha recebido de um ente superior alguma perfeição. Logo \textsc{Deus} é a sua vida. (C. \emph{Gentes}, l. I. c. 98).}.

82. A \textbf{vida} de \textbf{Deus} é \textbf{perfeitissima}.

a) Consistindo a vida na operação \emph{immanente}, é claro que

quanto mais \emph{immanente}, isto é, independente de um impulso

exterior, fôr a operação, tanto mais perfeita será a vida. Ora

a operação de \textsc{Deus}, identificando-se com a sua essencia, é \pg{511}  absolutamente independente de todo e qualquer impulso exterior. Logo a vida de \textsc{Deus} é \emph{perfeitissima}.

b) A vida de \textsc{Deus} é uma vida \emph{intellectual}, e por isso será tanto mais perfeita, quanto mais immunes do influxo de principios extranhos forem os actos da sua intelligencia e da sua vontade. Ora os actos da intelligencia e da vontade de \textsc{Deus} são inteiramente immunes de todo o influxo extranho; porque \textsc{Deus} é a propria Verdade e Bondade, e os seus actos, como tambem as faculdades, são uma e a mesma coisa com a essencia. Logo a vida de \textsc{Deus} é \emph{perfeitissima} \footnote{A vida \emph{intellectual} é propria das \emph{substancias espirituaes}, mas só em \textsc{Deus} ella attinge o mais alto grau de \emph{immanencia}, e por isso de perfeição. Diz S. Thomaz: « Cum vivere dicantur aliqua secundum quod operantur ex seipsis, et non quasi ab aliis mota; quanto perfectius competit hoc alicui, tanto perfectius in eo invenitur vita » (\emph{Sum}. \emph{Th}., p. I, q. 18, a. 3). --- Consideremos a vida \emph{intellectual} do \emph{homem}. A nossa intelligencia move-se a si mesma, forma a especie intelligivel, que é o \emph{principio} do conhecimento, e produz o \emph{verbo}, que é a imagem do objecto entendido e o \emph{termo} do acto intellectual. Tem, pois, em si o \emph{principio} e o \emph{termo} da sua operação, e esta operação é propriamente \emph{immanente}. Mas, se reflectirmos --- em que a \emph{especie intelligivel} é formada pela intelligencia humana depois de esta receber, por meios dos sentidos, um impulso da parte dos objectos externos, --- e em que o \emph{verbo mental}, embora unido com a intelligencia, é comtudo uma coisa accidental, distincto d'esta faculdade, devemos concluir que o nosso acto intellectual não é tão immanente, não é tão intimo, que não se possa conceber outro acto mais perfeito. --- Acto mais perfeito, infinitamente mais perfeito é o de \textsc{Deus}, porque elle, e só elle, é infinitamente immanente. \textsc{Deus}, absolutamente independente no ser, é absolutamente independente na operação. A especie intelligivel, que, segundo o nosso modo de entender, move e determina a intelligencia divina ao acto do entendimento, não depende, nem proxima nem remotamente, do impulso de um agente externo, porque é a propria essencia divina; como tambem o \emph{verbo}, que é o termo do acto do entendimento, não é uma coisa accidental, distincta da intelligencia, mas é a propria intelligencia, é a propria essencia divina. Mais immanente, mais intimo não pode ser o acto da intelligencia divina, e por isso não pode ser mais perfeito. Acto perfeitissimo, vida perfeitissima. --- O que dissemos com relação ao acto da intelligencia divina, deve applicar-se ao acto da vontade. A nossa vontade move-se a si mesma para um bem particular; mas, antes de se mover para um bem particular, ella recebe um impulso e é movida pelo ultimo fim, pela felicidade. Portanto o acto da nossa vontade é \emph{immanente}, mas não é perfeitamente immanente. Pelo contrario, a vontade de \textsc{Deus}, tendendo para a felicidade, não é movida por uma coisa extranha, porque Elle é a sua propria felicidade, o seu ultimo fim. Diz S. Thomaz: « Perfectior modus vivendi est eorum, quae habent intellectum. Haec enim perfectius movent seipsa... Sed quamvis intellectus noster ad aliqua se agat, tamen aliqua sunt ei praestituta a natura; sicut sunt prima principia, circa quae non potest aliter se babere, et ultimus finis, quem non potest non velle.\pgn{512}Unde, licet quantum ad aliquid moveat se, tamen oportet quod quantum ad aliqua ab alio moveatur. Illud igitur, cujus sua natura est ipsum ejus intelligere, et cui id, quod naturaliter habet, non determinatur ab alio, hoc est quod obtinet summum gradum vitae. Tale autem est Deus. Unde in Deo maxime est vita. Unde Philosophus, ostenso quod Deus sit intelligens, concludit quod habeat vitam perfectissimam et sempiternam, quia intellectus ejus est perfectissimus et semper in actu » (\emph{Sum}. \emph{Th}., p. I, q. 18, a. 3).}.

\pg{512} 83. \textbf{A vida de Deus é sempiterna.}

a) Um ser deixa de viver, quando d'elle se separa a vida. Ora a vida não pode separar-se de \textsc{Deus}; porque, em \textsc{Deus}, a vida é o seu ser, e uma coisa não pode separar-se de si mesma, pois a separação consiste na divisão de uma coisa de outra. Logo a vida de \textsc{Deus} é sempiterna.

b) Um ser, cuja vida não é sempiterna, mas tem o seu começo e o seu termo, e está sujeita á successão, é mutavel, composto e dependente de uma causa superior. Ora \textsc{Deus} é immutavel, simplicissimo, causa suprema de todas as coisas. Logo a vida de \textsc{Deus} é sempiterna \footnote{Estes e outros argumentos são de S. Thomaz (\emph{C}. \emph{Gentes}, 1. I, c. 99).}.

84. \textbf{Todas as coisas creadas são vida em Deus.} --- O viver de \textsc{Deus} é o seu entender. Ora, em \textsc{Deus}, o entender, o objecto entendido e a faculdade intellectual são uma e a mesma coisa. Portanto tudo o que existe em \textsc{Deus} como objecto entendido é o viver de \textsc{Deus}, é a sua vida. Ora todas as coisas creadas, sendo effeitos de \textsc{Deus}, existem nelle como objectos entendidos. Logo todas as coisas creadas são vida em \textsc{Deus} \footnote{O argumento, adduzido no texto, é tambem do Angelico Doutor (\emph{Sum}. \emph{Th}., p. I, q. 18, a. 4). --- Para que evitemos equivocos, S. Thomaz faz as seguintes advertencias. --- 1.º) As creaturas são ou existem em \textsc{Deus} por dois modos: emquanto são conservadas pelo poder divino (assim dizemos que existe em nós o que está em nosso poder), e emquanto são conhecidas pela intelligencia divina (assim dizemos que existe em nós uma pedra, cuja especie está na nossa intelligencia). Emquanto são conservadas pelo poder divino, ellas existem em \textsc{Deus} \emph{como existem nas suas proprias naturezas}. E' neste sentido que se devem tomar as palavras de S. Paulo, quando diz: \emph{In ipso} (\emph{Deo}) \emph{vivimus, movemur et sumus}; porque o nosso viver, o nosso ser e o nosso movimento são causados por \textsc{Deus}. Emquanto são conhecidas pela intelligencia divina, as creaturas são em \textsc{Deus} pelos exemplares archetypos, que as representam como são em si mesmas e que são a propria essencia divina. Por isso as coisas, emquanto existem em \textsc{Deus} por este segundo modo, são a essencia divina e a vida divina (\emph{Ib}. ad 1.) --- 2.º) Uma coisa feita segundo um modelo deve conformar-se com este, emquanto este é um \emph{modelo}, e não emquanto tem \emph{um certo modo de ser}. Porque o \emph{ser} do \emph{modelo} pode ser e é muitas vezes diverso do \emph{ser} da \emph{coisa modelada}; assim o modelo de uma casa, emquanto existe na intelligencia do artifice, tem um ser immaterial e intelligivel, ao passo que, se o considerarmos applicado \pgn{513}na propria casa, que existe fóra da intelligencia, o modelo tem um ser material e sensivel. Portanto as ideas exemplares divinas, se as considerarmos nas coisas creadas, ás quaes são applicadas, não vivem, ou podem não viver; mas, se as considerarmos na intelligencia divina, são a propria vida divina, porque tem o ser divino (\emph{Ib}. ad 2).}.

\pg{513} CAPITULO QUARTO

\textbf{OPERAÇÕES TRANSEUNTES DE DEUS}

%%>>> art-44

\section*{Artigo I.\\ Creação, sua natureza, principio e fim.}
\addcontentsline{toc}{section}{Artigo I. --- Creação, sua natureza, principio e fim.}

85. \textbf{Creação}. --- \emph{Creação} é o acto, pelo qual \textsc{Deus} tirou do \emph{nada} todas as coisas. --- Tirar do \emph{nada} o ente é produzil-o na sua \emph{totalidade}, isto é, não só emquanto á \emph{forma}, mas tambem emquanto á \emph{materia} \footnote{Tratamos, neste capitulo, dos principaes actos do amor \emph{effectivo} de \textsc{Deus} para com as creaturas, que são: a \emph{creação}, a conservação, o concurso e a \emph{providencia}. --- Neste artigo occupamo-nos da \emph{creação}. (\emph{Cf}. \emph{Sum}. \emph{Th}., p. I, q. 44, 45; \emph{c}. \emph{Gent}., l. II, c. 15 e seg).\par A \emph{creação} é a \emph{producção de uma coisa do nada}. Para que, pois, a producção de uma coisa possa chamar-se \emph{creação}, são necessarios e sufficientes dois requisitos: 1.º) que \emph{nada preexista} da coisa, que se diz \emph{creada}, a não ser a não-repugnancia para existir, ou a absoluta possibilidade, ou a potencia objectiva, e por isso a creação differe de todas as outras especies de producção, como são a geração, a alteração, etc.; --- 2.º) que a coisa, que se diz creada, passe do não-ser para o ser, não emquanto ella é alguma coisa no termo \emph{a quo}, que é o \emph{não-ser}, mas emquanto, pelo influxo de uma causa proporcionada, recebe o ser, que antes não tinha.\par A \emph{creação} diz-se \emph{activa} ou \emph{passiva}, conforme se considera em \textsc{Deus} ou na creatura. --- Considerada em \textsc{Deus}, a \emph{creação} é o acto, pelo qual \textsc{Deus} tira do \emph{nada} as coisas. E' a definição apresentada no texto. Este acto de \textsc{Deus} é a sua propria essencia, emquanto se refere ás creaturas. Para a comprehensão d'esta verdade, é necessario distinguir um duplice acto --- um \emph{immanente} e outro \emph{transeunte}. O \emph{immanente} é o que permanece no agente, e é uma perfeição do proprio agente, como entender, querer, etc.; o \emph{transeunte} é o que passa para o effeito externo, e é uma perfeição do proprio effeito, que existe por elle, como aquecer, fabricar, etc. Este duplice acto compete a \textsc{Deus}. Compete-lhe o primeiro acto, porque \textsc{Deus} entende, quer, ama, gosa; compete-lhe o segundo, emquanto produz as coisas, conserva-as e\pgn{514}rege-as. E porque o primeiro acto é perfeição do agente, e o segundo é perfeição do effeito, e o agente é anterior ao effeito e é causa d'elle, é evidente que o primeiro acto é a razão do segundo e naturalmente o procede (\emph{C}. \emph{Gent}., l. II, c. 1). --- Mas o \emph{acto} não se attribue a \textsc{Deus} e ás creaturas num sentido \emph{univoco}, mas só num sentido \emph{analogico}; e por isso do acto de \textsc{Deus} devemos afastar toda a imperfeição, que se encontra no acto das creaturas. O acto das creaturas sahe do agente e é recebido no paciente, porque é uma coisa \emph{accidental}, e por isso é formalmente \emph{transeunte}; ao passo que o acto divino não sahe de \textsc{Deus} para as creaturas, porque o acto de \textsc{Deus} é a sua propria \emph{substancia}, e é simplesmente \emph{immanente}. --- Todavia o acto divino não é tão \emph{immanente}, que se não possa chamar, de algum modo, \emph{transeunte}; porque convem com o acto \emph{transeunte} das creaturas em que produz um effeito fóra do agente. Por isso, o acto divino chama-se \emph{formalmente} immanente e \emph{eminente} ou \emph{virtualmente} transeunte, emquanto, ficando no agente com que se identifica, produz o que o agente creado não pode produzir senão por um acto, que d'elle passa para o paciente. Emquanto \textsc{Deus} entende e quer, produz o effeito. --- Considerada nas creaturas, a \emph{creação} é uma \emph{relação das creaturas com o Creador, como com o principio do proprio ser.} Esta relação é \emph{não-mutua;} porque, se as creaturas se referem \emph{realmente} a \textsc{Deus}, \textsc{Deus} só \emph{logicamente} refere-se ás creaturas; pois o \emph{fundamento} da relação, que é \emph{o ser tirado do nada}, só é \emph{real} nas creaturas. Comtudo, não devemos d'ahi concluir que \textsc{Deus} só \emph{logicamente} é Senhor das creaturas; porque, para \textsc{Deus} ser \emph{realmente} Senhor das creaturas, basta que estas \emph{realmente} derivem e dependam de \textsc{Deus}. --- Outras noções, relativas á \emph{creação}, foram expostas na \emph{Cosmologia}, tratando da \emph{origem do mundo}. Ahi consideramos a creação mais no aspecto \emph{passivo}, e aqui mais no aspecto \emph{activo}.}.

\pg{514} \textbf{86. A creação é o acto proprio de Deus.}

a) A \emph{creação}, por ser um acto \emph{absoluto}, isto é, independente de toda a materia, deve ser exclusivamente proprio da causa, que possue uma existencia \emph{absoluta}, isto é, independente de todo o ser; porque o modo de operar é proporcionado ao modo de existir. Ora só \textsc{Deus} possue uma existencia \emph{absoluta}. Logo a \emph{creação} é o acto proprio de \textsc{Deus}.

b) Todo o effeito deve ser proporcionado e semelhante á forma do agente; porque todo o agente opera pela sua forma. Ora a forma, por que \textsc{Deus} opera, é a sua essencia, e esta consiste em que \textsc{Deus} é o proprio \emph{Ser} subsistente, contendo em si toda a plenitude e perfeição da entidade. Logo o effeito proporcionado e semelhante a \textsc{Deus} é o proprio \emph{ser} das coisas. Mas o ser das coisas, pelo facto de abranger \emph{todos} os elementos constitutivos, só pode ser produzido por \emph{creação}. Logo a \emph{creação} é o acto proprio de \textsc{Deus} \footnote{Desenvolvamus e esclareçamos este ponto. A Causa primeira, que communica a outras coisas todo o ser, que possuem, e por isso produz o ente emquanto é ente, opera pela \emph{intelligencia} com o concurso da \emph{vontade} e da \emph{potencia operativa}. Para isso, para entendermos, de algum modo, a acto creador\pgn{515} de \textsc{Deus}, devemos fallar da sua \emph{potencia}, ou \emph{energia operativa}, tendo, no capitulo antecedente, fallado da sua intelligencia e vontade.\par A \emph{potencia}, sendo uma perfeição \emph{absoluta}, deve encontrar-se em \textsc{Deus} no seu conceito \emph{formal} e em toda a extensão da sua perfeição, excluindo todas as imperfeições e limitações, de que se acha cercada nas creaturas. Ora a \emph{potencia}, no seu conceito \emph{formal}, é \emph{um principio efficiente do effeito}. Este conceito, como se vê, importa simplesmente o \emph{acto}, e exclue toda a imperfeição da \emph{potencialidade}. --- Mas a \emph{potencia} attribue-se só \emph{analogicamente} a \textsc{Deus} e ás creaturas. Por isso, nas creaturas, a \emph{potencia} distingue-se \emph{realmente} tanto da \emph{essencia}, como do \emph{acto}, \emph{é principio efficiente} do \emph{acto} e do \emph{effeito}, e está sempre unida com a potencialidade, porque suppõe a passagem da potencia para o acto; ao passo que, em \textsc{Deus}, no qual a \emph{essencia}, a \emph{potencia} e o \emph{acto} são uma e a mesma coisa (pois \textsc{Deus} é o proprio Ser subsistente), a \emph{potencia operativa}, emquanto ás operações \emph{transeuntes}, não é principio efficiente do \emph{acto}, mas unicamente do \emph{effeito}, e, d'este modo, encontra-se em \textsc{Deus} em toda a sua perfeição. (\emph{Cf. Sum. Th.}, p. I, q. 25, a. 1 ad 3).\par Emquanto ás relações da \emph{potencia operativa}, ou \emph{activa}, de \textsc{Deus} com a sua \emph{intelligencia} e \emph{vontade}, dizemos que, na producção e governo das creaturas, a intelligencia \emph{dirige}, a vontade \emph{manda} e a potencia \emph{executa}. Por quanto, \textsc{Deus}, operando fóra da sua essencia não por necessidade da sua natureza, mas pela liberdade do seu juizo, não produz todos os effeitos, que a sua potencia é capaz de produzir; mas, como convém a um agente independente, que livremente se determina a produzir um effeito preconcebido pela intelligencia, produz, sem alguma mudança de si mesmo, os effeitos externos, e os produz no modo, na ordem e no tempo, que o eterno juizo da sua vontade livremente determinou. Por isso, todos os effeitos attribuem-se, como á causa efficiente, ao imperio da vontade divina, porque a potencia só produz os effeitos, que a vontade imperou. Com razão diz o Psalmista (\emph{Ps}. 148): « \emph{Ipse dixit et facta sunt}; \emph{Ipse mandavit et creata sunt} » (\emph{Cf. S. Th.}, 2 \emph{Sent}. Dist. I, q. 1, a. 5 ad 11).\par A potencia divina, identificando-se com a \emph{essencia}, é simplesmente \emph{infinita} (\emph{Sum. Th.}, p. I, q. 25, a. 2). --- E pode chamar-se \emph{infinita} tambem com relação aos \emph{effeitos}, --- não emquanto a potencia divina produza uma multidão actualmente infinita, ou produza um effeito infinitamente perfeito (o que repugna á razão de effeito), --- mas emquanto não pode exgottar-se na producção dos effeitos, e a efficacia da sua acção nas creaturas é multiplice e illimitada. Diz S. Thomaz: « Numquam (Deus) tot effectus facit, quin plures facere possit, nec umquam ita intense operatur, quin intensius operari possit. Intensio autem in operatione divina non est attendenda secundum quod operatio est in operante, quia sic semper est infinita, cum operatio sit divina essentia; sed attendenda est secundum quod attingit effectum. Sic enim a Deo moventur quaedam efficacius, quaedam minus efficaciter » (\emph{De pot.},\pgn{516}q. 1, a 2). --- D'onde se segue que \textsc{Deus} é \emph{omnipotente}, emquanto pode fazer todo o que é \emph{absolutamente possivel}. \emph{O absolutamente impossivel} é uma \emph{contradicção}, e a contradicção não pode ser feita. (\emph{Sum. Th.}, p. I, q. 25, a. 3, --- \emph{c. Gent.}, l. II, c. 25, --- \emph{De pot.}, q. 1, a. 3).\par Uma potencia \emph{infinita} deve, na sua operação, ser independente de toda e qualquer coisa extranha, e por isso, para produzir o effeito, não exige que alguma coisa d'esse effeito preexista ou se presupponha. O acto, que não presuppõe nada da coisa, que por elle é produzida, diz-se \emph{creação}. --- A \emph{creação} é o acto proprio, especificamente proprio de \textsc{Deus}; porque só \textsc{Deus}, por ser a Causa universal e o primeiro Principio das coisas, pode, na sua operação, ser independente de toda e qualquer causa anterior, material ou efficiente. Por isso é absurdo dizer que a creatura possa, por si, tirar do \emph{nada} as coisas; porque é absurdo que a creatura possa ser o primeiro principio das coisas.\par Faz-se a questão se \textsc{Deus} possa communicar á creatura a virtude ou força de crear. Os escholasticos respondem que essa virtude ou força não pode ser communicada á creatura, e dão as seguintes razões. --- A creação exige uma potencia \emph{infinita}, que opere pelo \emph{ser}, como tal, e não pelo ser \emph{determinado} a uma certa especie; porque o termo da creação é o \emph{ser}, e a causa deve ser proporcionada ao effeito. Ora nenhuma creatura pode receber uma qualidade infinita, porque uma qualidade recebida deve accommodar-se á capacidade do sujeito, que a recebe. --- Além d'isso, o acto da creatura é uma coisa \emph{accidental}, que deve ser necessariamente recebido num sujeito para que possa produzir o effeito, e por isso a creatura não pode fazer nada senão transformando a materia preexistente: o que é diametralmente opposto ao acto \emph{creador}. Portanto a creação exige uma força, que se identifique com a essencia do agente, que seja acto puro e não seja distincto da propria operação. (\emph{De pot.}, q. 3, a 4). --- S. Thomaz accrescenta que a creatura não pode intervir, nem como causa \emph{instrumental}, no acto \emph{creador} de \textsc{Deus}. De facto, o \emph{instrumento} serve para dispór o sujeito, de modo que este possa receber em si o acto do agente principal; ora, como a creação não presuppõe algum sujeito, a creatura não pode intervir, nem como \emph{instrumento}, no acto creador de \textsc{Deus}. Além d'isso, o acto de uma causa, embora \emph{instrumental}, deve derivar da potencia d'ella; ora, sendo impossivel que da potencia de uma causa \emph{finita} derive um acto de uma força infinita, qual é o acto necessario para a creação, é tambem impossivel que a creatura seja causa \emph{instrumental} do ente \emph{creado} (\emph{Cf. De pot.}, q. a. 3, 4; \emph{Sum. Th.}, p. I, q. 45, a. 5).\par Dissemos que, na producção das creaturas, a intelligencia de \textsc{Deus} \emph{dirige}, a sua vontade \emph{manda} e a sua potencia \emph{executa}. Para entendermos, de algum modo, o acto da intelligencia divina nessa obra, lembremo-nos da nobilissima theoria do Angelico Doutor. A intelligencia divina, diz S. Thomaz, conhece perfeitamente a essencia divina --- acto purissimo e oceano infinito de toda a perfeição, --- e conhece-a não só emquanto é tal essencia, mas tambem emquanto pode ser imitada e participada por infinitos modos, embora todos infinitamente inadequados. Conhecendo assim a essencia divina, a intelligencia forma innumeraveis conceitos, ou archetypos, de coisas possi\pgn{517}veis, cada um dos quaes imita, mais ou menos, a mesma essencia, e, em conformidade com estes archetypos, \textsc{Deus} produz, numa certa differença de tempo, tudo o que Elle quer produzir. E S. Agostinho já tinha ensinado que todas as creaturas existiam, desde a eternidade, na arte e na sapiencia da intelligencia divina, pela qual foram depois creadas: « \emph{Sapientia Dei, per quam facta sunt omnia, secundum artem continet omnia, antequam fabricet omnia} » (\emph{In Joan.}, c. I, n. 17; \emph{cf. Sum. Th.}, p. I, q. 15, a. 4).}.

\pg{515} \textbf{87. Deus creou todas as coisas.}

a) \emph{Crear} é tirar do \emph{nada}. Ora \textsc{Deus} tirou do \emph{nada} todas as coisas; porque todas as coisas são \emph{contingentes} na \emph{totalidade} da sua substancia, isto é, tanto em relação á \emph{materia} como em relação á \emph{forma}, e todo o ente \emph{contingente} é produzido, em ultima analyse, pelo ente \emph{necessario}. Logo \textsc{Deus} \emph{creou} todas as coisas.

b) Se o acto da Causa primeira, na producção das coisas, devesse presuppór alguma materia não-produzida, esta materia possuiria o ser \emph{por essencia}; porque o ente, que não possue o ser \emph{por participação}, possue-o \emph{por essencia}. Ora é impossivel \pg{516} \pg{517}  que a materia possua o ser \emph{por essencia}; porque o ente que possue o ser \emph{por essencia é acto purissimo}, emquanto a materia é \emph{mera potencia}. Logo o acto da Causa primeira, na producção das coisas, não presuppõe alguma materia não-produzida, e por isso é um acto \emph{creador}. Logo \textsc{Deus} \emph{creou} todas as coisas \footnote{Outros argumentos encontram-se em S. Thomaz (\emph{C. Gentes}, l. II, c. 15, --- \emph{De pot.}, q. a. 5). --- Contra esta these fazem-se muitas objecções, algumas das quaes foram expostas e resolvidas na \emph{Cosmologia}. Para a solução de outras, que se referem á Causa creadora, expomos os seguintes principios: --- 1º) \textsc{Deus} opera pelas ideas da sua intelligencia, e por isso, assim como não tem propriamente a idea da materia, separada da forma, mas só a idea do composto subsistente, assim tambem não crea a materia sem a forma, mas crea a materia formada, isto é, o composto subsistente. Todavia, pode dizer-se que \textsc{Deus} tem uma idea da materia, emquanto a materia imita, de algum modo, a essencia divina (\emph{De pot.}, q. 3, a. 1 ad 13). --- 2.º) O acto de \textsc{Deus}, na producção das coisas, passa do agente para o paciente --- emquanto, permanecendo na essencia divina, com a qual se identifica, produz fôra de \textsc{Deus} toda a substancia do paciente, --- mas não emquanto, derivando da essencia divina, é recebido no paciente; porque, como dissemos, o acto divino não se distingue da essencia, e esta produz o effeito sem produzir o acto. --- 3.º) A causa universal, que é \textsc{Deus}, produz o ente \emph{emquanto ente} --- não no sentido que produz a razão abstracta do ente communissimo e analogico (que não pode existir na realidade), --- mas no sentido que produz o ente na sua \emph{totalidade}, isto é, segundo \emph{todos} os elementos, de que o ente é constituido e que são a materia e a forma.\par Não nos demoramos neste assumpto, pois que foi sufficientemente tratado na \emph{Cosmologia}. --- Mas não podemos deixar de fazer uma breve critica do systema de A. \emph{Rosmini} acerca da \emph{creação}.\par Este escriptor, na sua \emph{Theosophia}, explica a \emph{creação} por um triplice acto da intelligencia divina: \emph{abstracção}, \emph{imaginação} e \emph{synthese}. --- A Pessôa divina do \emph{Pae}, contemplando o \emph{Verbo}, \emph{abstrahe} d'Elle o \emph{ser inicial}, em que vê contido virtualmente todo o ser finito. E' a primeira \emph{creação}, que dá origem á primeira creatura, isto é, ao \emph{ser inicial}. Este ser não tem uma existencia \emph{real}, mas somente \emph{logica}, e reside primeiramente na intelligencia de \textsc{Deus} e depois na das creaturas. --- A esta \emph{abstracção} segue-se a \emph{imaginação}, pela qual \textsc{Deus}-Pae concebe varios \emph{termos} reaes e limitados, de que se compõe a realidade do universo. --- Finalmente, realiza-se a \emph{synthese}, pela qual \textsc{Deus} une os dois elementos, que são o \emph{ser inicial} e os diversos \emph{termos} d'elle, isto é, os entes \emph{reaes finitos} (\emph{Teos.}, v. 1, l. 2, sez. 4. c. 6, a. 7).\par O systema de \emph{Rosmini} é falso pelas seguintes razões:\par 1.º) Depois de attribuir só á Pessoa do \emph{Pae} a creação, que (como toda a acção de \textsc{Deus} \emph{ad extra}) \emph{é commum} ás tres divinas Pessoas, \emph{Rosmini} reduz\pgn{518}a creação a um puro \emph{idealismo}. Na verdade, as tres operações, de que ella seria constituida e que são a \emph{abstracção}, a \emph{imaginação} e a \emph{synthese}, são productos da intelligencia e não dão origem a nenhuma realidade. Porquanto, o proprio \emph{Rosmini} ensina que a primeira creatura --- o \emph{ser inicial} --- não tem existencia \emph{real}, mas só \emph{logica}. Se não tem existencia \emph{real} o producto da \emph{abstracção}, nem a podem ter os productos da \emph{imaginação} e da \emph{synthese}. D'este modo, os entes creados não passam de entes \emph{ideaes}.\par 2.º) \emph{Rosmini} admitte como \emph{increado} o ser das creaturas. De facto, o ser das creaturas, na opinião d'esse escriptor, é constituido pelo \emph{ser inicial}, e este \emph{ser inicial} é unico, simplicissimo e eterno, porque é o proprio \emph{ser} do \emph{Verbo}, abstrahido somente da subsistencia. Ora esta opinião, além de conter o \emph{pantheismo}, destroe o conceito da \emph{creação}; porque a creação, como dissemos, tem por termo o proprio \emph{ser}, pois produz a \emph{totalidade} da substancia. Assim ensina S. Thomaz: « \emph{Creatio terminatur ad esse, tamquam ad proprium effectum} » (\emph{De pot.}, q. 3, a. 16 ad 21). E noutro logar: « \emph{Cum Deus sit ipsum esse per suam essentiam, oportet quod esse creatum sit proprius effectus ejus, sicut ignire est proprius effectus ignis} » (\emph{Sum. Th.}, p. I, q. 8, a. 1). --- Coerente com esse principio, \emph{Rosmini} diz que --- o ser das creaturas não é \emph{produzido}, mas é \emph{posto}, --- o \emph{ser das creaturas} não se distingue do \emph{ser divino}, senão por causa dos \emph{termos}, --- só \textsc{Deus} pode crear, porque o \emph{ser inicial}, que dá á creatura a forma de ente, é \emph{alguma coisa} de \textsc{Deus}, e só \textsc{Deus} pode dispór do que lhe pertence (\emph{Teos.}, v. 1, p. 369).}.

\pg{518} \textbf{88. Deus foi livre na creação do mundo.} --- \textsc{Deus} foi \emph{livre} na creação do mundo, se a creação não era exigida --- nem pela essencia divina, --- nem pela essencia do mundo. Ora a creação do mundo não era exigida --- nem pela essencia divina, porque esta é um oceano infinito de perfeições, e, para ser felicissima, não precisa de algum bem creado, --- nem pela essencia do mundo, porque o mundo é contingente, e todo o ente contingente é por si indifferente para a existencia e para a não-existencia. Logo \textsc{Deus} foi \emph{livre} na creação do mundo \footnote{Não insistimos neste ponto, que foi sufficientemente desenvolvido, tratando da \emph{vontade} de \textsc{Deus}. \textsc{Deus} é causa do mundo, --- não \emph{natural}, isto é, que não opera por necessidade da sua natureza, --- mas \emph{livre}, isto é, que opera pelo livre juizo da sua vontade. Por isso, não só não foi necessario que \textsc{Deus} creasse o mundo \emph{desde a eternidade}, mas nem foi necessario que o creasse \emph{no tempo}. Diz S. Thomaz: « \emph{Voluntas Dei est causa rerum. Sic ergo aliqua necesse est esse, sicut necesse est Deum velle illa: cum necessitas effectus ex necessitate causae dependeat. Ostensum est autem quod, absolute loquendo, non est necesse Deum velle aliquid, nisi seipsum. Non est ergo necessarium Deum velle quod mundus fuerit semper. Sed eatenus mundus est, quatenus Deus vult illum esse, cum esse mundi ex voluntate Dei dependeat, sicut ex sua causa. Non est \emph{igitur} necessarium mundum semper esse} » (\emph{Sum. Th.}, p. I, q. 46, a. 1). --- Lembramos que, quando se diz que \textsc{Deus} creou o mundo \emph{no tempo}, --- não entendemos que o tempo tenha preexistido ao mundo, ou que o tempo seja a medida de toda a duração, mesmo da duração divina (o que é absurdo), --- mas entendemos que \textsc{Deus}, immovel na sua eternidade, que abrange e excede toda a duração\pgn{519}temporanea, creou o mundo \emph{com o tempo}, emquanto, creado o mundo, começou o \emph{movimento}, de que o tempo é medida. O tempo, que se pensa ter existido antes do mundo, é \emph{imaginario}, não \emph{real} (\emph{Cf. De pot.}, q. 3, a. 17).}.

\pg{519} 89. \textbf{Deus creou o mundo para a sua gloria.} --- \textsc{Deus}, sendo infinitamente intelligente, devia propor-se um fim na creação do mundo. Este fim não pode ser senão á communicação e por isso a manifestação das perfeições divinas, --- communicação e manifestação, de que resulta a \emph{gloria} de \textsc{Deus}. Logo \textsc{Deus} creou o mundo para a sua \emph{gloria} \footnote{A \emph{gloria} é o \emph{conhecimento} claro da \emph{bondade} de um ser, a qual é digna de louvor. Por isso define-se: \emph{clara notitia cum laude}. D'onde se segue que a \emph{gloria} é maior ou menor, conforme maior ou menor é a \emph{bondade}, de que um ser é dotado, e conforme maior ou menor é o \emph{conhecimento}, que se tem d'essa bondade. A gloria das creaturas é imperfeita, porque a sua bondade é finita e participada. Só a gloria de \textsc{Deus} é infinitamente perfeita, porque só a sua bondade é essencial e infinita. --- Mas, como a bondade divina só por \textsc{Deus} pode ser conhecida adequadamente e por isso louvada como merece, só a gloria, que \textsc{Deus} recebe de si mesmo, é digna de \textsc{Deus}. A gloria, que a creatura pode e deve dar a \textsc{Deus}, é sempre infinitamente inferior ao merito; porque infinitamente deficiente é o conhecimento, que a creatura tem da bondade divina. --- Ainda que infinitamente inferiores ao merito, os nossos louvores dão ao bom \textsc{Deus} uma satisfação immensa e procuram á nossa alma uma felicidade eterna.\par Esta proposição --- que \textsc{Deus} creou o mundo para a \emph{sua gloria} --- não é contraria á outra proposição --- que \textsc{Deus} creou o mundo para communicar a sua bondade. Communicando a sua bondade, \textsc{Deus} entende a sua gloria. Como tambem, dando \emph{gloria} a \textsc{Deus}, as creaturas participam da felicidade divina. \textsc{Deus} não quiz separar do nosso bem a sua gloria; o nosso bem é a sua gloria, e a sua gloria é o nosso bem. Felizes as almas, que reconhecem os beneficios de \textsc{Deus} e Lhe tributam a homenagem da gratidão, fidelidade e da obediencia!}.

%%>>> art-45

\section*{Artigo II.\\ Conservação, sua divisão e necessidade.}
\addcontentsline{toc}{section}{Artigo II. --- Conservação, sua divisão e necessidade.}

90. \textbf{Conservação}. --- \emph{Conservação} é o acto da divina vontade, pelo qual as coisas, uma vez creadas, perseveram na existencia. --- Diz-se --- acto da divina \emph{vontade}; porque \textsc{Deus} não produz as operações \emph{transeuntes} por necessidade da sua \emph{natureza}, mas pelo \emph{livre} arbitrio da sua \emph{vontade} \footnote{S. Thomaz trata da \emph{conservação} das coisas creadas --- na \emph{Sum. Th.}, p. I, q. 104, a. 1, --- na \emph{C. Gentes}, 1. III, c. 65, --- \emph{De pot.}, q. 5, a 1.}.

91. \textbf{Conservação} positiva e negativa. --- A conservação \emph{positiva}, ou \emph{directa}, é a que suppõe e exige, da parte de \textsc{Deus}, um \pg{520}  continuo e real influxo, pelo qual as coisas, uma vez creadas, possam perseverar na existencia. A \emph{negativa}, ou \emph{indirecta}, é a que suppõe e exige que \textsc{Deus} --- não anniquile as coisas creadas, --- ou afaste d'ellas tudo o que as possa corromper \footnote{Todos os philosophos admittem que \textsc{Deus} conserva as coisas creadas; mas alguns, como \emph{Crousat}, \emph{Bayle}, \emph{Galluppi}, sustentam que essa conservação é apenas \emph{negativa}. \textsc{Deus}, na opinião d'elles, conserva o mundo, \emph{emquanto o não destroe}. Mas os escholasticos ensinam que \textsc{Deus} conserva o mundo de um modo \emph{positivo}, e que por isso as coisas não poderiam subsistir, nem por um instante, se \textsc{Deus} lhes subtrahisse o seu influxo e o seu apoio.}.

92. \textbf{Deus conserva todas as coisas de um modo positivo.}

a) O ser de todas as coisas, mesmo depois de creadas, não deixou de ser \emph{contingente}, isto é, indifferente ou indeterminado tanto para a existencia como para a não-existencia; porque esse ser é recebido e limitado na essencia, e não é a propria essencia, mas distingue-se realmente d'ella. Se o ser de todas as coisas, mesmo depois de creadas, é, por si, indifferente ou indeterminado tanto para a existencia como para a não-existencia, não tem em si a razão sufficiente da sua continuação na existencia e por isso precisa do influxo \emph{positivo}, ou \emph{directo}, de uma causa, que constantemente o tire da indifferença e o conserve. Esta causa, que assim conserva o ser de todas as coisas, é a mesma que o tirou do nada e que é \textsc{Deus}; porque o ser das coisas finitas depende sempre e exclusivamente d'Aquelle, que é o \emph{proprio Ser subsistente}. Logo \textsc{Deus} conserva todas as coisas de um modo \emph{positivo} \footnote{Este argumento costuma propor-se por differentes modos; mas funda-se sempre sobre o mesmo \emph{termo medio} --- o caracter de \emph{contingencia} e de \emph{participação}, que revelam todas as coisas creadas. --- Todo o effeito, \emph{emquanto tal}, depende sempre da causa, que o produziu, de modo que o effeito dura tanto, quanto dura o influxo da causa; porque a razão da existencia do effeito não se encontra no proprio effeito, mas na sua causa. Assim a vida do nosso corpo, derivando da alma, depende da propria alma e dura tanto, quanto dura a união da alma com o corpo. Ora o \emph{ser} de todas as coisas é um \emph{effeito} exclusivo de \textsc{Deus}. E' um \emph{effeito}, porque não é o ser de \textsc{Deus}: e por isso é um ser contingente, participado. Sendo \emph{contingente}, o \emph{ser creado} deve ser sempre determinado, para a existencia, pelo Ente \emph{necessario}; e, sendo \emph{participado}, não é a sua essencia e por isso não pode subsistir por si mesmo, mas deve subsistir por virtude do Ente, que é o proprio Ser subsistente. --- S. Thomaz compara a acção de \textsc{Deus} nas coisas creadas á acção do sol no ar. Toda a creatura, diz o Santo Doutor, refere-se a \textsc{Deus}, como o ar se refere ao sol, que o illumina. Assim como o sol é luminoso por sua natureza, e o ar se torna luminoso pela participação da \emph{luz} do sol, e não pela participação da \emph{natureza} do sol, assim tambem \textsc{Deus}, e só Elle, é o Ser pela sua essencia, porque a sua essencia é o seu ser, e toda a creatura\pgn{521} é ser por participação, porque a sua essencia não é o seu ser. Por isso, assim como o ar se torna luminoso ao apparecer do sol, e se cobre de trevas ao desapparecer d'elle, assim tambem a creatura existe emquanto \textsc{Deus} a conserva, e deixaria de existir e voltaria para o seu nada, se \textsc{Deus} lhe subtrahisse o influxo da sua potencia (\emph{Sum. Th.}, p. I, q. 104, a 1).}.

\pg{521} b) Se as coisas creadas se conservassem na existencia pela sua propria virtude, e não por um influxo positivo e directo do Creador, só poderiam ser \emph{anniquiladas}, isto é, reduzidas ao nada, por um acto \emph{positivo} da divina vontade. Ora isto é simplesmente absurdo; porque um acto \emph{positivo} de \textsc{Deus} não pode ter por termo o \emph{nada}. Logo o \emph{anniquilação} das creaturas, que é \emph{possivel}, só pode conceber-se por um acto \emph{negativo} de \textsc{Deus}, emquanto \textsc{Deus} lhes \emph{subtrahe} o influxo \emph{positivo} e \emph{directo}, que as conserva na existencia. Logo \textsc{Deus} conserva todas as coisas de um modo \emph{positivo} \footnote{Outros argumentos adduz S. Thomaz nos logares citados. --- Agora expomos e resolvemos as principaes objecções contra a nossa these.\par Obj. 1.ª --- \textsc{Deus} é mais poderoso do que todo o agente creado. Ora o agente creado communica ao effeito o poder de se conservar no ser, mesmo depois de cessado o acto productor. Logo, com maior razão, \textsc{Deus} pode conferir ás coisas creadas o poder de se conservarem no ser, sem um positivo e continuo influxo da sua infinita energia.\par Resposta. A razão, por que os effeitos do agente creado se conservam no ser, mesmo depois de cessar o acto productor, não se deve attribuir á perfeição e efficacia do mesmo agente, mas unicamente á potencia da Causa primeira, da qual o ser de todas as coisas é effeito proprio e especifico. Por quanto, o agente creado não pode operar senão transformando a materia, e por isso o seu acto presuppõe o sujeito constituido no seu ser. D'este modo, a conservação do effeito de um agente creado no proprio ser, longe de revelar a perfeição d'esse agente, revela, pelo contrario, a sua imperfeição, emquanto demonstra --- que o agente creado não pode, na sua operação, prescindir de um sujeito preexistente, --- e que é absolutamente necessaria a acção de \textsc{Deus}, que produza e conserve o ser das coisas. --- A difficuldade teria algum valor, se o agente creado produzisse por tal modo o effeito, que, ao cessar do acto productor, o effeito permanecesse no proprio ser, independentemente de toda a causa: ora isto é falso. (\emph{De pot.}, q. 5, a. 1 ad 4).\par Obj. 2.ª --- O agente creado produz a \emph{forma}, e por isso o \emph{ser} das coisas; porque o \emph{ser} deriva da \emph{forma}. Ora, acabado o acto do agente creado, o effeito persevera no ser. Logo tambem, acabado o acto do Creador, que é a causa do \emph{ser} das coisas, estas perseveram no ser.\par Resposta. O agente creado, como dissemos, não opera senão transformando a materia, e a sua efficacia extende-se unicamente ás coisas, que, de algum modo, dependem da materia ou estão unidas com a materia. Por isso, se a efficacia do agente creado chega a produzir a \emph{forma substancial}, esta forma não a produz independentemente de todo o sujeito, como se ella fosse subsistente, mas tira-a da potencia da materia, dependentemente da materia, e por isso o ser, que se produz e que é proprio do composto (embora\pgn{522}derive da forma para o composto), não depende, pelo menos exclusivamente, do agente creado, e pode permanecer e permanece, embora acabe a acção d'esse agente. Mas \textsc{Deus} é a causa total e proporcionada do \emph{ser} das coisas, e por isso esse ser não poderia continuar a existir, se \textsc{Deus} não lhe continuasse o influxo do seu poder.\par Obj. 3.ª --- Os substancias \emph{materiaes} são compostas de materia e de forma; ora a materia sustenta a forma, e a forma é o principio do ser. Do mesmo modo, as substancias \emph{immateriaes} são \emph{naturalmente indestructiveis}. Logo todas as substancias creadas continuam a existir pelos seus elementos \emph{intrinsecos}, e não pelo accão \emph{positiva} de \textsc{Deus}.\par Resposta. As substancias \emph{materiaes} são compostas de materia e de forma. Ninguem o contesta. Mas esses elementos, sendo \emph{contingentes} e \emph{participados}, não podem existir por si mesmos, por virtude propria; mas exigem necessariamente o auxilio de um agente extrinseco, de uma Causa superior, que os conserve. --- Diga-se o mesmo com relação ás substancias \emph{immateriaes}. Estas são \emph{naturalmente} indestructiveis, mas na supposição que \textsc{Deus} conserve os principios essenciaes, de que são compostas. Diz S. Thomaz: « Quamdiu principia essentialia rerum sunt, tamdiu res conservantur in esse; sed et ipsa rerum principia esse desinerent, divina actione cessante » (\emph{De pot.}, q. 5, a. 1 ad 3).\par Obj. 4.ª --- Todas as coisas creadas, sendo obras de \textsc{Deus}, são perfeitas. Ora não seriam perfeitas, se, para a sua conservação, precisassem do influxo de um agente exterior; porque perfeito é o ente, a que nada falta. Logo as coisas creadas não precisam do influxo divino para a sua conservação.\par Resposta. Todas as coisas creadas são perfeitas --- emquanto têm tudo o que convem á sua natureza, --- mas não emquanto são dotadas de \emph{todas} as perfeições e independentemente de \textsc{Deus}. Ora uma das perfeições, que convem á natureza das coisas creadas, é que sejam conservadas por \textsc{Deus}, para que não voltem para o nada, d'onde foram tiradas. Escreve S. Agostinho « Creatoris potentia, et Omnipotentis virtus causa subsistendi est omni creaturae; quae virtus ab eis, quae creata sunt regendis, si aliquando cessaret, simul illorum cessaret species, omnisque natura concideret » (\emph{De Gen. ad lit.}, l. IV, c. 12). Por isso S. Thomaz responde á difficuldade com estas palavras: « Creaturae Dei sunt perfectae in sua natura et in suo ordine; sed inter alia, quae ad eorum perfectionem requiruntur, hoc etiam est, quod a Deo contineantur in esse » (\emph{De pot.}, q. 5, a. 1 ad 1).}.

\pg{522} 93. \textbf{Deus conserva todas as coisas pela continuação do acto creador.} --- O ser das coisas finitas, em todos os instantes da sua duração, é o mesmo que era no primeiro instante, e por isso, para continuar na existencia, exige aquella mesma força, que foi necessaria para o seu começo. Ora, para o começo de todas as coisas, foi necessario o acto \emph{creador} de \textsc{Deus}. Logo, para que as coisas continuem na existencia, é necessaria a continuação do mesmo acto. Logo \textsc{Deus} conserva todas as coisas pela continuação do acto \emph{creador} \footnote{O acto, pelo qual \textsc{Deus} conserva as coisas, não é diverso do acto, com que as creou, mas é a continuação do proprio acto creador. Esta verdade é certissima, --- não só porque em \textsc{Deus} não ha distincção ou diversidade de actos --- mas tambem porque, mesmo na esphera das causas finitas, quando\pgn{523}um effeito é conservado pela mesma causa, que o produziu, o acto, que conserva o effeito, não é \emph{realmente} distincto do acto, que o produziu, mas é uma continuação do mesmo; assim o acto, com que o sol continua a illuminar o ar, não é \emph{realmente} distincto do acto, com que o illuminou no primeiro instante, mas é uma \emph{continuação} d'esse primeiro acto. Sendo uma continuação do acto \emph{creador}, o acto \emph{conservador} não dá ás creaturas um novo ser, distincto do ser, que receberam pelo acto \emph{creador}, mas conserva aquelle mesmo ser, que as creaturas receberam na creação. Diz S. Thomaz: « Conservatio rerum a Deo non est per aliquam novam actionem, sed per continuationem actionis, qua dat esse; quae quidem actio est sine motu et tempore; sicut etiam conservatio luminis in aëre est per continuatum influxum solis » (\emph{Sum. Th.}, p. I, q. 104, a. 1 ad 4). --- Dissemos que o acto \emph{conservador} não é \emph{realmente} distincto do acto \emph{creador}; porque \emph{logicamente} um é distincto do outro, emquanto o acto \emph{creador} refere-se ao \emph{primeiro} instante da existencia das coisas, e o \emph{conservador} aos instantes \emph{posteriores}, --- o acto \emph{creador} não suppõe a existencia das coisas, o que suppõe o acto \emph{conservador}. --- Com razão, pois, a \emph{conservação} chama-se uma creação \emph{continua}, emquanto a mesma virtude infinita, que tirou do \emph{nada} as coisas, é necessaria em cada momento para as conservar na existencia, --- e não emquanto \textsc{Deus} anniquile e crie continuamente as coisas, como pensavam \emph{Bayle}, \emph{Bertaud} e outros. Sob o influxo da acção divina, as coisas perseveram no ser, sem que sejam anniquiladas em cada instante, como, sob o influxo do sol, o ar continua a ser luminoso, sem que, em cada instante, se torne escuro.}.

\pg{523} 94. \textbf{Deus não anniquila alguma coisa creada}. --- As intenções de \textsc{Deus} relativamente ás coisas revelam-se na propria natureza das coisas. Ora a natureza das coisas mostra que \textsc{Deus} não anniquila nada do que Elle creou. Por quanto, as coisas creadas --- ou são \emph{immateriaes}, --- ou são \emph{materiaes}. Se são \emph{immateriaes}, são por si indestructiveis, porque não ha nellas a potencia para a não-existencia; se são \emph{materiaes}, são indestructiveis emquanto á \emph{materia}. porque esta é o sujeito de toda a geração e corrupção. Logo \textsc{Deus} não anniquila alguma coisa creada \footnote{\textsc{Deus} creou as coisas, emquanto livremente lhes communicou o ser, e conserva-as, emquanto influe nellas constantemente o ser. Por isso, assim como podia não communicar-lhes o ser, antes que existissem, assim tambem pode subtrahir-lhes o seu influxo, depois de existirem. Subtrahindo ás coisas creadas o seu influxo, \textsc{Deus} reduziria as coisas ao \emph{nada}; porque, como dissemos, a reducção das coisas ao \emph{nada} só pode ser effeito de uma acção \emph{negativa} de \textsc{Deus}, e não de uma acção \emph{positiva}, que tem por termo o \emph{ser} e não o \emph{nada} (\emph{Sum. Th.}, p. I, q. 104, a. 3). --- Mas \textsc{Deus}, embora \emph{possa} reduzir ao nada as coisas creadas, comtudo é certo que não as reduzirá; e esta vontade divina revela-se nos proprios caracteres das coisas. Diz S. Thomaz: « Quae facturus est Deus secundum ordinem naturalem, rebus inditum, considerari possunt ex ipsis rerum naturis » (\emph{Ib.}, a. 4). A potencia e a bondade de \textsc{Deus} revelam-se mais na conservação do que na anniquilação das coisas. --- O homem deve corresponder a esta bondade de \textsc{Deus}. Se \textsc{Deus} conserva em cada instante a nossa vida, é justo que em cada instante a nossa vida seja um hymno de reconhecimento pelo nosso Bemfeitor.}.

%%>>> art-46

\section*{Artigo III.\\ Concurso, sua natureza e extensão.}
\addcontentsline{toc}{section}{Artigo III. --- Concurso, sua natureza e extensão.}

\pg{524} 95. \textbf{Concurso}. --- \emph{Concurso} é o acto, pelo qual \textsc{Deus}, como Causa primeira, move ou auxilia as forças \emph{physicas} das causas segundas, e assim influe em todos os actos \emph{naturaes} d'essas causas. --- O \emph{concurso}, pois, distingue-se da \emph{conservação}. Esta refere-se á \emph{existencia} das creaturas; aquelle, ás \emph{operações} \footnote{Fallamos, como se vê, do concurso \emph{physico}, com que \textsc{DEus}, como Causa primeira e Auctor da natureza, move ou auxilia as forças das causas segundas, --- e não do concurso \emph{moral}, com que \textsc{DEus} pode determinar a causa segunda ao acto por meio de \emph{conselhos}, \emph{exhortações}, etc. Esse concurso \emph{physico} pertence á ordem \emph{natural}, refere-se ás operações \emph{naturaes} das causas segundas, e por isso distingue-se do auxilio da \emph{graça}, com que \textsc{DEus}, emquanto Auctor da ordem \emph{sobrenatural}, influe nas operações da vida \emph{sobrenatural}. --- O \emph{concurso} costuma chamar-se tambem \emph{influxo}, moção, acção de \textsc{Deus} nas obras das creaturas. Devemos, comtudo, advertir que o vocabulo « \emph{concurso} » aqui se toma em sentido \emph{lato}; porque, no seu sentido \emph{rigoroso}, \emph{concurso} significa a acção de dois agentes, cada um dos quaes influe no effeito, mas de uma maneira \emph{parcial} (e por isso esse effeito não depende \emph{inteiramente} nem de um nem de outro), e portanto não pode significar o \emph{influxo} de \textsc{DEus} nas operações das causas segunda, visto que o effeito d'essas operações attribue-se \emph{inteiramente} a \textsc{DEus} e \emph{inteiramente} á causa segunda, ainda que de um modo \emph{diverso}, como veremos.}.

96. \textbf{Especies do concurso}. --- O \emph{concurso} pode ser \emph{immediato} ou \emph{mediato}, \emph{previo} ou \emph{simultaneo}.

\textbf{a) Concurso immediato ou mediato.} --- E' \emph{immediato}, se \textsc{Deus}, como Causa primeira e universal, opera juntamente com as creaturas, causas segundas e particulares, e com ellas realmente produz o mesmo effeito. --- E' \emph{mediato}, se \textsc{Deus} influe no effeito, produzido pelas creaturas, só emquanto \emph{conserva} nellas as faculdades operativas, de que as dotou.

\textbf{b) Concurso previo ou simultaneo.} --- E' \emph{previo}, se \textsc{Deus} move e applica as creaturas á operação, de modo que o concurso de \textsc{Deus} \emph{prevém} ou \emph{precede} a determinação das creaturas. --- E' \emph{simultaneo}, se \textsc{Deus} não prevém, mas limita-se a \emph{acompanhar} a operação das creaturas \footnote{São as especies \emph{principaes}. Porquanto o \emph{concurso} pode ser tambem \emph{especial} ou \emph{geral}, \emph{intrinseca} ou \emph{extrinsecamente efficaz}. --- E' \emph{especial}, se \textsc{Deus} para cada operação da creatura concede um determinado auxilio, --- é \emph{geral}, se \textsc{Deus} concede o mesmo auxilio para todas as operações. --- E' \emph{intrinsecamente efficaz}, se o auxilio de \textsc{Deus} produz o seu effeito por si, independentemente\pgn{525} de toda a determinação da creatura, --- é \emph{extrinsecamente efficaz}, se o auxilio de \textsc{Deus} não produz o seu effeito por si, mas só dependentemente da previa determinação da creatura.}.

\pg{525} 97. \textbf{Deus concorre immediatamente em todas as operações das creaturas.}

a) Todas as operações das creaturas produzem um effeito; e este effeito contém o \emph{ser}, \emph{determinado e particularisado por algumas perfeições}. Ora, se as \emph{perfeições} que determinam e particularisam o \emph{ser}, sendo coisas \emph{particulares}, constituem o effeito \emph{proprio} das causas \emph{particulares}, que são as creaturas, o \emph{ser}, sendo uma coisa \emph{universal}, é o effeito \emph{proprio} da Causa \emph{universal}, que é \textsc{Deus}; porque a ordem dos effeitos deve corresponder á ordem das causas. Mas uma coisa, que está contida no effeito de uma causa \emph{particular}, não poderia ser produzida pela causa \emph{universal}, se a mesma causa \emph{particular} não operasse pela virtude ou energia da causa \emph{universal}, isto é, se a causa \emph{universal} não concorresse \emph{immediatamente} na operação da causa \emph{particular}. Logo \textsc{Deus} concorre \emph{immediatamente} em todas as operações das creaturas \footnote{O argumento é de S. Thomaz. Eis as suas palavras: « Secundum ordinem causarum est ordo effectuum. Primum autem in omnibus effectibus est \emph{esse}; nam omnia alia sunt determinationes ipsius. Igitur \emph{esse} est proprius effectus primi agentis, et omnia alia agunt ipsum, inquantum agunt in virtute primi agentis; secunda autem agentia, quae sunt \emph{quasi particulantia et determinantia actionem primi agentis}, agunt, sicut proprios effectus, alias perfectiones, quae determinant esse » (\emph{C. Gentes}, l. III, c. 66). --- Este argumento funda-se naquella verdade: um effeito \emph{commum} e \emph{universal} deve derivar de uma causa \emph{commum} e \emph{universal}, como do seu principio adequado e proporcionado. Por quanto, se uma causa \emph{creada} e \emph{particular} podesse produzir um effeito \emph{commum} a varios entes e á propria causa, o effeito seria mais nobre do que a causa, e esta seria causa de si mesma. E' por isso que uma causa, subordinada a uma certa especie, não pode por si produzir a propria especie. Posto este fundamento solidissimo, conclue-se facilmente --- que o \emph{ser}, que está contido em todos os effeitos, provenientes das causas particulares, deve attribuir-se á virtude ou energia da causa \emph{universal}, --- que as causas particulares operam pela virtude ou energia da propria causa \emph{universal}, --- e que por isso \textsc{Deus} concorre \emph{immediatamente} em todas as operações das creaturas. E' o que diz S. Thomaz em poucas palavras: « Universaliores effectus oportet reducere in universaliores et priores causas. Inter omnes autem effectus universalissimum est ipsum esse. Unde oportet quod sit proprius effectus primae et universalissimae causae, quae est Deus » (\emph{Sum. Th.}, p. I, q. 45, a. 3).}.

b) Todas as creaturas, quando operam, passam da potencia para o acto e adquirem uma nova perfeição. Ora nenhum ente pode passar da potencia para o acto, nem adquirir uma

\pg{526} nova perfeição, sem o concurso de outro ente, que esteja em acto e possua essa perfeição. Logo as creaturas, quando operam, precisam do concurso de um ente, que esteja em acto, para que possa determinar a potencialidade d'ellas, e que possua immensas perfeições, para que as possa communicar aos que não as têm. Esse ente, em ultima analyse, é \textsc{Deus}; porque só Elle é superior a todas as creaturas, e é acto purissimo e infinitamente perfeito. Logo \textsc{Deus} concorre \emph{immediatamente} em todas as operações das creaturas \footnote{Estes e outros argumentos encontram-se em S. Thomaz (\emph{Sum}. \emph{Th}., p. I, q. 105, a. 5, --- \emph{c}. \emph{Gentes}, l. III, c. 66, 67, --- \emph{De pot}., q. 3, a. 7). Ainda que apresentados sob varios aspectos, todos esses argumentos reduzem-se a este: \textsc{Deus} é o ente \emph{por essencia} (cuja essencia é o seu ser) e por isso de \textsc{Deus} deve derivar todo o ente \emph{participado}. Ora toda a operação da creatura é uma realidade \emph{participada}, é um novo e limitado grau de entidade. Logo toda a operação da creatura deve derivar de \textsc{Deus}, como uma imperfeita participação da realidade infinita. Mas não poderia essa operação derivar de \textsc{Deus}, se \textsc{Deus} não concorresse \emph{immediatamente} na operação da creatura. Logo \textsc{Deus} concorre \emph{immediatamente} em toda a operação da creatura. --- E', pois, necessario admittir, contra \emph{Durando} e outros poucos, que \textsc{Deus}, além do concurso \emph{mediato}, emquanto concede ás creaturas e lhes conserva a faculdade ou energia de operar, presta tambem um concurso \emph{immediato}, emquanto applica essa faculdade ou energia á operação e emquanto toda a energia creada opera pela propria energia increada. --- De resto, o concurso \emph{immediato} é um corollario da conservação \emph{positiva}. Por quanto, se todo o ente finito e participado depende, na \emph{continuação} da sua existencia, de um concurso \emph{positivo} de \textsc{Deus}, tambem deve depender do mesmo concurso \emph{positivo} ou \emph{immediato} de \textsc{Deus} no seu \emph{começo}; porque, sendo a \emph{continuação} na existencia uma \emph{extensão do proprio começo}, a energia que \emph{conserva} o ser das coisas é a mesma que \emph{produziu}. Logo \textsc{Deus} produz todo o ente finito e participado. E como o effeito da operação das creaturas é um ente finito, deve concluir-se que \textsc{Deus} concorre em todas as operações das creaturas.\par Demonstrada a \emph{existencia} do \emph{concurso immediato} de \textsc{Deus} em todas as operações das creaturas, devemos fazer algumas considerações, que esclarecem a \emph{natureza} e a \emph{condição} d'esse concurso.\par I. Não ha nada na operação e no effeito da creatura, que não derive total e immediatamente de \textsc{Deus}. Na verdade, se houvesse nessa operação e nesse effeito alguma coisa, que não derivasse \emph{immediatamente} de \textsc{Deus}, o concurso de \textsc{Deus} seria \emph{mediato}. --- D'onde se segue que a mesma unica e indivisivel operação da creatura procede \emph{inteiramente} e \emph{immediatamente} de \textsc{Deus}. Assim como uma e a mesma é a operação da causa instrumental e da causa principal, e não pode determinar-se uma parte da operação, que seja produzida unicamente pela causa instrumental sem o influxo da principal; assim tambem toda a operação das creaturas deriva da energia natural d'ellas e da energia da Causa primeira, e não ha uma parte só d'essa operação que possa attribuir-se exclusivamente ás causas segundas. \textsc{Deus}, pois, e a creatura são causas \emph{totaes}, cada uma na sua ordem. Diz S. Thomaz: « Patet quod non sic idem effectus causae naturali et divinae virtuti\pgn{527}attribuitur, quasi \emph{partim} a Deo et \emph{partim} a naturali agente fiat, sed \emph{totus} ab \emph{utroque} secundum \emph{alium modum}; sicut idem effectus totus attribuitur instrumento, et principali agenti etiam totus » (\emph{C}. \emph{Gentes}, l. III, c. 70). E noutro logar: « In operatione, qua Deus operatur movendo naturam, non operatur natura; sed ipsa naturae operatio est etiam operatio virtutis divinae, sicut operatio instrumenti est per virtutem agentis principalis. Nec impeditur quin natura et Deus ad idem operentur, propter ordinem qui est inter Deum et naturam » (\emph{De pot}., q. 3, a. 7 ad 3). E noutro: « Una actio non procedit a duobus agentibus unius ordinis; sed nihil prohibet quin una et eadem actio procedat a primo et secundo agente » (\emph{Sum}. \emph{Th}., p. I, q. 105, a. 5 ad 2).\par II. A acção das creaturas é necessaria --- não pela insufficiencia da energia de \textsc{Deus}, --- mas pela immensidade da sua bondade. Por quanto, \textsc{Deus} teria podido produzir, só por si, o effeito \emph{todo}, não só emquanto ao elemento universal, que é o \emph{ser}, mas tambem emquanto aos elementos particulares, que \emph{determinam} o ser. Todavia, sendo infinitamente bom, quiz communicar ás creaturas a propria semelhança, pela qual não só fossem bôas, mas podessem, operando, communicar aos outros a bondade. Mas, como a creatura não pode ter uma perfeita semelhança com o Creador, nem a energia participada egualar a energia infinita, segue-se --- que a creatura, operando, não pode attingir o effeito \emph{universal}, que é proprio da energia infinita e que é o \emph{ser}, --- e que por isso, na producção de um effeito, que contem esse ser, a mesma creatura precisa do concurso \emph{immediato} de \textsc{Deus}, isto é, não pode ella operar senão pela energia mesma de \textsc{Deus}. Diz S. Thomaz: « Neque est superfluum, si Deus per se ipsum potest omnes effectus naturales producere, quod per quasdam alias causas producantur. Non enim hoc est ex insufficientia divinae virtutis, sed ex immensitate bonitatis ipsius, per quam suam similitudinem rebus communicare voluit, non solum quantum ad hoc quod essent, sed etiam quantum ad hoc quod aliorum causae essent. His enim duobus modis creaturae communiter omnes divinam similitudinem consequuntur... Per hoc etiam decor ordinis in rebus creatis apparet » (\emph{C}. \emph{Gentes}, l. III, c. 70)... « Creatura habet aliquam Dei similitudinem participando bonitatem ipsius, in quantum est et agit, non tamen ita quod per similitudinis perfectionem ad aequalitatem perveniat; et ideo sicut imperfectum indiget perfecto, ita virtus naturae in agendo indiget operatione divina » (\emph{De pot}., q. 3, a. 7 ad 10).\par III. Toda a causa creada, \emph{ainda que seja completa na sua ordem}, comtudo, absolutamente fallando, é por si incompleta e insufficiente para a producção de todo e qualquer effeito. Na verdade, se a causa segunda não pode operar sem o concurso \emph{immediato} da Causa primeira, é claro que ella é incompleta e insufficiente com relação ao effeito; visto que uma causa absolutamente completa e adequada pode só por si produzir o effeito. Portanto a operação da creatura, que é tambem a operação de \textsc{Deus}, deve conceber-se como uma entidade, dotada de uma duplice relação de dependencia, --- uma com a propria causa creada, e outra com a Causa primeira e universal, porque de uma e de outra depende de um modo indivisivel. --- Dissemos que a causa \emph{creada pode ser completa} na sua \emph{ordem}; porque, entre essas causas,\pgn{528}ha umas que podem operar sem o auxilio de outra causa creada e por isso são completas na sua ordem, ao passo outras, para operarem, precisam do auxilio de outra causa e por isso são incompletas e parciaes mesmo na ordem de causas segundas. --- Advirta-se que essa \emph{insuficiencia} das causas segundas não significa que as suas faculdades operativas sejam \emph{imperfeitas}. Essas faculdades são perfeitas, muito perfeitas, mas \emph{na sua ordem}, e a ordem d'ellas exige que os seus actos sejam exercidos dependentemente da Causa primeira. Assim como \textsc{Deus}, quando produz as coisas, não pode fazer com que as coisas existam independentemente d'Elle, assim tambem, quando concede ás creaturas as faculdades operativas, não pode fazer com que essas faculdades, embora perfeitas, se exerçem independentemente d'Elle; porque \textsc{Deus} não pode fazer com que o ser participado seja ser por essencia. (\emph{De pot}., q. 3, a. 7 ad 5).\par IV. O effeito depende mais da Causa primeira do que da causa segunda. Na verdade, a causa segunda, sendo effeito da Causa primeira, recebe da Causa primeira a propria substancia. Recebendo da Causa primeira a substancia, a causa segunda recebe d'ella tambem a potencia ou energia operativa; porque a potencia é uma consequencia da substancia. Logo, se a causa segunda produz, pela sua potencia, um determinado effeito, esta causalidade deriva da Causa primeira. Portanto, a razão de causa, com relação a um certo effeito, compete mais á Causa primeira, do que á causa segunda, e por isso o effeito depende mais d'aquella do que d'esta. (S. Th., \emph{De causis}, l. I, par. 6). --- Além d'isso, emquanto a energia da creatura não pode extender-se até á producção da materia ou do sujeito, porque a sua actividade presuppõe sempre a materia ou o sujeito, a energia de \textsc{Deus} extende-se até á materia, creando-a, e assim o effeito procede de \textsc{Deus} segundo \emph{toda} a sua physica entidade. (\emph{De pot}., q. 3, a. 4. 7). --- D'ahi se segue que toda a operação da causa segunda se funda necessariamente na operação da Causa primeira, não só porque a creatura não pode operar sem o concurso divino, mas tambem porque toda a operação da creatura presuppõe o sujeito ou a materia, creada por \textsc{Deus}. (\emph{De pot}., q. 3, a. 4).\par V. As causas segundas são para a Causa primeira, como o instrumento é para o agente principal. E' a opinião de S. Thomaz e de muitos outros. Por quanto, assim como o instrumento depende, na sua operação, do agente principal, nem pode operar senão pela força do mesmo agente, de modo que é uma causa essencialmente deficiente e imperfeita com relação ao effeito; assim tambem as causas segundas não podem operar sem o concurso da Causa primeira. --- Assim como una é a acção do agente principal e do instrumento e uno é o effeito; assim tambem una é acção e uno é o effeito da Causa primeira e da causa segunda. --- Assim como o instrumento não produz o effeito, se não fôr applicado e movido pelo agente principal; assim tambem a causa segunda não produz o effeito, se não fôr applicada e movida pela Causa primeira. --- Assim como o agente principal se serve de instrumentos diversos, conforme a variedade dos actos e dos effeitos; assim tambem \textsc{Deus} se serve de muitas e diversas causas segundas para a producção de muitos e diversos effeitos. --- Todavia, as creaturas são \emph{instrumentos} --- não\pgn{529}num sentido \emph{rigoroso}, mas só num sentido \emph{analogico} ou \emph{metaphorico}. Por quanto, o instrumento não torna o effeito semelhante a si mesmo, mas ao agente principal, porque o instrumento não opera por virtude ou energia propria, mas por virtude do agente principal, que, operando pela propria virtude e forma, transfunde no effeito essa forma, e portanto a propria semelhança; ao passo que as causas segundas, embora operem sob o influxo de \textsc{Deus}, tornam o effeito semelhante a si mesmas, e não a \textsc{Deus}, porque operam pela propria forma. Além d'isso, o instrumento opera por virtude do agente principal, sem receber d'este um principio \emph{intrinseco} do seu movimento; ao passo que as causas segundas são dotadas de um principio \emph{intrinseco} do seu movimento, e algumas são por tal maneira movidas por \textsc{Deus}, que conservam o dominio dos proprios actos, a liberdade dos seus movimentos. (\emph{De verit}. q. 24, a. 1 ad 5).}.

\pg{527} 98. \textbf{Deus concorre immediatamente} em \textbf{todas as operações da vontade humana}.

a) A volição actual é uma realidade mais nobre e mais perfeita do que a propria faculdade volitiva; porque o acto é \pg{528}  mais perfeito do que a potencia. Ora nenhuma realidade, nenhuma perfeição pode encontrar-se nas creaturas, que não derive de \textsc{Deus}, fonte de toda a perfeição e causa universal de toda a realidade participada. Logo \textsc{Deus} é a causa --- não \pg{529}  só da vontade, que é o principio proximo do acto volitivo, --- mas tambem do acto volitivo, e por isso concorre \emph{immediatamente} em todas as operações da vontade humana \footnote{E' necessario distinguir, sempre \emph{da parte do effeito}, uma acção divina que tem por termo a \emph{nossa vontade}, e uma outra acção que tem por termo o \emph{acto volitivo}. A primeira por tal modo produz e conserva a nossa faculdade volitiva, que a creatura não coopera com a acção divina, que lhe confere tal faculdade. A segunda, pelo contrario, produz por tal modo na nossa vontade o acto volitivo, que a mesma vontade coopera com a acção divina como verdadeira causa, embora \emph{particular}, que não pode operar, senão em virtude da causa \emph{universal}. (\emph{De pot.}, q. 3, a. 7. ad 7).}.

b) \textsc{Deus} por tal modo confere a todos os agentes finitos, quer naturaes quer voluntarios, a faculdade ou potencia de operar, que nenhum d'elles pode servir-se d'essa faculdade, se não é determinado ou estimulado pela virtude ou energia do proprio Creador; porque todo o agente potencial e por participação, como é o agente finito, não pode operar senão pela virtude ou energia do agente, que é acto e opera pela sua essencia. Ora o ente, cuja virtude ou energia determina o acto de um agente, quer natural quer voluntario, é causa não só da faculdade operativa, mas tambem do proprio acto d'esse agente. Logo \textsc{Deus} não só é causa da nossa vontade, mas tambem do nosso acto volitivo, e por isso concorre \emph{immediatamente} em todas as operações da vontade humana \footnote{E' o argumento de S. Thomaz, na Sum. \emph{Contra Gentes}, l. III, c. 89, n. 4. --- Como nas operações de todas as creaturas, \textsc{Deus} influe nas operações da vontade humana por differentes modos. Primeiramente, \textsc{Deus}, emquanto primeiro motor immovel e causa universal de todas as coisas, move a nossa vontade --- \emph{necessariamente} para o bem \emph{universal}, como para o objecto proporcionado d'essa faculdade e \emph{contingentemente} para os bens \emph{particulares}, que não têm um nexo necessario com o bem universal; porque a inclinação natural de uma faculdade para um termo não pode provir senão da Causa, de que deriva essa faculdade (\emph{Sum. Th.}, 1-2, q. 9, a. 6). \textsc{Deus}, portanto, move\pgn{530}a nossa vontade, emquanto nos deu uma certa \emph{forma}, de que derivam certas faculdades e tendencias. --- Em segundo logar, \textsc{Deus} move a nossa vontade emquanto ao \emph{objecto}. Com effeito, a vontade é movida pelo \emph{bem}, e de um modo \emph{necessario} ou \emph{contingente}, conforme o bem é \emph{universal} ou \emph{particular}. Ora \textsc{Deus} é o bem \emph{universal}, e todo o bem \emph{particular} não pode mover a vontade senão emquanto é uma participação da bondade divina, a qual é o ultimo fim do homem, como é o seu primeiro principio. \textsc{Deus}, portanto, move a vontade humana emquanto esta tende para o objecto, que é o \emph{bem}. (\emph{De malo}, q. 6, a. un, --- \emph{Sum. Th.}, p. I, q. 105, a. 4). --- Finalmente, \textsc{Deus} move a nossa vontade, emquanto influe para que a vontade \emph{queira actualmente}, e assim é \emph{principio efficiente} do \emph{acto volitivo}. E, como \textsc{Deus} não corrompe mas conserva a natureza creada por Elle, segue-se que a acção divina move \emph{diversamante} para a operação os agentes \emph{naturaes} e os agentes \emph{livres}. Move os agentes \emph{naturaes} emquanto lhes communica a energia necessaria para actuarem nos outros seres; move os agentes \emph{livres} emquanto lhes communica a energia necessaria para se determinarem ao movimento e para se moverem de modo, que sejam senhores dos seus actos. Por isso, o influxo divino, embora, considerado no seu \emph{principio}, seja um e o mesmo (porque a acção de \textsc{Deus} é a sua essencia), --- comtudo, considerado no seu \emph{termo}, que é a creatura, é multiplice e diverso, emquanto \textsc{Deus} influe nas causas segundas por um modo proporcionado á natureza d'ellas. Diz S. Thomaz: « Ad providentiam divinam non pertinet naturam rerum corrumpere, sed servare. Unde omnia movet \emph{secundum eorum conditionem}; ita quod ex \emph{causis necessariis} per motionem divinam sequantur \emph{effectus ex necessitate}, ex \emph{causis} autem \emph{contingentibus} sequantur \emph{effectus contingentes}. Quia igitur voluntas est activum principium \emph{non determinatum ad unum}, sed indifferenter se habens ad multa; sic Deus ipsam movet, quod non ex necessitate ad unum determinat, sed remanet motus ejus \emph{contingens} et \emph{non necessarius}, nisi in his ad quae \emph{naturaliter} movetur » (\emph{Sum. Th.}, 1-2, q. 10, a. 4). --- E' claro que \textsc{Deus} concorre \emph{immediatamente} nas operações da nossa vontade emquanto é o principio \emph{efficiente} do acto volitivo.\par Para a intelligencia d'esta verdade e para a solução das objecções, devemos fazer as advertencias seguintes:\par I. A nossa vontade não pode ser coagida; já o provámos. O acto elicito da vontade escapa a todo o movimento \emph{violento}. O movimento de um ente é \emph{violento}, quando esse ente é determinado ao movimento por um principio \emph{extrinseco}, que não só não communica ao ente nenhuma inclinação\pgn{531}\emph{natural}, mas contraria essa inclinação. \textsc{Deus}, porém, sendo o Auctor de todas as naturezas, confere a todas a inclinação conveniente e move-as por tal modo, que o movimento está em perfeita harmonia com as inclinações d'ellas. Por isso o movimento, por que \textsc{Deus} actua nos entes creados, não é violento, mas é \emph{natural}, isto é, proporcionado á natureza de cada um. Diz S. Thomaz: « Illud quod movetur ab altero dicitur cogi, si moveatur contra inclinationem propriam; sed si moveatur ab alio, \emph{quod ipsi dat propriam inclinationem}, non dicitur cogi... Sic igitur Deus, movendo voluntatem, non cogit ipsam; quia dat ei ejus propriam inclinationem » (\emph{Sum. Th.}, p. I, q. 105, a. 4 ad 1). --- Advertimos que o movimento, que Deus imprime na nossa vontade, é e diz-se \emph{natural}, --- não emquanto a vontade seja por isso obrigada a operar \emph{necessariamente}, --- mas emquanto esse movimento é proporcionado á \emph{natureza} da vontade, a qual exige que \textsc{Deus} a mova e a mova de modo que lhe conserve a liberdade dos seus actos.\par II. O movimento de um ente vivo deriva de um principio intrinseco, e, se esse ente é dotado de vontade, é indeterminado, emquanto aos actos voluntarios, tanto para uma coisa como para outra. Mas, se o ente vivo é causa \emph{segunda}, isto é, se é vivo por participação e não por essencia, então não só não exclue, mas exige a acção de um extrinseco principio superior, que é Deus; e como recebeu de \textsc{Deus}, --- Causa primeira e universal, vida por essencia e acto purissimo, --- o ser e a potencia operativa, assim tambem deve receber \emph{actualmente} d'Elle a energia necessaria para realizar a sua operação. Portanto o facto de um ser vivo se mover a si mesmo e, se é livre, de se determinar a si mesmo para uma coisa ou para outra, não exclue o influxo da Causa primeira, que é o principio universal de toda a perfeição e entidade participada, e que faz com que o ser vivo e livre se mova a si mesmo e se mova livremente. Diz S. Thomaz: « Cum aliquid movet seipsum, non excluditur, quin ab alio moveatur, \emph{a quo habet hoc ipsum, quod seipsum movet}; et sic non repugnat libertati, quod Deus est causa actus liberi arbitrii » (\emph{De malo}, q. 3, a. 2 ad 4).\par III. Quando um ente é movido por outro e não se move de modo algum por si mesmo, então aquelle movimento não é imputavel ao mesmo ente; porque, para um acto ser imputavel, é preciso que o agente tenha o poder de se determinar a si mesmo, isto é, é preciso que seja livre. Pelo contrario, se o ente, que é movido por \textsc{Deus}, é livre, então o seu acto é imputavel. --- Aqui advertimos que a Causa primeira e universal communica ás causas segundas e particulares, dotadas de liberdade, toda a perfeição, que se encontra no acto, e não a deformidade ou imperfeição; de modo que toda a perfeição do acto deve attribuir-se á Causa primeira como ao agente primario, que influe no effeito mais efficazmente do que o agente secundario, e toda a imperfeição ou todo o defeito deve attribuir-se á causa segunda, que pode faltar e desviar-se do recto caminho. Diz S. Thomaz: « Quia causa\pgn{532}prima magis influit in effectum quam secunda, ideo quidquid perfectionis est in effectu, \emph{principaliter} reducitur ad primam causam; quod autem est de defectu, reducendum est in causam secundam, quae non ita efficaciter operatur, sicut causa prima » (\emph{De pot.}, q. 3, a. 7 ad 15).\par IV. A Causa primeira e universal, assim como produziu todas as coisas pela intelligencia e pela vontade, assim tambem pela intelligencia e pela vontade as governa e lhes concede tudo o que lhes é necessario para operarem actualmente conforme a natureza e a condição de cada uma. Por isso \textsc{Deus} não só não impõe a necessidade á nossa vontade, que é um principio activo, indifferente para muitas coisas, mas actua por tal modo nesta faculdade, que ella possa determinar-se mais para uma do que para outra coisa. A moção divina accommoda-se á indole do agente creado. Diz S. Thomaz: « Voluntas divina non solum se extendit ut aliquid fiat per rem quam movet, sed ut etiam eo modo fiat, quo congruit naturae ipsius. Et ideo magis repugnaret divinae motioni, si voluntas ex necessitate moveretur, quod suae naturae non competit, quam si moveretur libere, prout competit suae naturae » (\emph{Sum. Th.}, 1-2, q. 10, a. 4 ad 1).\par V. A vontade humana tem o dominio dos seus actos. Esse dominio --- se exclue todo o influxo de uma causa exterior, que determine necessariamente a vontade a praticar um certo acto, --- não exclue o influxo da Causa primeira, que não só não determina necessariamente o acto, mas faz com que o proprio acto seja livre. Diz S. Thomaz: « Voluntas dicitur habere dominium sui actus non per exclusionem causae primae, sed quia causa prima non ita agit in voluntate, ut eam de necessitate ad unum determinet, sicut determinat naturam; et ideo determinatio actus relinquitur in potestate rationis et voluntatis » (\emph{De pot.}, q. 3, a. 7 ad 13; --- Cf. c. \emph{Gentes}, I, 68).}.

\pg{530} 99. O \textbf{concurso de Deus} é \textbf{tambem previo ás operações das creaturas}.

a) As creaturas, como dissemos, não podem por si mesmas passar da potencia para o acto, mas devem ser reduzidas ao acto por \textsc{Deus}, emquanto \textsc{Deus} influe na potencialidade das causas segundas e as tira da indifferença. Ora este influxo não pode deixar de ser previo ou anterior, na ordem da natureza e da causalidade, ás operações das creaturas; porque o acto do motor deve ser previo ao acto do movel. Logo o concurso de \textsc{Deus} é tambem \emph{previo} ás operações das creaturas.

b) Se o concurso de \textsc{Deus} nas operações das creaturas é exclusivamente \emph{simultaneo}, e não é tambem \emph{previo}, a \emph{determinação} d'essas operações é independente de \textsc{Deus}, e por isso não é \textsc{Deus} quem determina o acto da creatura, mas, pelo contrario, é a creatura que determina o concurso de Deus. Ora estas consequencias são falsas; porque não pode haver, no \pg{531} \pg{532}  mundo, a mais simples entidade, que escape á universal causalidade de \textsc{Deus}, e porque a creatura, não podendo produzir acto algum sem o concurso de \textsc{Deus}, não pode determinar esse concurso. Logo o concurso de \textsc{Deus} é tambem \emph{previo} ás operações das creaturas \footnote{O concurso de \textsc{Deus} nas operações das creaturas pode ser, como dissemos, \emph{previo} ou \emph{simultaneo}. --- E' \emph{previo}, quando precede, pelo menos na ordem da natureza e da causalidade, a operação do agente, ao qual se presta, e por isso é recebido na potencialidade, ou faculdade, do proprio agente, e a determina para o acto. Tal é concurso, ou influxo, do esculptor, quando applica o cinzel e o move para reproduzir no marmore uma determinada imagem. --- E' \emph{simultaneo}, quando não precede a operação do agente nem é recebido na potencialidade d'este, mas é apenas collateral, parallelo e extrinseco á propria operação do agente, e é recebido no \emph{effeito}, que, d'este modo, se attribue \emph{parcialmente} a um e outro agente. Tal é o concurso, que um cavallo dá a outro na tracção de um carro; porque um cavallo não influe na causalidade do outro, e o movimento do carro não se attribue \emph{totalmente}, mas só \emph{parcialmente} a um e ao outro. --- Advertimos que o concurso \emph{previo} não exclue o \emph{simultaneo}; este pode ser uma \emph{continuação} d'aquelle. Assim, o esculptor applica o cinzel ao marmore, continua a propria acção previa e primitiva com a acção do instrumento; e, emquanto o cinzel produz com a sua acção o effeito proprio, que é \emph{cortar}, o esculptor tambem produz, por\pgn{533}meio da acção do cinzel, o effeito proprio, que é \emph{reproduzir a imagem}; e, d'este modo, o esculptor e o cinzel são causa do mesmo effeito, embora numa ordem differente.\par Dada esta explicação, perguntamos: Deverá admittir-se o concurso \emph{previo} de \textsc{Deus} nas operações das creaturas, ou bastará o concurso exclusivamente \emph{simultaneo}? --- A esta pergunta têm respondido duas celebres escholas: a dos \emph{thomistas} e a dos \emph{molinistas}, sustentando aquella o concurso \emph{previo}, e esta o exclusivamente \emph{simultaneo}. --- Não temos a presumpção de dar uma solução definitiva a esta controversia, que tem apaixonado tantos sabios, sobretudo nos seculos XVI e XVII. Limitamo-nos a expor, com muita serenidade e imparcialidade, as theorias d'essas duas escholas, citando os argumentos, em que se firmam, e as objecções, com que os adversarios se combatem.\par \emph{A}) \emph{Theoria do concurso simultaneo}. --- Foi um Jesuita hespanhol, \emph{Luiz Molina}, quem propoz esta theoria. Diz assim: « A acção de \textsc{Deus} e a acção da creatura podem comparar-se á acção de dois homens, que puxam uma barca. O movimento todo deriva de um e de outro, embora nenhum dos dois seja a causa total d'elle » (\emph{Concord.}, q. 14, a. 13, dis. 26). Todavia, a acção de \textsc{Deus} e a acção da creatura não são duas acções distinctas. A acção é uma só, produzida contemporaneamente por duas causas \emph{parciaes}, que, unidas, formam uma causa \emph{completa}.\par Como se vê, nesta theoria, \textsc{Deus} não opera \emph{na} creatura, mas com a creatura; não é \emph{motor}, é \emph{cooperador}; não \emph{dirige} o acto da creatura, limita-se a \emph{acompanhal-o}; não é a causa \emph{primeira} do acto, mas \emph{concorre} para elle de um modo \emph{parallelo}; não dá ás causas segundas o impulso ao acto, mas dá, em união com ellas, o ser ao effeito; não é o senhor absoluto dos agentes creados, é o servo sempre prompto a prestar-lhes auxilio. Eis o fraco da theoria: dizem os adversarios. --- Por quanto, se \textsc{Deus} e a creatura não são agentes \emph{subordinados}, mas apenas \emph{unidos} ou \emph{associados}, sem que um dependa de outro, a acção da creatura escapa á causalidade divina e não se poderá dizer que toda a entidade, a qual se encontra na creatura, deriva de \textsc{Deus}. --- Ha mais: A acção de \textsc{Deus} subordina-se, nesta theoria, á acção da creatura. Com effeito, o concurso de \textsc{Deus}, que se suppõe indifferente para este ou aquelle acto, recebe da creatura a propria determinação, sem a qual \textsc{Deus} não poderia operar. Diz \emph{Molina}: « Para nós o concurso divino não determina a vontade ao consenso. Pelo contrario, é o influxo particular do livre arbitrio, que determina o concurso divino ao acto, emquanto a vontade se resolve a querer ou a não querer, a querer isto e não aquillo » (\emph{Concord.}, q. 23, a. 4 e 5, dis. 1). \textsc{Deus}, pois, para operar, deve esperar que o homem o determine para este ou aquelle acto. Mas, como poderá a creatura determinar o concurso divino, se nada pode ella operar sem esse concurso? D'este modo, a liberdade humana, que se pretende salvar, está reduzida á impotencia.\par Em vista d'estas difficuldades, alguns sequazes de \emph{Molina} abandonaram a theoria do concurso \emph{simultaneo}, embora lhe conservassem o nome, e apresentaram uma theoria, que pode chamar-se --- a theoria da \emph{premoção indeterminada}. Á comparação de dois homens, que puxam uma barca, substituiram a de um navio impellido pelo vento para um porto indeterminado.\pgn{534}O seu raciocinio é o seguinte. O vento impelle o navio, por. ex., para o levante, mas não determina o porto. Este é determinado pelo piloto, que governa e dirige o navio. Ora o vento é a imagem da moção divina, que prevem a acção das creaturas e as move para um termo geral; o piloto é a vontade livre, que determina este ou aquelle acto. Fundam esta theoria na auctoridade de S. Thomaz, quando diz --- que \textsc{Deus}, como motor universal, move a vontade do homem para o objecto universal da vontade, que é o bem, --- e que o homem, pela razão, se determina para este ou aquelle bem particular, verdadeiro ou apparente (\emph{Sum. Th.}, 1-2, q. 9, a. 6 ad 3; \emph{De pot.}, q. 1, a. 4 ad 3).\par A theoria da \emph{premoção indeterminada} é superior á do \emph{concurso simultaneo}. Admittindo uma premoção da nossa vontade pela acção divina, ella explica a passagem d'aquella faculdade da potencia para o acto. Todavia, nem esta theoria reconhece uma \emph{total} dependencia do acto da creatura com relação ao acto de \textsc{Deus}. O vento explica o movimento do navio, mas não a sua direcção: esta depende exclusivamente do piloto. Do mesmo modo, se o movimento da vontade humana para o bem em geral depende da moção divina, a \emph{direcção} d'esse movimento para este ou aquelle bem em particular depende da propria vontade. Portanto admitte-se uma realidade positiva, uma verdadeira perfeição (tal é essa direcção), que escapa á causalidade divina, que não deriva da fonte de todo o bem e de toda a perfeição. --- Alem d'isso, tambem esta theoria admitte que o concurso divino, indeterminado por si mesmo, é determinado e especificado pela activa influencia do livre arbitrio humano. Mas como poderá a creatura determinar e especificar com o seu acto o concurso divino, se ella precisa do concurso divino para praticar todo e qualquer acto? --- Por estas e outras razões, muitos philosophos e theologos acceitam a theoria dos \emph{thomistas}.\par \emph{B}) \emph{Theoria do concurso previo}. --- Esta theoria é mais \emph{synthetica} e grandiosa. \textsc{Deus} é o motor primeiro e universal, e nada pode escapar á sua acção e causalidade soberana. Diz S. Thomaz: « Quantumcumque natura corporalis vel spiritualis ponatur perfecta, non potest in suum actum procedere, nisi moveatur a Deo » (\emph{Sum. Th.}, 1-2, q. 109, a. 1). --- « Deus comparatur ad omnia sicut movens ad motum, et activum ad id quod est in potentia » (\emph{C. Gent.}, III, 100). Deus não move, como movem as causas segundas, por um modo exterior e violento; a sua moção realiza-se no intimo dos seres, e adapta-se á natureza e ás disposições de cada um. Diz o mesmo Angelico: « Deus cuilibet agenti adest interius quasi in ipso agens, dum ipsum ad agendum movet » (\emph{Comp. Th.}, c. 120). Por isso a moção divina, longe de ser determinada e especificada pelo concurso activo das creaturas, determina-se e especifica-se por si mesma, conforme a natureza e as disposições do ser, em que é recebida, fazendo com que o ser necessario produza um acto necessario e o ser livre um acto livre. Na efficacia universal e transcendente da causalidade divina encontra-se a razão sufficiente de todos os actos das creaturas, mesmo dos actos livres. Assim como, emquanto é o Ser primeiro, \textsc{Deus} influe continuamente em todas as coisas o \emph{ser}, que é a ultima actualidade na linha do ser, e por este influxo physico e preveniente confere a\pgn{535}existencia, assim tambem, emquanto é a Causa primeira, Elle influe em todos os agentes o \emph{acto}, que é a ultima actualidade na linha da acção, e applica activamente as faculdades operativas aos seus termos connaturaes. Por isso \textsc{Deus} não é apenas o \emph{companheiro} dos agentes finitos; é o \emph{senhor} e \emph{moderador}, que se serve d'elles, como de instrumentos, para a execução das obras da sua providencia; é a Causa primeira, que actua immediatamente, não nos effeitos, mas nas causas segundas, infundindo-lhes força ou vigor, completando-as e applicando-as ao acto; é o agente universal, cuja influencia abrange a natureza inteira, todos os agentes de cada especie, todas as energias de cada agente. --- Esta influencia da Causa primeira, que determina e applica ao acto as causas segundas, é chamada \emph{premoção physica}. E' uma \emph{premoção}, porque a moção da Causa primeira \emph{precede}, pela sua natureza e causalidade, o acto das causas segundas, --- e é uma premoção \emph{physica}, porque é o acto de uma Causa efficiente, que move e applica activamente a causa segunda ao acto connatural. A \emph{premoção} chama-se tambem \emph{predeterminação}, no mesmo sentido e pelo mesmo motivo.\par Os \emph{thomistas} provam a sua theoria com muitos argumentos, baseados nos principios do Angelico Doutor. Expomos os principaes.\par \emph{a}) \textsc{Deus} é o primeiro motor. As creaturas não podem operar senão emquanto são movidas e applicadas por \textsc{Deus}; não só porque toda a passagem da potencia para o acto deve ter e tem a sua razão sufficiente na Causa primeira, mas tambem porque um motor segundo deve ser necessariamente movido pelo motor primeiro. Esta moção de \textsc{Deus} --- é \emph{physica}, porque deriva d'Elle, como da primeira Causa efficiente, e influe physicamente em todas as creaturas, --- e é \emph{preveniente}, isto é, anterior ao acto da creatura, porque a moção do motor precede, na ordem da natureza e da causalidade, a moção do movel: « \emph{Motio moventis praecedit motum mobilis, ratione et causa} » (\emph{C. Gent.}, III, 149). Ora uma moção physica e preveniente é uma \emph{premoção physica}. Logo deve admittir-se esta \emph{premoção}. --- Accrescente-se que, se a causa segunda não é applicada ao acto por uma moção actual e preveniente, mas só recebe de \textsc{Deus} o auxilio de um concurso \emph{simultaneo}, nem ella nem o seu effeito dependem de \textsc{Deus}, como da Causa verdadeiramente primeira. Uma causa não é nem pode dizer-se \emph{primeira} com relação a outra, se não emquanto influe actualmente nella.\par \emph{b}) As causas segundas podem operar e influir no effeito, mas não operam e não influem sempre, e, quando operam, passam da potencia para o acto. Este acto, ou este influxo actual, é uma entidade, uma perfeição, distincta da potencia. Sendo uma entidade, uma perfeição, o acto da creatura deve derivar de \textsc{Deus}, fonte unica de toda a entidade e perfeição. Mas \textsc{Deus} não communica ás causas segundas o poder de influir actualmente no effeito, senão emquanto Elle exerce nessas causas um influxo proprio. Este influxo divino deve preceder naturalmente o influxo das creaturas, como a causa precede o effeito, e por isso é uma verdadeira \emph{premoção physica}. --- Nem se diga que o influxo actual, que as causas segundas exercem no effeito, deriva de \textsc{Deus}, emquanto \textsc{Deus} concedeu a essas causas as faculdades, de que o influxo deriva. Porquanto, \textsc{Deus}, dando ás causas segundas\pgn{536}as faculdades ou energias operativas, deu-lhes o \emph{poder} de operar, mas não lhes deu o \emph{exercicio}, o \emph{influxo actual} d'esse poder. Quer Elle que, em relação a esse \emph{exercicio}, a causa segunda dependa da Causa primeira, como o ser segundo depende do Ser primeiro.\par \emph{c}) E' certo que uma grande parte das causas creadas são por si indifferentes e contingentes, como são os seres dotados de liberdade. Ora uma causa indifferente e contingente, que pode produzir varios actos mas que por si não é determinada para nenhum d'elles, ficará sempre no seu estado de potencialidade e não produzirá acto algum, se não fôr determinada. Esta determinação --- não pode provir da propria \emph{causa}, que é por si indifferente, --- nem do \emph{acto}, que é um effeito, e um effeito não pode determinar a causa, --- nem do \emph{objecto}, porque este, quando se refere a uma causa \emph{livre}, excerce uma moção moral e não physica, deixa a causa na sua indifferença, e, se \emph{especifica} o acto, não determina o \emph{exercicio} d'esse acto. Logo a determinação não pode provir, em ultima analyse, senão da Causa primeira por uma moção anterior ao acto da causa segunda, isto é, por uma \emph{premoção physica}. (\emph{Sum. Th.}, 1-2, q. 9, a. 4).\par \emph{d}) Em todas as faculdades operativas devemos distinguir dois estados --- um \emph{habitual} e outro \emph{actual}. Encontra-se a faculdade no estado \emph{habitual}, ou em \emph{acto primeiro}, quando \emph{pode} ou é \emph{capaz} de operar, mas actualmente não opera; --- encontra-se no estado \emph{actual}, ou em \emph{acto segundo}, quando \emph{actualmente} opera. Logo uma faculdade operativa, para que possa actualmente operar, deve passar do estado habitual para o actual, do acto primeiro para o segundo. Mas não pode passar por si mesma, pela propria energia; porque está em potencia para o acto, e um ente em potencia só pode ser movido por um ente em acto, e porque o estado habitual é menos perfeito que o actual, e nenhum ente pode dar a si mesmo uma perfeição superior á que possue. Logo a faculdade deve passar da potencia para o acto pela energia de um motor, que seja acto purissimo, porque só este pode dar ás forças creadas a ultima actualidade. Diz S. Thomaz: « Omne quod quandoque est agens in actu, et quandoque in potentia, indiget moveri ab aliquo movente » (\emph{Sum. Th.}, 1-2, q. 9, a. 4). --- « Oportet enim ut id quod est in potentia reducatur in actum per aliquid quod est in actu » (\emph{l. c.}, a. 1). Este motor, que completa as faculdades das causas segundas e que é acto purissimo, é o primeiro agente, é \textsc{Deus}. Diz o mesmo S. Doutor: « Complementum virtutis agentis secundi est ex virtute agentis primi » (\emph{C. Gent.}, III, 66). Este complemento, esta ultima actualidade é communicada por \textsc{Deus} por uma moção \emph{transitoria}, que só nElle tem uma existencia fixa (\emph{De pot.}, q. 3, a. 7). Tal moção \emph{transitoria} não pode ser um concurso \emph{simultaneo}; porque o concurso \emph{simultaneo} suppõe o agente no estado \emph{actual}, e não no \emph{habitual}; logo deve ser \emph{previa}, deve ser uma \emph{premoção}.\par \emph{e}) Os defensores do \emph{concurso previo} appellam tambem para a auctoridade do \emph{Catecismo do Concilio de Trento}, que diz: « Non solum Deus universa quae sunt providentia sua tuetur atque administrat, verum etiam \emph{quae moventur et agunt aliquid}, intima virtute, ad motum atque operationem ita \emph{impellit}, ut, quamvis causarum secundarum efficientiam \emph{non impediat}, \emph{praeveniat} tamen; cum ejus \emph{occultissima} vis ad singula pertineat, et,\pgn{537}quemadmodum Sapiens testatur (\emph{Sap.} VIII, 1), \emph{attingat a fine usque ad finem fortiter et omnia disponat suaviter} » (P. I, in I art. Symb., n. 22). --- Trata-se --- de \emph{premoção}, porque \textsc{Deus} \emph{impelle} ao acto e \emph{prevem} a efficiencia das causas segundas, --- e de premoção \emph{physica}, e não de premoção \emph{moral}, porque o \emph{Catecismo} falla de uma \emph{moção} commum a \emph{todos} os agentes, e de uma moção que deriva da omnipotencia divina e consiste numa \emph{virtude occulta} ou \emph{energia intima}, ora a premoção \emph{moral} --- não é \emph{commum} a \emph{todos} os agentes (porque só convem aos agentes \emph{livres}), --- não é attributo proprio da omnipotencia divina, porque o anjo e o homem podem mover moralmente a nossa vontade, --- nem é uma energia occulta ou intima, como é evidente. --- Os \emph{thomistas} adduzem a auctoridade do \emph{Catecismo}, não para elevarem a \emph{premoção physica} a um artigo de Fé, mas para mostrarem a solidez da sua theoria, recebida e considerada pela Egreja como a melhor explicação da omnipotencia divina e para refutarem os que disseram que a \emph{premoção physica} tinha sido condemnada no Concilio de Trento.\par Os defensores do concurso exclusivamente \emph{simultaneo} adduzem varios trechos de S. Thomaz, com que pretendem demonstrar que o Angelico Doutor é contrario á \emph{premoção physica}. --- Os \emph{thomistas} respondem que esses trechos, rectamente interpretados, nada provam contra a \emph{premoção}, e, para essa recta interpretação, apresentam os seguintes criterios: --- 1.º) Quando S. Thomaz affirma que \textsc{Deus} \emph{não predetermina} os nossos actos livres, ou que os actos da nossa vontade \emph{não derivam} de um principio externo \emph{determinante}, o S. Doutor refere-se á predeterminação \emph{natural}, que impõe ás coisas um certo e determinado modo de operar e exclue toda e qualquer liberdade de escolha, mas não exclue a predeterminação propria de um ser livre, que livremente é applicado ao seu acto (\emph{Cf. Sum. Th.}, p. I, q. 32, a. 1 ad 1). --- 2.º) Quando S. Thomaz diz que « \emph{a nossa vontade se move a si mesma, que se determina e é senhora dos seus actos}, etc., » attribue á vontade o que é proprio d'ella, como \emph{agente proximo}, mas não exclue, antes suppõe e exige a moção e a applicação do \emph{primeiro agente} (\emph{Sum. Th.}, p. I, q. 83, a. 1 ad 3). --- 3.º) Quando S. Thomaz ensina que \textsc{Deus} move a vontade, sem deixar nella impressão alguma, não exclue com isso a moção divina, mas apenas um novo estado habitual, uma nova qualidade permanente: porque, quando \textsc{Deus} influe o vigor, a energia na creatura, para ella operar actualmente, a acção divina na creatura é \emph{transitoria}, e por isso não produz um novo estado ou uma nova qualidade (\emph{De pot.}, q. 3, a. 7; --- \emph{De verit.}, q. 22, a. 8). --- 4.º) Quando S. Thomaz sustenta que \textsc{Deus} move a vontade para o bem \emph{universal} e que a propria vontade se move por si mesma para este ou aquelle bem \emph{particular}, refere-se o S. Doutor áquella \emph{primitiva moção especial}, pela qual \textsc{Deus} imprime na vontade humana uma tendencia \emph{necessaria} para o bem \emph{universal}, para a \emph{felicidade}, mas não exclue a divina moção \emph{geral}, que dá á nossa vontade e a todos os agentes o vigor, a energia indispensavel para o exercicio do acto; porque ensina expressamente, em varios logares das suas obras, que \emph{todo} o acto da nossa vontade deriva --- não só da propria vontade, como do agente \emph{immediato} --- mas tambem, e sobretudo, de \textsc{Deus}, como do primeiro agente (\emph{Sum. Th.}, 1-2, q. 109, a. 1; --- \emph{De verit.}, q. 22, a. 8).}.

100. O \textbf{concurso previo de Deus deixa intacta} a \textbf{liberdade da vontade humana}. --- Sendo o homem causa segunda, como o é todo e qualquer agente creado, para que a liberdade da sua vontade \pg{533} \pg{534}  fique intacta, não se requer que a vontade não seja movida de modo algum pela Causa primeira, mas requer-se e basta que essa moção não tire á vontade a sua indifferença, o dominio \pg{535}  de seus actos. Ora a moção divina, sendo \emph{natural} e \emph{intrinseca} ás faculdades das creaturas, não só não destroe o que é proprio d'essas faculdades, mas presta-lhes o auxilio, para que \pg{536}  possam passar da potencia para o acto e operar em conformidade com a sua natureza; de modo que, movendo uma \pg{537}  faculdade \emph{necessaria}, \textsc{Deus} faz com que esta produza um acto \emph{necessario}, e, movendo uma faculdade \emph{livre}, faz com que esta \pg{538}  produza um acto \emph{livre}. Logo o concurso \emph{previo} de \textsc{Deus} deixa intacta a liberdade da vontade humana \footnote{E' o argumento de S. Thomaz, que diz: « Liberum arbitrium est causa sui motus, quia homo per liberum arbitrium seipsum movet ad agendum. Non tamen hoc est de necessitate libertatis quod sit prima causa sui id quod liberum est: sicut nec ad hoc quod aliquid sit causa alterius, requiritur quod sit prima causa ejus. Deus igitur est prima causa movens et \emph{naturales} causas, et \emph{voluntarias}. Et sicut \emph{naturalibus} causis, movendo eas, non aufert quin actus earum sint \emph{naturales}; ita movendo causas \emph{voluntarias}, non aufert quin actiones earum sint \emph{voluntariae, sed potius hoc in eis operatur}; \emph{operatur enim in unoquoque secundum ejus proprietatem} » (\emph{Sum}. \emph{Th}., p. I, q. 83, a. 1 ad 3).\par Esta conclusão está contida nas conclusões precedentes. E' certo que \textsc{Deus} \emph{premove} a nossa vontade; é certo que a nossa vontade é \emph{livre}; logo deve ser tambem certo que a premoção divina deixa intacta a liberdade da nossa vontade. --- Alguns philosophos e theologos reconhecem que \textsc{Deus} \emph{move} a vontade humana, e que a vontade, movida por \textsc{Deus}, permanece livre; mas não admittem a \emph{premoção}, porque dizem que esta é incompativel com a liberdade. Mas não têm razão. Já vimos que, segundo S. Thomaz, toda a \emph{moção} é uma \emph{premoção}; porque a moção do motor \emph{precede} necessariamente, na ordem da natureza e da causalidade, a moção do movel. Logo, se a \emph{premoção} é incompativel com a liberdade, tambem é incompativel a \emph{moção}; e, se não é incompativel esta, nem o pode ser aquella. O concurso \emph{simultaneo} não é verdadeira \emph{moção}; quem corre com outro, nem por isso o move. --- Os adversarios da \emph{premoção} esforçam-se para demonstrar que a nossa vontade, sob o impulso \emph{previo} de \textsc{Deus}, não pode ser \emph{livre}; mas todas as razões, que adduzem, fundam-se no supposto que \textsc{Deus} opere, como operam as creaturas. Porque uma causa finita e extrinseca, predeterminando physicamente a nossa vontade, tira-lhe a indifferença e por isso a liberdade, concluem que se \textsc{Deus} predetermina physicamente ao acto a nossa vontade, tira-lhe tambem a indifferença e a liberdade. Ora tal supposto é falso. Entre a causalidade creada e a causalidade divina poderá haver e ha uma \emph{analogia}, mas não ha uma \emph{egualdade}. A creatura não pode mover \emph{naturalmente} a nossa vontade (porque esta, sendo faculdade de uma substancia espiritual, creada por \textsc{Deus}, escapa ao dominio das causas segundas), e por isso pode exercer nella uma violencia, e tirar-lhe de algum modo a liberdade; mas \textsc{Deus}, sendo o Creador e o Senhor da nossa vontade, pode mover-a \emph{naturalmente}, dando-lhe \emph{uma inclinação} para a coisa, para a qual a move. D'este modo, a moção divina é preveniente e infallivel, e comtudo a vontade humana conserva sempre a sua liberdade (\emph{Sum}. \emph{Th}., p. 1, q. 105, a. 4 ad 1; 1-2, q. 9, a. 6).\par Algumas difficuldades foram já resolvidas. Vamos resolver ainda uma ou outra das que se apresentam com mais força.\par \emph{a}) Dizem: Se a nossa vontade é \emph{premovida} por \textsc{Deus}, ella não se determina a si mesma ao acto. Ora a essencia da liberdade consiste precisamente nesta auto-determinação. Logo a premoção é incompativel com a liberdade.\par Resposta. A premoção divina não só não é contraria a esta auto-determinação, mas é ella mesma que faz com que a vontade se determine a si mesma. Sem a premoção divina, a nossa vontade não poderia exercer nenhum\pgn{539} acto. E' a premoção, que, como o decreto divino, determina que uma causa \emph{livre} exerça, num certo tempo, um certo acto e o exerça \emph{livremente}. \textsc{Deus} é a causa universalissima e efficacissima. A nossa vontade, como diz S. Agostinho (\emph{De corrept. et gratia} c. 14, n. 45), está mais nas mãos de \textsc{Deus} do que nas nossas; e por isso, continua S. Thomaz (\emph{De ver.}, a. 22 a. 8), se a nossa vontade pode, pela virtude propria, passar de um acto para outro diverso, e até contrario, muito mais pode realizar essa passagem pela virtude divina, de que deriva e a que está subordinada toda a virtude creada. Mas \textsc{Deus}, repetimos, actua na nossa vontade, respeitando a dignidade d'ella e inclinando-a intrinsecamente para a coisa, para a qual Elle a premove, de modo que a nossa vontade produz o seu acto, porque assim quer.\par \emph{b}) Continuam: Se \textsc{Deus} \emph{premove} a nossa vontade, esta deve operar \emph{necessariamente}. Ora a necessidade oppõe-se á liberdade. Logo a premoção é incompativel com a liberdade.\par Resposta. Se Deus premove a nossa vontade, a operação d'esta faculdade é \emph{infallivel}, mas não é \emph{necessaria}. E' \emph{infallivel}; porque a premoção de \textsc{Deus}, sendo efficacissima, é seguida pela operação da nossa vontade, e a vontade, quando opera, não pode deixar de operar. Mas não é \emph{necessaria}; porque, como dissemos, a premoção determina a vontade em harmonia com a natureza d'esta faculdade, e por isso não só não lhe tira a liberdade, mas faz com que ella possa operar \emph{livremente}, como se fosse inteiramente independente de todo e qualquer concurso. Diz S. Thomaz: « Deus movet quidem voluntatem \emph{immutabiliter} propter \emph{efficaciam virtutis moventis}, quae deficere non potest; sed propter naturam voluntatis motae, quae indifferenter se habet ad diversa, \emph{non inducitur necessitas, sed manet libertas}; sicut etiam in omnibus providentia divina \emph{infallibiliter} operatur; et tamen a causis contingentibus proveniunt effectus contingenter, inquantum Deus omnia movet \emph{proportionabiliter}, \emph{unumquodque secundum suum modum} » (\emph{De malo}, q. 6, a. u. ad 3).\par c) Replicam: Se a premoção de \textsc{Deus} e a operação da nossa vontade estão infallivelmente unidas, é evidente que a nossa vontade, assim como não tem dominio sobre a premoção divina, assim tambem não o tem sobre a operação propria. Logo a premoção e a liberdade são incompativeis.\par Resposta. E' certo que a nossa vontade não tem dominio sobre o \emph{principio} da premoção, que é o proprio \textsc{Deus}, mas é falso que o não tenha sobre o \emph{termo}, que é o acto da mesma vontade. A nossa vontade, sob a premoção divina, pode suspender, interromper, revogar o proprio acto, sem prejuizo da efficacia da mesma premoção; porque o proprio poder de revogar ou\pgn{540} interromper o acto, que a premoção tinha começado, não pode deixar de ser effeito da premoção. E', pois, sob a causalidade de \textsc{Deus}, motor primeiro, que a vontade passa livremente de um acto, ou de um objecto, para outro acto, ou outro objecto diverso e até contrario. --- Nem isto se oppõe á infallibilidade da Providencia divina. Quando a nossa vontade, dispondo livremente de si mesma, interrompe um acto, para o qual tinha sido determinada pela premoção, o novo acto não se realiza sem uma moção preveniente de \textsc{Deus}, á qual nenhum acto das creaturas, embora livre, pode escapar; de modo que as causas livres concorrem livremente para a execução dos divinos decretos.\par \emph{d}) Insistem: Se \textsc{Deus} premove a nossa vontade para todos e cada um dos actos, deve premovel-a tambem para o peccado, e assim \textsc{Deus} é o auctor do peccado. Ora isto é absurdo e impio.\par Já respondemos a esta difficuldade. O peccado é um \emph{acto defeituoso}. Emquanto \emph{acto}, é uma perfeição e por isso deriva de \textsc{Deus}, como do agente primeiro. Emquanto \emph{defeituoso}, é uma imperfeição, que deve attribuir-se á deficiencia das causas segundas. Assim, o \emph{andar torto} de um coxo, --- emquanto é \emph{andar}, é uma perfeição e deriva da alma, fonte de toda a energia vital, --- emquanto é \emph{torto}, é imperfeição e deriva do defeito da perna. Diz S. Thomaz: « Deus est causa omnis actionis, in quantum est actio. Peccatum nominat \emph{ens} et \emph{actionem} cum quodam \emph{defectu}; defectus autem ille est ex causa creata, scilicet libero arbitrio, in quantum deficit ab ordine primi agentis, scilicet Dei. Unde defectus non reducitur in Deum, sicut in causam, sed in liberum arbitrium; sicut defectus claudicationis reducitur in tibiam curvam, sicut in causam, non autem in virtutem motivam, a qua tamen causatur quidquid est motionis in claudicatione. Et secundum hoc Deus est causa \emph{actus} peccati, non tamen est causa \emph{peccati}, quia non est causa hujus, quod actus sit cum defectu » (\emph{Sum}. \emph{Th}., 1-2, q. 79, a. 2). --- Do resto, esta difficuldade pode oppor-se tambem contra a theoria do concurso \emph{simultaneo}, porque tambem essa theoria admitte que \textsc{Deus} \emph{concorre immediatamente} em todos os nossos actos, e, para o caso, pouco importa que esse concurso seja \emph{previo} ou \emph{simultaneo}.\par Concluimos que, embora a theoria do concurso \emph{simultaneo pareça} explicar ou salvar a liberdade humana melhor do que a do concurso \emph{previo}, comtudo é certo, e todos o reconhecem, que a theoria do concurso \emph{previo} favorece e salva a universalidade e a efficacia da causalidade divina melhor que a outra. Dissemos --- embora \emph{pareça}, porque, na \emph{realidade}, a theoria do concurso \emph{previo}, mais do que a do concurso \emph{simultaneo}, favorece não só os direitos de Deus, mas tambem os da nossa liberdade. Por quanto, a perfeição da creatura consiste na sua união com \textsc{Deus}: e quanto mais intima fôr esta união, tanto mais elevada será a sua perfeição. Ora a união nossa com \textsc{Deus} em nenhum systema se affirma mais intima, como no systema de concurso \emph{previo}, o qual propugna que a nossa vontade não passa da liberdade habitual para a liberdade actual, se não porque ou emquanto é movida pela primeira e increada liberdade.}.

%%>>> art-47

\section*{Artigo IV.}
\addcontentsline{toc}{section}{Artigo IV.}

\pg{539} \textbf{Providencia}, \textbf{sua existencia}, \textbf{extensão}

e \textbf{propriedades}.

101. \textbf{Providencia}. --- \emph{Providencia} é o acto, pelo qual \textsc{Deus} \emph{ordena} e \emph{dirige} todas as coisas para o fim. --- A Providencia, como se vê, importa duas coisas: a \emph{ordenação}, ou \emph{disposição}, \pg{540}  das coisas para o fim, e a \emph{execução} d'essa ordenação. A \emph{ordenação}, por ser acto da intelligencia divina, é \emph{eterna}; a \emph{execução}, por ser realisada no tempo, é \emph{temporanea} \footnote{\textsc{Deus}, como vimos, \emph{creou} todas as coisas; depois de as ter creado, \emph{conserva-as} na existencia, e \emph{concorre} em todos os actos, que ellas produzem.\pgn{541}Mas não basta que \textsc{Deus} concorra em todos os actos das creaturas; é necessario tambem que \emph{regule} esses actos e os \emph{dirija} para o fim, que estabeleceu ás mesmas creaturas. Esta direcção é obra da \emph{providencia}. --- A providencia, como dizemos no texto, abrange duas coisas: a \emph{ordenação} ou \emph{disposição} de todas as creaturas para o fim, e a \emph{execução} d'essa ordenação. A \emph{ordenação}, ou \emph{disposição}, é um acto da \emph{intelligencia} divina, a qual concebe a \emph{idea} ou a razão da ordem das coisas para o fim; e esta ordenação chama-se propriamente \emph{providencia}, que foi por isso definida por \emph{Boecio}: « \emph{divina ratio in summo omnium Principe constituta, quae cuncta disponit} » (\emph{De cons}., l. IV, pr. 6). A \emph{execução} accrescenta ao acto da intelligencia, que estabelece a ordem, o acto da \emph{vontade}, que manda a execução da ordem, e o acto da \emph{potencia}, que executa o mandato da vontade e diz-se \emph{governo}, \emph{direcção}. --- Embora a Providencia divina, considerada no seu complexo, importe uma acção concorde da intelligencia, da vontade e da potencia, comtudo, considerada no seu conceito \emph{formal}, é um acto proprio da \emph{intelligencia}; porque \emph{governar} é acto de \emph{prudencia} e de \emph{sabedoria}, e estas qualidades residem na intelligencia, como no seu sujeito. De facto, é proprio da intelligencia conhecer a relação entre os meios e o fim, e determinar os meios mais efficazes para a consecução do fim. --- A Providencia suppõe nas creaturas uma absoluta dependencia de \textsc{Deus}, um fim proprio de cada uma, e um fim commum e unico, para o qual estão todas ordenadas. --- Emquanto a \emph{ordenação} é uma coisa \emph{eterna}, porque eterno é acto da intelligencia que concebeu a ordem, a \emph{execução} ou o \emph{governo} é uma coisa \emph{temporanea}, porque temporaneas são as creaturas, que Deus dirige para o fim. --- S. Thomaz resumiu estas noções nas seguintes palavras: « \emph{Ad providentiae} curam duo pertinent: scilicet \emph{ratio ordinis}, quae dicitur \emph{providentia} et \emph{dispositio}; et \emph{executio ordinis}, quae dicitur \emph{gubernatio}. Quorum \emph{primum} est \emph{aeternum}; \emph{secundum}, \emph{temporale} » (\emph{Sum}. \emph{Th}., p. I, q. 22, a 1 ad 2). --- Quando o \emph{fim} e os \emph{meios} para o fim \emph{excedem} as forças e as exigencias das creaturas, a providencia é \emph{sobrenatural}; quando o fim e os meios são \emph{proporcionados} ás forças e ás exigencias das creaturas, a providencia é \emph{natural}. Da providencia \emph{sobrenatural} trata a Theologia revelada; aqui occupamo-nos da \emph{natural}. --- S. Thomaz desenvolve esta questão --- na \emph{Sum}. \emph{Th}., p. I, q. 22, --- \emph{c}. \emph{Gent}., l. III, c. 64, --- \emph{De ver}., q. 5.}.

\pg{541} \textbf{102. Deus ordena todas as coisas para um fim.}

a) Todo o agente, dotado de intelligencia, não só opera para um fim, mas tambem dispõe ou ordena todas as suas obras para esse fim. Ora \textsc{Deus} é o agente dotado de intelligencia infinita. Logo \textsc{Deus} ordena todas as coisas para um fim \footnote{O fim, que levou \textsc{Deus} a crear a mundo, é a sua bondade, e o fim, que \textsc{Deus} estabeleceu ás suas creaturas, é tambem a sua bondade. --- As creaturas têm tambem um fim \emph{proximo}, que é a perfeição ou a bondade de cada uma; mas o seu fim \emph{supremo} ou \emph{remoto} é sempre a bondade divina. --- Tendendo para o fim \emph{proximo}, as creaturas tendem para o fim \emph{supremo}.}.

b) Conforme demonstram a experiencia e a razão, cada uma das coisas creadas está ordenada ou disposta para um fim particular e proximo, que é a propria perfeição, e todas ellas estão ordenadas ou dispostas para um fim universal e supremo, \pg{542}  que é a participação e a manifestação da perfeição divina. Ora esta ordenação ou disposição, sendo um bem e um bem das creaturas, deve derivar de \textsc{Deus}, fonte de todo o bem e creador de todas as coisas, e por isso deve existir na intelligencia divina a razão, o plano d'essa ordenação ou disposição. Logo \textsc{Deus} ordena todas as coisas para um fim \footnote{Diz S. Thomaz: « Omne bonum, quod est in rebus, a Deo creatum est. In rebus autem creatis invenitur bonum, non solum quantum ad substantiam rerum, sed etiam quantum ad ordinem earum in finem, et praecipue in finem ultimum, qui est bonitas divina. Hoc igitur bonum ordinis in rebus creatis existens a Deo creatum est. Cum autem Deus sit causa rerum per suum intellectum, et sic cujuslibet sui effectus, necesse est quod ratio ordinis rerum in finem in mente divina praeexistat. Ratio autem ordinandorum in finem proprie \emph{providentia} est » (\emph{Sum}. \emph{Th}., p. I, q. 22, a. 1). --- Não nos demoramos neste ponto, que foi tratado e sufficientemente desenvolvido na \emph{Cosmologia geral} e no I capitulo da \emph{Theodicéa}.}.

103. \textbf{Deus dirige todas as coisas para o fim.}

a) Sendo proprio de um agente perfeito produzir obras relativamente perfeitas, convem á infinita Bondade de \textsc{Deus} governar todas as coisas e dirigil-as para o mais alto grau de perfeição, a que podem chegar. Ora o mais alto grau de perfeição, a que podem chegar as coisas creadas, consiste na consecução do fim. Logo \textsc{Deus} dirige todas as coisas para o fim \footnote{Diz S. Thomaz: « Cum optimi sit optima producere, non convenit summae Dei bonitati quod res productas ad perfectum non perducat. Ultima autem perfectio uniuscujusque est in consecutione finis. Unde ad divinam bonitatem pertinet, ut sicut produxit res in esse, ita etiam eas ad finem perducat: quod est \emph{gubernare} » (\emph{Sum}. \emph{Th}., p. I, q. 103, a. 1).}.

b) Se considerarmos as coisas creadas, veremos que todas ellas, embora diversas, não só estão ordenadas para um fim commum, mas tambem tendem para esse fim e o conseguem pelo uso dos meios mais simples e efficazes. Ora a razão sufficiente d'essa tendencia e d'essa consecução não pode encontrar-se nas energias naturaes das proprias coisas, que são geralmente cegas e se encontram todas no estado de potencialidade. Logo tal razão deve encontrar-se na suprema vontade de \textsc{Deus}, que, executando a ordem por Elle estabelecida, governa e dirige todas as coisas para o fim \footnote{Tambem este argumento encontra-se em S. Thomaz, que diz: « Impossibile est aliqua contraria et dissonantia in unum ordinem concordare semper vel pluries, nisi alicuius gubernatione, ex qua omnibus et singulis tribuitur, ut ad certum finem tendant. Sed in mundo videmus res diversarum naturarum in unum ordinem concordare, non ut raro et a casu, sed ut semper vel in maiori parte. Oportet ergo esse aliquid, cujus providentia mundus gubernatur; et hoc dicimus Deum. (C. \emph{Gent}., l. I, c. 13).}.

\pg{543} \textbf{104. A Providencia extende-se a cada uma das coisas.}

a) Governando \textsc{Deus} as coisas, emquanto foram por Elle creadas, é claro que a sua Providencia se deve extender até onde se extende a sua causalidade. Ora a causalidade divina extende-se a cada uma das coisas, --- não só ás universaes mas tambem ás particulares, não só ás necessarias mas tambem ás contingentes, não só ás corporeas mas tambem ás incorporeas, não só ás substancias mas tambem aos accidentes, --- de modo que não ha coisa alguma, que não seja produzida por \textsc{Deus}. Logo a Providencia extende-se a cada uma das coisas, de modo que não ha e não pode haver coisa, que escape á disposição e ao governo de \textsc{Deus}\footnote{Diz S. Thomaz: « Cum omne agens agat propter finem, tantum se extendit ordinatio effectuum in finem, quantum se extendit causalitas primi agentis... Causalitas autem Dei, qui est primum agens, se extendit usque ad omnia entia, non solum quantum ad principia speciei, sed etiam quantum ad individualia principia, non solum incorruptibilium, sed etiam corruptibilium. Unde necesse est omnia, quae habent quocumque modo esse, ordinata esse a Deo in finem. Cum ergo nihil aliud sit Dei providentia quam ratio ordinis rerum in finem, necesse est omnia, inquantum participant esse, intantum subdi divinae providentiae » (\emph{Sum}. \emph{Th}., p. I, q. 22, a. 2).}.

b) Se a Providencia não se extendesse a cada uma das coisas, isto sería --- ou porque \textsc{Deus} \emph{não pode}, --- ou porque, embora possa, \emph{não quer}. Ora \textsc{Deus} --- \emph{pode,} porque é omnipotente e é capaz de regular e governar tudo com um acto simplicissimo da sua vontade, --- e \emph{quer}, porque um artifice infinitamente bom, como é \textsc{Deus}, não pode deixar de ordenar as suas obras e dirigil-as para o fim, que lhes estabeleceu e que ellas, abandonadas a si mesmas, não podem conseguir. Logo a Providencia extende-se a cada uma das coisas\footnote{O cuidado, que \textsc{Deus} tem de todas as coisas, ainda das que nos parecem mais vis, não é incompativel com a dignidade e grandeza de \textsc{Deus}. Todas as coisas são relativamente boas e perfeitas, porque todas participam da bondade e perfeição divina; e se \textsc{Deus} nada perdeu na sua dignidade, quando as creou, nada pode perder, quando as governa. Antes o contrario. A divina bondade não se revela nunca tão esplendidamente, como quando extende os seus cuidados aos seres mais pequeninos e imperceptiveis da terra. Diz Jesus: « Respicite volatilia coeli, quoniam non serunt, neque metunt, neque congregant in horrea: et Pater vester coelestis pascit illa. Nonne vos magis pluris estis illis? » (\emph{Matth}., VI, 26; cf. \emph{ib}. X, 29-31; \emph{Luc}. XII, 27-29).}.

\textbf{105. A Providencia tem um cuidado especial do homem.}

a) A Providencia de \textsc{Deus}, sendo sapientissima e rectissima, deve ordenar e dirigir todas as creaturas em harmonia --- com a natureza, de que são dotadas, --- e com a ordem dos fins, \pg{544}  que as une e conserva unidas. Ora, o homem é mais nobre do que todas as creaturas racionaes, e, na ordem dos fins, é o fim \emph{proximo} d'essas creaturas; porque as creaturas irracionaes estão ordenadas ou dispostas para a perfeição do homem, e para o seu bem temporal e eterno. Logo a Providencia tem um cuidado especial do homem \footnote{A Providencia é proporcionada á natureza e dignidade dos seres, que lhe estão sujeitos. Nem admira. Um bom pae de familia tem mais cuidado dos filhos do que dos servos, mais dos servos do que dos animaes. Ora o homem é mais nobre do que todas as creaturas visiveis, e é o fim \emph{proximo} de todas ellas. Logo \textsc{Deus} deve ter e tem do homem um cuidado especial. --- E' o que diz o mesmo \textsc{Jesus} nos logares citados.}.

b) O homem, por ser dotado de livre arbitrio, produz actos moralmente bons ou maos, e por isso imputaveis e dignos de premio ou de castigo. Ora a indole de todo o governo deve corresponder á indole da coisa governada. Logo \textsc{Deus} deve governar o homem de modo que este se sirva rectamente do livre arbitrio, praticando o bem e fugindo do mal, e assim alcance o seu ultimo fim. --- De facto, \textsc{Deus} gravou no animo do homem a lei moral, que o dirija em todos os seus actos, e, para a devida sancção d'essa lei, prometteu premios e ameaçou castigos. Logo \textsc{Deus} tem um cuidado especial do homem \footnote{Que \textsc{Deus} dirija e governe todas e cada uma das coisas, sobretudo o homem, com uma sapientissima e efficacissima providencia, prova-o tambem o unanime consenso do genero humano. A Religião, que se encontra em todos os povos, funda-se principalmente na convicção de que existe a Providencia divina. Por isso disse S. Agostinho: « Si Dei providentia non praesidet rebus humanis, nihil est de religione satagendum » (\emph{De util}. \emph{cred}., c. 16). Na verdade, se \textsc{Deus} abandonou todas as coisas creadas ao acaso, não têm razão de ser os juramentos, os sacrificios, os votos, as orações. --- Attesta o mesmo o instincto natural, que, na ameaça de um grande perigo, leva-nos a levantar os olhos, as mãos e a voz para o ceu e a invocar o auxilio divino. --- Demonstram tambem esta verdade as blasphemias dos impios contra \textsc{Deus}; porque os impios não as pronunciariam, se não estivessem convencidos de que nada acontece no mundo, que \textsc{Deus} não queira ou não permitta.\par Concluimos, pois, --- que \textsc{Deus}, principio primeiro e fim ultimo de toda a creação, governa todas e cada uma das creaturas, quer materiaes quer espirituaes, não só em geral e emquanto ás especies, mas tambem em particular e emquanto aos individuos de cada especie; --- que os individuos materiaes e corruptiveis estão ordenados e servem para o bem e a conservação da especie, que se perpetua no meio de uma continua geração e corrupção dos individuos; --- que as creaturas racionaes, como é o homem, são\pgn{545}governadas com uma providencia especial, proporcionada á natureza e á dignidade d'ellas, e por isso por meio de leis, preceitos, conselhos, exhortações, promessas de premios e ameaças de castigos.\par *\par São muitos os adversarios da Providencia divina, não só entre os \emph{antigos}, mas tambem entre os \emph{modernos}. --- Na antiguidade, negaram a existencia da providencia --- os \emph{epicureus}, que chamavam a providencia um jugo insupportavel para o homem, e indecoroso para \textsc{Deus}, --- os \emph{estoicos}, ensinando que todas as coisas, mesmo as nossas vontades e o proprio \textsc{Deus}, estão sujeitas ás leis do \emph{fato}, --- alguns sequazes de \emph{Platão}, que, embora admittissem uma providencia \emph{geral} com relação aos generos, ás especies e ás causas universaes, negavam uma providencia \emph{particular} com relação aos individuos e ás causas particulares.\par Nos tempos modernos negaram a providencia --- \emph{indirectamente} os \emph{pessimistas} e os \emph{fatalistas}, --- \emph{directamente} os \emph{deistas}.\par Entre os \emph{pessimistas} destacam-se \emph{Schopenhauer}, \emph{Hartmann} e os \emph{Neo-Budhistas}. Na opinião de \emph{Schopenhauer}, o mundo é a evolução de uma unica \emph{vontade} cega e maligna, que tudo move de uma maneira absurda. Na opinião de \emph{Hartmann}, o principio da evolução universal é o \emph{ente-inconscio}, cuja maligna vontade é temperada por outro principio, isto é, pelo desejo da anniquilação, e por isso o mundo, embora essencialmente mau e peior que o nada, comtudo não é absolutamente pessimo. Na opinião dos \emph{Neo-Budhistas}, a vida presente, cheia de miserias, é intolleravel, e os homens só podem encontrar a felicidade em \emph{Nirvana}, isto é, na total anniquilação da propria substancia. --- Os \emph{fatalistas} sustentam que todas as coisas estão sujeitas a um \emph{fato} indeclinavel. Patrono do moderno \emph{fatalismo é Taine}.\par Os \emph{deistas}, embora admittam a existencia de \textsc{Deus}, rejeitam, em tudo ou em parte, o seu governo no mundo. Por isso, o seu systema, chamado \emph{deismo}, \emph{é absoluto} e \emph{relativo}. --- Os sequazes do deismo \emph{absoluto} rejeitam totalmente a providencia, porque incompativel com a justiça, bondade e santidade de \textsc{Deus}. Assim dizem \emph{Bayle}, \emph{Voltaire}, \emph{Rousseau}, \emph{Diderot}, \emph{d'Alembert}. --- Os sequazes do deismo \emph{relativo} dizem --- que \textsc{Deus} estabeleceu leis universaes, pelas quaes todos os seres são governados, --- mas que a sua providencia não se extende a todos e a cada um dos seres creados, porque isso repugna á immutabilidade divina e á liberdade humana, --- e que por isso toda a religião revelada é impossivel e a oração, emquanto deprecativa, é inutil e injuriosa para \textsc{Deus}. Estas ideas encontram-se, mais ou menos explicitas, nas obras de \emph{Thiers}, \emph{Cousin}, \emph{J}. \emph{Simon}, \emph{Saisset}, \emph{Damiron}, \emph{G}. \emph{Séailles}.\par Tendo provado, no texto, a \emph{existencia} e a \emph{universalidade} da providencia divina, resta-nos demonstrar a falsidade do \emph{pessimismo} e a inefficacia dos argumentos dos \emph{deistas}.\par O \emph{pessimismo} --- é absurdo nos seus principios, --- é contrario ao senso commum e á natureza, --- e leva ás mais perniciosas consequencias. --- E' absurdo nos seus principios. Por quanto, a unica vontade cega e maligna de \emph{Schopenhauer}, o \emph{ente-inconscio} de \emph{Hartmann}, como o \emph{eu-puro} de \emph{Fichte},\pgn{546}o \emph{absoluto} de \emph{Schelling}, a \emph{idea-ente} de \emph{Hegel}, encerram o mais grosseiro \emph{pantheismo}. Além d'isso, a vida não é puro soffrimento. A vida é, por si, boa: o soffrimento é \emph{accidental}. --- E' contrario ao senso commum; por quanto, todos os homens desejam a vida e temem a morte, e, commummente uma vida miseravel prefere-se á anniquilação da propria substancia. --- E' tambem contrario á natureza; visto que todos os homens, opprimidos pela dor, sentem um desejo innato de uma felicidade futura. --- Leva ás mais perniciosas consequencias; porque o homem, que soffre, se nega a existencia de uma Providencia superior e perde a esperança numa vida futura, não pode ver diante de si senão o desespero e o suicidio, e não pode desejar senão a extincção total do genero humano e a absorpção dos individuos pelo famoso \emph{Nirvana}, isto é, pelo \emph{nada}. Os \emph{pessimistas} admittem estas consequencias.\par Os \emph{deistas} adduzem muitas e varias difficuldades contra a Providencia divina. Expomos e resolvemos as principaes.\par \emph{a}) Dizem: Se todas as coisas fossem governadas pela divina Providencia, não poderia produzir-se nenhum effeito \emph{fortuito} ou \emph{casual}. Ora produzem-se muitos effeitos \emph{fortuitos} ou \emph{casuaes}. Logo nem todas as coisas estão sujeitas á divina Providencia.\par Resposta. Concedemos que haja effeitos \emph{fortuitos} ou \emph{casuaes} com relação ás causas \emph{segundas} e \emph{particulares}, mas negamos que os haja ou possa haver com relação á Causa \emph{primeira} e \emph{universal}. Na verdade, \emph{fortuito} ou \emph{casual} diz-se o effeito que se produz \emph{fóra da intenção do agente}. Ora nenhum effeito pode escapar e escapa á ordem e ao influxo da Causa primeira e universal, que tudo conhece e cuja acção não pode ser estorvada por nenhum outro agente; e por isso, com relação a \textsc{Deus}, nada pode chamar-se fortuito ou casual, pois tudo acontece por divina disposição. Pelo contrario, o effeito de uma causa segunda e particular pode ser estorvado pela acção contraria de outros agentes; e por isso, pode produzir-se um effeito, que é fortuito relativamente á causa particular, porque produz-se fóra da ordem e da intenção d'essa causa, ainda que não se produza fóra da ordem e da intenção da causa universal, porque todos os actos das causas particulares estão contidos na ordem da causa universal, e tudo acontece conforme a disposição d'essa causa universal. (\emph{Sum}. \emph{Th}., p. I, q. 22, a. 2 ad 1).\par \emph{b}) Continuam: Um governador sapiente e bom afasta, por quanto é possivel, das coisas, que rege, todo o mal. Ora no mundo ha muitos males. Logo o mundo não é governado por Deus.\par Resposta. Um governador sapiente e bom afasta das coisas, que rege, o mal da \emph{communidade}, mas não deve afastar o mal do \emph{individuo}, quando este concorre para o bem da mesma communidade. Com effeito, o que constitue o maximo bem do mundo é a ordem ou disposição de todos os seres para o seu fim commum. Este bem não pode, ás vezes, alcançar-se senão com o prejuizo ou damno de um ser particular, que é como um membro da grande totalidade cosmica. N'este caso, o governador sapiente, que olha para o bem da sua communidade, permitte o mal do ser particular, porque e emquanto é ordenado ou concorre para o bem commum. Um governador bom e sapiente só é obrigado a afastar de um ser particular o defeito ou o mal, quando só tem ao seu cuidado esse ser particular. Ora \textsc{Deus} é o\pgn{547}governador \emph{universal} de todas as creaturas, e por isso deve \emph{permittir} o defeito ou o mal de qualquer ser particular, para que não soffra nem seja estorvado o bem perfeito e total do universo. Porquanto, se \textsc{Deus} não permittisse nenhum mal, muito bem deixaria de se dar ou de existir. Por isso dizia S. Agostinho (\emph{Enchir}., c. XI): « Deus omnipotens nullo modo sineret malum aliquod esse in operibus suis, nisi usque adeo esset omnipotens et bonus, ut bene faceret etiam de malo » (\emph{Cf}. \emph{Sum}. \emph{Th}., \emph{ib}., ad 2).\par \emph{c}) Replicam: Se existisse a Providencia divina, não existiriam no mundo os males \emph{physicos}, como a \emph{dór}, a \emph{pobreza}, as \emph{tempestades}, os \emph{terremotos}, etc. Ora estes males existem para a infelicidade do genero humano.\par Resposta. \textsc{Deus}, como dissemos, não quer \emph{directamente} os males \emph{physicos}; pois o mal \emph{physico} é a privação de uma entidade ou perfeição physica, e o termo directo da vontade divina não pode ser uma \emph{privação}. Logo \textsc{Deus} quer esses males \emph{indirectamente}, emquanto os \emph{permitte}. --- Esta permissão não só não exclue, mas até suppõe a providencia divina. Por quanto, o mal physico \emph{particular} --- não só serve para a consecução do bem physico \emph{geral}, --- mas tambem procura ao homem um bem de uma ordem mais elevada, assim a dor purifica e fortifica a nossa alma e é occasião de muitos meritos para a vida eterna.\par \emph{d}) Insistem: No mundo encontra-se o mal \emph{moral}, que é \emph{o peccado}. Ora o peccado é um mal tão grande, que por um governador sapiente e santo, como é \textsc{Deus}, não pode, de maneira nenhuma, ser permittido ou tolerado.\par Resposta. O peccado, embora seja um mal tão grande, que chega a ter uma malicia quasi infinita, porque offensa de uma Bondade infinita, comtudo pode ser permittido. Por quanto, o poder peccar, ainda que não constitua a essencia da liberdade, comtudo é uma consequencia da liberdade, de que é dotada a creatura racional. Ora é conveniente que a Providencia divina governe as creaturas por um modo proporcionado á essencia d'ellas e que as deixe operar em conformidade com a natureza de cada uma, e por isso é conveniente que deixe ao homem o poder de praticar o mal, isto é, de commetter o peccado. --- Além d'isso, Deus não poderia supprimir no homem a faculdade de peccar, sem destruir, ao mesmo tempo, a liberdade. Ora repugna que a Providencia divina possa destruir as qualidades, de que dotou as creaturas. --- Finalmente, esta \emph{permissão} do mal \emph{moral} não repugna nem á \emph{santidade}, nem á \emph{sabedoria}, nem á \emph{bondade} de \textsc{Deus}. Não repugna á \emph{santidade}; porque esta exige que \textsc{Deus} \emph{odeie} o peccado, e não que \emph{impeça} todo o peccado, e porque \textsc{Deus} concede ao homem muitos auxilios para que possa evitar o peccado e exercer a virtude. Não repugna á \emph{sabedoria}; porque o peccado não destroe a ordem \emph{universal}, que consiste na \emph{gloria divina}, mas só destroe a ordem \emph{particular}, que consiste na subordinação de uma creatura ao seu Creador, e esta destruição da ordem \emph{particular} serve para a conservação e perfeição da ordem \emph{universal}, porque tambem os peccadores devem dar gloria a \textsc{Deus}, exaltando a sua \emph{justiça}, se não quizerem exaltar a sua \emph{misericordia}. Não repugna á \emph{bondade}; porque a liberdade, de que os peccadores abusam, é um dom excellente da munificencia divina, é principio de \emph{merito} para a vida presente e de um \emph{premio eterno} na vida futura. O proprio \emph{Rousseau} diz: « Como?! para impedir que o homem seja mau, devia\pgn{548}elle ser rebaixado até á condição dos brutos? Não, Deus da minha alma, eu nunca me queixarei por me teres feito á tua imagem » (\emph{Emile}, l. IV, c. 61).\par \emph{e}) Accrescentam: Embora Deus permitta o peccado, comtudo não devia ser indifferente para com os bons e para com os maus. Ora parece e é indifferente; porque os peccadores levam sempre uma vida feliz, e os justos uma vida cheia de miserias e de amarguras. Esta desegualdade de tratamento mostra que não ha uma \emph{providencia}; pois, se houvesse, seria injusta. --- E' o raciocinio de \emph{Bayle}, \emph{Quinet}, \emph{Havet}, \emph{Michelet}.\par Resposta. E' falso que os peccadores levem sempre uma vida feliz e os justos uma vida miseravel; muitas vezes acontece o contrario. Por isso a pergunta não deve formular-se nestes termos: porque é que os justos soffrem e os peccadores gosam? mas nos seguintes: porque é que os justos e os peccadores estão, promiscuamente, sujeitos aos bens e aos males de vida presente, quando os justos deveriam estar isentos de todos os males, e os peccadores privados de todos os bens? --- Formulada nestes termos, a pergunta pode ter muitas respostas. Antes de tudo, esta sujeição dos justos e dos peccadores aos bens e aos males da vida presente é uma consequencia da condição humana, exposta á acção diversa e contraria das causas segundas. Deus poderia, com certeza, fazer com que os justos não recebessem algum mal e os peccadores não recebessem algum bem; mas é conveniente que o não faça, que não impeça a actividade das causas segundas, sobretudo das livres, e não perturbe a constancia das leis geraes. --- Alem d'isso, se \textsc{Deus} desse, immediatamente, aos justos o premio e aos peccadores o castigo, ficaria destruida toda a ordem moral; porque, então, os homens praticariam o bem não por amor da virtude, mas pela esperança de uma felicidade temporal, e evitariam o mal unicamente pelo medo do castigo. Accrescente-se que o bem e o mal d'esta vida não podem constituir a recompensa adequada da virtude e do vicio. --- Ha mais: essa desegualdade de tratamento, que ás vezes se encontra no mundo, poderia subministrar aos adversarios uma arma efficaz contra a providencia, se esta se extendesse e se limitasse á vida \emph{presente}. Mas não é assim. Porquanto \textsc{Deus} reserva, na vida \emph{futura}, a recompensa \emph{adequada} para a virtude, e o castigo \emph{adequado} para o vicio. Por isso a desegualdade de tratamento na vida presente é o signal mais certo da existencia da providencia. --- De resto, os \emph{males}, de que são affligidos os bons, servem para os purificar dos seus defeitos, para os afastar dos perigos e seducções do mundo, e para lhes subministrar occasião de meritos, --- e os \emph{bens}, que os peccadores gosam ou possuem, são a recompensa de algumas boas obras, que elles não deixam de praticar, e são avisos da bondade divina, que não os quer eternamente perdidos.}.

106. Propriedades da Providencia divina. --- A Providencia divina é \emph{universal}, \emph{suave}, \emph{infallivel}, \emph{immediata}.

\pg{545} a) E' \emph{universal}. Por quanto, como dissemos, a extensão da Providencia de \textsc{Deus} está na razão directa da sua causalidade, e por isso extende-se a tudo o que tem razão de ente.

b) E' \emph{suave}. \textsc{Deus}, sendo o Auctor da propria natureza, governa as coisas conforme a inclinação de cada uma; e por \pg{546} \pg{547} isso cada uma tende, não de um modo violento, mas espontaneo, para o seu fim, para o qual a dirige o impulso divino.

\pg{548} c) E' \emph{infallivel}. Sendo \textsc{Deus} o Agente absoluto e a Causa primeira de todas as coisas, só pode acontecer o que \textsc{Deus} decretou e no modo que \textsc{Deus} decretou, sem prejuizo do livre arbitrio, de que algumas creaturas são dotadas \footnote{« In hoc est \emph{immobilis} et \emph{certus} divinae providentiae ordo, quod ea quae ab ipso providentur, cuncta eveniunt eo modo quo ipse providet, sive \emph{necessario}, sive \emph{contingenter} » (\emph{Sum}. \emph{Th}., p. I, q. 32, a. 4 ad 2). --- E se perguntarmos a S. Thomaz a razão por que tudo acontece conforme os decretos da providencia divina, responde-nos o Angelico, dizendo que toda a inclinação dos seres, quer \emph{natural} quer \emph{voluntaria}, não é outra coisa senão uma\pgn{549}impressão do primeiro motor, como a direcção de uma flecha para o alvo não é senão o effeito da impulsão do archeiro, e que por isso os actos de todos os agentes, quer naturaes quer voluntarios, tendem. e chegam quasi espontaneamente ao termo, que a divina Providencia lhes determinou. Eis as suas palavras: « Omnis inclinatio alicujus rei vel naturalis vel voluntaria, nihil est aliud quam quaedam impressio a primo movente; sicut inclinatio sagittae ad signum determinatum nihil aliud est quam quaedam impressio a sagittante. Unde omnia quae aguntur vel \emph{naturaliter} vel \emph{voluntarie}, quasi \emph{propria sponte} perveniunt in id, ad quod divinitus ordinantur. Et ideo dicitur Deus \emph{omnia disponere suaviter} » (\emph{Sum}. \emph{Th}., p. I, q. 103, a. 8).}.

Resposta. E' falso que os peccadores levem sempre uma vida feliz e os justos uma vida miseravel; muitas vezes acontece o contrario. Por isso a pergunta não deve formular-se nestes termos: porque é que os justos sofrem e os peccadores gosam? mas nos seguintes: porque é que os justos e os peccadores estão, promiscuamente, sujeitos aos bens e aos males de vida presente, quando os justos deveriam estar isentos de todos os males, e os peccadores privados de todos os bens? --- Formulada nestes termos, a pergunta pode ter muitas respostas. Antes de tudo, esta sujeição dos justos e dos peccadores aos bens e aos males da vida presente é uma consequencia da condição humana, exposta á acção diversa e contraria das causas segundas. \textsc{Deus} poderia, com certeza, fazer com que os justos não recebessem algum mal e os peccadores não recebessem algum bem; mas é conveniente que o não faça, que não impeça a actividade das causas segundas, sobretudo das livres, e não perturbe a constancia das leis geraes. --- Alem d'isso, se \textsc{Deus} desse, immediatamente, aos justos o premio e aos peccadores o castigo, ficaria destruida toda a ordem moral; porque, então, os homens praticariam o bem não por amor da virtude, mas pela esperança de uma felicidade temporal, e evitariam o mal unicamente pelo medo do castigo. Accrescente-se que o bem e o mal d'esta vida não podem constituir a recompensa adequada da virtude e do vicio. --- Ha mais: essa desigualdade de tratamento, que ás vezes se encontra no mundo, poderia subministrar aos adversarios uma arma efficaz contra a providencia, se esta se extendesse e se limitasse á vida presente. Mas não é assim. Porquanto \textsc{Deus} reserva, na vida futura, a recompensa adequada para a virtude, e o castigo adequado para o vicio. Por isso a desigualdade de tratamento na vida presente é o signal mais certo da existencia da providencia. --- De resto, os males, de que são affligidos os bons, servem para os purificar dos seus defeitos, para os afastar dos perigos e seducções do mundo, e para lhes subministrar occasião de meritos, --- e os bens, que os peccadores gosam ou possuem, são a recompensa de algumas boas obras, que elles não deixam de praticar, e são avisos da bondade divina, que não os quer eternamente perdidos.

\pg{549} d) E' \emph{immediata}. Por quanto, a \emph{ordem} da Providencia foi estabelecida e predeterminada por \textsc{Deus}, desde a eternidade, sem o concurso das causas segundas; embora, na \emph{execução} d'essa ordem, \textsc{Deus} possa servir-se e se sirva d'essas causas. --- \textsc{Deus} serve-se das causas segundas, não porque lhe falte o poder, mas porque é tão bom, que deseja communicar a dignidade da causalidade tambem ás creaturas \footnote{Diz S. Thomaz: « Ad providentiam duo pertinent: scilicet \emph{ratio ordinis} rerum provisarum in finem, et \emph{executio} hujus ordinis, quae \emph{gubernatio} dicitur. --- Quantum igitur ad \emph{primum} horum, Deus \emph{immediate} omnibus providet; quia in suo intellectu habet \emph{rationem} omnium etiam minimorum, et quascumque causas aliquibus effectibus praefecit, dedit eis virtutem ad illos effectus producendos. Unde oportet quod \emph{ordinem} illorum effectuum \emph{in sua ratione praehabuerit}. --- Quantum autem ad \emph{secundum}, sunt aliqua \emph{media} divinae providentiae; quia inferiora gubernat per superiora, non propter defectum suae virtutis, sed propter abundantiam suae bonitatis, ut dignitatem causalitatis etiam creaturis communicet » (\emph{Sum}. \emph{Th}., p. I, q. 22. a. 3). --- Servindo-se \textsc{Deus} das causas segundas na \emph{execução} da ordem, que Elle, só por si, estabeleceu, é claro que os defeitos e as imperfeições, que se encontram no mundo, não se podem attribuir a \textsc{Deus}, mas unicamente a essas causas segundas. --- \textsc{Deus}, mesmo na \emph{execução} da sua ordem, pode prescindir e prescinde, algumas vezes, do concurso das causas segundas e opera immediatamente por si e fóra da ordem d'essas causas. E' quando faz um \emph{milagre}. Mas, neste caso, \textsc{Deus} opera de um modo \emph{extraordinario}. --- Do \emph{milagre} tratámos, com o devido desenvolvimento, na \emph{Cosmologia}.}.

\textbf{107. A Providencia de Deus não exclue a cooperação do homem.} --- A Providencia de \textsc{Deus} excluiria a cooperação do homem, se \textsc{Deus} fizesse, só por si, as nossas acções. Mas não é assim. \textsc{Deus} não faz, só por si, as nossas acções, mas presta o seu concurso aos que operam, e concede aos que pedem os auxilios necessarios para a consecução do fim. Logo a Providencia de \textsc{Deus} não exclue a cooperação do homem. --- D'ahi a utilidade e a necessidade da \emph{oração} \footnote{A \emph{providencia} de \textsc{Deus} não exclue, mas auxilia os cuidados humanos; assim como o seu \emph{concurso} não destroe, mas auxilia os actos da nossa vontade. Por isso, querendo alcançar um fim, devemos empregar todos os\pgn{550}meios licitos, que estiverem ao nosso alcance, e depois pedir a \textsc{Deus} pelo bom exito das nossas emprezas. A experiencia mostra que, muitas vezes, homens ignorantes e fracos, mas apoiados em \textsc{Deus}, levam ao fim negocios, que homens espertos e poderosos, mas confiados nas proprias forças, não poderam concluir. --- A \emph{immutabilidade} da providencia não exclue a utilidade da \emph{oração}, como não impõe a necessidades ás coisas, por ella governadas. A \emph{oração}, mesmo quando \emph{efficaz}, não muda a disposição eterna da providencia; porque o proprio acto, com que \textsc{Deus} decreta a concessão de uma graça a quem a pede, pertence á mesma ordem da providencia. Diz S. Thomaz: « Orationes apud Deum efficaces sunt, non tamen ordinem immutabilem divinae providentiae solvunt, quia et hoc ipsum quod tali petenti hoc concedatur sub ordine divinae providentiae cadit ». (\emph{C. Gent.}, 1. III, c. 96).}.

\pg{550} \textbf{108. Corollario}. --- Creador e Conservador de todas as coisas. \textsc{Deus} tem sobre ellas um dominio absoluto, inalienavel. Ao dominio de \textsc{Deus} deve corresponder, da parte das creaturas, uma sujeição affectuosa, completa. Nesta sujeição \textsc{Deus} encontra a sua gloria, e a creatura uma felicidade eterna \footnote{Os theologos apontam outro titulo para o dominio de \textsc{Deus} sobre todos os homens: o titulo da \emph{Redempção}. Os homens eram escravos do \emph{peccado}, do \emph{demonio} e do \emph{inferno}. D'esta triplice escravidão resgatou-nos o \textsc{HOMEM-DEUS}, levado unicamente pelo seu amor, e resgatou-nos offerecendo um preço infinito --- o seu \textsc{SANGUE}! Os resgatados da escravidão de um tyranno passam naturalmente para o dominio de quem deu o preço do resgate. Somos todos de \textsc{JESUS}. Conquistado pelo amor, o dominio de \textsc{JESUS} exerce-se pelo amor e só exige uma sujeição de amor. E' o reino do Coração de Jesus! A este dominio e a esta sujeição referem-se as seguintes palavras de S. Paulo: « \emph{Non estis vestri}; \emph{empti enim estis pretio magno}. \emph{Glorificate et portate} \textsc{DEUM} \emph{in corpore vestro} » (I. \emph{Cor}., VI, 19-20).}.

\pg{551} 1. \textbf{Moral}. --- \emph{Moral} é a \emph{sciencia que}, \emph{por meio das regras de moralidade}, \emph{dirige os actos humanos ao nosso ultimo fim}. --- A \emph{Moral} --- é \emph{sciencia}, porque é um conjuncto ordenado de conclusões e principios relativos a um objecto determinado; --- \emph{que dirige os actos humanos}, isto é, os actos que derivam da nossa \emph{intelligencia} e \emph{livre vontade}; --- \emph{ao ultimo fim do homem}, porque a \emph{Moral} não se occupa dos fins \emph{proximos} e \emph{secundarios} dos nossos actos, mas do fim \emph{ultimo}; --- \emph{por meio das regras de moralidade}, porque os actos humanos só podem alcançar o nosso ultimo fim, quando \emph{moralmente} bons e honestos \footnote{A \emph{Moral} é a sciencia dos \emph{costumes} (\emph{mores}). --- \emph{Costume}, em sentido \emph{rigoroso}, é a repetição frequente e uniforme de certos actos, \emph{causada pelo livre arbitrio}. D'ahi a denominação de \emph{moral} applicada ao \emph{acto}, cuja repetição produz o \emph{costume}. Por isso acto \emph{moral} é o que deriva da vontade \emph{livre} do homem. --- A existencia da \emph{liberdade} dos nossos actos é, pois, uma condição necessaria para a existencia da \emph{Moral}.\par D'aqui se segue que a \emph{Moral} dos \emph{pantheistas}, \emph{fatalistas}, \emph{deterministas}, \emph{positivistas}, e, em geral, de todos os que não reconhecem no homem o \emph{livre arbitrio}, tudo poderá ser, menos a sciencia dos \emph{costumes}.\par a) Os \emph{pantheistas} admittem que o homem, com os actos que pratica, e o mundo, com as substancias que o compõem, são emanações \emph{necessarias} da substancia divina. Sendo assim, não ha acções livres no homem, e, não havendo acções livres, não é possivel uma direcção \emph{moral}.\par b) Os \emph{fatalistas}, sustentando que todos os acontecimentos, que se realizam no homem e no mundo, estão sujeitos a uma necessidade fatal e inevitavel, não reconhecem acções \emph{livres}, e por isso não podem deixar de desvirtuar a natureza da \emph{Moral}.\par c) Os \emph{deterministas}, negando de \emph{facto} a \emph{liberdade} da nossa vontade, como provámos na \emph{Anthropologia}, devem tambem negar o objecto da \emph{Moral}.\par d) Os \emph{positivistas} pertencem á classe dos \emph{deterministas}. Na verdade, os \emph{positivistas} admittem --- que a sciencia \emph{moral} é uma parte da \emph{physica social}, que é tambem uma especie da \emph{physica organica}, sujeita a leis \emph{biologicas}, --- que todos os phenomenos humanos são o producto do organismo sob a acção do meio, --- que entre o homem e o animal não existe nenhuma differença\pgn{552}\emph{essencial}. D'estes principios deriva necessariamente a negação do \emph{livre arbitrio}, e por isso a negação da verdadeira Philosophia \emph{moral}. --- Não negamos que os \emph{positivistas} tambem estabelecem uma sciencia \emph{moral}; mas para elles a \emph{Moral} é a \emph{sciencia das condições do aperfeiçoamento das sociedades humanas}, de modo que a \emph{Moral}, no sentido d'esses escriptores, tem uma esphera muito mais extensa do que a que lhe marcou o bom senso, abrangendo não só a educação, a instrucção, a penalidade, mas tambem as bellas artes, as bellas lettras, isto é, abrangendo tudo o que não devia abranger e não abrangendo o que devia abranger, que é a \emph{livre} actividade do homem em ordem ao seu fim. --- Mas, ao menos neste ponto, os \emph{positivistas} são coherentes. Se a alma humana é uma coisa \emph{ficticia}, se \textsc{Deus} é uma coisa \emph{imaginaria}, ou, pelo menos, \emph{incognoscivel}, como é possivel formar a verdadeira sciencia \emph{moral}, que dirija os actos \emph{livres} do homem para \textsc{Deus}, que é o nosso ultimo fim? --- Teremos muitas occasiões de fallar nas doutrinas \emph{positivistas}, e sobretudo na 2.ª parte, em que daremos um resumo da \emph{physica social} ou \emph{sociologia}.}.

\pg{552} 2. \textbf{Objecto da Moral}. --- O \emph{objecto} da \emph{Moral} é constituido pelos actos \emph{humanos}, emquanto, pelas regras da \emph{moralidade}, se hão de dirigir para o nosso \emph{ultimo} fim. --- Por isso a \emph{Moral} --- é uma sciencia \emph{pratica}, porque applica os conhecimentos ás operações, --- e distingue-se das outras duas partes da Philosophia, que são a \emph{Logica} e a \emph{Metaphysica}, porque a \emph{Moral} regula os actos da \emph{vontade}, ao passo que a \emph{Logica} dirige os actos de \emph{intelligencia}, e a \emph{Metaphysica} estuda os entes \emph{reaes} \footnote{A \emph{Moral} distingue-se das outras duas partes da Philosophia, que são a \emph{Logica} e \emph{Metaphysica}. Na verdade, a Philosophia é altissima sabedoria. Ora é proprio da sabedoria estabelecer a ordem entre as coisas. As coisas, que se podem ordenar, --- ou são os actos da nossa \emph{intelligencia}, --- ou são os actos da nossa \emph{vontade}, --- ou são os objectos que compõem o universo. Os actos da nossa \emph{intelligencia} constituem o ente \emph{logico}, que é o objecto da \emph{Logica}, pois esta dirige a nossa intelligencia para a \emph{verdade}; --- os actos da \emph{vontade} formam o ente \emph{moral}, que é o objecto da \emph{Moral}, porque a \emph{Moral} dirige a nossa vontade para o \emph{bem}; --- os objectos, que compõem o universo, constituem o ente \emph{real}, que é o objecto da \emph{Metaphysica}; porque a \emph{Metaphysica} trata dos entes que existem independentemente dos nossos actos. Se, pois, o objecto da \emph{Moral} é distincto do objecto da \emph{Logica} e da \emph{Metaphysica}, é evidente que a primeira é uma sciencia distincta d'estas duas ultimas.\par A \emph{Moral} não só se distingue da \emph{Logica} e da \emph{Metaphysica}, mas tambem da \emph{Theologia moral}. Na verdade, o \emph{sujeito} da \emph{Theologia moral} é o \emph{christão}, o da philosophia \emph{moral} é o \emph{homem}; --- os \emph{principios} da \emph{Theologia moral} são revelados ou deduzidos das verdades reveladas, os da \emph{philosophia moral} são descobertos pela luz natural da razão; --- a \emph{origem} dos preceitos da \emph{Theologia moral} é \emph{divina}, porque derivam da Revelação sobrenatural, ao passo que a origem dos preceitos da philosophia \emph{moral} é \emph{natural}, porque esses preceitos estão baseados na lei natural. --- D'onde se vê o erro de \emph{Kant}, De Wette, \emph{Bruch}, que identificam a philosophia \emph{moral} com a Theologia \emph{moral}.\par Todavia a Theologia \emph{moral} auxilia e illumina a philosophia \emph{moral}, afastando-a dos perigos e guiando-a para um caminho seguro. A historia da philosophia attesta que o homem, abandonado ás suas forças naturaes, tem\pgn{553}descoberto algumas verdades, mas só depois de muito tempo e misturadas com muitos erros. Sirvam de exemplo os escriptos dos maiores entre os philosophos pagãos --- \emph{Socrates}, \emph{Platão}, \emph{Aristoteles}. \emph{Kant} diz: « A maravilhosa religião do Christianismo, na sua grande singeleza, enriquece a philosophia \emph{moral} com ideas moraes mais exactas e mais puras do que as que a mesma philosophia tinha subministrado até essa epocha ». (\emph{Crit. Judic.}, pag. 442).}.

\pg{553} 3. \textbf{Excellencia da Moral}. --- A \emph{Moral} não só é mais excellente que todas as sciencias racionaes e experimentaes, mas tambem é mais nobre que a \emph{Logica} e a \emph{Metaphysica}. --- E' mais nobre que a \emph{Logica}; porque a \emph{Logica} ensina a raciocinar, e a \emph{Moral} a bem viver. --- E' mais nobre que a \emph{Metaphysica}; porque, se esta descobre á nossa intelligencia a natureza do mundo e da nossa alma, a existencia de um \textsc{Deus} infinitamente sabio, poderoso e bom; a \emph{Moral}, applicando á pratica esses conhecimentos, dirige a nossa alma á consecução da felicidade eterna, que consiste na posse do proprio \textsc{Deus} \footnote{D'onde se vê a utilidade, que do estudo da philosophia \emph{moral} deriva para o individuo e para a sociedade. Por quanto, a felicidade do individuo e da sociedade consiste na consecução dos seus fins respectivos. Ora é a \emph{Moral} que aponta aos individuos e á sociedade o \emph{fim}, que se deve conseguir, e os \emph{meios}, que se devem empregar. --- E não só a \emph{Moral} é util para a vida \emph{pratica} do homem e da sociedade, mas tambem para o estudo das outras sciencias praticas, sobretudo do \emph{direito}, da \emph{politica} e da \emph{economia politica}. E' util para o estudo do \emph{direito}; porque as leis \emph{positivas}, de que se occupa o \emph{direito}, nunca poderão ser conhecidas perfeitamente, se não forem relacionadas com a lei \emph{natural}, em que se devem fundar e que forma o objecto da \emph{Moral}. E' util para o estudo da \emph{politica}; porque é a \emph{Moral} que determina o fim do homem e da sociedade, e os deveres dos governantes e dos governados. E' util para a \emph{economia politica}; porque a \emph{Moral} ensina a subordinação do bem \emph{util}, que é o objecto de \emph{economia politica}, ao bem \emph{honesto}, e, definindo as relações entre patrões e operarios, entre amos e servos, lança os fundamentos do bem estar social.} (1).

4. \textbf{Methodo da Moral}. --- O \emph{methodo} da \emph{Moral} --- a) não póde ser exclusivamente \emph{synthetico}, --- b) nem exclusivamente \emph{analytico}, --- c) mas deve ser \emph{synthetico-analytico}.

a) Não póde ser exclusivamente \emph{synthetico}, isto é, \emph{racional}; porque a \emph{Moral} deve considerar os actos \emph{humanos} não em \emph{abstracto} mas em \emph{concreto}, para se conhecer o modo de applicar os principios absolutos e universaes da \emph{moralidade} ás diversas circumstancias da vida humana \footnote{Exaggeram a força da \emph{razão}, no estudo da \emph{Moral}, os sequazes do \emph{direito natural puro}, e os \emph{racionalistas}.\par \emph{a}) Os sequazes do \emph{direito natural puro}, como \emph{Rousseau}, sem se importar da \emph{historia} e da \emph{experiencia}, pretendem deduzir da natureza humana, considerada em \emph{abstracto}, todas as regras de \emph{moral}. --- Mas enganam-se. Se no estudo da \emph{Moral} não tivermos presentes a historia do nosso espirito e o\pgn{554}movimento da epocha actual, perderemos o tempo em questões inuteis ou applicaremos sem tacto e criterio as leis geraes e os principios universaes ás circumstancias individuaes e concretas.\par \emph{b}) Os \emph{racionalistas}, como \emph{Kant}, prescindem inteiramente do auxilio da divina Revelação. --- Tambem estes se enganam. A razão humana, abandonada a si mesma, tem sempre errado, sobre tudo em assumptos moraes e religiosos; e por isso, se quizer encontrar a verdade, deve acceitar, com reconhecimento e veneração, o auxilio de Revelação divina, não para provar as suas proposições, mas para conhecer o caminho da verdade.}.

\pg{554} b) Não póde ser exclusivamente \emph{analytico}, isto é, \emph{experimental}; porque, para dirigir os actos humanos para o seu fim ultimo, é necessario conhecer esse fim, e o fim de um ser conhece-se pelo estudo da sua essencia, e o estudo das essencias exige uma operação da \emph{razão}, que das propriedades e operações se eleve, pelo raciocinio, ao descobrimento do principio, de que derivam \footnote{Attenuam ou rejeitam a \emph{força} da \emph{razão}, no estudo da \emph{Moral}, os \emph{tradicionalistas}, os \emph{historicos}, os \emph{positivistas}, os \emph{scepticos}.\par a) Os \emph{tradicionalistas} dizem que a razão individual não póde conhecer, com \emph{certeza}, nenhuma verdade, sobre tudo de ordem \emph{moral}, e que por isso todas as verdades devem ser ensinadas pela auctoridade, divina ou humana. --- Não nos detemos aqui em refutar o \emph{tradicionalismo}; porque já foi refutado na \emph{Anthropologia}.\par b) Os \emph{historicos} reconhecem a \emph{historia} como a unica fonte das regras de \emph{moralidade}. --- O exagero é manifesto. A \emph{historia}, embora seja de grande auxilio no tratado da \emph{Moral}, não póde comtudo ser a \emph{unica} fonte das regras moraes; porque, se não for estudada á luz de principios sãos, leva facilmente a muitos erros. Não são os principios que se hão de accommodar aos costumes, mas são estes que se devem conformar com aquelles.\par c) Os \emph{positivistas} affirmam que devemos estudar a \emph{Moral}, independentemente da \emph{Metaphysica}, isto é, de todo o \emph{supposto transcendental}, apoiando-nos unica e exclusivamente na \emph{experiencia}. A vida \emph{intellectual} e \emph{moral} do homem, segundo elles, está sujeita, como todas as coisas, a uma \emph{evolução} constante; e por isso a \emph{Moral} é uma parte da \emph{Mechanica} e tira os seus principios da \emph{Physica}, da \emph{Biologia}, da \emph{Psychologia} e da \emph{Sociologia}, que explicam as leis universaes da evolução. D'este modo o munus da \emph{Moral} consiste em applicar essas leis aos phenomenos da vida moral, mostrando --- não o que os homens devem fazer, pois os homens, segundo os \emph{positivistas}, são destituidos de liberdade, --- mas o que \emph{farão}. --- O erro d'estes escriptores é evidente. A ordem \emph{especulativa} não pode deixar de ser o fundamento da ordem \emph{moral}. Se não conhecermos ou se desvirtuarmos as grandes verdades, que se ensinam na \emph{Metaphysica} e que são o principio de \emph{contradicção} e de \emph{causalidade}, a existencia de \textsc{Deus}, a espiritualidade, liberdade e immortalidade da alma humana, o fim ultimo da nossa vida, não é possivel compor um tratado de \emph{Moral}, sem errarmos, em cada passo, ácerca da \emph{lei natural}, da sua obrigação e sancção, ácerca da consciencia, da honestidade, do direito. A \emph{Moral} dos \emph{positivistas} não passa de uma \emph{historia natural}. De resto, na \emph{Anthropologia} refutámos o \emph{determinismo}, que é a base de \emph{Moral} dos \emph{positivistas} e a negação da liberdade, e mais adeante veremos que a nossa vida \emph{moral} não está nem póde estar sujeita a uma evolução fatal e irresistivel.\par d) Os \emph{scepticos} negam que as leis moraes sejam universaes e immutaveis para todos os homens, em todos os tempos e em todos os logares. --- Mais tarde refutaremos esta opinião.}.

\pg{555} c) Mas deve ser \emph{analytico-synthetico}, isto é, \emph{experimental} e \emph{racional}. --- Deve ser \emph{analytico} ou \emph{experimental}; porque a \emph{Moral} dirige os actos humanos, que, por serem \emph{contingentes}, só podem ser conhecidos pela experiencia. --- Deve ser \emph{synthetico} ou \emph{racional}; porque a \emph{Moral} dirige os actos humanos por meio de principios universaes e absolutos, que só podem ser percebidos pela razão \footnote{O methodo \emph{analytico-synthetico} foi seguido por \emph{Socrates}, \emph{Platão}, \emph{Aristoteles}, \emph{S. Agostinho}, \emph{S. Thomaz} e por todos os escholasticos e escriptores de bom senso, que, nas suas especulações, se servem dos principios da razão e da experiencia.}.

5. \textbf{Divisão da Moral}. --- A \emph{Moral} divide-se em duas partes: \emph{geral} e \emph{especial}. --- A Moral \emph{geral} trata dos \emph{actos humanos}, considerando-os no seu \emph{fim}, \emph{essencia}, \emph{moralidade} e \emph{regras}. A Moral \emph{especial} tem por objecto os \emph{direitos} e os \emph{deveres}, que são proprios do homem emquanto é \emph{individuo} e \emph{membro da sociedade}. --- A Moral \emph{geral} costuma chamar-se \emph{Ethica}; --- a Moral \emph{especial} denomina-se commummente \emph{Direito natural} \footnote{Advirta-se que a palavra \emph{ethica} (\emph{etos-costume}) tem a mesma significação que a palavra \emph{moral} (\emph{mos-costume}), e que por isso não é muito rigoroso dar a uma \emph{parte} da divisão o nome que significa o \emph{todo}. Mas é questão de palavras, e accommodamo-nos aos mais auctorizados escriptores.\par Os \emph{philosophos} e os \emph{juristas} discutem se o \emph{direito natural} é, ou não, inseparavel da \emph{ethica}. As opiniões divergem; uns identificam as duas partes, outros separam-nas. --- A questão não póde ser resolvida aqui; visto que, para esta solução, é necessario conhecer as definições da \emph{ethica} e do \emph{direito natural}, e as relações que existem entre uma parte e outra, --- definições e relações, que só serão conhecidas no decurso d'este tratado.}.
